% !TeX program = xelatex
% Standalone English volume. Use XeLaTeX until the references stabilize.
\documentclass[a4paper,11pt]{article}
\usepackage[T1]{fontenc}
\usepackage{lmodern}
\usepackage[english]{babel}
\usepackage[left=26mm,right=26mm,top=25mm,bottom=25mm,headheight=15pt,headsep=7mm,footskip=11mm]{geometry}
\usepackage{amsmath,amssymb,mathtools,amsthm}
\newtheorem{proposition}{Proposition}[section]
\newtheorem{lemma}[proposition]{Lemma}
\newtheorem{theorem}[proposition]{Theorem}
\newtheorem{corollary}[proposition]{Corollary}
\newtheorem{maintheorem}[proposition]{Main Theorem}
\theoremstyle{definition}
\newtheorem{definition}[proposition]{Definition}
\theoremstyle{remark}
\newtheorem{remark}[proposition]{Remark}
\makeatletter
\renewcommand*\l@section{\@dottedtocline{1}{0em}{2.4em}}
\renewcommand*\l@subsection{\@dottedtocline{2}{1.45em}{3.45em}}
\renewcommand*\l@subsubsection{\@dottedtocline{3}{3.8em}{4.55em}}
\makeatother
\usepackage{xcolor,fancyhdr,titlesec,needspace,xurl}
\usepackage{longtable,multicol}
\usepackage{hyperref}
\definecolor{ink}{HTML}{1A3542}
\definecolor{muted}{HTML}{4F606A}
\hypersetup{unicode=true,colorlinks=true,linkcolor=ink,filecolor=ink,urlcolor=ink,citecolor=ink,bookmarksnumbered=true,
 pdfdisplaydoctitle=true,pdflang={en},
 pdftitle={The generalized Riemann hypothesis for Dirichlet L-functions. Edition 930, Volume XXVI},
 pdfsubject={A natural-language reconstruction of the proposed proof},
 pdfauthor={A mathematical reconstruction}}
\titleformat{\section}{\Large\bfseries\color{ink}}{\thesection}{.65em}{}
\titlespacing*{\section}{0pt}{.8em}{.65em}
\setlength{\parindent}{1.2em}
\linespread{1.045}
\setlength{\parskip}{.22em plus .05em minus .03em}
\setlength{\emergencystretch}{2em}
\setcounter{topnumber}{2}
\setcounter{tocdepth}{3}\setcounter{secnumdepth}{3}
\widowpenalty=4500
\clubpenalty=4500
\pagestyle{fancy}\fancyhf{}
\fancyhead[L]{\small\color{muted}Dirichlet GRH \textperiodcentered\ Volume I}
\fancyhead[R]{\small\color{muted}Dirichlet zeros and power intervals}
\fancyfoot[C]{\small\thepage}
\renewcommand{\headrulewidth}{.3pt}
\newcommand{\lead}[1]{\par\smallskip\noindent{\bfseries\color{ink}#1.}\enspace}
\newcommand{\sour}[2]{\unskip\penalty150\hspace{.35em}{\small\hyperlink{ref#1}{[#2]}}}
\titleformat{\subsection}{\large\bfseries}{\thesubsection}{.6em}{}
\titlespacing*{\subsection}{0pt}{1.15em plus .15em}{.4em}
\setlength{\abovedisplayskip}{9pt plus 2pt minus 3pt}
\setlength{\belowdisplayskip}{9pt plus 2pt minus 3pt}
\newcommand{\RR}{\mathbb R}\newcommand{\CC}{\mathbb C}\newcommand{\QQ}{\mathbb Q}
\newcommand{\ZZ}{\mathbb Z}\newcommand{\ord}{\operatorname{ord}}
\newcommand{\Frob}{\operatorname{Frob}}\newcommand{\Gal}{\operatorname{Gal}}
\newcommand{\len}{\operatorname{length}}\newcommand{\Res}{\operatorname*{Res}}
\newcommand{\wmean}{\operatorname{av}}\newcommand{\Ker}{\operatorname{ker}}
\newcommand{\Lh}{\mathcal L}\newcommand{\Gg}{\mathcal G}
\newcommand{\real}{\operatorname{Re}}\newcommand{\imag}{\operatorname{Im}}
\newcommand{\abs}[1]{\left|#1\right|}\newcommand{\eps}{\varepsilon}
\newcommand{\anchoritem}[2]{\par\Needspace{2.8\baselineskip}\hypertarget{ref#1}{}\noindent{\bfseries\color{ink}#2.}\enspace}
\newcommand{\srcloc}[2]{{\small\nolinkurl{#1.lean}}, line~#2}
\numberwithin{equation}{section}
\raggedbottom

\newcommand{\Sel}{\mathcal S}
\newcommand{\Cl}{\operatorname{Cl}}
\newcommand{\Aut}{\operatorname{Aut}}
\newcommand{\GL}{\operatorname{GL}}
\newcommand{\OO}{\mathcal O}
\hypersetup{bookmarksdepth=3}


\addto\captionsenglish{\renewcommand{\partname}{Volume}}


\makeatletter
\expandafter\def\csname r@cor:fieldrelations712\endcsname{{12.4}{32}{}{proposition.12.4}{Dirichlet-GRH-v930-Volume-II.pdf}}
\expandafter\def\csname r@cor:integratedmark873\endcsname{{134.13}{98}{}{proposition.134.13}{Dirichlet-GRH-v930-Volume-XXV.pdf}}
\expandafter\def\csname r@cor:oddwindow677\endcsname{{11.30}{27}{}{proposition.11.30}{Dirichlet-GRH-v930-Volume-II.pdf}}
\expandafter\def\csname r@cor:realmarked682\endcsname{{131.2}{16}{}{proposition.131.2}{Dirichlet-GRH-v930-Volume-XXV.pdf}}
\expandafter\def\csname r@cor:squaredamplitudetail700\endcsname{{132.18}{64}{}{proposition.132.18}{Dirichlet-GRH-v930-Volume-XXV.pdf}}
\expandafter\def\csname r@eq:linkenergy869\endcsname{{134.6}{83}{}{equation.134.6}{Dirichlet-GRH-v930-Volume-XXV.pdf}}
\expandafter\def\csname r@eq:orbitfactors871\endcsname{{134.10}{89}{}{equation.134.10}{Dirichlet-GRH-v930-Volume-XXV.pdf}}
\expandafter\def\csname r@eq:spheremean872\endcsname{{134.13}{94}{}{equation.134.13}{Dirichlet-GRH-v930-Volume-XXV.pdf}}
\expandafter\def\csname r@lem:calibratedlevel703\endcsname{{133.8}{69}{}{proposition.133.8}{Dirichlet-GRH-v930-Volume-XXV.pdf}}
\expandafter\def\csname r@lem:coefficientdecay697\endcsname{{132.8}{57}{}{proposition.132.8}{Dirichlet-GRH-v930-Volume-XXV.pdf}}
\expandafter\def\csname r@lem:cosetpacking699\endcsname{{132.13}{60}{}{proposition.132.13}{Dirichlet-GRH-v930-Volume-XXV.pdf}}
\expandafter\def\csname r@lem:massscale869\endcsname{{134.3}{82}{}{proposition.134.3}{Dirichlet-GRH-v930-Volume-XXV.pdf}}
\expandafter\def\csname r@lem:pivots871\endcsname{{134.7}{88}{}{proposition.134.7}{Dirichlet-GRH-v930-Volume-XXV.pdf}}
\expandafter\def\csname r@lem:unipotentvolume703\endcsname{{133.9}{70}{}{proposition.133.9}{Dirichlet-GRH-v930-Volume-XXV.pdf}}
\expandafter\def\csname r@nav:marked-open-closed\endcsname{{131}{14}{}{section.131}{Dirichlet-GRH-v930-Volume-XXV.pdf}}
\expandafter\def\csname r@nav:poincare-spanning\endcsname{{133}{64}{}{section.133}{Dirichlet-GRH-v930-Volume-XXV.pdf}}
\expandafter\def\csname r@nav:spanning-and-sphere\endcsname{{134}{79}{}{section.134}{Dirichlet-GRH-v930-Volume-XXV.pdf}}
\expandafter\def\csname r@nav:stationary-cutoff\endcsname{{132}{51}{}{section.132}{Dirichlet-GRH-v930-Volume-XXV.pdf}}
\expandafter\def\csname r@nav:two-mark-flags\endcsname{{130}{3}{}{section.130}{Dirichlet-GRH-v930-Volume-XXV.pdf}}
\expandafter\def\csname r@prop:averagedschur684\endcsname{{131.6}{19}{}{proposition.131.6}{Dirichlet-GRH-v930-Volume-XXV.pdf}}
\expandafter\def\csname r@prop:localpivot873\endcsname{{134.12}{97}{}{proposition.134.12}{Dirichlet-GRH-v930-Volume-XXV.pdf}}
\expandafter\def\csname r@prop:stationaryweights695\endcsname{{132.1}{51}{}{proposition.132.1}{Dirichlet-GRH-v930-Volume-XXV.pdf}}
\expandafter\def\csname r@sec:bothmarks873\endcsname{{134.6}{96}{}{subsection.134.6}{Dirichlet-GRH-v930-Volume-XXV.pdf}}
\expandafter\def\csname r@sec:discriminant681\endcsname{{130.4}{11}{}{subsection.130.4}{Dirichlet-GRH-v930-Volume-XXV.pdf}}
\expandafter\def\csname r@sec:gramtrace683\endcsname{{131.2}{16}{}{subsection.131.2}{Dirichlet-GRH-v930-Volume-XXV.pdf}}
\expandafter\def\csname r@sec:missing\endcsname{{6.6}{32}{}{subsection.6.6}{Dirichlet-GRH-v930-Volume-I.pdf}}
\expandafter\def\csname r@sec:orbitfactors871\endcsname{{134.4}{88}{}{subsection.134.4}{Dirichlet-GRH-v930-Volume-XXV.pdf}}
\expandafter\def\csname r@sec:rationalflags680\endcsname{{130.3}{8}{}{subsection.130.3}{Dirichlet-GRH-v930-Volume-XXV.pdf}}
\expandafter\def\csname r@sec:realgram682\endcsname{{131.1}{14}{}{subsection.131.1}{Dirichlet-GRH-v930-Volume-XXV.pdf}}
\expandafter\def\csname r@seriesguide\endcsname{{1.1}{4}{}{subsection.1.1}{Dirichlet-GRH-v930-Volume-I.pdf}}
\expandafter\def\csname r@thm:bothmarks873\endcsname{{134.11}{96}{}{proposition.134.11}{Dirichlet-GRH-v930-Volume-XXV.pdf}}
\expandafter\def\csname r@thm:globalamplitudetail700\endcsname{{132.17}{63}{}{proposition.132.17}{Dirichlet-GRH-v930-Volume-XXV.pdf}}
\expandafter\def\csname r@thm:norm91\endcsname{{94.9}{44}{}{proposition.94.9}{Dirichlet-GRH-v930-Volume-XVIII.pdf}}
\expandafter\def\csname r@thm:normalizedcut699\endcsname{{132.14}{61}{}{proposition.132.14}{Dirichlet-GRH-v930-Volume-XXV.pdf}}
\expandafter\def\csname r@thm:orbitfactors871\endcsname{{134.8}{89}{}{proposition.134.8}{Dirichlet-GRH-v930-Volume-XXV.pdf}}
\expandafter\def\csname r@thm:realgram682\endcsname{{131.1}{14}{}{proposition.131.1}{Dirichlet-GRH-v930-Volume-XXV.pdf}}
\expandafter\def\csname r@thm:spheremean872\endcsname{{134.10}{94}{}{proposition.134.10}{Dirichlet-GRH-v930-Volume-XXV.pdf}}
\makeatother
\begin{document}
\setcounter{part}{25}
\fancyhead[L]{\small\color{muted}Dirichlet GRH \textperiodcentered\ Volume XXVI}
\fancyhead[R]{\small\color{muted}Local Bessel correlations}
\thispagestyle{plain}\begin{center}
{\LARGE\bfseries The generalized Riemann hypothesis\\[2pt]for Dirichlet \(L\)-functions}\\[6pt]
{\small 27 November 2026 \quad\textperiodcentered\quad Edition 930}\end{center}
\part{Local Bessel correlations and relative Gram frequencies}\label{part:XXVI}
\begin{abstract}\noindent We compute position integrals for marked vectors and bound the first relative Gram transform at split and inert finite places. The rational sum puts these estimates on the number fields and intervals of the projected period. Its real factor must be estimated at the corresponding archimedean frequencies.\end{abstract}
\setcounter{section}{135}
\noindent The \hyperref[seriesguide]{guide in Volume~I} lists all twenty-seven volumes.
\begingroup\small\setlength{\parskip}{0pt}\tableofcontents\endgroup
\clearpage

\section{Fixed vectors and local position integrals}\label{nav:marked-orbital-energy}\label{nav:volume-XXVI}

\subsection{A fixed frame vector at a split mark with unit pivots}
\label{sec:selectedopen874}

Use the fields, the shorter interval, the two old primes and the measures of
Theorem~\ref{thm:bothmarks873}. At a marked place \(w\), let \(O_w\) be the
set of residues with all the arm pivots units in both components. For
\(J\subseteq\Xi\), let \(E(\Xi;J)\) be the energy in
\eqref{eq:linkenergy869} with the prescribed vector fixed at each place of
\(\Xi\), and with \(s_w\in O_w\) at every unselected place and at every
place of \(J\). Here and below the constraint on \(z(s)\) has been dropped.
All these energies are sums of nonnegative squared norms indexed by
residues and links. Write \(E(\Xi)=E(\Xi;\varnothing)\). We use the notation
\(\Xi\setminus J\) also for the induced choice of vectors at the places left
in \(\Xi\).

\begin{theorem}[Integrated fixed-vector factor on the open cell]\leavevmode
\label{thm:selectedopen874}
Fix \(w\in J\), put \(q=q_w\), and write the chosen vector in its prescribed
residue component as \(d=\ell+c z\ne0\), with \(\ell\in L\) in that component
and \(c\in\mathbb F_q\). Let \(\theta\) be the terminal entry in the sphere
formula \eqref{eq:spheremean872} and put \(\beta=(\theta\bar\theta-N^2)/N\).
If the unrestricted local factor of Corollary~\ref{cor:integratedmark873}
vanishes for a pair of heads, the restricted factor also vanishes. Otherwise
\(\theta,\beta\) are integral at \(w\). The restricted factor is the
unrestricted one multiplied by
\begin{equation}
 R_w(d,\theta)=
 \begin{cases}
 q^{-2a+1-k}\,\delta_w(\beta/q)/\delta_w(\beta),
 &c\ne0,\quad \theta_{\rm sel}\equiv-N\pmod q,\quad
                         \beta\in q\mathcal O_{F_w},\\
 0,&\text{otherwise}.
 \end{cases}
 \label{eq:fixedratio874}
\end{equation}
Here \(\theta_{\rm sel}\) is the component of \(\theta\) in the same
copy of the residue field as \(d\). In particular \(0\le R_w\le
q^{1-k-2a}\), uniformly in the terminal entry and both heads.
\end{theorem}
\begin{proof}
In the open cell the two contributing heads must be units with the
same reductions, and the middle matrix \(x\) must be integral.
For such heads the two character factors cancel on each residue by
Proposition~\ref{prop:localpivot873}. With \(\bar x\) and the pivot
reductions fixed, the possible residues are parametrized, each once, by
the two independent unipotents in
\(s=\bar u\,w_{\bar{\mathbf p}}\bar x\,\bar v^{-1}\).
Consequently the restriction amounts to the probability that \(s d=d\)
for independent uniform \(\bar u,\bar v\).

Work in the selected component, using the paired dual frame as the
basis if it is the second component. The split group in this basis is
a general linear group with the flag of \(L\). The radical preserves
\(L\), acts trivially on the middle quotient, and contains every shear
from the middle space to \(L\). In the other component the inverse
adjoint supplies the paired action, so the use of either component gives
the full group of shears. For \(c=0\), the nonzero vector
\(w_{\bar{\mathbf p}}\bar x\,\bar v^{-1}d\) belongs to the dual arm,
whereas \(\bar u^{-1}d\) belongs to \(L\). Equality is impossible.
For \(c\ne0\), both \(\bar u^{-1}d\) and \(\bar v^{-1}d\) are uniform
on \(cz+L\), independently. Indeed the shears act transitively on
this affine space, so each point has the same number of preimages.
Write these points as \(cz+\ell_1,cz+\ell_2\). The equation
\(cz+\ell_1=w_{\bar{\mathbf p}}\bar x(cz+\ell_2)\) equates a vector in
\(L\oplus W\) to one in \(W\oplus L^\vee\). It forces
\(\ell_1=\ell_2=0\) and \(\bar x z=z\) in the chosen component. These
three conditions are also sufficient. The residue count is thus the
unrestricted count times \(q^{-2a}{\bf1}_{\bar x z=z}\).

On the integral fiber in \eqref{eq:spheremean872},
\(xz=\omega-(\theta/N)z\). The extra middle condition is constant on
the fiber and says that \(\omega_{\rm sel}\equiv0\) and
\(\theta_{\rm sel}\equiv-N\). The volume of the fiber of integral
isometries remains \(|\mathrm{GL}_k(\mathbb F_q)|\). It remains to
compute the measure of integral \(\omega\) with the indicated
congruence. The split, unimodular middle lattice has coordinate
vectors \(v,\lambda\in\mathcal O_{F_w}^{\,k}\) in dual bases, so
\(Q(\omega)=v\cdot\lambda\). Use \(v\) for the selected component,
interchanging the names in the other orientation. If
\(v\in q\mathcal O_{F_w}^{\,k}\), the pairing belongs to
\(q\mathcal O_{F_w}\). For \(\beta\in q\mathcal O_{F_w}\), the change
\(v=qv'\) multiplies the ambient measure by \(q^{-k}\) and the
target measure by \(q^{-1}\). Equivalently one may count preimages
of \(\beta+q^h\mathcal O_{F_w}\) and divide by \(q^{-h}\) for
\(h>\operatorname{ord}_w\beta\). The restricted sphere measure is
therefore \(q^{1-k}\delta_w(\beta/q)\), giving
\eqref{eq:fixedratio874}.

For completeness, if \(j=\operatorname{ord}_w\beta\ge0\), then
\[
 \delta_w(\beta)=(1-q^{-k})\sum_{i=0}^{j}q^{-(k-1)i}.
\]
In fact the set of \(v\) of minimum coordinate valuation \(i\) has
measure \(q^{-ki}(1-q^{-k})\). For each such \(v\), the map
\(\lambda\mapsto v\cdot\lambda\) sends probability Haar measure to
probability Haar measure on \(q^i\mathcal O_{F_w}\), which has
density \(q^i\) on that ideal. This proves the formula and
\(\delta_w(\beta/q)\le\delta_w(\beta)\) when \(j\ge1\). The terminal
spheres in question have \(\beta\ne0\) by Lemma~\ref{lem:pivots871}.
\end{proof}

\begin{corollary}[Reduction to singular minors at every selected mark]
\leavevmode\label{cor:singularselected874}
Let \(E_{\rm sing}(\Xi)\) be \(E(\Xi)\) restricted to residues
\(s_w\notin O_w\) at every selected \(w\), and set \(B_0=b^{6k}\).
For the choices \(|\Xi|<f/4\) at each sufficiently large node,
\begin{align}
 E(\Xi;J)&\le\prod_{w\in J}q_w^{1-k-2a}\,E(\Xi\setminus J),
 \label{eq:opensave874}\\
 \max_\Xi B_0^{|\Xi|} E(\Xi)
 &\le (1-2f b^{-4k})^{-1}
                 \max_\Xi B_0^{|\Xi|}E_{\rm sing}(\Xi).
 \label{eq:singreduce874}
\end{align}
In particular the right side of \eqref{eq:singreduce874} bounds
\(\max_\Xi B_0^{|\Xi|}\mathcal E_\Xi^\sharp\) for the original
cutoff in Theorem~\ref{thm:bothmarks873}.
\end{corollary}
\begin{proof}
Fix the terminal entry in \eqref{eq:spheremean872}. At every place
of \(J\), both the fixed and the unrestricted residue sums impose
the same unit and congruence requirements on the two heads.
The ratio \eqref{eq:fixedratio874} is independent of the heads.
For a nonintegral terminal entry at any of these places both
contributions are zero. Otherwise the entire contribution for the
restricted mask at this terminal entry equals the contribution for
\(\Xi\setminus J\) times \(\prod_{w\in J}R_w(d_w,\theta)\).
This is an identity for the full quadratic expression. The
unrestricted expression is nonnegative for each terminal entry by
its sphere integral of squared norms. We may therefore apply the
uniform upper bounds for the ratios and then sum in \(\theta\).
This proves \eqref{eq:opensave874}.

In \(E(\Xi)\), a residue outside the all-singular part belongs to
\(O_w\) for at least one selected \(w\). Nonnegativity before the
expansion in heads bounds its total contribution by
\(\sum_{w\in\Xi} E(\Xi;\{w\})\). Since \(q_w\ge b^{10}\) and \(a\ge1\),
we have \(B_0q_w^{1-k-2a}\le b^{-4k}\). After weighting the
energies by \(B_0^{|\Xi|}\), each term in this sum is therefore at
most \(b^{-4k}\max_{\Xi'} B_0^{|\Xi'|}E(\Xi')\).
There are at most \(2f\) summands, and \(2f b^{-4k}<1\) on the stated
interval. Taking a maximum and rearranging gives
\eqref{eq:singreduce874}. Restoring the discarded condition on \(z(s)\)
can only decrease the sum of squared norms.
\end{proof}

For the empty choice, \(E_{\rm sing}(\varnothing)\) includes every
residue with unit pivots at both marks. The remaining maximal energy
thus includes this term as well as the terms having a singular
minor at each selected place.
\clearpage

\subsection{Summation of rational frequencies on the shorter interval}
\label{sec:freqnode875}

At a node let \(E\) be the CM field, \(f=[E:\mathbb Q]/2\), and let
\(\Phi_v\) range over one embedding from each conjugate pair. For a finite
set \(T\) of primes of \(E\), write \(\mathcal O_{E,T}\) for the ring of
elements integral outside \(T\). If \(y\in E^\times\), define
\[
 I_T(y)=\prod_{\mathfrak p\in T}
      (\mathrm N\mathfrak p)^{\operatorname{ord}_{\mathfrak p}y},
 \qquad
 D_T(y)=\prod_{\mathfrak p\in T}
      (\mathrm N\mathfrak p)^{\max(0,-\operatorname{ord}_{\mathfrak p}y)}.
\]
For a nonzero vector of scales \(c=(c_v)_v\) with every \(c_v\ne0\),
let \(g=(\prod_v|c_v|)^{1/f}\). With \(K\ge0\) define the nonnegative
weight on nonzero \(T\)-integers
\begin{equation}
 W_{T,c,K}(y)=
 \frac{I_T(y)^{K/2}}{D_T(y)^3}
 \prod_v |c_v\Phi_v y|^{-K}
                    \min\{1,|c_v\Phi_v y|^{-6}\}.
 \label{eq:freqweight875}
\end{equation}

\begin{lemma}[Summability with unequal archimedean scales]\leavevmode
\label{lem:freqheight875}
There is an absolute \(C\ge1\) such that, for every \(E,T,c,K\) as above,
\[
 \sum_{0\ne y\in\mathcal O_{E,T}} W_{T,c,K}(y)
       \le C^f\,g^{-Kf}\max\{1,g^{-1}\}^{6f}.
\]
\end{lemma}
\begin{proof}
Write \((y)=\mathfrak n/\mathfrak d\) with coprime integral ideals.
Since \(y\) is integral outside \(T\), \(D_T(y)=\mathrm N\mathfrak d\).
The factors in the norm of \((y)\) at primes outside \(T\) are at least
one, so \(I_T(y)\le\prod_v|\Phi_v y|^2\). This cancels the power \(K\)
of the embeddings and leaves \(g^{-Kf}\).
Put \(u_v=|c_v|/g\). Then \(\prod_vu_v=1\), and
\[
 \min\{1,|c_v\Phi_v y|^{-6}\}
   \le \max\{1,g^{-1}\}^6\max\{1,|u_v\Phi_v y|\}^{-6}.
\]
It suffices to sum the product on the right together with
\((\mathrm N\mathfrak d)^{-3}\).

For fixed \(\mathfrak d\), all the points lie in
\(\mathfrak d^{-1}\), which is a union of \(\mathrm N\mathfrak d\)
cosets of \(\mathcal O_E\). The index follows by factoring \(\mathfrak d\)
into prime ideals and taking local indices in the corresponding discrete
valuation rings. In each coset embed \(y\) as \((u_v\Phi_v y)_v\).
A product of half-open squares of side \(1/2\) in these complex planes
contains at most one point of the coset. Indeed the difference of two
distinct points comes from a nonzero algebraic integer, so the product
of its squared absolute coordinates is an integer at least one, even
after the rescaling. Two points in the same box would have each
coordinate difference of modulus at most \(1/\sqrt2\).

The sum over all squares of the supremum of \(\max(1,|z|)^{-6}\)
on the square is some absolute constant \(C_0\). On the square with
lower corner \((h/2,j/2)\), where \(t=\max(|h|,|j|)\ge3\), the
supremum is at most \((2/(t-1))^6\), and there are \(8t\) squares
of that order. Taking the product of the sums gives \(C_0^f\) per
coset. Summing over all denominators costs at most
\(\sum_{\mathfrak d}(\mathrm N\mathfrak d)^{-2}
 \le\zeta(2)^{2f}\).
For the inequality, each rational prime has at most \(2f\) prime ideals
above it and each has norm at least that prime. Unique ideal
factorization then bounds each Euler factor. The nonnegative series
may be summed in any order. This proves the lemma with
\(C=C_0\zeta(2)^2\).
\end{proof}

\begin{theorem}[The two frequency scales at an actual node]\leavevmode
\label{thm:freqnode875}
Use \(m=2n\), \(K=k-n^7\), \(M=p_0p_1\), and the different generator
\(\delta_E\) of Corollary~\ref{cor:realmarked682}. Put
\[
 \kappa_v=\frac{M|\Phi_v\delta_E|}{2\pi},\qquad
 b_{j,v}=\frac{\kappa_v^{2(n-j)}}{N}\quad(1\le j\le n),\qquad
 Q_E=\frac{M\rho_E}{2\pi}.
\]
For either of the choices
\[
   c_v=\frac{M}{b_{j,v}}
   \qquad\hbox{or}\qquad
   c_v=\frac{2\pi i}{\Phi_v(\delta_E)b_{j,v}},
\]
for every \(j\) and every finite \(T\), one has
\begin{equation}
 \sum_{0\ne y\in\mathcal O_{E,T}} W_{T,c,K}(y)
 \ \le\ C^f
       \left(\frac{b^{15/2}}{M w^m}\right)^{K f}
 \ \le\ b^{-12kf}
 \label{eq:actualfreq875}
\end{equation}
at all sufficiently large nodes of the same shorter interval.
Here \(w\) is the geometric mean of the weights in
Lemma~\ref{lem:calibratedlevel703}. The constants and the threshold
are independent of \(T,j\) and the particular choice of the prime
in that lemma's interval.
\end{theorem}
\begin{proof}
The norm of the different generator is \(|\operatorname{disc}E|\),
so \((\prod_v\kappa_v)^{1/f}=Q_E\). This quantity is at least one at
the stated nodes. The respective values of \(g^{-1}\) are
\[
       \frac{Q_E^{\,2(n-j)}}{MN},
       \qquad
       \frac{Q_E^{\,2(n-j)+1}}{MN}.
\]
Each exponent is at most \(m-1\). The lower endpoint for \(N\) is
\(w^m Q_E^{m-1}b^{-15/2}\). Thus both values are at most
\(b^{15/2}/(M w^m)\le b^{-25/2}\), since \(M\ge b^{20}\)
and \(w\ge1\). Lemma~\ref{lem:freqheight875} gives the first inequality.
For the second, \(n\le\sqrt x\) and \(k\ge x^5-2\sqrt x\) imply
\(n^7/k\longrightarrow0\). Hence
\((25/2)(k-n^7)-\log C/\log b\ge12k\) eventually.

The first scale pairs the finite argument
\(y=(p^{-1}-(p')^{-1})/M\) with the difference of the inverses of
the normalized real entries \(b_{j,v}\Phi_v p\) and
\(b_{j,v}\Phi_v p'\). The second arises when sampling a common entry
with the additive character defined by the trace divided by
\(\delta_E\). Indeed at a complex place that trace contributes
\(2\Re(z\,\Phi_v y/\Phi_v\delta_E)\). In the rescaled variable
\(b_{j,v}z\), the wave \(u\mapsto\exp(u-\bar u)\) has coefficient
\(2\pi i\,\Phi_v y/(\Phi_v(\delta_E)b_{j,v})\), up to the choice
of the sign in the Fourier transform. At a finite prime the
annihilator of the integers is the integers for this pairing,
by the definition of the inverse different.
\end{proof}

\begin{corollary}[Matrix bound for unequal last pivots]\leavevmode
\label{cor:pivotoperator875}
Let \(P\) be any finite set of distinct possible nonzero last
pivots at a fixed arm position \(j\). Choose \(T\) such that
\(y_{p,p'}=(p^{-1}-(p')^{-1})/M\) is integral outside \(T\) for
the pairs under consideration. Form the matrix on \(P\) with
entries \(W_{T,c,K}(y_{p,p'})\) for those unequal pairs and
zero elsewhere, using \(c_v=M/b_{j,v}\). Its operator norm on
\(\ell^2(P)\) is at most \(b^{-12kf}\). This bound is uniform in
the number of pivots and holds also after deleting any set of entries.
For the pairs occurring together at the unused marks of
Proposition~\ref{prop:localpivot873}, one may exclude the primes
over those marks from \(T\).
\end{corollary}
\begin{proof}
For each fixed \(p\), distinct \(p'\ne p\) give distinct nonzero
\(y_{p,p'}\). The same injectivity holds in a fixed column.
Thus all row sums and all column sums of the nonnegative matrix are
at most the sum in \eqref{eq:actualfreq875}. If \(A\) is its matrix
and this bound is \(e\), Cauchy--Schwarz gives
\(|(Ah)_p|^2\le e\sum_{p'} A_{p,p'}|h_{p'}|^2\).
Summation in \(p\) gives \(\|Ah\|_2^2\le e^2\|h\|_2^2\).
Deleting entries decreases the row and column sums. Finally the two
pivots at an unused mark are congruent units in each component,
so their inverses are congruent there as well. At each prime over
that mark \(M\) has valuation one, since both old primes split
completely. Division by \(M\) therefore leaves a local integer.
\end{proof}

Estimate \eqref{eq:actualfreq875} controls the arithmetic weight.
Applying it to the signed integrals in \eqref{eq:spheremean872}
requires estimates for the Fourier coefficients of those integrals
at these same frequencies.
\clearpage

\subsection{Mixed Bessel periods and the discrete-series mass}
\label{sec:mixedperiod876}

There are two tests in the outer product of an orbital integral.
At every real place retain the full discrete-series vector \(F_v\) in the
first position and the cut \(F_v^\flat\) in the second. Let \(G_v\) be the real unitary group
at the embedding \(v\) of the totally real field \(F\), and set
\[
 I_v^{\rm mix}(g)=
       \int_{G_v} F_v(y)F_v^\flat(g^{-1}y)^*\,dy .
\]
Both test vectors and the Haar measure are the ones in
Theorem~\ref{thm:normalizedcut699}.

\begin{lemma}[Absolute convergence on the real open cell]\leavevmode
\label{lem:mixedintegral876}
The integral defining \(I_v^{\rm mix}(g)\) converges absolutely on
the whole cell with nonzero leading arm minors, at the nodes of the
shorter interval. Its absolute integral is locally bounded there.
\end{lemma}
\begin{proof}
First take \(g=w_{\mathbf p}\sigma\) in the normal form of
Lemma~\ref{lem:pivots871}. In \(y=g u c\), integrate first over
\(u\in U_v\) and then over \(\dot c\in U_v\backslash G_v\).
On the support of the second test we may take \(c\) in a fixed
compact set. Indeed the support has bounded
\(\inf_{s\in S_v}\|s c\|_{\rm op}\), the subgroup \(S_v/U_v\)
is compact, and a bounded set in this unitary matrix group has
compact closure. Its image in the quotient has finite measure.

We bound the measure of \(u\) for which the length
\(\ell(g u c)\le R\), \(R\ge1\). For any such \(u\) there is
\(u'\in U_v\) with \(\|u' g u\|_{\rm op}\le C_1(1+R)\).
This follows by writing the element in the infimum as a compact
factor times a unipotent and using boundedness of \(c^{-1}\).
Order the basis in upper flag order, reversing the order on the dual
arm. Then \(g^{-1}u'g\) is lower block unitriangular, and \(u\) is
upper block unitriangular. Both remain unitriangular when the
coordinates of the middle block are listed separately. In
\(A=g^{-1}(u' g u)\), the entry \(u_{ij}\) for \(i\le j\) equals
the minor with rows \(1,\ldots,i\) and columns \(1,\ldots,i-1,j\).
Indeed the first \(i\) rows of the lower factor only use its first
\(i\) columns, whose leading determinant is one. The indicated
minor of the upper factor is \(u_{ij}\). Thus
\(\|u\|_{\rm op}\le C_g(1+R)^r\), with the fixed changes of basis
absorbed in \(C_g\).

The radical has \(2a+1\) diagonal blocks, so \(\log u\) is a
polynomial of degree at most \(2a\) in \(u-1\). Haar measure in
these logarithm coordinates is a constant times Lebesgue measure
on its real Lie algebra. To see this, choose coordinates adapted
to the lower central series. By the finite
Baker--Campbell--Hausdorff formula, translation adds a constant
in the first quotient and has identity diagonal blocks in each
later quotient, with polynomial additions from earlier quotients.
Its Jacobian is one. Since the real dimension is at most \(r^2\),
the measure under consideration is \(O_g((1+R)^{2a r^3})\).
The bounds just obtained are uniform for \(g=w_{\mathbf p}\sigma\)
in a compact part of the normal-form parameters.

The full vector is bounded by a constant at the node times
\(\max(1,\ell/b^{6n})^{-3r^3 f}\), by
Theorem~\ref{thm:normalizedcut699}. The interval has \(f\ge n=a+1\).
Summation over dyadic lengths now converges because \(3r^3f>2ar^3\).
The second test is bounded and its norm is constant on the unipotent
fibers. This proves the asserted integrability. General elements of
the cell differ from their normal forms by left and right
unipotents. Equivariance and Haar invariance leave the absolute
integral unchanged, and the normal-form parameters vary continuously
on the cell. This proves the local bound.
\end{proof}

\begin{theorem}[One full test in the arithmetic link sum]\leavevmode
\label{thm:mixedperiod876}
Use the same finite set \(\mathcal Q^a\) of heads and the same
terminal set \(\mathcal T\) as in \eqref{eq:orbitfactors871}
and \eqref{eq:spheremean872}. For any of the masks in
Sections~\ref{sec:bothmarks873} and~\ref{sec:selectedopen874}, let
\(E_{\rm cut}\) denote its cut energy and form \(E_{\rm mix}\)
by replacing \(I_v\) by \(I_v^{\rm mix}\) at every real place in
these formulas. On the same fields and interval,
\begin{equation}
  B_0^{|\Xi|/2}
  \left|\sqrt{E_{\rm mix}}-\sqrt{E_{\rm cut}}\right|
       \le b^{-k^3 f/3}\sqrt{\mathcal M}
       \qquad (|\Xi|<f/4)
 \label{eq:mixedperiod876}
\end{equation}
at all sufficiently large nodes. The finite orbital factors,
the sphere densities and the fixed-vector ratios
\eqref{eq:fixedratio874} are the same for the mixed expression.
\end{theorem}
\begin{proof}
Let \(F\) be the product of the local full tests and write
\(F^\flat(y)=\alpha(y)F(y)\), \(0\le\alpha\le1\), for the
product cut. At \(z=\gamma^{-1}y\), the difference between the
mixed outer product and the cut outer product is
\((1-\alpha(y))\alpha(z)F(y)F(z)^*\).
Its norm is at most the full-to-cut error norm, since
\((1-\alpha(y))\alpha(z)\le1-\alpha(y)\alpha(z)\).
Theorem~\ref{thm:globalamplitudetail700} estimates the sum of
these absolute errors with the two nonnegative class weights.

Restrict the rational terms to the displayed head and terminal
sets, and group them within each \((c,s,c')\) by their middle
isometry. These restrictions are invariant under the two rational
unipotent actions used to unfold the sum, by
Lemma~\ref{lem:pivots871}. Every nonzero cut term is retained.
Groups in which the complete cut integral is zero may also be
retained, assigning it value zero. In a block the sum over all
groups of their error norms is at most \(E_{csc'}/v_U\) in the
notation of Theorem~\ref{thm:globalamplitudetail700}. The sum of
their squared norms is at most its square. Hence the difference
of the two square roots is at most
\(\sqrt{|\mathcal C|}\,b^{-k^3 f/2}\) by
Corollary~\ref{cor:squaredamplitudetail700}. This argument applies
to every residue mask in the statement, since it only deletes
terms from a sum of nonnegative errors.
Now \(B_0^{|\Xi|/2}\le b^{3kf/4}\) and
\(\mathcal M\ge D_{\rm id}\ge|\mathcal C|\) by
Lemma~\ref{lem:massscale869} at large nodes. This proves
\eqref{eq:mixedperiod876}.

For the identification of the mixed expression, unfold each
rational double orbit exactly as in Theorem~\ref{thm:orbitfactors871}.
The absolute row bound just used, together with absolute convergence
of the cut row, justifies removing the two partitions and regrouping.
Both tests transform by the same unipotent character, so these
changes of variables give the same finite character factors.
Lemma~\ref{lem:mixedintegral876} permits the individual real
integrals and their products. Finite levels and Haar measures are
as before. Left and right translations by the compact middle frame
fixer act on the two real factors by the same unitary actions
as for the cut tests. Consequently the Gram products of the
archimedean matrices still depend on the terminal entry and the
two heads. The finite sphere disintegration of
Theorem~\ref{thm:spheremean872}, applied to this finite list of
matrices, gives precisely the claimed mixed expression.
\end{proof}

At one real place let \(R(x)\) act by right translation in the
vector-valued quotient Hilbert space of the unipotent character,
and put \(C_Z(x)=\langle Z,R(x)Z\rangle\), linear in the second
argument. Write \(d_{\pi,v}\) for the formal degree in the specified
real Haar measure and \(d_v=\|F_v\|^2\).

\begin{theorem}[The mixed position integral and its tail]\leavevmode
\label{thm:mixedmass876}
Set \(D_\pi=\prod_vd_{\pi,v}\), \(d_B=\prod_vd_v\), and let
\(\mathcal B=\{(x_v):\max_v\|x_v\|_{\rm op}\le b\}\).
The integral
\[
 J_{\rm mix}
   =\int_{\prod_vG_v}\prod_v
            \overline{C_{F_v}(x_v)}\,C_{F_v^\flat}(x_v)\,dx
\]
is absolutely convergent and satisfies
\begin{align}
  (1-b^{-r^3f})^{2f}\,\frac{d_B^2}{D_\pi}
       &\le J_{\rm mix}\le\frac{d_B^2}{D_\pi},
     \label{eq:mixedmass876}\\
 \int_{\mathcal B^c}\prod_v
          |C_{F_v}(x_v)C_{F_v^\flat}(x_v)|\,dx
       &\le b^{-r^8/2}\frac{d_B^2}{D_\pi}
     \label{eq:mixedtail876}
\end{align}
for all sufficiently large nodes.
\end{theorem}
\begin{proof}
Write \(c_v(x)=\langle w_v,\pi_v(x)w_v\rangle\) for the unit
discrete-series coefficient defining the full test; here \(w_v\)
denotes a vector. Proposition~\ref{prop:averagedschur684} and
normalization give \(C_{F_v}(x)=d_v c_v(x)\). In the quotient
representation set
\[
             P_v=d_{\pi,v}\int_{G_v}\overline{c_v(x)}R(x)\,dx .
\]
The coefficient is integrable by that proposition. More explicitly
this integral defines a bounded operator by integration on each
vector. It is an orthogonal projection. Indeed
\(p=d_{\pi,v}\overline{c_v}\) satisfies \(p^*=p\) since
\(c_v(x^{-1})=\overline{c_v(x)}\), and \(p*p=p\) by Schur
orthogonality: integrating \(c_v(h)c_v(h^{-1}x)\) in \(h\) gives
\(c_v(x)/d_{\pi,v}\). Conjugating gives the stated convolution
identity. Integration in any unitary representation respects
convolution and adjoints by absolute integrability.
The reproduction identity in Theorem~\ref{thm:normalizedcut699}
gives \(P_vF_v=F_v\). Thus
\[
 \int_{G_v}\overline{C_{F_v}(x)}C_{F_v^\flat}(x)\,dx
       =\frac{d_v}{d_{\pi,v}}\,\|P_vF_v^\flat\|^2.
\]
The norm on the right is at most \(\sqrt{d_v}\) before squaring;
it is at least \((1-b^{-r^3f})\sqrt{d_v}\) by contraction and
the local cut error. Absolute convergence follows also from
\(|C_{F_v^\flat}|\le d_v\). Multiplying proves
\eqref{eq:mixedmass876}.

Put \(P=9r^2\). Lemmas~\ref{lem:coefficientdecay697}
and~\ref{lem:cosetpacking699} give
\(|c_v(x)|\le\min(1,(2/\|x\|_{\rm op})^{r^8})\) and
\(d_{\pi,v}\int (2/\|x\|_{\rm op})^Pdx\le b^{38r^2}\).
Hence \(d_{\pi,v}\int |c_v|\le b^{38r^2}\); on \(\|x\|_{\rm op}>b\)
the latter bound acquires the additional factor \((2/b)^{r^8-P}\).
Cover the complement of \(\mathcal B\) by the \(f\) events where
one coordinate exceeds \(b\). After dividing by \(d_B^2/D_\pi\)
the result is at most \(f b^{38r^2 f}(2/b)^{r^8-9r^2}\).
Here \(f\ll_{q_0,d}x^{12}\), \(r\ge x^5\), and \(b\to\infty\).
In particular \(fr^2=o(r^8)\) and \(2/b\le b^{-3/4}\) eventually.
This proves \eqref{eq:mixedtail876}.
\end{proof}

The position integral is now evaluated for tests that also occur
in the arithmetic energy. Passing from the discrete head sums and
their finite orbital factors to this integral still requires control
of the intermediate correlations.
\clearpage

\subsection{Fourier coefficients of the unramified terminal density}
\label{sec:terminalfourier877}

Fix a finite place \(w\) of the totally real field. Write \(O\) for
its integers, \(\pi\) for a uniformizer, \(q=|O/\pi O|\), and \(A\)
for the integers of the quadratic algebra \(E\otimes_F F_w\).
Assume that \(N\) is a unit, the middle lattice is unimodular
and \(A\) is split or unramified. Put \(\eta=1\) in the split case
and \(\eta=-1\) in the unramified field case, and set \(e=\eta^k\).
Give \(A\) and \(O\) additive volume one. Let \(\delta_w\) be
the integral density of the rank \(k\) middle norm \(Q\), with
target measure of volume one on \(O\), as in
Theorem~\ref{thm:spheremean872}. As a function of the terminal
entry on the compact additive group \(A\), consider
\[
        H_w(t)=\delta_w\big((t\bar t-N^2)/N\big),\qquad t\in A .
\]
The extension by zero to the quadratic local algebra incorporates
the integrality of \(t\). This is the terminal factor at the
standard places outside \(S\) in \eqref{eq:spheremean872}.
At the unused marked places it is the entire dependence on \(t\)
in Corollary~\ref{cor:integratedmark873} for compatible heads.

\begin{theorem}[Fourier expansion at a standard finite place]\leavevmode
\label{thm:terminalfourier877}
For \(k\ge4\), \(H_w\) is continuous, and its coefficient at the
trivial character of \(A\) is
\begin{equation}
         \widehat H_w(0)=1-\eta^{k+1}q^{-(k+1)} .
 \label{eq:terminalmean877}
\end{equation}
For a nontrivial additive character \(\xi\) of \(A\), let \(L\)
be the least integer with \(\xi|_{\pi^L A}=1\). Then
\begin{equation}
       |\widehat H_w(\xi)|\le 4q^{-kL},\qquad
       \sum_{\xi\ne0}|\widehat H_w(\xi)|\le8q^{2-k}.
 \label{eq:terminalfourier877}
\end{equation}
The coefficients here are integrals against \(\overline{\xi(t)}\).
These formulas apply also in residue characteristic two.
\end{theorem}
\begin{proof}
For \(s\in O\) the density is
\begin{equation}
 \delta_w(s)=(1-e q^{-k})\sum_{h\ge0}
            e^h q^{-(k-1)h}{\bf1}_{s\in\pi^hO}.
 \label{eq:terminalseries877}
\end{equation}
We verify this including the value at zero. Modulo \(\pi\) the norm
has \(q^{2k-1}-e q^{k-1}\) solutions at a unit target and
\(q^{2k-1}+e(q-1)q^{k-1}\) at zero. Indeed a nonsingular residue
Hermitian form splits as an orthogonal sum of unit lines. In the
field case this follows by choosing a vector of nonzero norm
and passing to its orthogonal complement. A form all of whose
norms vanish has all its pairings zero, by varying the scalar in
one coordinate and using the nondegenerate trace pairing. In the
split case take dual bases in the two components. For a nontrivial
residue additive character the sum over a single line is \(\eta q\):
in the split case sum over the second variable in \(uv\), and in
the field case each nonzero norm has \(q+1\) preimages. Additive
orthogonality gives the indicated counts.

Each nonzero residue vector is nonsingular for the scalar norm
equation, since the trace pairing is perfect. Each lift of a
prescribed next target digit then has \(q^{2k-1}\) choices of
the next vector digit. For a unit \(s\) the density is
\(1-e q^{-k}\). For \(s\in\pi O\), the primitive part has density
\(1+e(q-1)q^{-k}-q^{1-2k}\). The vectors in \(\pi A^k\) contribute
only for \(s\in\pi^2O\), and there they contribute
\(q^{2-2k}\delta_w(s/\pi^2)\) by scaling the vector and target.
The first two terms of the series \eqref{eq:terminalseries877}
for \(s\in\pi O\) sum to this primitive part, since \(e^2=1\).
The terms from the third onward are its stated recursive
remainder. This proves the series for each nonzero \(s\).
At zero one can use the finite congruence densities
\(q^j\Pr(Q(x)\in\pi^jO)\). Their primitive part is the same for
\(j\ge1\), and for \(j\ge2\) their remaining part is \(q^{2-2k}\)
times the density at precision \(j-2\). The two starting densities
are finite and the ratio is less than one. Taking the limit gives
the same sum at zero. The series converges uniformly and its
terms are locally constant.

For \(h\ge1\) let
\(V_h=\{t\in A:t\bar t\equiv N^2\bmod\pi^h\}\).
Its measure is \((1-\eta/q)q^{-h}\). In the split case there
are \(q-1\) pairs modulo \(\pi\) with the prescribed nonzero
product, and in the field case there are \(q+1\) such elements.
Each solution has \(q\) lifts for a specified next norm digit.
Indeed the differential of the norm at a unit is its unit norm
times the trace after dividing the increment by that unit.
The residue trace is surjective for either algebra. Dividing
these counts by \(|A/\pi^h A|=q^{2h}\) gives the measure.
Substitute \(s=(t\bar t-N^2)/N\) in
\eqref{eq:terminalseries877} and integrate over \(A\). One obtains
\[
 (1-e q^{-k})
   \left(1+(1-\eta/q)\sum_{h\ge1}(e q^{-k})^h\right)
       =1-\eta e q^{-(k+1)}.
\]
For nontrivial \(\xi\), its integral against the constant term or
against \({\bf1}_{V_h}\) for \(h<L\) vanishes by translation on
\(\pi^h A\). For \(h\ge L\) the absolute value of the latter
integral is at most the measure of \(V_h\). Thus the coefficient
is at most
\[
    |1-e q^{-k}|(1-\eta/q)\frac{q^{-k L}}{1-q^{-k}}
       \le 4q^{-k L}.
\]
At most \(q^{2L}\) characters have conductor exponent \(L\), since
they descend to \(A/\pi^L A\). Summing over \(L\ge1\) proves
\eqref{eq:terminalfourier877}. Fourier inversion here can be seen
without a limiting inversion hypothesis: average \(H_w\) on cosets
of \(\pi^jA\) and use finite character orthogonality. These averages
converge uniformly to \(H_w\), and their finite expansions exhaust
the absolutely summable coefficients.
\end{proof}

\begin{corollary}[The product and the denominator of a terminal frequency]
\leavevmode\label{cor:terminalproduct877}
For any set \(\mathcal V\) of the standard places just described,
the product of \(H_w\) on \(\prod_{w\in\mathcal V}A_w\) has an
absolutely convergent character expansion. Its constant coefficient
is \(\prod_{w\in\mathcal V}(1-\eta_w^{k+1}q_w^{-(k+1)})\). Relative
to this coefficient, the sum of the absolute values of all other
coefficients is \(O(f2^{2-k})\), uniformly on the shorter interval.
For a character given by a nonzero rational frequency \(y\) in the
trace pairing divided by \(\delta_E\), let \(T\) consist of the
primes of \(E\) over the finite set of places where this character
on \(A_w\) is nontrivial in \(\mathcal V\). The absolute value of
its coefficient divided by the constant coefficient is at most
\begin{equation}
             I_T(y)^{(k-n^7)/2}D_T(y)^{-3}
 \label{eq:terminaldenom877}
\end{equation}
at sufficiently large nodes, with the notation of
Section~\ref{sec:freqnode875}.
\end{corollary}
\begin{proof}
The local mean is positive and at least \(1-2^{-k-1}\). After
division by the means the nonconstant coefficients have sum
\(O(q_w^{2-k})\). At most \(f\) places lie over any rational
prime. For \(k\ge5\), summing \(m^{2-k}\) over integers \(m\ge2\)
costs \(O(2^{2-k})\), by comparison with the integral on
\([2,\infty)\) for the terms after \(m=2\). Hence
\(\sum_w q_w^{2-k}=O(f2^{2-k})=o(1)\) at the given nodes.
The product of one plus each local bound is at most the
exponential of their sum. Expanding finite products and taking
the uniformly absolutely convergent limit gives the asserted
relative sum and constant coefficient.

Locally the trace pairing divided by the different generator has
\(A_w\) as its own annihilator. Thus the exponent \(L\) is the least
integer for which \(\pi^L y\in A_w\). If \(D_w\) denotes the
product of the denominator norms over this place, then
\(q_w\le D_w\le q_w^{2L}\) for a nontrivial character. In the split
case one adds the two nonnegative denominator exponents. In the
unramified field case the residue norm is \(q_w^2\). By
\eqref{eq:terminalfourier877} the normalized coefficient is at most
\(8q_w^{-kL}\le8D_w^{-k/2}\). Since \(D_w\ge2\) and \(n^7\ge12\)
at large nodes, this is at most \(D_w^{-(k-n^7)/2-3}\).
The signed index at this place is at least \(D_w^{-1}\), directly
from the valuations. Multiplication over the nontrivial places
proves \eqref{eq:terminaldenom877}.
\end{proof}

For a compatible pair of heads, \eqref{eq:terminaldenom877}
supplies the finite part of the frequency weight
\eqref{eq:freqweight875} at all these places. At a selected mark
with a singular minor the test also contains the signed unipotent
orbital integrals, so their Fourier coefficients have to be
controlled in addition to this density.
\clearpage

\subsection{Propagation of a fixed vector at a singular mark}
\label{sec:singularflag878}

Fix a selected split mark \(w\) of residue order \(q\). Work in the
component where the fixed vector was prescribed. Its group is
\(\mathrm{GL}(V)\) over \(\mathbb F_q\), with the paired inverse adjoint
determining the other component. Use the frame \(f_i=e_i\) for \(i\le a\),
\(f_n=z\), \(n=a+1\), and \(T_i=\langle f_1,\ldots,f_i\rangle\),
\(T_0=0\). Thus \(T_a=L\). Let \(\bar U\) be the residue unipotent group,
and let \(\bar\chi\) be its level character. Call a residue \(s\)
\emph{compatible} if
\[
 \bar\chi(u)=\bar\chi(s^{-1}u s)
           \quad(u\in\bar U\cap s\bar U s^{-1}).
\]

\begin{lemma}[Character compatibility on the residue double orbit]
\leavevmode\label{lem:residueshift878}
If \(K_w(s;g)\ne0\) for any local orbital parameter \(g\), then \(s\)
is compatible. For any compatible \(s\) there is a linear map \(A\)
on \(V\) such that, for \(i<n\),
\[
 A f_i-f_{i+1}\in T_i,\qquad A s f_i-s f_{i+1}\in sT_i .
\]
This holds even if some of the leading arm minors of \(s\) vanish.
\end{lemma}
\begin{proof}
In the local integral of Theorem~\ref{thm:orbitfactors871}, change
the two unipotent variables by integral lifts of \(u_0,v_0\in\bar U\).
Surjectivity of reduction follows from the split upper block matrices.
The principal kernel has trivial character, so this gives
\[
 K_w(u_0s v_0^{-1};g)
      =\overline{\bar\chi(u_0)}\,\bar\chi(v_0)\,K_w(s;g).
\]
Taking \(u_0=s v_0s^{-1}\) gives the necessity of compatibility.

Write \(\bar U=1+\mathfrak n\), where \(\mathfrak n\) consists of the
strictly upper blocks in the chosen component, including the strict
upper entries on the arm lines. The character of \(1+X\) is of the
form \(\psi(\Lambda(X))\) for a nontrivial additive character
\(\psi\) of \(\mathbb F_q\) and a linear functional \(\Lambda\).
Here is its relevant part. The frame character in
Section~\ref{sec:rationalflags680} sums the consecutive frame entries
in both components under the trace. Its coefficient in the selected
component, denoted by \(c_0\), is nonzero: \(M\) has valuation one
and the trace divided by the different has conductor the integer
ring. The contribution to \(\Lambda(X)\) from an entry taking a
later block into the line \(e_i\) is \(c_0\) times the coefficient
of \(e_i\) in \(X f_{i+1}\). Indeed the other component is the
inverse adjoint. On the first off-diagonal blocks it reads minus the
adjoint of \(X\), all higher powers having zero such blocks.
Its frame uses entries on the opposite arm and hence has zero
coefficient on the entries under consideration. These statements
also apply with the two components interchanged.

For \(X\in\mathfrak n\cap s\mathfrak n s^{-1}\), apply compatibility
to \(1+tX\) for all \(t\in\mathbb F_q\). Since \(\psi\) is nontrivial,
\(\Lambda(X)=\Lambda(s^{-1}Xs)\). The two functionals therefore
extend to one linear functional on the sum of the two matrix spaces,
and a basis extension makes it a functional on all endomorphisms.
Write it as \(X\mapsto\operatorname{tr}(BX)\) by the perfect matrix
trace pairing. Testing maps from each later block into \(e_i\)
determines the part of \(B e_i\) modulo \(T_i\) as \(c_0 f_{i+1}\).
Testing their conjugates by \(s\) gives the same assertion for
\(s^{-1}B s\). Taking \(A=c_0^{-1}B\) proves the claim.
\end{proof}

\begin{theorem}[A chain of common vectors and the resulting valuations]
\leavevmode\label{thm:singularflag878}
Suppose \(s\) is compatible and maps a vector of
\(\lambda f_j+T_{j-1}\) into \(\lambda f_j+T_{j-1}\), where
\(\lambda\ne0\) and \(1\le j\le n\). In particular this holds if
\(s\) fixes the prescribed vector of top index \(j\).
For every \(j\le i\le n\), the two affine sets
\(\lambda f_i+T_{i-1}\) and \(s(\lambda f_i+T_{i-1})\) intersect.
The resulting common vectors at the successive levels are linearly
independent modulo \(T_{j-1}\).

Let \(\mathfrak l\) be the prime of \(E\) for the selected component.
For any integral local lift with residue \(s\), its integral leading
frame determinants satisfy
\begin{equation}
 \operatorname{ord}_{\mathfrak l}D_i\ge i-j+1\quad(j\le i\le a),
 \qquad
 \operatorname{ord}_{\mathfrak l}(D_n+ND_a)\ge n-j+1.
 \label{eq:singdepth878}
\end{equation}
We put \(\operatorname{ord}(0)=+\infty\). In particular the actual
head and terminal entry of any contributing orbital factor at this
residue satisfy
\(\operatorname{ord}_{\mathfrak l}(D_a(\theta+N))\ge n-j+1\).
\end{theorem}
\begin{proof}
Start with a common vector in the two stated affine sets at \(j\).
The map of Lemma~\ref{lem:residueshift878} carries any such common
vector at \(i<n\) to one at \(i+1\): the lower terms on either
side map into the preceding levels of the next affine sets, and
the top coefficient is still \(\lambda\). Iteration produces the
vectors claimed. Distinct top indices give independence modulo
\(T_{j-1}\).

Reduce the entire matrix \(C=(h(f_i,g f_l))_{i,l\le n}\) of an
integral lift \(g\) in the chosen component. If \(j\le i\le a\),
the source coordinates of the common vectors of levels
\(j,\ldots,i\) give independent vectors in the first \(i\)
columns. Their images under \(s\) lie in \(L\), so pair to zero
with every frame row. Thus the reduction of the leading
\(i\)-square has nullity at least \(i-j+1\). For the last common
vector, the image is \(\lambda z+\ell\), \(\ell\in L\).
Pairing with the frame rows gives zero apart from the last
entry, which is \(-N\lambda\). Add \(N\) to the last diagonal
entry of \(C\). The modified residue matrix therefore kills
the source coordinate vectors of all \(n-j+1\) levels,
including the last one with terminal coefficient \(\lambda\).
Its determinant is \(D_n+N D_a\) by expansion in that entry.
An integral square matrix with an \(h\)-dimensional residual
kernel has determinant divisible by \(\mathfrak l^h\):
extend a basis of that kernel to a residue basis and lift the
basis change to an invertible integral matrix. The \(h\)
corresponding columns then belong to \(\mathfrak l\). This
proves \eqref{eq:singdepth878}. In a contributing local
integral some unipotent translates give an integral lift of
the residue. They preserve the leading determinants, and
\(D_n=D_a\theta\) by Lemma~\ref{lem:pivots871}.
\end{proof}

\begin{corollary}[Terminal fixers suffice at the selected singular marks]
\leavevmode\label{cor:terminalfixer878}
Let \(j_w\) be the top index of the chosen vector at each selected
mark, and let \(\Xi_z\) replace it by \(z\) in the same component.
For the cut or mixed energies restricted to a singular arm minor
at each selected place,
\begin{equation}
             E_{\rm sing}(\Xi)
       \le \left(\prod_{w\in\Xi}q_w^{\,n-j_w}\right)
                         E_{\rm sing}(\Xi_z).
 \label{eq:terminalfixer878}
\end{equation}
In particular the product costs at most \(b^{11a|\Xi|}\) at the same
nodes. The empty choice incurs factor one.
\end{corollary}
\begin{proof}
Within a fixed double orbit of \(\bar U\), the squared norm in the
sphere formula is constant as a function of the residue, with all
other variables fixed. The reason is the first identity in the proof
of Lemma~\ref{lem:residueshift878}: its scalar has modulus one and is
common to every head in the sum. A double orbit that contributes
is compatible. The singular-minor requirement is constant on the
orbit because the leading minors are unipotent invariants.

A uniform element of the double orbit of \(s\) is obtained as
\(u s v^{-1}\) for independent uniform \(u,v\in\bar U\). All fibers
have equal cardinality by transitivity. For a fixed vector \(d\)
of top index \(j\) and coefficient \(\lambda\), its \(\bar U\)-orbit
is exactly \(D=\lambda f_j+T_{j-1}\). For an arm line this follows
from the independent upper entries in its column, and for the
terminal vector from the free middle-to-arm shears. The probability
of fixing \(d\) on this double orbit equals
\(|D\cap sD|/|D|^2\). If this is nonzero, put
\(h_i=\dim(T_i\cap sT_i)\). The intersection of affine spaces is a
translate of \(T_{j-1}\cap sT_{j-1}\); the probability is thus
\(q^{h_{j-1}-2(j-1)}\). By Theorem~\ref{thm:singularflag878}
there is also a match at \(z+L\), after division by \(\lambda\).
Its corresponding probability is \(q^{h_a-2a}\). For \(j\le a\),
the independent common vectors of levels \(j,\ldots,a\) give
\(h_a\ge h_{j-1}+a-j+1\). The ratio of the two probabilities
is therefore at most \(q^{a-j+1}=q^{n-j}\). For \(j=n\) the two
probabilities are equal by the affine orbit formula itself.

The inequality holds separately on each double orbit of positive
weight, and trivially on those with no fixed \(d\). Sum the
nonnegative contributions, then apply the same comparison successively
at all selected places. Each replacement preserves the requirements
at the other places. Finally \(q_w\le b^{11}\) and \(n-j_w\le a\).
\end{proof}

Thus a bound for terminal fixers can be applied to all prescribed
frame vectors with the explicit loss in \eqref{eq:terminalfixer878}.
The signed orbital factor at a singular terminal fixer remains in
this bound.
\clearpage

\subsection{Transversality of the full isotropic arm}
\label{sec:armopen879}

At a split mark \(w\) write \(A=\mathcal O_{F_w}\times\mathcal O_{F_w}\)
and \(\bar A=\mathbb F_q\times\mathbb F_q\).
Let \(P_w\) be the set of residues for which the full \(a\)-by-\(a\)
arm determinant \(D_a\) is a unit in both components. Thus \(P_w\)
contains \(O_w\) of Section~\ref{sec:selectedopen874}. We calculate
the integrated factor on \(P_w\) while allowing the leading minors
of orders strictly less than \(a\) to vanish modulo the mark.

Let \(V_0\) be the radical for the three blocks \(L,W,L^\vee\),
and let \(U_A\) be the subgroup of \(U\) acting within the two arms
and trivially on \(W\). Thus \(U=V_0\rtimes U_A\). For an invertible
arm matrix \(B\) over the paired algebra, let \(w_B\) carry \(L\)
to \(L^\vee\) by \(B\) in their dual bases, carry \(L^\vee\) back to
\(L\) by the inverse adjoint, and act identically on \(W\).
For \(h,h'\in U_A(F_w)\) write \(B(h,h';\mathbf p)\) for the
arm matrix determined by \(h w_{\mathbf p}(h')^{-1}=w_{B(h,h';\mathbf p)}\).
Define
\[
 L_w(\bar B,\mathbf p)=\int_{U_A(F_w)^2}
 {\bf1}_{\,B(h,h';\mathbf p)\in{\rm GL}_a(A),\
             B(h,h';\mathbf p)\bmod\mathfrak p_w=\bar B}
      \overline{\chi_w(h)}\chi_w(h')\,dh\,dh' .
\]
The arm measures give their principal integral kernels volume one.

\begin{theorem}[The marked factor with transverse arms]\leavevmode
\label{thm:armopen879}
Let \(U_w(\mathbf p,\mathbf p',\theta)\) denote the entire local
factor of the expanded sphere expression \eqref{eq:spheremean872}
when \(s_w\in P_w\) and no vector is fixed at \(w\). With \(H_w\)
from Theorem~\ref{thm:terminalfourier877}, extended by zero off \(A\),
one has
\begin{align}
 U_w(\mathbf p,\mathbf p',\theta)
  ={}&|\bar V_0|^2\,|\mathrm{GL}_k(\mathbb F_q)|\,H_w(\theta)
        \sum_{\bar B\in\mathrm{GL}_a(\bar A)}
         L_w(\bar B,\mathbf p)\overline{L_w(\bar B,\mathbf p')}.
 \label{eq:armorbital879}
\end{align}
Imposing a fixed frame vector \(d\) multiplies this expression
by \(R_w(d,\theta)\) of \eqref{eq:fixedratio874}, with value zero
if \(\theta\notin A\). The multiplier is independent of both heads.
Both formulas continue to hold with \(\bar B\) restricted to any
fixed subset, using that same subset in the sum.
\end{theorem}
\begin{proof}
We first verify the three-block decomposition used over the local
ring. An integral isometry whose residue lies in \(P_w\) factors
uniquely as
\[
              v_1\, w_B x\,v_2^{-1}
       \quad(v_i\in V_0(\mathcal O_{F_w}),\
        B\in {\rm GL}_a(A),\ x\in\mathrm U(W)(\mathcal O_{F_w})).
\]
To see this, identify the chosen component with a general linear
group and put \(J=w_{1_a}\). In \(J^{-1}g\) the leading diagonal
block of size \(a\) is invertible: its determinant is the given
arm pairing in that component up to a unit from the bases.
The leading square of size \(a+\dim W\) is also invertible.
By the complementary-minor identity, its determinant is, up to
a unit, the determinant of the final \(a\)-square of \(g^{-1}J\).
The latter is the arm pairing in the paired component up to a
unit, since the paired action is inverse adjoint. The middle
Gram matrix is unimodular here because \(N\) is a unit.
Block elimination now gives a lower unitriangular factor, a
block diagonal factor, and an upper unitriangular factor, all
integral with integral inverses. In detail, clear the first
block column and row with the inverse of the first pivot block.
The determinant of the next pivot block is the ratio of the
two invertible leading determinants. Clear the second column
and row in the same way. All choices are forced. Multiplying
back by \(J\) gives the displayed factorization, including its
paired component by inverse adjunction. Exactly the same proof
works over the residue field, and the factors of any integral
lift reduce to its unique residue factors.

Write the two unipotents in the integral defining \(K_w\) in
the form \(v_1h,v_2h'\). The product measures give the prescribed
Haar normalization: unipotent conjugation by \(U_A\) preserves
Haar on \(V_0\) (its determinant on the Lie algebra is one), and
the principal subgroup is the semidirect product of the two
principal subgroups. Since \(h,h'\) act trivially on \(W\),
the middle of the resulting product is \(w_{B(h,h';\mathbf p)}x\).
Write the residue decomposition as
\(\bar v_{1,s} w_{\bar B_s}\bar x_s \bar v_{2,s}^{-1}\).
By the uniqueness just proved, \(v_1,v_2\) must belong to the
integral cosets with these reductions, each of volume one.
The middle factor must be integral and reduce to \(\bar x_s\).
Conversely these conditions, together with the prescribed
reduction of \(B\), suffice. This yields
\[
 K_w(s;w_{\mathbf p}x)
   =\overline{\bar\chi(\bar v_{1,s})}\bar\chi(\bar v_{2,s})
          L_w(\bar B_s,\mathbf p)\,
       {\bf1}_{\,x\ {\rm integral},\ \bar x=\bar x_s}.
\]
The integral defining \(L_w\) is finite for a fixed nonzero head.
Indeed the full flag decomposition over the local field at
\(w_{\mathbf p}\) has continuous inverse on the locus where its
leading determinants are nonzero. In the compact set of allowed
arm matrices these determinants have fixed nonzero values on the
orbit, so restricting to that closed level set bounds both arm
unipotents. The character is constant on each principal coset.

For fixed \(\bar B,\bar x\) there are \(|\bar V_0|^2\) distinct
residues, given by freely choosing the two outer radicals.
Their two character factors cancel in the product for the paired
heads. The integral over the fiber of \(xz=\omega-(\theta/N)z\)
has volume \(|\mathrm{GL}_k(\mathbb F_q)|\) for integral
\(\theta,\omega\) and is zero otherwise, as proved in
Corollary~\ref{cor:integratedmark873}. Integrating the norm sphere
therefore gives \(H_w(\theta)\), proving \eqref{eq:armorbital879}.

For \(d=\ell+c z\), \(c=0\) gives no fixed vector: the two
inverse radical translates lie in \(L\) and the invertible arm
matrix sends a nonzero vector of \(L\) into \(L^\vee\). For
\(c\ne0\), the two inverse radical translates independently
range uniformly over \(cz+L\), since \(V_0\) itself contains
the middle-to-arm shears. The equation between them through
\(w_{\bar B}\bar x\) forces both arm components to vanish and
\(\bar xz=z\) in the selected component. Its probability is
\(q^{-2a}{\bf1}_{\bar xz=z}\), independently of \(\bar B\).
The same change of variable on the split norm sphere as in
Theorem~\ref{thm:selectedopen874} now gives precisely
\(R_w(d,\theta)\). This proof was for each \(\bar B\) separately.
\end{proof}

\begin{corollary}[Separation of the places with a defective full arm]
\leavevmode\label{cor:armdefect879}
Retain the original choice \(\Xi\), \(|\Xi|<f/4\). For
\(J\subseteq\Xi\), let \(E_{\rm def}(\Xi\setminus J;J)\) be the
following cut or mixed energy. At the original unselected places
impose all unit pivots. At the places of \(J\) impose only
\(s_w\in P_w\). At every place of \(\Xi\setminus J\) retain the
fixed vector and impose \(s_w\notin P_w\). In the energy \(E(\Xi)\)
of Section~\ref{sec:selectedopen874}, or in its mixed version, one
then has
\begin{align}
 B_0^{|\Xi|}E(\Xi)
 \le \sum_{J\subseteq\Xi}
     \left(\prod_{w\in J} B_0q_w^{1-k-2a}\right)
       B_0^{|\Xi\setminus J|}
       E_{\rm def}(\Xi\setminus J;J).
 \label{eq:armdefect879}
\end{align}
Each factor in parentheses per place is at most \(b^{-4k}\).
The sum of these products over nonempty \(J\) is
\(O(f b^{-4k})\), uniformly on the interval. All the terms still
have the original set of more than \(7f/4\) places with full
pivot congruences from Proposition~\ref{prop:localpivot873}.
\end{corollary}
\begin{proof}
Partition the actual residues by the subset \(J\) of selected
places at which both full arm determinants are units. Each part
has nonnegative energy. In its sphere integral keep the terminal
entry fixed. At every \(w\in J\), Theorem~\ref{thm:armopen879}
identifies the complete local factor, for every pair of heads,
with the one omitting the vector condition times \(R_w(d_w,\theta)\).
The product of these scalars can be taken outside the head sum.
The resulting contribution with no vector restrictions on \(J\)
is nonnegative at each terminal entry, by the squared norm in
the sphere formula. Bound the product using
\(R_w\le q_w^{1-k-2a}\) and sum the entries and the disjoint
parts. This proves \eqref{eq:armdefect879}. As \(q_w\ge b^{10}\),
\(B_0=b^{6k}\) and \(a\ge1\), each weighted factor is at most
\(b^{-4k}\). Their nonempty subset sum is at most
\((1+b^{-4k})^{|\Xi|}-1=O(f b^{-4k})\), since \(f b^{-4k}\to0\).
The original unselected conditions were retained in every part.
\end{proof}

The term with \(J=\varnothing\) now has a noninvertible full arm
pairing in at least one component at every selected mark. The
integral on such residues, as well as the correlations of the
real orbital factors at unequal heads, still enter the arithmetic
second moment.
\clearpage

\clearpage
\subsection{Integral positions at a marked terminal stabilizer}
\label{sec:integralmass881}

Fix one of the original marked places in the interval of
Theorem~\ref{thm:bothmarks873}. Work in the component in which the
vector is specified. In a split basis put \(G=\mathrm{GL}_r(F_w)\),
\(K=\mathrm{GL}_r(\mathcal O_w)\), and let \(J\) be the kernel of
reduction from \(K\) to \(\bar G=\mathrm{GL}(V)\), where
\(V=\mathbb F_q^r\). The residual flag has blocks of dimensions
\(1^a,d,1^a\) with \(d=k+1\), \(r=2a+d\) and \(n=a+1\).
Write \(V=L\oplus W\oplus L^\vee\), with \(z\in W\) the terminal
vector. It is nonzero since its norm is a unit. Let \(Z\) be the
inverse image of its stabilizer in \(\bar G\). Use the unipotent
group \(U\) and character \(\chi=\chi_w\) of
Theorem~\ref{thm:orbitfactors871}. Give \(J\) and \(U\cap J\)
volume one in their respective groups. The Haar measure on
\(U\backslash G\) is the quotient of these two measures. Both
groups are unimodular.

For \(y=u j\in UJ\) put \(\phi(y)=\chi(u)\) and extend \(\phi\)
by zero. The character is trivial on \(U\cap J\), so this
definition is independent of the decomposition. For a subgroup
\(J\subset H\subset K\) write \(\bar H=H/J\), choose lifts
\(k_d\in H\) of its elements \(d\), and set
\[
 \Phi_d(y)=\phi(y k_d^{-1}),\qquad
 C_H(x)=\int_H\phi(hxh^{-1})\,dh .
\]
In subscripts of \(\Phi\) the letter \(d\) denotes a residue class.
The functions are independent of the lifts, since \(J\) is normal
in \(K\). Each is right \(J\)-invariant, transforms by \(\chi\)
on the left under \(U\), and has support compact modulo \(U\).

\begin{proposition}[Two correlations of the marked functions]
\leavevmode\label{prop:twocorrelations881}
For an open rational head and middle element \(g=w_{\mathbf p}\sigma\)
of Theorem~\ref{thm:orbitfactors871} in the class-removed local
coordinates, put
\[
 I(g)_{d,e}=\int_G \Phi_d(y)\overline{\Phi_e(g^{-1}y)}\,dy
        \qquad(d,e\in\bar H).
\]
Then \(I(g)_{d,e}=K_w(de^{-1};g)\). In particular, for two
such parameters \(g,g'\),
\[
 \frac{1}{|\bar H|}\sum_{d,e\in\bar H}
 I(g)_{d,e}\overline{I(g')_{d,e}}
   =\sum_{s\in\bar H}K_w(s;g)\overline{K_w(s;g')}.
\]
For every \(x\in G\) the right translates of the same functions
satisfy
\begin{equation}
 C_H(x)=\int_{U\backslash G}\sum_{d\in\bar H}
                \overline{\Phi_d(y)}\Phi_d(yx)\,d\dot y .
 \label{eq:rightposition881}
\end{equation}
\end{proposition}
\begin{proof}
The identity
\(\Phi_d(y)=\int_U\chi(u){\bf1}_{u k_dJ}(y)\,du\)
follows since the contributing set of \(u\)'s is a coset of
\(U\cap J\) and \(k_d\) normalizes \(J\). Expand the two factors
of \(I(g)_{d,e}\). The left cosets \(u k_dJ\) and
\(g v k_eJ\) intersect precisely when
\(u^{-1}g v\in k_d J k_e^{-1}\). Their intersection then has
volume one. For fixed nonzero rational pivots the preimage of
this compact coset on the double orbit has compact unipotent
coordinates, by the inverse flag elimination in the proof of
Theorem~\ref{thm:orbitfactors871}. Thus the calculation also
converges with both characters replaced by absolute values.
Invert \(u\) and \(v\). Haar measure is preserved, and the
characters become \(\overline{\chi(u)}\chi(v)\), as in the
definition of \(K_w(de^{-1};g)\). Each value of \(de^{-1}\)
occurs \(|\bar H|\) times in the double sum.

For the right correlation fix \(d\). The first factor restricts
\(\dot y\) to the image of \(J k_d\). Two elements \(j k_d\)
represent the same point precisely when their \(j\)'s differ on
the left by \(U\cap J\). The quotient measure on this image is
the pushforward of \(dj\) on \(J\), since \(U\cap J\) has
volume one. The product in \eqref{eq:rightposition881} is left
\(U\)-invariant. At \(y=j k_d\) its first factor is one, and its
second is \(\phi(j k_d x k_d^{-1})=\phi((j k_d)x(j k_d)^{-1})\)
by right \(J\)-invariance. Sum the integrals over the disjoint
cosets \(J k_d\) of \(H\).
\end{proof}

The integral part of the position mass at \(H\) will be written
\[
 A_H=\frac{\int_K |C_H(x)|^2\,dx}
                  {\operatorname{vol}(J)\operatorname{vol}(H)^2}.
\]
Let \(\mathfrak n\) be the space of strictly upper block matrices
of the above flag, with the dual arm in reverse order. Write
\(\bar\chi(1+X)=\psi(\Lambda(X))\) on \(\mathfrak n\), as in
Lemma~\ref{lem:residueshift878}.

\begin{lemma}[An exact sum on the residue space]
\leavevmode\label{lem:finiteposition881}
For \(H=K\) or \(H=Z\) one has
\begin{equation}
 A_H=\sum_{\mathcal C}
   \frac{1}{|\mathcal C|}
       \left|\sum_{X\in\mathcal C\cap\mathfrak n}
                    \psi(\Lambda(X))\right|^2 .
 \label{eq:finiteposition881}
\end{equation}
Here \(\mathcal C\) ranges over the nilpotent conjugacy orbits of
\(\bar H\) on \(\operatorname{End}(V)\). The zero orbit contributes
one. A contributing orbit has rank at most \(2a\).
\end{lemma}
\begin{proof}
For \(y\in K\), membership in \(UJ\) is equivalent to reduction
into \(1+\mathfrak n\). Indeed \(U\cap K\) surjects onto this group
by lifting its upper block entries. Conversely, \(y=u j\in K\)
forces \(u\in K\). Hence \(\phi(y)\) on \(K\) is the indicated
residue character extended by zero. On the residue class of
\(x\in K\), \(C_H(x)/\operatorname{vol}(H)\) is the average of
this function on the \(\bar H\)-conjugates of \(\bar x\).
Subtracting the identity translates an orbit of unipotent
elements into an orbit of nilpotent endomorphisms. Every point
of the orbit occurs equally often in the average. Each residue
class of \(K\) has volume \(\operatorname{vol}(J)\). Squaring
and summing therefore gives \eqref{eq:finiteposition881}.
For \(X\in\mathfrak n\) one has \(X(L+W)\subset L\). Its image
on the remaining \(a\) dimensions adds at most \(a\) more
dimensions.
\end{proof}

\begin{theorem}[Integral mass for the terminal vector]
\leavevmode\label{thm:integralmass881}
Uniformly at the two marks of the same fields and intervals, for
all sufficiently large nodes,
\begin{align}
 A_Z &\le q^{20n^2},\label{eq:integralmass881}\\
 \frac{\int_K |C_Z(x)|^2\,dx}
      {\operatorname{vol}(J)\operatorname{vol}(Z)\operatorname{vol}(K)}
   &=\frac{A_Z}{q^r-1}\le 2q^{-r+20n^2}.\label{eq:fixernormalized881}
\end{align}
The estimate includes matrices of all residual arm ranks. Even
after multiplication by \(q^a B_0\), the last bound is at most
\(b^{-3k}\), where \(B_0=b^{6k}\).
\end{theorem}
\begin{proof}
Group the nonzero nilpotent matrices by their rank \(\ell\) and
by \(h={\bf1}_{z\in XV}\), both invariant under \(\bar Z\). We
first count the possible \(X\) in \(\mathfrak n\). Let \(S=XV\)
and \(s=\dim(S\cap L)\). A basis of \(S\) can be chosen by taking
\(s\) vectors in \(L\), including \(z\) as an additional vector
if \(h=1\), and then choosing \(\ell-s-h\) vectors in \(V\).
This is possible since \(z\notin L\). It gives at most
\(q^{as+r(\ell-s-h)}\) choices of \(S\), allowing repetitions
of a given space. For fixed \(S\), the images of \(L+W\)
belong to \(S\cap L\) and the remaining \(a\) basis vectors have
images in \(S\). This gives at most \(q^{(a+d)s+a\ell}\) maps.
Since \(r=2a+d\), the exponent of \(s\) cancels. Thus the total
number \(N_{\ell,h}\) of such \(X\) satisfies
\[
                   N_{\ell,h}\le(\ell+1)q^{r(\ell-h)+a\ell}.
\]

We next bound every relevant orbit from below. For a nilpotent
\(X\) with Jordan block lengths \(t_1,\ldots,t_m\), its
commuting algebra has dimension \(\sum_{i,j}\min(t_i,t_j)\).
Indeed a map between two cyclic Jordan blocks is determined by
the image of a generator, which must be killed by the
corresponding power of \(X\). Writing \(b_j=\#\{i:t_i\ge j\}\)
turns this dimension into \(\sum_{j\ge1}b_j^2\). Here
\(b_1=r-\ell\) and \(\sum_{j\ge2}b_j=\ell\), so the dimension
is at most \((r-\ell)^2+\ell^2\). The group of invertible
commuting maps is contained in this algebra. Also
\(|\bar G|=q^{r^2}\prod_{i=1}^r(1-q^{-i})\ge q^{r^2}/4\).
For the last inequality use \(q\ge2\), take the first two factors
at \(q=2\), and bound the remaining product below by
\(1-\sum_{i\ge3}2^{-i}=3/4\). Orbit--stabilizer gives
\(|\bar G\cdot X|\ge q^{2\ell(r-\ell)}/4\).

Set \(D=C_{\bar G}(X)\). Applying orbit--stabilizer also to \(z\)
under \(\bar G\) and \(D\) yields
\[
            |\bar Z\cdot X|
       =|\bar G\cdot X|\frac{|D z|}{q^r-1}.
\]
If \(h=0\), split \(V=V_1\oplus V_0\) into the blocks of length at
least two and the blocks of length one. Each length \(t\ge2\)
uses \(t\le2(t-1)\) dimensions. Thus \(\dim V_0\ge r-2\ell\).
If \(z\) has nonzero \(V_0\) component, \(\mathrm{GL}(V_0)\)
fixing \(V_1\) pointwise already gives at least \(q^{r-2\ell}-1\)
values of \(Dz\). Otherwise \(z\in V_1\setminus XV\). Take a
linear functional \(\lambda\) which is one on \(z\) and vanishes
on \(XV+V_0\). For \(v\in V_0\) the maps
\(1+v\lambda\) are invertible, commute with \(X\), and send \(z\)
to the distinct vectors \(z+v\). This proves the same lower
bound. At the nodes in question \(r-2\ell\ge1\). When \(h=1\)
we use \(|Dz|\ge1\). We have obtained the two orbit bounds
\[
 |\bar Z\cdot X|\ge
 \begin{cases}
 q^{2r\ell-2\ell^2-2\ell}/8,&h=0,\\
 q^{2r\ell-2\ell^2-r}/4,&h=1 .
 \end{cases}
\]

In \eqref{eq:finiteposition881}, the sum of the squared
intersection sizes for given \((\ell,h)\) is at most
\(N_{\ell,h}^2\). The contribution with \(h=0\) is therefore at
most \(8(\ell+1)^2q^{2a\ell+2\ell^2+2\ell}\). For \(h=1\) it
is at most \(4(\ell+1)^2q^{-r+2a\ell+2\ell^2}\). Include the
zero orbit and sum over \(1\le\ell\le2a\). The result is at most
\(1+24a(2a+1)^2q^{12a^2+4a}\), which is at most \(q^{20n^2}\)
at all sufficiently large nodes. Finally \([K:Z]=q^r-1\) since
\(\bar G\) acts transitively on the nonzero vectors.
This proves \eqref{eq:fixernormalized881}. On our interval
\(q\ge b^{10}\), \(n\le\sqrt{x}\), and
\(k\ge x^5-2\sqrt{x}\). In particular \(20n^2+a=o(k)\).
Thus \(2q^{-r+20n^2+a}\le b^{-9k}\) for all sufficiently
large nodes. Multiplication by \(B_0\) proves the last assertion.
\end{proof}

Corollary~\ref{cor:terminalfixer878} costs at most \(q^a\) in
passing from a selected frame vector to \(z\) on the defective
arm locus. The same proof applies with the mask \(s_w\notin P_w\)
of Corollary~\ref{cor:armdefect879}, since the full arm minor
is invariant under both residual unipotents. At a terminal fixer
one may enlarge this residue set to all of \(\bar Z\) in the sum
of squared norms before expanding the squares. Proposition~%
\ref{prop:twocorrelations881} identifies the resulting local
residue pairing and its companion position correlation.
Theorem~\ref{thm:integralmass881} estimates the latter on \(K\).
To bound the original arithmetic second moment by this method
one must also treat \(G\setminus K\) and justify the comparison
with the sum over rational heads and the archimedean factors.
\clearpage

\subsection{The identity term on integral positions at an unrestricted mark}
\label{sec:unrestrictedmass882}

Continue with \(V=L\oplus W\oplus R\), where \(R=L^\vee\) has
its basis in reverse dual order. Write \(L_i\) for the span of the
first \(i\) basis vectors of \(L\). In the notation of
Lemma~\ref{lem:finiteposition881}, every \(X\in\mathfrak n\) has
the form
\[
 X=\begin{pmatrix}A&B&C\\0&0&D\\0&0&E\end{pmatrix}.
\]
Here \(A,E\) are strictly upper triangular \(a\)-squares,
\(C:R\to L\), \(B:W\to L\), and \(D:R\to W\). Denote the
rows of \(B\) by \(\beta_i\) and the first column of \(D\) by
\(\delta\). The actual residue character has phase
\begin{equation}
 \Lambda(X)=\sum_{i<a}c_i A_{i,i+1}
       +c_a\beta_a(z)+\eta(\delta)
       +\sum_{i<a}e_i E_{i,i+1},
 \label{eq:actualphase882}
\end{equation}
where the \(c_i,e_i\) are nonzero and \(\eta\in W^*\) satisfies
\(\eta(z)\ne0\). Indeed the first split component of the frame
character of Section~\ref{sec:rationalflags680} reads the
successive left entries and \(\beta_a(z)\). In the other component
the isometry is the inverse adjoint. On consecutive blocks its
off-diagonal part is minus the adjoint of the corresponding
block, since every term of higher degree skips a block. In the
reverse dual basis this reads the consecutive right entries and
the pairing of the first column of \(D\) with the paired terminal
vector. The two scalar coefficients coming from the conductor are
units by Lemma~\ref{lem:residueshift878}, and the pairing of the
two terminal vectors is the unit \(-N\). This proves
\eqref{eq:actualphase882} also with the components interchanged.

For a \(\bar G\)-conjugacy class \(\mathcal C\) of nilpotent
matrices set \(S(\mathcal C)=\sum_{X\in\mathcal C\cap\mathfrak n}
\psi(\Lambda(X))\). The next estimate keeps cancellation inside
each class.

\begin{proposition}[One middle dimension saved in the class sums]
\leavevmode\label{prop:middlecancellation882}
For \(1\le\ell\le2a\), \(d=\dim W=k+1\), one has
\begin{equation}
 \sum_{\operatorname{rank}\mathcal C=\ell}|S(\mathcal C)|
   \ \le\ 10q(a+1)^2q^{4a^2+d(\ell-1)} .
 \label{eq:middlecancellation882}
\end{equation}
\end{proposition}
\begin{proof}
Put \(b'=\operatorname{rank}B\), \(c'=\operatorname{rank}D\)
and \(T=(a+1)^2q^{4a^2}\). The image on \(W\) has dimension
\(b'\) in \(L\), and projection of the full image modulo \(L\)
has dimension at least \(c'\). Hence \(b'+c'\le\ell\). For
prescribed ranks the number of pairs \(B,D\) is at most
\(q^{(d+a)(b'+c')}\), by factoring each map through a space of
its rank. The choices for \(A,C,E\) number
\(q^{2a^2-a}\). There are at most \((a+1)^2\) rank pairs, and
\(a(b'+c')\le2a^2\). Thus the total count at rank \(\ell\) is
at most \(Tq^{d\ell}\), and the count with \(b'+c'<\ell\) is
at most \(Tq^{d(\ell-1)}\).

There is another set of size at most \(Tq^{d(\ell-1)}\): those
matrices with \(\eta\) in the row space of \(B\). In that case
\(b'\ge1\). A basis for the row space can be chosen starting with
\(\eta\), so \(b'-1\) arbitrary covectors and at most \(ab'\)
scalar coefficients suffice to specify \(B\). The same bound
for \(D,A,C,E\) as above gives the asserted count.

First consider all matrices with \(\beta_a\ne0\).
Conjugation by \(\mathrm{GL}(W)\), acting as the identity on the
two arms, preserves this set and each class. Average the phase
over these conjugations. The middle part becomes
\(c_a\beta_a(w)+\mu(\delta)\), with \((w,\mu)\) uniform on
\[
    w\ne0,\qquad \mu(w)=\eta(z).
\]
Indeed \((w,\mu)=(g^{-1}z,\eta g)\), and a basis of \(\ker\mu\)
together with \(w\) shows transitivity, since \(\eta(z)\ne0\).
When \(\delta=0\), the average of the additive character has
absolute value \(1/(q^d-1)\), by summing a nontrivial linear
character over nonzero \(w\). When \(\delta\ne0\), average first
over all \(\mu\) with fixed value on \(w\). This gives zero unless
\(\delta\in\mathbb F_qw\), an event of probability
\((q-1)/(q^d-1)\). Both averages are bounded in absolute value
by \(2q^{1-d}\). The other entries of the phase are unchanged.
By the triangle inequality the sum of the absolute values of
the class contributions from this set is at most
\(2qTq^{d(\ell-1)}\).

Next suppose \(B\ne0\) and its last nonzero row is \(j<a\).
Then \(B(W)\subset L_j\). If \(X(L)\not\subset L_j\), the
image of \(L+W\) has dimension at least \(b'+1\), which implies
\(b'+c'<\ell\). These matrices have the smaller count already
proved. If \(\beta_j\) is proportional to \(\eta\), it is a
nonzero multiple, so the second smaller count applies.

On the rest of this set, fix \(B\) and a class. Let \(\xi\) be
the coordinate of the \((j+1)\)-st left basis vector, extended
by zero on \(W+R\). For \(u\in\ker\eta\) conjugate by
\(1+u\xi\), with inverse on the left. Here \(Xu=Bu\),
\(\xi(Xu)=0\), and \(\xi X\) vanishes on \(L+W\). The new
matrix is
\[
             X'=X+(Bu)\xi-u(\xi X).
\]
It belongs to \(\mathfrak n\), preserves the value of \(B\), and
still has \(X'(L)\subset L_j\). Indeed \((Bu)\xi\) takes the
\((j+1)\)-st line into \(L_j\), and the second change sends only
\(R\) into \(W\). These conjugations permute the set under
consideration. The phase changes by \(c_j\beta_j(u)\), since the
change involving the first right column is killed by \(\eta(u)=0\).
The restriction of \(\beta_j\) to \(\ker\eta\) is nonzero.
Averaging over \(u\) makes the class sum on this set vanish.

It remains to treat \(B=0\). If \(\ell>c'\), the first smaller
count applies. Otherwise \(\ell=c'>0\). Transpose the matrix and
reverse the two complete outer flags, moving the old right arm
to the left. This bijection carries the matrices to the same
block form on the dual space, now with \(B\) of rank \(\ell\)
and \(D=0\). It carries conjugacy classes to conjugacy classes,
and interchanges the two nonzero terminal data in
\eqref{eq:actualphase882}, up to nonzero scalars. The pairing of
the new terminal data is still nonzero. The preceding averaging
arguments apply on this restricted set: the middle basis
changes preserve \(D=0\), and in the case \(j<a\) the whole
image of \(X\) equals \(B(W)\subset L_j\), since both have
dimension \(\ell\). Consequently \(\xi X=0\) and the shears
also preserve \(D=0\). The exceptional row-space count holds on
this subset as well.

All cancellations and basis averages were taken within a single
class. Summing their bounds over classes uses the total counts
above. The two middle basis averages cost at most
\(4qTq^{d(\ell-1)}\) together. The smaller counts from the
remaining cases, including the transposed row-space condition,
cost at most \(4Tq^{d(\ell-1)}\). This proves
\eqref{eq:middlecancellation882}.
\end{proof}

\begin{theorem}[Integral mass with its identity term]
\leavevmode\label{thm:unrestrictedmass882}
At every original marked place, for all sufficiently large nodes,
the full integral position mass with \(H=K\) satisfies
\begin{equation}
 0\le A_K-1
   =\frac{\int_{K\setminus J}|C_K(x)|^2\,dx}
              {\operatorname{vol}(J)\operatorname{vol}(K)^2}
       \le q^{-2d+20n^2}.
 \label{eq:unrestrictedmass882}
\end{equation}
\end{theorem}
\begin{proof}
On \(J\) the conjugated function in \(C_K\) is identically one
because \(J\) is normal in \(K\). The zero orbit in
\eqref{eq:finiteposition881} accounts for exactly this integral.
At rank \(\ell>0\), every \(\bar G\)-orbit has size at least
\(q^{2r\ell-2\ell^2}/4\), by the commuting-algebra calculation in
the proof of Theorem~\ref{thm:integralmass881}. The sum of the
squared absolute class sums is bounded by the square of
\eqref{eq:middlecancellation882}. Since \(r=d+2a\) and
\(2\ell^2-4a\ell\le0\) for \(\ell\le2a\), the nonzero ranks
together contribute at most
\[
        800a(a+1)^4 q^{-2d+8a^2+2}.
\]
The factor apart from \(q^{-2d}\) is at most \(q^{20n^2}\) at
all sufficiently large nodes. This proves the stated bound.
\end{proof}

\begin{corollary}[The product of integral masses at the original marks]
\leavevmode\label{cor:integralproduct882}
Let \(\mathcal W\) be the \(2f\) marked places of
Theorem~\ref{thm:bothmarks873}. At \(w\in\Xi\) use the stabilizer
\(H_w=Z_w\) of \(z\) in the prescribed component, and at the
other places take \(H_w=K_w^{\rm int}=\mathrm{GL}_r(\mathcal O_w)\).
Here \(K_w^{\rm int}\) denotes the integral group at any marked place.
Write
\(M_w(H)=A_H\,\operatorname{vol}(H)/\operatorname{vol}(K_w^{\rm int})\).
For the top indices \(j_w\) of the original selected frame vectors,
\[
 B_0^{|\Xi|}
 \left(\prod_{w\in\Xi}q_w^{n-j_w}\right)
 \prod_{w\in\mathcal W} M_w(H_w)
      \le b^{-3k|\Xi|}\bigl(1+O(f b^{-19k})\bigr).
\]
\end{corollary}
\begin{proof}
At a selected place apply \eqref{eq:fixernormalized881} and
\(n-j_w\le a\). At every other place
\eqref{eq:unrestrictedmass882} gives \(M_w(K_w^{\rm int})\le
1+b^{-19k}\), since \(d=k+1\), \(20n^2=o(k)\) and \(q_w\ge b^{10}\).
The product of at most \(2f\) such factors is bounded by
\(\exp(2f b^{-19k})=1+O(f b^{-19k})\), using the polynomial
bound for \(f\) on the interval.
\end{proof}

The product here integrates independent integral positions at the
original marked completions. The common rational heads and the real
integrals in Theorem~\ref{thm:orbitfactors871} also enter the
projected period. An estimate for that second moment must still
account for the nonintegral positions and for those correlations.
\clearpage

\subsection{The two middle root coordinates on nonintegral positions}
\label{sec:allpositions883}

Fix an original split mark. Write \(O=\mathcal O_w\), choose a
uniformizer \(\pi\), and take an additive character \(\Psi\) of \(F_w\)
of conductor \(O\). In a split integral basis the contribution to
\(\chi(u)\) of the two simple blocks adjacent to \(W\) is
\[
             \Psi\big((Pz+\eta v)/\pi\big).
\]
Here \(P\) is the last row of the block \(W\to L\), multiplied
by its character unit, and \(v\) is the first column of the
block \(R\to W\). The primitive covector \(\eta\) includes the
unit multiplying the opposite terminal root, and \(c=\eta z\in
O^\times\). These statements over \(O\) follow from the frame
formula of Section~\ref{sec:rationalflags680} by inverse adjunction
on consecutive blocks, just as for \eqref{eq:actualphase882}.
The remaining simple entries on the two arms will be denoted by
\(s_i\), including their character units. If \(y=u j\in UJ\),
any two choices of \(u\) differ on the right by \(U\cap J\).
Consecutive blocks add under multiplication in \(U\), so the
classes of \(s_i,P,v\) modulo the corresponding \(\pi O\)-lattices
depend only on \(y\).

Put \(t=0\) if both \(P\) and \(v\) lie in their \(\pi O\)-lattices.
Otherwise let \(t\ge1\) be least with \(\pi^{t-1}P,\pi^{t-1}v\)
integral. Write
\[
 R_t=O/\pi^tO,\qquad p=\pi^{t-1}P\bmod\pi^t,\qquad
 b'=\pi^{t-1}v\bmod\pi^t,\qquad
 \psi_t(h)=\Psi(\pi^{-t}h).
\]
At least one of \(p,b'\) is primitive over \(R_t\), meaning that
its reduction modulo \(\pi\) is nonzero.

\begin{lemma}[The finite middle average at precision \(t\)]
\leavevmode\label{lem:middleaverage883}
Let \(F_{K,t}\) be the normalized average of the middle character
under changes of integral basis in \(W\), acting trivially on
both arms. Define \(F_{Z,t}\) with the additional condition that
the basis change fixes \(z\) modulo \(\pi\). At \(t=0\) both
averages equal one. For \(t\ge1\),
\begin{equation}
 F_{K,t}=\frac{\displaystyle\sum_{\alpha\in R_t^\times}
                 \psi_t(-c\alpha-\alpha^{-1}p(b'))}
                {q^{td}(1-q^{-d})},\qquad
 |F_{K,t}|\le 2q^{-t(d-1)}.
 \label{eq:fullmiddle883}
\end{equation}
If \(b'\) is primitive,
\begin{equation}
 F_{Z,t}=q^{-(t-1)d}
      \sum_{\substack{h\in R_t^\times\\
                      \bar h\,\bar z=\overline{b'}}}
                    \psi_t(ch+h^{-1}p(b')) .
 \label{eq:fixedmiddle883}
\end{equation}
For \(t\ge2\) the latter average is zero if \(b'\) is
nonprimitive. In all cases \(t\ge1\) one has
\(|F_{Z,t}|\le q^{-(t-1)(d-1)}\).
For either average to be nonzero when \(t\ge2\), the two
frequencies \(p,b'\) must both be primitive, and
\(\overline{p(b')}/\bar c\) must be a nonzero square in
\(\mathbb F_q\). For \(F_{Z,t}\) one further needs
\(\overline{b'}=\bar h\bar z\) with
\(\bar h^2=\overline{p(b')}/\bar c\).
\end{lemma}
\begin{proof}
The pair \((w,\mu)=(g^{-1}z,\eta g)\) for a uniform integral
basis change reduces to the uniform distribution on
\(\mu(w)=c\) over \(R_t\). To see transitivity, \(w\) is primitive
and \(R_t^d=R_t w\oplus\ker\mu\) because \(c\) is a unit.
A basis of its kernel is obtained by solving for a coordinate
with unit coefficient. Send \(w\) to \(z\) and a kernel basis
to a basis of \(\ker\eta\). This gives \(g\) over \(R_t\),
and any lift of its entries is invertible over \(O\). If
\(g\) is to fix \(\bar z\), precisely the pairs with
\(w\equiv z\bmod\pi\) occur, by the same construction. A
transitive action of a finite group gives uniform measure on its
orbit. For each primitive \(w\) there are \(q^{t(d-1)}\)
covectors with \(\mu(w)=c\).

For the full average insert
\(q^{-t}\sum_{\alpha\in R_t}\psi_t(\alpha(\mu(w)-c))\)
in the unrestricted sum of \(\psi_t(pw+\mu b')\) over \(w,\mu\).
The insertion is the indicator of \(\mu(w)=c\): the additive
pairing of \(R_t\) with itself is nondegenerate by the
conductor of \(\Psi\). If \(b'\) is primitive, summing over
\(\mu\) requires \(\alpha w=-b'\), so that \(\alpha\) is a unit
and \(w=-\alpha^{-1}b'\). The surviving factor from the
\(\mu\)-sum is \(q^{td}\). If \(p\) is primitive instead,
sum first over \(w\) to obtain the same expression. Dividing by
\(q^{td}(1-q^{-d})q^{t(d-1)}\), the number of pairs, gives
\eqref{eq:fullmiddle883}. The unit sum has at most \(q^t\) terms.

For the stabilizer average fix \(w\equiv z\bmod\pi\). The
average of \(\psi_t(\mu b')\) over \(\mu(w)=c\) is zero unless
\(b'=h w\) for some \(h\in R_t\); this follows by extending \(w\)
to a basis and varying the remaining coordinates of \(\mu\).
When it exists \(h\) is unique and the average equals
\(\psi_t(ch)\). If \(b'\) is primitive, \(h\) is a unit and
\(w=h^{-1}b'\). There are \(q^{(t-1)d}\) possible \(w\)'s before
this restriction. This gives \eqref{eq:fixedmiddle883}, with at
most \(q^{t-1}\) terms. If \(b'\) is nonprimitive and \(t\ge2\),
any surviving \(h\) lies in \(\pi R_t\), and \(p\) is primitive.
For every residue vector \(l\) the replacement
\(w\mapsto w+\pi^{t-1}l\) preserves \(b'=hw\), the class of \(w\)
modulo \(\pi\), and \(h\). The factor \(\psi_t(pw)\) averages to
zero over these replacements. At \(t=1\) the elementary bound
one suffices.

For \(t\ge2\) translate a unit summation variable \(h\) by
\(\pi^{t-1}s\), \(s\in\mathbb F_q\). Modulo \(\pi^t\) its inverse
changes by \(-\pi^{t-1}s h^{-2}\), since \(2(t-1)\ge t\).
Summing this last digit in either displayed formula gives zero
unless \(\bar c=\overline{p(b')}\bar h^{-2}\). This forces both
primitive conditions and the asserted square condition. The
calculation also holds in residue characteristic two.
\end{proof}

\begin{theorem}[Middle and arm averages in the actual position correlation]
\leavevmode\label{thm:allpositions883}
For \(H=K\) or \(H=Z\), define on \(UJ\)
\[
 D_H(y)=F_{H,t}\prod_i R(s_i),\qquad
 R(s)=
 \begin{cases}
 1&s\in\pi O,\\
 -1/(q-1)&s\in O^\times,\\
 0&s\notin O .
 \end{cases}
\]
Extend \(D_H\) by zero off \(UJ\). For every \(x\in G\),
\begin{equation}
             C_H(x)=\int_H D_H(hxh^{-1})\,dh .
 \label{eq:allpositions883}
\end{equation}
In particular, retain in this last integral only contributing
conjugates of depth \(t\ge T\), where \(T\ge2\). Denote by
\(\sigma_H(x,T)\) the Haar probability in \(H\) that
\(hxh^{-1}\in UJ\) has \(t\ge T\), integral outer simple entries,
and satisfies the necessary middle conditions of
Lemma~\ref{lem:middleaverage883}. The absolute value of this part
of \(C_H(x)/\operatorname{vol}(H)\) is at most
\[
 \begin{cases}
      2q^{-Tk}\sigma_K(x,T),&H=K,\\
      q^{-(T-1)k}\sigma_Z(x,T),&H=Z,
 \end{cases}
 \qquad k=d-1 .
\]
\end{theorem}
\begin{proof}
The middle basis changes used in the lemma are contained in \(H\)
and normalize \(U\) and \(J\). The same holds for diagonal units
on the outer lines, with the left line next to \(W\) and the
first right line both held fixed. The successive outer diagonal
ratios run independently through \(O^\times\): their values
recover the outer units recursively from the two fixed values.
These changes scale the respective simple entries independently
and preserve \(P,v\). The middle changes preserve all the outer
simple entries. Averaging a factor \(\Psi(su/\pi)\) on
\(u\in O^\times\) gives \(R(s)\), by subtracting its integral
on \(\pi O\) from that on \(O\), then dividing by \(1-q^{-1}\).
Multiplicative Haar probability on the units equals normalized
additive measure there, since multiplication by a unit preserves
additive measure.

For fixed \(x\), average \(\phi(lhxh^{-1}l^{-1})\) first over
the product of these two compact groups of basis changes \(l\).
Membership in \(UJ\) is preserved. On its support the middle
and outer averages just computed give \(D_H(hxh^{-1})\).
Averaging in \(h\) gives \(C_H(x)\) by left invariance on \(H\).
All integrands are bounded on compact domains. This argument
also preserves \(t\) because the basis changes are invertible
over \(O\). Now apply Lemma~\ref{lem:middleaverage883} and
\(|R(s)|\le1\) on the indicated set.
\end{proof}

This calculation applies to every position outside \(K\) as well.
It detects large denominators in the two simple blocks adjacent
to \(W\). A large entry in a block farther above the diagonal may
occur with both of those coordinates in \(\pi O\). Controlling
the measure of such positions, and integrating the depth bounds
against the remaining local factors, are still required for the
full position mass.
\clearpage

\subsection{Summing over nonintegral positions at the marked places}
\label{sec:nonspheretail884}

Use the split integral lattice \(\mathcal L=O^r\) and the matrix
norm for which the norm of an integral unit matrix is one and
\(|\pi|=q^{-1}\). For \(x\in G\) let \(m\) be the integer given by
\(q^m=\max(\|x\|,\|x^{-1}\|)\), where the matrix norm is the
maximum of the absolute values of the entries. For an \(F_w\)-linear
map \(T\) on \(\mathcal L\otimes F_w\) put
\[
     p(T)=\operatorname{length}_O
                \big((T\mathcal L+\pi\mathcal L)/\pi\mathcal L\big).
\]
This module is finite. A nonzero Smith entry of valuation \(e\)
contributes \(\max(1-e,0)\) to its length, by diagonalizing over
\(O\) after clearing denominators. Put \(p_+=p(x-1)\),
\(p_-=p(x^{-1}-1)\). Write \(\mathcal S(m,P,Q)\) for the set
on which \((m,p_+,p_-)=(m,P,Q)\).

These sets are invariant under conjugation by \(K\) and
multiplication on either side by \(J\). For the norm this follows
by multiplying back by an integral inverse. For the length use
\(xj-1=(x-1)j+(j-1)\) and \(jx-1=j(x-1)+(j-1)\), with
\((j-1)\mathcal L\subset\pi\mathcal L\). Apply the same calculation
to the inverse. Each shell contains finitely many cosets modulo \(J\).
Indeed the bounds on the matrix and its inverse give a compact set:
the pairs of bounded matrices \(x,y\) satisfying \(xy=yx=1\)
form a closed subset of a compact set, and projection to \(x\)
is continuous. The cosets modulo \(J\) are open. Nonintegral
positions have \(m\ge1\).

\begin{lemma}[Counting cosets and bounding their conjugacy orbits]
\leavevmode\label{lem:shellorbits884}
The number \(N(m,P,Q)\) of cosets \(xJ\) in \(\mathcal S(m,P,Q)\)
meeting \(U\) is at most
\begin{equation}
         q^{4a^2(m+1)+d\min(P,Q)} .
 \label{eq:shellcount884}
\end{equation}
If this number is positive then \(P,Q\le2a(m+1)\). Every
\(K\)-conjugacy orbit on \(G/J\) meeting \(U\) in this shell
has cardinality at least
\begin{equation}
              \tfrac14 q^{(d-2a)(P+Q)}.
 \label{eq:shellorbit884}
\end{equation}
\end{lemma}
\begin{proof}
Write \(u=1+X\in U\) in the three blocks of
Section~\ref{sec:unrestrictedmass882}. Both \(u-1\) and \(u^{-1}-1\)
have rank at most \(2a\), by the image calculation of
Lemma~\ref{lem:finiteposition881}, which works over \(F_w\).
All their entries have valuation at least \(-m\). The assertion
on \(P,Q\) follows from the Smith formula.

The cosets counted in \(N(m,P,Q)\) correspond to the cosets of
\(U\cap J\) in the unipotent shell, hence their count equals its
unipotent Haar measure under our normalization. In upper block
coordinates this measure is the product of additive measures
giving \(\pi O\) volume one: left multiplication is triangular
with diagonal one if the coordinates are ordered by increasing
block separation, and the product assigns volume one to \(U\cap J\).

Set \(h=m+1\). Each coordinate in \(\pi^{-m}O\), reduced modulo
\(\pi O\), has \(q^h\) possibilities of measure one. The two
triangular arm blocks and the full block \(C:R\to L\) together
cost at most \(q^{2a^2h}\). Let \(s_B,s_D\) be the lengths
of the images of \(B,D\) on their integral source lattices,
modulo \(\pi\) times their target lattices. Then
\(s_B+s_D\le P\): the image of \(X\) on \(W\) gives the
\(B\)-image in \(L\), and projection to \(W\) gives the
\(D\)-image. These two contributions enter the kernel and the
image, respectively, of that projection on
\((X\mathcal L+\pi\mathcal L)/\pi\mathcal L\).

The columns of \(B\) take values in a submodule of
\((\pi^{-m}O/\pi O)^a\) of length \(s_B\). The rows of \(D\) take
values in such a submodule of length \(s_D\), since transposition
preserves Smith entries. There are at most \(q^{a^2h}\) choices
for either submodule: its inverse image in \((\pi^{-m}O)^a\)
is a free lattice of rank \(a\), so it has \(a\) generators
after reduction. Once the two submodules are fixed, their \(d\)
columns and \(d\) rows together give at most \(q^{d(s_B+s_D)}\)
choices. This bounds the shell measure by \(q^{4a^2h+dP}\).
Inversion preserves Haar measure on the unimodular group \(U\)
and interchanges \(P,Q\), giving \eqref{eq:shellcount884}.

For the orbit estimate take a representative \(u\in U\). Inside
the middle lattice choose a direct summand containing \(D(R)\)
over \(F_w\) of rank \(\operatorname{rank}D\), by taking the
saturation of its image after clearing denominators. On an
integral complement split off the kernel of \(B\); it is a
direct summand since the quotient is torsion free over a discrete
valuation ring. We obtain a summand \(W_0\) of rank \(d_0\ge d-2a\)
on which \(B=0\), with \(D(R)\) in the complementary vector
space. Consequently in an integral basis the matrix of \(u\)
is \(Y\oplus1_{d_0}\), on a complementary summand and \(W_0\).
This integral change of basis preserves both lengths.

An element \(g\) of \(K\) stabilizes \(uJ\) under conjugation
precisely when \(u^{-1}g u-g\in\pi M_r(O)\), by multiplying the
congruence by \(g\) on the right. Each column \(c_0\) in the
block from \(W_0\) to the complement must satisfy
\((Y^{-1}-1)c_0\in\pi O^{r-d_0}\), an index \(q^Q\) condition
by the definition of the image length. Each row \(r_0\) in the
opposite block must satisfy \(r_0(Y-1)\in\pi O^{r-d_0}\), an
index \(q^P\) condition by transposition. In additive probability
on \(M_r(O)\) these conditions are independent across the \(d_0\)
columns and \(d_0\) rows. The measure of the stabilizer is thus
at most \(q^{-d_0(P+Q)}\). On \(K\), additive measure restricted
from the matrix ring is proportional to group Haar: multiplication
by an integral invertible matrix preserves it. The additive
measure of \(K\) is at least \(1/4\), as in the calculation of
\(|\bar G|\) in Theorem~\ref{thm:integralmass881}.
Dividing the two measures proves \eqref{eq:shellorbit884}.
\end{proof}

On every shell contributing to the position mass one has \(P=Q\),
by Lemma~\ref{lem:detlength885} below. The case of separated
lengths in the following estimate consequently has zero mass.

For \(m\ge1\) put \(l_m=\lceil(m+1)/(2a+1)\rceil\). On a shell with
\(|P-Q|\ge l_m\), use the full position correlation \(C_H\). On any
other nonintegral shell restrict the integral \eqref{eq:allpositions883}
to the conjugates of middle depth at least \(l_m\) if \(H=K\),
and at least \(l_m+1\) if \(H=Z\). Denote the resulting function
on \(G\setminus K\) by \(C_H^{*}\).

\begin{theorem}[An integrated bound on nonintegral shells]
\leavevmode\label{thm:nonspheretail884}
For \(H=K\) or \(H=Z\), on the fields and intervals of
Theorem~\ref{thm:bothmarks873}, all sufficiently large nodes satisfy
\begin{equation}
 \frac{\int_{G\setminus K}|C_H^{*}(x)|^2\,dx}
 {\operatorname{vol}(J)\operatorname{vol}(H)\operatorname{vol}(K)}
                 \le q^{-9k/10}.
 \label{eq:nonspheretail884}
\end{equation}
At a selected place this still gives a bound \(b^{-2k}\) after
multiplication by \(B_0q^a\).
\end{theorem}
\begin{proof}
On the shell in question let \(F\) be \(D_H\) of
Theorem~\ref{thm:allpositions883} with the indicated depth
restriction, extended by zero off \(UJ\). This function is
right \(J\)-invariant. The shell integral in
\eqref{eq:nonspheretail884}, with its normalization, equals
\begin{equation}
  \frac1{[K:H]}\sum_{\mathcal C}
       \frac{\left|\sum_{yJ\in\mathcal C}F(y)\right|^2}{|\mathcal C|},
 \label{eq:cosetnorm884}
\end{equation}
where \(\mathcal C\) ranges over \(H\)-conjugacy orbits on the
cosets in that shell. Indeed \(H\subset K\) normalizes \(J\),
so its conjugation average on each such orbit is the ordinary
finite average. Each coset has volume \(\operatorname{vol}(J)\).
For any \(yJ\), its \(H\)-orbit has size at least the size of its
\(K\)-orbit divided by \([K:H]\), by orbit--stabilizer. Thus the
index in \eqref{eq:cosetnorm884} absorbs the possible loss from
passing to \(H\).

Let \(\epsilon\) bound \(|F|\) on this shell. Its support uses at
most \(N(m,P,Q)\) cosets across all orbits. Bound the sum of the
squared orbit counts by the square of the total count.
Lemma~\ref{lem:shellorbits884} bounds \eqref{eq:cosetnorm884} by
\[
 4\epsilon^2 q^{8a^2(m+1)+2d\min(P,Q)-(d-2a)(P+Q)}
       \le 4\epsilon^2 q^{16a^2(m+1)-d|P-Q|},
\]
since \(P+Q\le4a(m+1)\). When \(|P-Q|\ge l_m\) we may take
\(\epsilon=1\). On all other shells Lemma~\ref{lem:middleaverage883}
gives \(\epsilon\le2q^{-k l_m}\), including \(l_m=1\).

There are at most \(2a+1\) possible \(m\)'s with a fixed \(l_m=l\)
and at most \((2a(2a+1)l+1)^2\) length pairs for each.
Uniformly on the interval \(16a^2(2a+1)\le k/32\) eventually,
since \(a<\sqrt{x}\) and \(k\ge x^5-2\sqrt x\).
For either kind of shell the last displayed bound is at most
\(16q^{-31kl/32}\). Sum first over the lengths and then over
\(l\ge1\). The sum of \(l^2 q^{-31k(l-1)/32}\) is bounded by
an absolute constant for large \(k\); the polynomial factor in
\(a\) is absorbed by \(q^{(31/32-9/10)k}\). This proves
\eqref{eq:nonspheretail884}. Finally \(q\ge b^{10}\),
\(B_0=b^{6k}\), and \(a=o(k)\), so
\(B_0q^{a-9k/10}\le b^{-2k}\) eventually.
\end{proof}

What is left in \(C_H-C_H^*\) has \(|P-Q|<l_m\) and middle depth
below the specified threshold. This difference is a correlation
on each shell and can interfere with \(C_H^*\) there. The square
root of \eqref{eq:nonspheretail884} bounds the contribution of
\(C_H^*\) in the triangle inequality for the full \(L^2\) norm.
An estimate for the remaining correlation is therefore still
needed, as is the comparison with rational heads in the arithmetic
second moment.
\clearpage

\subsection{Positions for the original open residues}
\label{sec:mixedopen885}

There is an exact relation between the lengths used in
Section~\ref{sec:nonspheretail884}. There is also a stronger
restriction on the position when the original residue test at an
unselected place is retained. Both statements hold at each of the
two marked primes.

\begin{lemma}[The two image lengths and the determinant]
\leavevmode\label{lem:detlength885}
For \(x\in\mathrm{GL}_r(F_w)\) and the lengths \(P=p(x-1)\),
\(Q=p(x^{-1}-1)\) defined above,
\[
                        Q-P=\operatorname{ord}_w\det x .
\]
In particular every shell contributing to \(C_H\) has \(P=Q\).
\end{lemma}
\begin{proof}
Let \(\mathcal L_x=\{v\in\mathcal L:(x-1)v\in\pi\mathcal L\}\).
The index \([\mathcal L:\mathcal L_x]\) is \(q^P\) by the
homomorphism to \((F_w^r)/\pi\mathcal L\). If \(v\in\mathcal L_x\),
then \(xv=v+(x-1)v\) is integral and
\((x^{-1}-1)xv=-(x-1)v\in\pi\mathcal L\). Applying the same argument
to \(x^{-1}\) proves \(x\mathcal L_x=\mathcal L_{x^{-1}}\).
With \(\operatorname{vol}(\mathcal L)=1\), comparison of volumes
gives \(q^{-Q}=|\det x|q^{-P}\). Every conjugate of an element
of \(UJ\) has determinant a unit, proving the last assertion.
\end{proof}

Put \(\bar G=\mathrm{GL}_r(\mathbb F_q)\). As before \(O_w\)
denotes the residues with all arm pivots units in both split
components. On the open double orbit write a canonical representative
as \(g=w_{\mathbf p}\sigma\), with both arm parameters and the middle
matrix retained. For \(\alpha,\beta\in\bar G\) set
\(s=\alpha\beta^{-1}\) and use the single-residue functions
\(\Phi_\alpha,\Phi_\beta\) of Proposition~\ref{prop:twocorrelations881}
with \(H=K\). Their product at the pair of arguments is
\[
   F_{\alpha,\beta}(g;y)
          =\Phi_\alpha(y)\overline{\Phi_\beta(g^{-1}y)} .
\]
Its \(y\)-integral is the original local factor \(K_w(s;g)\).

\begin{lemma}[Integral elimination on opposite flags]
\leavevmode\label{lem:oppositecosets885}
A lift \(h\in K\) of any \(s\in O_w\) has a unique decomposition
\[
             h=u_L\,g_s(h)\,u_R^{-1}
\]
in which \(u_L,u_R\in U(O)\) and \(g_s(h)\) is an integral canonical
representative. Its three reductions are determined by \(s\).
In addition,
\[
                        UJ\cap h^{-1}UJh=J .
\]
\end{lemma}
\begin{proof}
Use the block ordering \(1^a,d,1^a\) with the dual arm reversed.
Let \(w\) reverse the two arms, acting on the middle by the identity,
so that \(w^{-1}Uw=U^-\) is the lower block unipotent. The canonical
representatives have the form \(w\ell\) with \(\ell\) block diagonal.
For \(w^{-1}h\), the initial determinants at the ends of the first
\(a\) blocks are, up to basis units, the arm minors in the chosen
component. The remaining proper initial block determinants are
units by the complementary-minor formula: the complementary
minors of \((w^{-1}h)^{-1}\) are the arm minors in the paired
inverse-adjoint component, again up to units. Thus block Gaussian
elimination gives a lower unitriangular matrix, an invertible
block diagonal matrix, and an upper unitriangular matrix. At
each stage the next diagonal block is invertible over \(O\),
since its determinant is the ratio of two successive unit
determinants. Subtracting the corresponding multiples of rows
and columns constructs the factors over \(O\). The elimination
is forced and commutes with reduction. Multiplying back by \(w\)
gives the asserted decomposition. The same elimination over \(F_w\)
proves uniqueness even if the unipotents were originally allowed
to be nonintegral.

Every factor in the displayed decomposition normalizes \(J\).
The outer unipotents in \(U(O)\) normalize \(UJ\), and the
middle representative conjugates \(U\) to \(U^-\). Thus the
intersection assertion reduces, up to integral unipotent
conjugation, to \(UJ\cap U^-J=J\). If \(u j_1=v j_2\) with
\(u\in U,\ v\in U^-\) and \(j_i\in J\), then \(v^{-1}u\in J\).
Its lower-times-upper decomposition is unique over \(F_w\).
Elimination on a matrix congruent to the identity gives integral
factors reducing to the identity. Consequently \(u,v\in J\),
as required.
\end{proof}

\begin{theorem}[The position correlation for two possibly unequal heads]
\leavevmode\label{thm:mixedopen885}
Suppose \(s=\alpha\beta^{-1}\in O_w\). For every pair of canonical
parameters \(g,g'\), and every \(x\in G\),
\begin{align}
 &\int_G F_{\alpha,\beta}(g;y)
       \overline{F_{\alpha,\beta}(g';yx)}\,dy \notag\\
 &\hspace{2em}=
       K_w(s;g)\overline{K_w(s;g')}\,{\bf1}_J(x).
 \label{eq:pairposition885}
\end{align}
The measures here give \(J\) and \(U\cap J\) volume one.
In particular integrating in \(x\) recovers the actual product
of orbital factors, with no restriction on the difference of
the two heads beyond the local residue test in those factors.
\end{theorem}
\begin{proof}
Write the unique residue decomposition as
\(s=\bar u_L\,\bar g_s\,\bar u_R^{-1}\). Suppose
\(F_{\alpha,\beta}(g;y)\ne0\). Then \(y\in u k_\alpha J\) and
\(g^{-1}y\in v k_\beta J\) for \(u,v\in U\), so that
\(u^{-1}g v\in k_\alpha J k_\beta^{-1}\) is an integral lift of
\(s\). By Lemma~\ref{lem:oppositecosets885}, \(g,u,v\) are integral,
\(\bar g=\bar g_s\), \(\bar u=\bar u_L^{-1}\) and
\(\bar v=\bar u_R^{-1}\). The only possible coset of \(y\) is
therefore the coset \(y_sJ\) with
\(\bar y_s=\bar u_L^{-1}\alpha\).

Conversely, if \(g\) is integral and \(\bar g=\bar g_s\),
then on \(y_sJ\) one has
\(\overline{g^{-1}y}=\bar u_R^{-1}\beta\).
Lifting the two residual unipotents shows that both factors are
nonzero there. Their product is the unit complex number
\(\bar\chi(\bar u_L^{-1})\overline{\bar\chi(\bar u_R^{-1})}\),
independent of \(g\). In all cases, including incompatible heads,
this proves
\[
 F_{\alpha,\beta}(g;y)=K_w(s;g)\,{\bf1}_{y_sJ}(y),
\]
since the coset has volume one. For \(y\in y_sJ\), \(yx\) belongs
to \(y_sJ\) exactly when \(x\in J\). Multiplying the two
formulas proves \eqref{eq:pairposition885}.
\end{proof}

The completed local orbital integral has an equivalent normalization.
Give the canonical parameters the measure obtained from Haar measure
on the regular set of \(G\) by taking out the left and right
unipotent measures. This is well defined in the elimination
coordinates \(u_1 g u_2^{-1}\): Haar invariance makes the measures
in \(u_1,u_2\) their Haar measures, with the remaining factor a
measure \(d\nu(g)\). Equivalently one may pull back the matrix Haar
density through the invertible elimination coordinates and use
translation invariance in each unipotent variable. For \(s\in O_w\),
\begin{equation}
                   \int |K_w(s;g)|^2\,d\nu(g)=1.
 \label{eq:openisometry885}
\end{equation}
Indeed integrate the indicator of the coset of \(s\) in \(K\),
whose volume is one, in these coordinates. The two unipotents
are each restricted to one principal coset when \(g\) is integral
with residue \(\bar g_s\), and there are no other contributions.
Thus the unsigned double integral is precisely \(|K_w(s;g)|^2\)
at each \(g\).

\begin{corollary}[All unselected positions on the same fields]
\leavevmode\label{cor:unusedposition885}
For every term of the actual marked energy of
Theorem~\ref{thm:bothmarks873}, one may insert in its product of
paired orbital factors the position integrals
\eqref{eq:pairposition885} at all unselected marks simultaneously.
Their joint position dependence is exactly
\(\prod_{w\in\mathcal W\setminus\Xi}{\bf1}_{J_w}(x_w)\).
This applies before summing the rational heads, including in the
presence of conditions linking the residues at different places.
In the independent completed average with uniform residues, the
local open factor relative to all residues is
\[
 \rho_w=\frac{|O_w|}{|\bar G|}
       =\frac{(1-q^{-1})^{2a}}{\prod_{i=d+1}^{r}(1-q^{-i})},
 \qquad
 1-\frac{4af}{q_*}\ \le\
     \prod_{w\in\mathcal W\setminus\Xi}\rho_w\ \le 1 .
\]
\end{corollary}
\begin{proof}
At each unselected mark the actual residue belongs to \(O_w\).
Apply the theorem separately to each pair of heads in the
expanded finite sum in the sphere formula. The factor on the
right of \eqref{eq:pairposition885} depends on \(x\) only through
its indicated principal subgroup. The normalization of each
subgroup makes the new integral identical to the original
summand.

For the last statement use \eqref{eq:openisometry885} and count
the residue decompositions. The two unipotent factors each have
\(|\bar U|\) choices and the diagonal part has
\((q-1)^{2a}|\mathrm{GL}_d(\mathbb F_q)|\) choices. They are all
distinct by uniqueness. Dividing this product by
\(|\mathrm{GL}_r(\mathbb F_q)|\) and using the count by independent
columns gives the stated ratio: the powers of \(q\) cancel since
\(2\dim U+2a+d^2=r^2\). It lies between
\((1-q^{-1})^{2a}\ge1-2a/q\) and one. There are at most \(2f\)
places in the product and each \(q\ge q_*\).
\end{proof}

Thus nonintegral positions at the unselected marks disappear when
one keeps their original open residue tests. The remaining marked
places can have defective flags. The continuous local calculation
does not replace the rational summation and its archimedean
weights in the arithmetic second moment.
\clearpage

\subsection{Terminal stabilizers through a growing range of denominators}
\label{sec:terminalrange886}

At a selected split mark let \(Z\) still denote the stabilizer
of \(\bar z\) in \(K\), where the bar indicates the fixed residue
of the terminal vector in the chosen component. Write
\[
  C_Z^+(x)=\int_Z {\bf1}_{UJ}(hxh^{-1})\,dh .
\]
This bounds \(|C_Z(x)|\). We estimate its full second moment on
the first nonintegral shells. The middle root coordinates can
have arbitrary depth on the shells in this estimate.

\begin{lemma}[A moving terminal vector in the stabilizer of a coset]
\leavevmode\label{lem:terminalorbit886}
Take \(u\in U\), written in the three blocks of
Lemma~\ref{lem:shellorbits884}. Denote by \(D:R\to W\) its
right-to-middle block and let
\[
  T_D=D(R)\cap W(O),
\]
where \(D(R)\) is the image over \(F_w\). This is a saturated
sublattice of rank at most \(a\). Let \(\mathcal D\subset K\)
be the group stabilizing \(uJ\) under conjugation. If
\(\bar z\notin T_D/\pi T_D\), viewed inside \(W(O)/\pi W(O)\),
then
\[
                  |\mathcal D\bar z|\ge q^{d-2a}-1 .
\]
For \(u\) distributed by unipotent Haar measure on any nonempty
shell \(\mathcal S(m,P,P)\cap U\), the fraction for which
\(\bar z\) belongs to this residue subspace is at most
\((q^a-1)/(q^d-1)\).
\end{lemma}
\begin{proof}
Choose an integral complement \(V'\) to \(T_D\) in \(W(O)\).
On \(V'\) let \(W_0\) be the kernel of the other middle block
\(B:W\to L\). The kernel is saturated, so it has an integral
complement in \(V'\), and its rank \(d_0\) is at least \(d-2a\).
By assumption the projection of \(\bar z\) to \(\overline{V'}\)
is nonzero. The complement to \(W_0\) in \(V'\) can be chosen
so that the \(W_0\)-component of \(\bar z\) is also nonzero.
In fact if this component already is nonzero there is nothing
to choose. Otherwise change a complementary basis vector having
a nonzero coefficient in the projection by adding a primitive
vector of \(W_0\). This is still an integral basis and makes
the \(W_0\)-component nonzero. We use \(d-2a\ge1\), as holds
at these nodes.

Now \(u\) is the identity on \(W_0\) and preserves its chosen
complement in the full space: \(B\) vanishes on \(W_0\) and
the image of \(D\) is in \(T_D\otimes F_w\). The copy of
\(\mathrm{GL}(W_0)(O)\) acting trivially on that complement
commutes with \(u\). Its action on the nonzero residue component
of \(\bar z\) gives \(q^{d_0}-1\) values.

For the fraction assertion conjugate \(u\) by integral changes
of basis in \(W\) acting trivially on the arms. This preserves
the shell, \(U\), and its Haar measure: the automorphism preserves
\(U\cap J\), which has volume one. It carries \(T_D\) to its
image under the basis change. A uniform integral change sends a
fixed nonzero residue vector uniformly over the \(q^d-1\)
nonzero vectors. A residue subspace of dimension at most \(a\)
contains at most \(q^a-1\) of them. Integrate this inequality
in \(u\) on the finite-measure shell.
\end{proof}

\begin{theorem}[An unsigned estimate for all positions of bounded denominator]
\leavevmode\label{thm:terminalrange886}
For a shell \(\mathcal S=\mathcal S(m,P,P)\) meeting \(UJ\)
one has, at every sufficiently large node,
\begin{equation}
 \frac{\int_{\mathcal S}C_Z^+(x)^2\,dx}
      {\operatorname{vol}(J)\operatorname{vol}(Z)\operatorname{vol}(K)}
          \le 24 q^{-d+2a+16a^2(m+1)} .
 \label{eq:terminaloneshell886}
\end{equation}
Let \(\mathcal B\) be the union of all shells with
\(m+1\le k/(1024(a+1)^2)\). Uniformly at either original mark
on the same interval,
\begin{equation}
 \frac{\int_{\mathcal B} C_Z^+(x)^2\,dx}
      {\operatorname{vol}(J)\operatorname{vol}(Z)\operatorname{vol}(K)}
                    \le q^{-9k/10}.
 \label{eq:terminalrange886}
\end{equation}
This includes every defective arm configuration in this range.
\end{theorem}
\begin{proof}
Use \eqref{eq:cosetnorm884} with \(H=Z\) and
\(F={\bf1}_{UJ}\). For \(yJ\) let \(\mathcal D_y\) be its
\(K\)-stabilizer and let \(\mathcal C_K(y),\mathcal C_Z(y)\) denote
its two conjugacy orbits. Orbit--stabilizer gives the exact formula
\[
  [K:Z]\,|\mathcal C_Z(y)|
          =|\mathcal C_K(y)|\,|\mathcal D_y\bar z|.
\]
Indeed the stabilizer in \(Z\) is
\(\mathcal D_y\cap Z\), and its index in \(\mathcal D_y\)
is the last orbit size. That size is constant on each
\(Z\)-orbit since \(Z\) fixes \(\bar z\).

Put \(N=N(m,P,P)\) and \(D_0=q^{2(d-2a)P}/4\).
By Lemma~\ref{lem:shellorbits884} the \(K\)-orbit of each
contributing coset has size at least \(D_0\). On the \(Z\)-orbits
with \(|\mathcal D_y\bar z|\ge q^{d-2a}/2\), the normalized sum
of squared intersection counts is at most
\(2q^{-d+2a}N^2/D_0\). On the other orbits every unipotent
representative has \(\bar z\in\overline{T_D}\), by the preceding
lemma. The union of their unipotent representatives is a union
of cosets of \(U\cap J\), each of volume one. Its total measure
is at most \(2q^{a-d}N\), by the fraction assertion of the
lemma applied to the full shell. This use of measure permits
\(T_D\) itself to vary with the representative inside a coset.
Their contribution to the squared sum is therefore at most
\(4q^{2a-2d}N^2/D_0\).

The count \eqref{eq:shellcount884} and \(P\le2a(m+1)\) imply
\(N^2/D_0\le4q^{16a^2(m+1)}\). Adding the bounds proves
\eqref{eq:terminaloneshell886}. The equality of the two lengths
on every contributing coset is Lemma~\ref{lem:detlength885}.
In \(\mathcal B\) the positive exponent \(16a^2(m+1)\) is
at most \(k/64\). There are at most \(2a(m+1)+1\) possible
\(P\)'s for each \(m\), and at most \(k\) such exponents \(m\).
Since \(d=k+1\), \(a=o(k)\), and both \(a\) and \(k\) are
polynomially bounded in \(x\), this polynomial factor and
\(q^{2a}\) are absorbed by \(q^{(1-1/64-9/10)k}\) eventually.
This proves \eqref{eq:terminalrange886}.
\end{proof}

\begin{corollary}[Absolute comparison for a subset of the fixer residues]
\leavevmode\label{cor:terminalmask886}
For \(S\subset\bar Z\), let
\(c_\alpha(x)=\int_{Jk_\alpha}\phi(hxh^{-1})\,dh\).
Then the original relative-residue test in the compact position
average satisfies
\[
 \frac{1}{\operatorname{vol}(J)|\bar G|^2}
 \int_{\mathcal B}
      \sum_{\alpha\beta^{-1}\in S}|c_\alpha(x)c_\beta(x)|\,dx
        \le q^{-9k/10}.
\]
Thus the same estimate includes the extra defective-arm test.
At a selected place the two factors \(B_0\) and \(q^{n-j_w}\)
together leave a bound \(b^{-2k}\).
\end{corollary}
\begin{proof}
Enlarge \(S\) to \(\bar Z\) in the sum of absolute values and use
\(|c_\alpha|\le\int_{J k_\alpha}{\bf1}_{UJ}(hxh^{-1})dh\).
The allowed pairs are those for which \(\alpha,\beta\) belong
to the same coset \(\bar Z\gamma\). The sum of the latter
integrals in one such coset is
\(C_Z^+(k_\gamma xk_\gamma^{-1})\), by the change of variable
\(h=z_0 k_\gamma\), \(z_0\in Z\). The sum over pairs is thus
bounded by the sum of these squares over the cosets.
Conjugation by \(K\) preserves \(\mathcal B\) and Haar measure.
As \(\operatorname{vol}(J)=1\), one has
\(|\bar G|=[K:Z]\operatorname{vol}(Z)=\operatorname{vol}(K)\).
The resulting normalized integral is precisely the one in
\eqref{eq:terminalrange886}. Finally \(n-j_w\le a=o(k)\),
\(q\ge b^{10}\), and \(B_0=b^{6k}\).
\end{proof}

For denominators beyond this range, the term \(16a^2(m+1)\)
in \eqref{eq:terminaloneshell886} can exceed \(k\).
There the growth of the denominator must enter the saving as
well. The exact calculation at unselected places in
Theorem~\ref{thm:mixedopen885} still holds at every depth.
\clearpage

\subsection{A change of variables for the interior rows of a flag}
\label{sec:armshear887}

We work directly with the class sums in the position integral.
For \(H=K\) or \(Z\), an \(H\)-conjugacy orbit on \(G/J\) will
also denote its inverse image \(\mathcal C\) in \(G\). This set
is invariant under multiplication by \(J\) on both sides.
Indeed it is right \(J\)-invariant and invariant under conjugation
by \(J\subset H\). Its signed intersection count on a fixed shell is
\[
              S_{\mathcal C}=\int_{U\cap\mathcal C}\chi(u)\,du .
\]
This is exactly the sum in \eqref{eq:cosetnorm884} for \(F=\phi\):
each coset meeting \(U\) contributes one coset of \(U\cap J\),
of measure one and constant character. We shall partition this
last integral in unipotent matrix entries.

For a scalar, row, column or matrix \(V\) put
\(e(V)=\min\{j\ge0:\pi^j V\text{ is integral}\}\).
Write the upper unipotent matrix as
\[
 u=\begin{pmatrix}T&B&C\\0&I&D\\0&0&S\end{pmatrix},
 \qquad T=I+A,\quad S=I+E .
\]
The rows of \(B\) are \(\beta_i\). Write \(b_i=e(\beta_i)\),
\(b_*=\max_i b_i\), \(e_D=e(D)\), and let \(s_B,s_D\) be the two
middle image lengths of Lemma~\ref{lem:shellorbits884}. The integer
\(\Delta=P-s_B-s_D\) is nonnegative. The following elementary
relations control the entries outside \(B,D\).

\begin{lemma}[Lengths seen in one left row]
\leavevmode\label{lem:rowlength887}
For every left coordinate covector \(\lambda\),
where \(b_\lambda=e(\lambda B)\),
\begin{align*}
 e(\lambda T)&\le b_\lambda+\Delta,\\
 e\big(\lambda(C-BD)\big)&\le e_D+1+\Delta+b_\lambda,\\
 m&\le 2a\big(b_*+e_D+1+\Delta\big).
\end{align*}
On \(\mathcal S(m,P,P)\) the unipotent measure of the set
\(\Delta\ge L\), for an integer \(L\ge1\), is at most
\(q^{4a^2(m+1)+d(P-L)}\).
\end{lemma}
\begin{proof}
Take the image \(M\) of \(u-1\) in \(F_w^r/\pi\mathcal L\)
on integral source vectors. Projection onto the middle coordinates
has image of length \(s_D\); its kernel contains the image \(M_B\)
on \(W(O)\), lying in the left coordinates and of length \(s_B\).
The quotient of the kernel by \(M_B\) has length \(\Delta\) and
is killed by \(\pi^\Delta\). In particular \(\pi^\Delta A\)
has columns belonging modulo \(\pi\) to \(M_B\), giving the first
assertion. Multiplication by \(\pi^{e_D+1}\) takes \(M\) into the
kernel. Thus the columns of \(\pi^{e_D+1+\Delta}C\) also belong
modulo \(\pi\) to \(M_B\). Applying \(\lambda\), and bounding the
row of \(BD\) by the norms of \(\lambda B\) and \(D\), proves
the second assertion. The same argument on all coordinates
shows that every entry of \(u-1\) has pole at most
\(b_*+e_D+1+\Delta\), because \(\pi^{b_*+1}M_B=0\).
The strictly upper block matrix \(u-1\) has \((2a+1)\)-st
power zero. The finite geometric series for \(u^{-1}\) bounds
its pole by \(2a\) times that of \(u-1\), giving the asserted
bound for \(m\). Finally repeat the submodule count in
Lemma~\ref{lem:shellorbits884} with \(s_B+s_D\le P-L\).
For each pair of permitted submodules the \(d\) columns and
\(d\) rows have at most \(q^{d(P-L)}\) choices, and all other
counts are as in that lemma.
\end{proof}

Fix \(j<a\) and let \(\lambda\) read left coordinate \(j+1\).
Set
\[
 R=\max_{i>j}b_i,\qquad \xi=\lambda T,\qquad
 \gamma=\lambda(C-BD).
\]
For a column \(b\in L\) write \(b=b_{\le j}+b_{>j}\) by its
coordinates. Use the primitive terminal covector \(\eta\) from
Section~\ref{sec:allpositions883}.

\begin{proposition}[A shear which tests the \(j\)-th row]
\leavevmode\label{prop:armshear887}
Let \(L\ge1\) and fix any stratum specified by the row pole
orders \(b_i\), the integer \(e_D\), the shell, and the tests
\[
 b_a<L,\qquad e(\xi)\le b_{j+1}+L,\qquad
 e(\gamma)\le e_D+b_{j+1}+L+1.
\]
The index \(j\) may be chosen as a function of the \(b_i\).
Put \(h=R+L+3\). In the integral of \(\chi\) over this stratum
in any fixed \(\mathcal C\), the contribution from
\begin{equation}
   \beta_j(\ker\eta\cap W(O))\not\subset\pi^{\,1-h}O
 \label{eq:cancelrow887}
\end{equation}
is zero. If \(b_j\ge R+6L\), the fraction of the stratum
satisfying the opposite inclusion, measured on the whole shell
before restricting to \(\mathcal C\), is at most \(2q^{-kL}\).
\end{proposition}
\begin{proof}
For \(w\in\pi^h(\ker\eta\cap W(O))\), put \(b=B w\),
\(v=\lambda b\), and
\(A_w=I-b_{>j}\lambda/(1+v)\) on the left block. The scalar
\(v\) lies in \(\pi O\). Multiply \(u\) on the right and left,
respectively, by the following matrices:
\[
 r_w=\begin{pmatrix}I&0&0\\ w\xi&I&0\\0&0&I\end{pmatrix},
 \qquad
 l_w=\begin{pmatrix}
       A_w&0&0\\ -w\lambda&I+w\lambda B&0\\0&0&I
      \end{pmatrix}.
\]
Both are in \(J\). In fact \(b_{>j}\in\pi^{h-R}L(O)\),
\(w\lambda B\) has entries in \(\pi^{h-b_{j+1}}O\), and
\(w\xi\) in \(\pi^{h-b_{j+1}-L}O\). Direct multiplication shows
that \(l_w u r_w\) belongs to \(U\). Its five blocks are
\begin{equation}
 T'=T+b_{\le j}\xi,\quad B'=A_w B,\quad C'=A_w C,\quad
 D'=D-w\gamma,\quad S'=S .
 \label{eq:shearblocks887}
\end{equation}
For example \(\lambda(T+b\xi)=(1+v)\xi\), which leaves only
\(b_{\le j}\xi\) in the first block. The lower-left block cancels
since \((I+w\lambda B)w=(1+v)w\). The change in \(T\) is
strictly upper triangular as \(\xi\) starts in coordinate \(j+1\).

This operation preserves the stratum. Rows \(i\le j\) of \(B\)
are fixed, and every other row is multiplied on the right by
the integral unit matrix \((I+w\lambda B)^{-1}\).
Thus the individual \(b_i\)'s are preserved. Also
\(\lambda T'=\xi\) and \(\lambda(C'-B'D')=\gamma\), since
\(\lambda A_w=\lambda/(1+v)\) and \(\lambda B w=v\).
The change in \(D\) has entries in \(\pi^{2-e_D}O\), which
preserves its pole order, including the case \(e_D=0\).
Both the coset orbit and the shell are preserved by the two
principal multipliers. The inverse operation uses \(-w\):
the new \(\lambda B'=\lambda B/(1+v)\) and the new
\(b_{>j}\) for the argument \(w\) is \(b_{>j}/(1+v)\).
These identities in \eqref{eq:shearblocks887} give back the
original blocks.

We check preservation of the unipotent Haar measure explicitly.
The row \(R_1=\lambda B\) changes to \(R_1/(1+R_1w)\). On the
ball \(e(R_1)\le R\) this is a bijective isometry, with inverse
\(R_1\mapsto R_1/(1-R_1w)\). Indeed in the difference of two
outputs the change of denominator contributes at most
\(q^{R-h}\) times the norm of the difference of the inputs;
the term with common denominator has exactly that norm. Thus
balls of each radius are taken to balls of the same radius
and additive measure is preserved. Given this row, each other
changed row of \(B\) is multiplied by an integral unit matrix.
For fixed \(B\), the transformation of \(T\) is left
multiplication by \(I+b_{\le j}\lambda\) in the triangular
unipotent group. The transformation of \(C\) has determinant
of absolute value one on each column. Given \(B,C\), the
transformation of \(D\) is multiplication by \(I+w\lambda B\)
followed by a translation. All preserve the relevant measures.
The product formula for the measure on \(U\) now gives the claim
also after restriction to the invariant stratum.

By \eqref{eq:actualphase882} and its integral lift, the change of
the phase before division by \(\pi\) is \(c_j\beta_j(w)\) modulo
\(\pi O\). Indeed the consecutive entries of \(T\) change only
at \(j\), where \(\xi_{j+1}=1\). The change in the last row evaluated
at \(z\) is \(-(\beta_a w)(\lambda Bz)/(1+v)\), which belongs
to \(\pi O\) since \(h-b_a-b_{j+1}\ge1\). The \(D\)-term is
unchanged because \(\eta(w)=0\). The row \(\beta_j\) itself
is unchanged. Average over \(w\) with additive probability.
The character \(w\mapsto\Psi(c_j\beta_j(w)/\pi)\) has integral
zero exactly when it is nontrivial, giving \eqref{eq:cancelrow887}.

For the fraction bound the stated stratum on the whole shell
is invariant under integral middle basis changes: \(BD,T,C\)
are unchanged and the indicated row and matrix pole orders
are preserved. If \(b_j\ge R+6L\), the covector
\(\pi^{b_j}\beta_j\) is primitive. The surviving inclusion
requires that it vanish on the integral kernel of \(\eta\)
modulo \(\pi^{b_j+1-h}\), hence modulo \(\pi^L\), since
\(b_j+1-h\ge5L-2\ge L\). It must therefore be a scalar multiple
of \(\eta\) modulo \(\pi^L\), by choosing a basis adapted to
this primitive covector. Under a uniform integral middle basis
change the primitive row is uniform amongst
\(q^{Ld}(1-q^{-d})\) primitive rows modulo \(\pi^L\), by extension
to an integral basis. There are at most \(q^L\) multiples of
\(\eta\). Integrating this probability over the stratum gives
the result, since \(d-1=k\).
\end{proof}

The change of variables uses the principal group on both sides.
It applies equally to orbits for the terminal stabilizer and for
the full integral group. In using the cancellation one keeps the
whole stratum specified in the proposition; cutting it first by
additional conditions that vary under \eqref{eq:shearblocks887}
could change the integral.
\clearpage

\subsection{The full position mass at a selected mark}
\label{sec:allfixermass888}

The loss in the elementary count of Lemma~\ref{lem:shellorbits884}
grows with the denominator exponent \(m\). We use the interior
row test to obtain a saving proportional to \(m\) on the large
shells. All integrals and characters are normalized as in
Section~\ref{sec:integralmass881}.

\begin{proposition}[Signed class sums on large shells]
\leavevmode\label{prop:largeshell888}
Suppose \(m\ge200n^3\), put \(L=\lfloor m/(200n^3)\rfloor\), and
fix a shell \(\mathcal S(m,P,P)\). For \(H=K\) and for \(H=Z\),
the signed class sums of Section~\ref{sec:armshear887} satisfy
\begin{equation}
 \sum_{\mathcal C}|S_{\mathcal C}|
        \le 10 q^{4a^2(m+1)+dP-kL}.
 \label{eq:largeshell888}
\end{equation}
Here \(\mathcal C\) runs over the \(H\)-orbits on \(G/J\) in
this shell.
\end{proposition}
\begin{proof}
Write \(N_0=q^{4a^2(m+1)+dP}\), an upper bound for the entire
unipotent measure on the shell. Split the integral into the
two parts \(b_*\ge e_D\) and \(b_*<e_D\). On the latter part
transpose the full matrix and reverse the flags to exchange the
arms. On the dual space the new block \(B\) is the transpose
of the old \(D\) in the reversed arm basis, and the new \(D\)
is the transpose of the old \(B\). This change preserves unipotent
Haar measure, the matrix height, the Smith lengths and the sum
\(s_B+s_D\). The phase is again \eqref{eq:actualphase882} in the
new coordinates with the primitive middle vector and covector
exchanged up to units. The orbit sets \(\mathcal C\) are carried
to the corresponding orbit sets for the group induced on the
dual space: the image of a right \(J\)-coset is a left \(J\)-coset,
and the whole conjugacy orbit is bi-\(J\)-invariant. This gives
a bijection of orbit sets by applying the inverse transpose.
In particular the use of principal multipliers in
Proposition~\ref{prop:armshear887} is valid in both orientations.
We may now use \(b_*\ge e_D\) on either part, retaining the strict
inequality on the transposed part.

In a given orientation first take \(b_*<8aL\). If \(\Delta<L\),
Lemma~\ref{lem:rowlength887} would give
\(m\le2a(16a+1)L<200n^3L\), contradicting the choice of \(L\).
Thus this part has the smaller count with \(\Delta\ge L\).

For \(b_*\ge8aL\), treat the case \(b_a\ge L\) by averaging
the middle basis. In original coordinates this event implies
that at least one of the two middle simple blocks has pole order
at least \(L\), including when the dominant arm has been
transposed. All pole comparisons made so far are invariant under
middle basis changes. For \(K\) use the full middle group, and
for \(Z\) use its subgroup fixing the original terminal vector
modulo \(\pi\). These averages preserve each original
\(\mathcal C\). The depth \(t\) in Lemma~\ref{lem:middleaverage883}
is at least \(L+1\). Thus, in the sum of the absolute class
contributions, the total from this part is at most
\(2q^{-kL}N_0\) per orientation.

There remains \(b_*\ge8aL,\ b_a<L\). At some \(j<a\) one has
\[
             b_j\ge R+6L,\qquad R=\max_{i>j} b_i .
\]
Indeed if every increase of the successive maxima, read from
the last row to the first, were less than \(6L\), the final
maximum would be at most \(b_a+6(a-1)L<8aL\).
Fix the row pole orders and \(e_D\), choosing the least \(j\)
with this property. If either bound on \(\xi,\gamma\) in
Proposition~\ref{prop:armshear887} fails, Lemma~\ref{lem:rowlength887}
implies \(\Delta\ge L\). Together with the first case these
exceptions have measure at most \(q^{-dL}N_0\) per orientation.

On the remaining strata apply Proposition~\ref{prop:armshear887}
inside each orbit. The operations in its proof preserve the row
orders, \(e_D\), and hence the orientation choice. After cancellation
only the stated congruence on \(\beta_j\) remains. Its fraction on
the whole stratum is at most \(2q^{-kL}\). Summing over the
disjoint choices of the row orders and \(e_D\) costs at most
\(2q^{-kL}N_0\) per orientation. This application uses the full
strata defined by the row, column and pole tests. We used
\(\Delta\ge L\) only to count the parts which failed one of
those tests. Adding the bounds for both orientations, with
\(d=k+1\), proves \eqref{eq:largeshell888}.
\end{proof}

\begin{theorem}[Complete local mass at the terminal stabilizer]
\leavevmode\label{thm:allfixermass888}
Uniformly at each original split mark of the shorter interval,
at all sufficiently large nodes,
\begin{align}
 \frac{\int_{\{m\ge200n^3\}}|C_H(x)|^2\,dx}
      {\operatorname{vol}(J)\operatorname{vol}(H)\operatorname{vol}(K)}
      &\le q^{-3k/2}\qquad(H=K,Z), \label{eq:alllargemass888}\\
 \frac{\int_G|C_Z(x)|^2\,dx}
      {\operatorname{vol}(J)\operatorname{vol}(Z)\operatorname{vol}(K)}
      &\le 2q^{-9k/10}. \label{eq:allfixermass888}
\end{align}
\end{theorem}
\begin{proof}
Apply \eqref{eq:cosetnorm884} with \(F=\phi\). Its numerator
on each orbit is \(|S_{\mathcal C}|^2\). When the orbit meets
\(U\), orbit--stabilizer and Lemma~\ref{lem:shellorbits884}
give \([K:H]|\mathcal C|\ge q^{2(d-2a)P}/4\).
The sum of squared absolute class sums is bounded by the square of
\eqref{eq:largeshell888}. Using \(P\le2a(m+1)\), this bounds the
normalized integral on one shell by
\[
                     400q^{16a^2(m+1)-2kL}.
\]
For fixed \(L\ge1\) there are at most \(200n^3\) values of \(m\);
each satisfies \(m+1\le400n^3 L\) and has at most
\(2a(m+1)+1\) choices of \(P\). Lemma~\ref{lem:detlength885}
accounts for all the shells with nonzero integrand. On the
given interval \(k\ge x^5-2\sqrt{x}\) and \(n\le\sqrt{x}\),
hence \(6400a^2n^3=o(k)\). For all large nodes the last bound
is at most \(400q^{-31kL/16}\). The sum of
\(Lq^{-31k(L-1)/16}\) is bounded by an absolute constant for
large \(k\). The remaining polynomial factor in \(a,n\) is
absorbed by \(q^{(31/16-3/2)k}\), proving
\eqref{eq:alllargemass888}.

For \(m<200n^3\) one has \(m+1\le k/(1024n^2)\) at sufficiently
large nodes, because \(n^5=o(k)\). The entire contribution for
\(Z\) on that range is at most \(q^{-9k/10}\) by
Theorem~\ref{thm:terminalrange886} and \(|C_Z|\le C_Z^+\).
Adding the first estimate proves \eqref{eq:allfixermass888}.
\end{proof}

\begin{corollary}[The complete marked position integral with the selected tests]
\leavevmode\label{cor:allmarks888}
Let \(\mathcal W,\Xi\) and \(j_w\) be as in
Theorem~\ref{thm:bothmarks873} and
Corollary~\ref{cor:terminalfixer878}. At an unselected place use
the relative-residue test \(S_w=O_w\), and at a selected place
use \(S_w=\bar Z_w\). With the \(c_\alpha\) of
Corollary~\ref{cor:terminalmask886}, put
\[
 \mathcal P_w(S_w)=
 \frac1{\operatorname{vol}(J_w)|\bar G_w|^2}
 \int_{G_w}\left(\sum_{\alpha\beta^{-1}\in S_w}
                  c_\alpha(x)\overline{c_\beta(x)}\right)dx .
\]
The finite sum is taken before integration. At an unselected place
\(\mathcal P_w=\rho_w\). At a selected place it is nonnegative
and at most \(2q_w^{-9k/10}\). In particular
\begin{equation}
 B_0^{|\Xi|}\left(\prod_{w\in\Xi}q_w^{n-j_w}\right)
      \prod_{w\in\mathcal W}\mathcal P_w(S_w)
   \le b^{-2k|\Xi|}\prod_{w\notin\Xi,\ w\in\mathcal W}\rho_w .
 \label{eq:allmarks888}
\end{equation}
When \(\Xi\) is empty the two products agree exactly.
\end{corollary}
\begin{proof}
At an unselected place a pair with \(s=\alpha\beta^{-1}\in O_w\)
can contribute to the integrand only if both conjugates of \(x\)
belong to \(UJ\) for some \(h_1\in Jk_\alpha,\ h_2\in Jk_\beta\).
The residue of \(h_2h_1^{-1}\) is \(s^{-1}\in O_w\), since inversion
conjugates the leading Hermitian Gram determinants. Apply
Lemma~\ref{lem:oppositecosets885} to \(h_1 xh_1^{-1}\) and
\(h_2h_1^{-1}\). It gives \(x\in J\). On \(J\) both single-coset
integrals are one. The normalized result is therefore \(\rho_w\)
by the count in Corollary~\ref{cor:unusedposition885}.

At a selected place group the pairs according to their common
coset \(\bar Z\gamma\). The sum at \(x\) is exactly
\(\sum_\gamma|C_Z(k_\gamma x k_\gamma^{-1})|^2\).
Invariance of Haar measure under \(K\)-conjugation shows that
the normalized integral equals the left side of
\eqref{eq:allfixermass888}. This proves convergence and the
stated estimate with its normalization. Finally \(q_w\ge b^{10}\),
\(n-j_w\le a=o(k)\), and \(B_0=b^{6k}\), so
\(2B_0q_w^{a-9k/10}\le b^{-2k}\) eventually. Multiply the
separate factors. The enlargement to the entire terminal
stabilizer is available in the original sum of squared norms,
as in Section~\ref{sec:integralmass881}.
\end{proof}

Equation~\eqref{eq:allmarks888} evaluates or bounds all the
marked position factors, on the same fields and without a
denominator cutoff. The original arithmetic energy still sums
over rational heads coupled to the real orbital integrals.
Passing from that sum to position integrals requires a global
comparison. The local estimate alone does not give a bound
for the integral projected theta periods or for strict cohomology.
\clearpage

\subsection{Sampling a common pivot at the unselected places}
\label{sec:commoncompletion889}

In passing from the rational sum in \eqref{eq:spheremean872} to
an integral, the equality of two rational pivots and the equality
of their residues play different roles. We keep the original
congruence for unequal rational pivots and complete the additive
test on the actual diagonal. The following argument concerns the
enlarged product mask of Corollary~\ref{cor:integratedmark873}.

Write \(\mathcal V=\mathcal W\setminus\Xi\) for the unselected
places, \(A_w=O_w^{\rm int}\times O_w^{\rm int}\) for the paired
ring of integers there, and \(A_w^\times\) for its group of units.
Here \(O_w^{\rm int}=\mathcal O_{F_w}\). Give each additive ring
volume one. For \(p,p'\in E^\times\) put
\[
 I_w(p,p')={\bf1}_{p,p'\in A_w^\times,\,
                       p\equiv p'\bmod\pi A_w}.
\]
Up to the constant \(q_w^{2a(2k+2a+1)}
|\mathrm{GL}_k(\mathbb F_{q_w})|\), the factor from \(w\) for
two heads \(\mathbf p,\mathbf p'\) and terminal entry \(t\) is
\[
               H_w(t)\prod_{i=1}^a I_w(p_i,p_i'),
\]
by Corollary~\ref{cor:integratedmark873}. Put
\(h_w^*=(1-q_w^{-k})/(1-q_w^{1-k})\).

\begin{theorem}[A positive completion before additive sampling]
\leavevmode\label{thm:commoncompletion889}
In the actual cut or mixed squared energy one obtains an upper
bound by making both of the following replacements at the
unselected places:
\begin{align*}
 H_w(t)&\longmapsto h_w^*\,{\bf1}_{A_w}(t),\\
 I_w(p,p')&\longmapsto
 I_w(p,p')+{\bf1}_{p=p',\ p\in A_w\setminus A_w^\times}.
\end{align*}
The second replacement may be made for any subset of the arm
coordinates. At an unequal pair which survives, the two pivots
are congruent units in both components. At an equal pair
\(p=p'=c\) the completed factor is \({\bf1}_{A_w}(c)\).
Consequently both additive characters on this common entry must
be trivial on \(A_w\) for its local Fourier coefficient to be
nonzero. The same assertion holds for the completed terminal
entry.

Compared with the original local integrals over the common
entries, the ratio of the product of the new zero coefficients
to the old one, when all \(a\) arm entries and the terminal
entry are integrated, is
\begin{equation}
 \prod_{w\in\mathcal V}
 \frac{1-q_w^{-k}}{(1-q_w^{1-k})(1-q_w^{-(k+1)})}
                       (1-q_w^{-1})^{-2a}
                 =1+O(b^{-9}).
 \label{eq:completioncost889}
\end{equation}
This ratio is for the local zero coefficients in additive measure.
The one-sided bound for the coupled rational sum holds before
any zero coefficient is taken.
\end{theorem}
\begin{proof}
For fixed \(t\), the kernel on the finite list of heads supplied
by the remaining real and finite variables is positive
semidefinite. To verify this at the level of the original
integral, hold the other residues and middle elements fixed.
The matrix indexed by \(\mathbf p,\mathbf p'\) is the Gram
matrix of the real tensor \(A_{\mathbf p}(t)\) multiplied by
the corresponding finite scalar \(K(s;w_{\mathbf p}x)\) from
the remaining places. Summing residues or integrating with the
positive sphere and stabilizer measures preserves this property.
The same reasoning applies to the mixed tensors of
Theorem~\ref{thm:mixedperiod876}. The determinant weight and
the sphere-density constants outside these matrices are nonnegative.

By \eqref{eq:terminalseries877} with split sign, \(0\le H_w(t)
\le h_w^*{\bf1}_{A_w}(t)\). For each coordinate \(i\) the kernel
\(I_w\) on the list of heads is a sum of rank-one matrices: for
each unit residue take the outer product of the indicator of
that residue class with itself. The added kernel is likewise
positive semidefinite. Partition the list according to the
\emph{actual rational value} of \(p_i\) and take the rank-one
indicator matrices for the values in \(A_w\setminus A_w^\times\).
Hence every replacement adds a positive semidefinite matrix.
Entrywise products of positive semidefinite matrices are
positive semidefinite, as is seen by using tensor products of
their Gram vectors. Multiplying entrywise by the other kernels
and evaluating the quadratic form on the all-ones vector proves
the asserted upper bound one replacement at a time.

The equal and unequal evaluations follow directly from the formula.
In the unequal case congruence of the units also gives
\((p^{-1}-(p')^{-1})/M\in A_w\), since \(M\) has valuation one at
each component. On the diagonal, the additive Fourier transform
of \({\bf1}_{A_w}\) is its annihilator indicator. Under the
trace pairing divided by the different generator, the annihilator
is \(A_w\) itself. This is the defining duality of the different.
At the terminal entry the same calculation has the constant
\(h_w^*\).

On a common arm entry the original integral is
\((1-q_w^{-1})^2\) and the new one is one. On the terminal entry
the original integral is \(1-q_w^{-(k+1)}\) by
\eqref{eq:terminalmean877}. This proves the product formula.
The logarithm of each terminal ratio is \(O(q_w^{1-k})\) and
that of the arm ratio is \(O(a/q_w)\). Now
\(|\mathcal V|\le2f\) and \(12af/q_*\le b^{-9}\) by
Theorem~\ref{thm:bothmarks873}. Exponentiating proves
\eqref{eq:completioncost889}.
\end{proof}

We record the resulting rational sampling statement with its
covolume, since a loss per infinite place would enter the
leading constant. Use the additive characters of \(E\backslash
\mathbb A_E\) paired by \(\operatorname{Tr}(\,\cdot\,/\delta_E)\).
Take ordinary Lebesgue measure at the complex embeddings and
volume one on all finite integer rings. Let \(T\) be a finite
set of primes of \(E\) disjoint from the primes belonging to
\(\mathcal V\). If \(F\) is an integrable test on the adeles,
write \(\widehat F(y)\) for its transform at the rational
frequency \(y\), using the conjugate additive character.

\begin{proposition}[One-sided sampling after completion]
\leavevmode\label{prop:sampling889}
Suppose \(F\ge0\) is lower semicontinuous and integrable. Outside
\(T\) its finite part is the product of the additive integer-ring
indicators, up to a positive constant independent of the entry.
Inside \(T\) and at infinity the variables may be coupled.
If \(\sum_{y\in E}|\widehat F(y)|<\infty\), then for any finite
set \(D\subset E\),
\[
 \sum_{u\in D} F(u)
 \le \frac1{c_E}\left(\widehat F(0)+
         \sum_{0\ne y\in\mathcal O_{E,T}}|\widehat F(y)|\right),
 \qquad c_E=2^{-f}\sqrt{|\operatorname{disc}E|}.
\]
In particular if, at the scale \(c_v=2\pi i/
(\Phi_v(\delta_E)b_{j,v})\) of Theorem~\ref{thm:freqnode875},
the nonzero coefficients are at most \(B\,W_{T,c,k-n^7}(y)\)
for \(B\ge0\), their total contribution here is at most
\(B b^{-12kf}/c_E\).
\end{proposition}
\begin{proof}
The covolume of the diagonal \(E\) is \(c_E\). Indeed finite
approximation reduces a fundamental domain to the finite integer
rings times a fundamental parallelotope of the Minkowski lattice
of \(\mathcal O_E\), which has that volume. The characters of
the compact additive quotient are indexed by \(E\) for this
pairing, by additive adelic duality. Explicitly, on a quotient
by a finite open integer ideal the archimedean characters are
those of the dual fractional lattice. The trace divided by
\(\delta_E\) identifies the dual of any fractional ideal
\(\mathfrak a\) with \(\mathfrak a^{-1}\); taking the union over
the finite levels gives the rational characters just stated.
Compatible finite characters have their values on the diagonal
fixed by the archimedean character, by finite approximation.

Periodize \(F\) over the diagonal \(E\). Its periodization is
lower semicontinuous as the supremum of finite sums of
nonnegative translates. It descends to an integrable function
on the compact quotient, by Tonelli's theorem. In probability
Haar measure its Fourier coefficients are
\(c_E^{-1}\widehat F(y)\), by unfolding the fundamental domain.
Absolute summability gives a continuous Fourier series with
these coefficients. It agrees almost everywhere with the
periodization by uniqueness of Fourier coefficients on a compact
abelian group. This uniqueness also follows by density of finite
linear combinations of characters in the continuous functions.
At zero the lower semicontinuous periodization is at most the
value of the continuous representative: a strictly larger value
would persist on a nonempty open set of positive Haar measure.
Every summand on a finite rational subset is included there.
The integer-ring factors force \(y\) to be integral outside
\(T\). Taking absolute values of the nonzero terms yields the
first estimate. The final one follows from
\eqref{eq:actualfreq875}.
\end{proof}

Theorem~\ref{thm:commoncompletion889} verifies the required
integer-ring factors at the unselected marks when sampling a
common coordinate of the actual positive expression. For
unequal entries it retains the condition on the inverse
difference in Corollary~\ref{cor:pivotoperator875}. To apply the
quantitative sampling estimate to the full period, one must
still bound the transforms of the factors at the remaining
places and at infinity for the intervening arm coordinates.
\clearpage

\section{The first relative Gram transform at split places}\label{nav:split-gram-transforms}

\subsection{The first relative Gram frequency on the principal subgroup}
\label{sec:firstgram890}

We calculate one part of the Fourier integral left by the first
common arm coordinate. In the chosen split component let
\(e\) be the first left frame vector and \(\ell\) its primitive
paired frame covector. Thus \(\ell(e)=0\), and \(U\) fixes \(e\)
and \(\ell\). The two components of its Gram entry at a relative
element \(g\) are \(\ell g e\) and \(\ell g^{-1}e\), using the
dual action in the second component. The terminal vector \(z\)
is primitive with \(\ell(z)=0\) and \(\bar e,\bar z\) are
independent. We can choose an integral basis starting with
\(e,z\), in which \(\ell\) is the last coordinate. Write
\[
 \omega_{\alpha,\beta}(g)=
       \Psi\big(\alpha\ell g e+\beta\ell g^{-1}e\big)
       \qquad(\alpha,\beta\in F_w),
\]
with \(\Psi\) of conductor \(O\). This wave is invariant under
left and right \(U\).

\begin{lemma}[Fourier cancellation for an incident vector and covector]
\leavevmode\label{lem:flagwave890}
Let \(H=K\) or \(Z\), and let
\(t=\max(0,-\operatorname{ord}\alpha,-\operatorname{ord}\beta)\),
using \(\operatorname{ord}(0)=+\infty\). For \(t=0\) the average
of \(\omega_{\alpha,\beta}\) on \(H\) equals one. For \(t\ge1\)
its absolute value is at most
\begin{equation}
 6q^{-(r-1)t}\quad(H=K),\qquad
 8q^{\,1-(r-1)t}\quad(H=Z).
 \label{eq:flagwave890}
\end{equation}
\end{lemma}
\begin{proof}
By inversion we may suppose \(\operatorname{ord}\beta=-t\).
Put \(w=g e,\ \mu=\ell g^{-1}\). For \(g\) uniform on \(K\),
this is the uniform incident pair of a primitive vector and a
primitive covector satisfying \(\mu(w)=0\), in every finite
quotient \(O/\pi^tO\). Indeed extend \(w\) to a basis of
\(\ker\mu\), then add a vector on which \(\mu\) is one. This
gives the required integral change of basis. The basis extension
is valid over the finite quotient: a primitive covector has a
unit coefficient, and a primitive vector in its kernel extends
to a basis there by a unit coordinate. Conversely every basis
change has the prescribed incidence. Lifts of an invertible
matrix over the quotient are invertible over \(O\).

For \(Z\) the additional conditions are independence of
\(\bar w,\bar z\) and \(\bar\mu(\bar z)=0\). To verify
sufficiency in the basis construction, adjust \(z\) by a
multiple of \(\pi\) times a vector on which \(\mu\) is one, to
obtain \(z'\equiv z\bmod\pi\) in \(\ker\mu\). The independent
residues of \(w,z'\) allow them to begin a basis of the kernel.
Sending \(e,z\) to \(w,z'\) gives a change of basis in \(Z\).
Transitivity gives uniform probabilities, and the number of
allowed covectors for a fixed \(w\) is independent of \(w\).
Thus \(w\) is uniform among the primitive vectors for \(K\)
and among the vectors with \(\bar w\notin\mathbb F_q\bar z\)
for \(Z\).

Keep \(w\) fixed, lifting it integrally. Write
\(V_w=\{\mu\in\mathcal L^*:\mu(w)=0\}\), a free module of rank
\(r-1\). For \(K\) integrate \(\Psi(\beta\mu(e))\) over the
primitive part of \(V_w\). Subtract the integral on \(\pi V_w\)
from the integral on \(V_w\). Orthogonality gives the conditional
average
\[
 \frac{{\bf1}_{e\in O w+\pi^t\mathcal L}
       -q^{-(r-1)}{\bf1}_{e\in O w+\pi^{t-1}\mathcal L}}
      {1-q^{-(r-1)}}.
\]
All these tests depend only on reductions at precision \(t\).
For \(h\ge1\) the probability that \(e\in O w+\pi^h\mathcal L\)
is at most \(2q^{-(r-1)h}\): \(w\bmod\pi^h\) must be a unit
multiple of \(e\), out of \(q^{rh}(1-q^{-r})\) primitive vectors.
The same bound with \(h=0\) is immediate. The factor
\(\Psi(\alpha\ell w)\) has modulus one. Applying these bounds
in the preceding average proves the first estimate at the
ranks in question.

For \(Z\) start with \(V_w'=\{\mu\in V_w:\mu(z)\in\pi O\}\).
The excluded covectors are still \(\pi V_w\), now of relative
measure \(q^{-(r-2)}\). In a basis beginning with \(w,z\),
orthogonality changes the conditional average to
\[
 \frac{{\bf1}_{e\in O w+\pi^{t-1}O z+\pi^t\mathcal L}
       -q^{-(r-2)}{\bf1}_{e\in O w+\pi^{t-1}\mathcal L}}
      {1-q^{-(r-2)}} .
\]
In the first event \(e\equiv a w+\pi^{t-1}b z\bmod\pi^t\)
requires \(a\) to be a unit, even for \(t=1\), since
\(\bar e,\bar z\) are independent. There are at most \(q^{t+1}\)
possible \(w\bmod\pi^t\), out of
\(q^{rt}(1-q^{1-r})\) eligible vectors. The probability is at
most \(2q^{1-(r-1)t}\). The second event has probability at most
\(2q^{-(r-1)(t-1)}\), by counting possible unit multiples as
above when \(t>1\), and directly for \(t=1\).
Substitution proves the second estimate. At integral frequencies
the two arguments of \(\Psi\) are integral throughout \(H\).
\end{proof}

For \(H=K,Z\) consider the position pairing with the wave in
its relative variable. At \(x\in J\) its finite expression is
\[
 \sum_{\gamma,\delta\in\bar H}
 \int_{(U\backslash G)^2}
   \overline{\Phi_\gamma(y)}\Phi_\gamma(yx)
   \Phi_\delta(y')\overline{\Phi_\delta(y'x)}
   \,\omega_{\alpha,\beta}(y(y')^{-1})\,d\dot y\,d\dot y'.
\]
This uses exactly the two quotient positions in
\eqref{eq:rightposition881}. Inserting the Gram character on
their relative group element is well defined on the two
quotients by its \(U\)-invariance. Denote its integral over
\(x\in J\), divided by
\(\operatorname{vol}(J)\operatorname{vol}(H)\operatorname{vol}(K)\),
by \(B_H(\alpha,\beta)\).

\begin{theorem}[The principal positions in the first relative transform]
\leavevmode\label{thm:firstgram890}
One has
\begin{equation}
 B_H(\alpha,\beta)=\frac{\operatorname{vol}(H)}{\operatorname{vol}(K)}
              \int_H\omega_{\alpha,\beta}(g)\,
                     \frac{dg}{\operatorname{vol}(H)} .
 \label{eq:firstgram890}
\end{equation}
In particular \(B_H(0,0)=\operatorname{vol}(H)/\operatorname{vol}(K)\).
For a rational nonzero frequency with split components
\(\alpha,\beta\), set
\[
 I=q^{\operatorname{ord}\alpha+\operatorname{ord}\beta},\qquad
 D=q^{\max(0,-\operatorname{ord}\alpha)+
        \max(0,-\operatorname{ord}\beta)},\qquad K_0=k-n^7 .
\]
At sufficiently large nodes, the corresponding bound on the
same fields is
\[
             |B_H(\alpha,\beta)|\le B_H(0,0)\,
                          \min\{1,I^{K_0/2}D^{-3}\}.
\]
\end{theorem}
\begin{proof}
All \(\Phi_\gamma\) are right \(J\)-invariant. At \(x\in J\)
each pair of factors in one quotient variable is its squared
absolute value. Its support projects to \(J k_\gamma\), with
the quotient integral obtained from ordinary Haar on
\(J k_\gamma\) since \(U\cap J\) has volume one. This is the
quotient calculation in Proposition~\ref{prop:twocorrelations881}.
The wave is invariant under the excluded \(U\)-variables, as
the group fixes both \(e\) and \(\ell\). Summing the cosets gives
the integral of \(\omega_{\alpha,\beta}(h_1h_2^{-1})\) over
\(H^2\). One of these compact integrations gives
\(\operatorname{vol}(H)\), and the \(x\)-integration contributes
\(\operatorname{vol}(J)\), proving \eqref{eq:firstgram890}.

For \(t=0\) the weight is at least one before taking the minimum.
For \(t\ge1\), \(I\ge D^{-1}\), \(D\le q^{2t}\), so that
\(I^{K_0/2}D^{-3}\ge q^{-(K_0+6)t}\).
Since \(r-1=k+2a\), the bounds in \eqref{eq:flagwave890} are
at most \(q^{-(K_0+6)t}\) for all \(t\ge1\) at the stated
nodes. The estimate by one holds for every normalized average.
This proves the bound with the local weight occurring in
\eqref{eq:freqweight875}.
\end{proof}

The frequency in \eqref{eq:firstgram890} is applied to the
relative element \(y(y')^{-1}\). Restricting to the principal
positions made its distribution a compact group integral.
For \(x\notin J\) the two factors in each quotient variable
are generally different. Their unweighted sum is controlled
at the marks by Theorem~\ref{thm:allfixermass888} and
Corollary~\ref{cor:allmarks888}; the additional wave requires
an estimate for that weighted pairing.
\clearpage

\subsection{The tail of the first relative Gram transform}
\label{sec:gramtail891}

We now allow the position \(x\) in Theorem~\ref{thm:firstgram890}
to run over \(G\). Write \(v_D=\operatorname{vol}(D)\) for a compact
open subgroup \(D\). In the notation of that theorem the integrand,
including its normalization, is
\begin{equation}
 b_H(x;\alpha,\beta)=\frac{1}{v_Jv_Hv_K}
   \int_{H^2}\phi(h_1xh_1^{-1})\,
       \overline{\phi(h_2xh_2^{-1})}\,
       \omega_{\alpha,\beta}(h_1h_2^{-1})\,dh_1\,dh_2 .
 \label{eq:gramintegrand891}
\end{equation}
Indeed, for each of the two quotient positions in
\eqref{eq:rightposition881} the first factor restricts the
variable to \(U J k_\gamma\). Choose its representative \(h=j k_\gamma\).
As in Proposition~\ref{prop:twocorrelations881}, the quotient
measure on this part is the pushforward of \(dj\), and the factor
at \(hx\) is \(\phi(hxh^{-1})\). The wave is unchanged by any
of the choices in the two \(U\)-orbits. Summing over both
sets of residue cosets gives \eqref{eq:gramintegrand891}.
The compact integrals in this formula are taken before the
integration in \(x\).

For an integer \(t\ge1\) define
\[
 H_t=\{h\in H:he\equiv e,\ \ell h\equiv\ell\bmod\pi^t\},
 \qquad H_0=H,\qquad N_t=[H:H_t].
\]
Here \(e\) is the first left flag vector, and \(\ell\) is the
covector reading the last coordinate of the reversed right flag.
The left cosets \(H_t h\) are parametrized by the distinct pairs
\((h^{-1}e,\ell h)\bmod\pi^t\). Equality of two such pairs means
that the quotient \(h'h^{-1}\) fixes both initial coordinates to
the stated precision. In particular \(N_t\le q^{2rt}\).
For \(t=\max(0,-\operatorname{ord}\alpha,-\operatorname{ord}\beta)\)
the wave in \eqref{eq:gramintegrand891} is constant on each pair
of these cosets. In fact the two entries of the phase are obtained
by evaluating the covector of each pair on the vector of the other.

\begin{lemma}[Signed sums with both end coordinates fixed]
\leavevmode\label{lem:flagclass891}
For \(H=K,Z\), any \(t\ge0\), \(m\ge200n^3\), and
\(L=\lfloor m/(200n^3)\rfloor\), the \(H_t\)-orbits on the
right \(J\)-cosets in \(\mathcal S(m,P,P)\) satisfy
\[
 \sum_{\mathcal C}
    \left|\int_{U\cap\mathcal C}\chi(u)\,du\right|
       \le 10q^{4a^2(m+1)+dP-kL}.
\]
In the integral \(\mathcal C\) denotes the inverse image in \(G\).
\end{lemma}
\begin{proof}
Every element of \(H_t\) normalizes \(J\). Thus each indicated
orbit defines a right \(J\)-invariant set, stable under conjugation
by \(H_t\). It is consequently stable under left multiplication
by \(J\cap H_t\) as well. Its signed sum on \(G/J\) is the
displayed integral, since a coset meeting \(U\) cuts it in a
coset of \(U\cap J\), of measure one and constant character.

We check the transformations used in Proposition~\ref{prop:largeshell888}
for these smaller sets. All allowed middle basis changes in its
proof fix both ends of the flag exactly and belong to \(H_t\).
The row shear in Proposition~\ref{prop:armshear887} uses left
and right principal multipliers \(l_w,r_w\). They fix \(e\)
because \(\lambda(e)=0\) and \(\xi(e)=\lambda T e=0\), since
\(T e=e\). Their last row is unchanged from the identity, hence
they fix \(\ell\). Both therefore belong to \(J\cap H_t\) for
every \(t\), including for \(Z\) since they act trivially modulo
\(\pi\). These facts preserve each orbit set under the shear.

On the part with dominant opposite arm, transposition followed
by reversal exchanges the initial vector and the last covector.
The images of the orbit sets under this map remain disjoint.
In the reversed coordinates the two shear multipliers again fix
both ends exactly and are principal. Pulling them back to the
original coordinates gives elements of \(J\cap H_t\), with
the two sides exchanged. The transformed sets are consequently
stable under this operation as well. Middle averaging on this
part can be performed in the original coordinates, as in the
proof of Proposition~\ref{prop:largeshell888}.

For clarity, put \(V_0=q^{4a^2(m+1)+dP}\).
In each orientation the strata failing the height or row tests
there have total measure at most \(q^{-dL}V_0\), by
Lemma~\ref{lem:rowlength887}. On the strata with a deep terminal
middle coordinate the average of Lemma~\ref{lem:middleaverage883}
bounds the sum of the absolute contributions by \(2q^{-kL}V_0\).
On the remaining strata the row shear cancels within each set
unless the exceptional covector congruence holds. Its measure
on the whole strata is at most \(2q^{-kL}V_0\), by the fraction
bound of Proposition~\ref{prop:armshear887}. These measure bounds
apply to the sum because the sets are disjoint. Adding both
orientations and using \(d=k+1\) proves the assertion.
\end{proof}

\begin{theorem}[A conductor-dependent tail estimate]
\leavevmode\label{thm:gramtail891}
At each original split mark, for \(H=K,Z\) and all sufficiently
large nodes, \(\int_G |b_H(x;\alpha,\beta)|\,dx\) is finite.
For every pair of local frequencies, with \(t\) as above,
\begin{equation}
 \int_{\{m\ge1200n^3(t+1)\}}\!|b_H(x;\alpha,\beta)|\,dx
                         \le q^{-6k(t+1)} .
 \label{eq:gramtail891}
\end{equation}
\end{theorem}
\begin{proof}
Choose representatives \(k_i\) of \(H_t\backslash H\). Write
\(C_D(y)=\int_D\phi(hyh^{-1})\,dh\) and put
\(A_i(x)=C_{H_t}(k_i x k_i^{-1})/v_{H_t}\).
Let \(\omega_{ij}\) denote the common wave on the two cosets.
Then \eqref{eq:gramintegrand891} is
\[
 \frac{v_H}{v_Jv_K}\frac1{N_t^2}
       \sum_{i,j}\omega_{ij}A_i(x)\overline{A_j(x)} .
\]
The normalized sum is bounded in absolute value by
\((N_t^{-1}\sum_i |A_i|)^2\le N_t^{-1}\sum_i |A_i|^2\).
Conjugation by \(k_i\) preserves the shell and its measure.
On a shell \(\mathcal S\), integration gives at most
\[
 \frac{v_H}{v_K}\frac1{v_Jv_{H_t}^2}
                         \int_{\mathcal S}|C_{H_t}(x)|^2\,dx .
\]
The second factor is exactly
\(\sum_{\mathcal C}|S_{\mathcal C}|^2/|\mathcal C|\), where
\(S_{\mathcal C}=\int_{U\cap\mathcal C}\chi(u)\,du\).
Indeed on a right coset \(C_{H_t}/v_{H_t}\) is the average of
\(\phi\) over its orbit, and each coset has measure \(v_J\).
Lemma~\ref{lem:shellorbits884} gives
\([K:H_t]|\mathcal C|\ge|\mathcal C_K|\ge
 q^{2(d-2a)P}/4\) for any contributing orbit, with \(\mathcal C_K\)
its full \(K\)-orbit. The first inequality follows by comparing
the stabilizer in \(H_t\) with that in \(K\). Since
\((v_H/v_K)[K:H_t]=N_t\), the preceding lemma and the square
of its bound give
\[
 \int_{\mathcal S}|b_H(x;\alpha,\beta)|\,dx
 \le 400N_t q^{8a^2(m+1)+4aP-2kL}
 \le 400q^{2rt+16a^2(m+1)-2kL}.
\]
The only contributing shells have \(P=Q\) by
Lemma~\ref{lem:detlength885}. We used \(P\le 2a(m+1)\).

On the tail in the theorem \(L\ge6(t+1)\), and
\(m+1\le400n^3L\). On this interval \(6400a^2n^3\le k/16\)
and \(2r\le9k/4\) for all sufficiently large nodes. Thus
\(2rt\le3kL/8\) and the last bound is at most \(400q^{-25kL/16}\).
For fixed \(L\) there are at most \(200n^3(800an^3L+1)\) pairs
\((m,P)\). Extract \(q^{-6k(t+1)}\) using \(L\ge6(t+1)\).
The remaining sum is at most
\[
 400\cdot200n^3\sum_{L\ge1}(800an^3L+1)q^{-9kL/16},
\]
which is at most one for all large nodes, uniformly in the
frequency and the mark. This follows by summing the geometric
series and its first moment, since \(q\ge2\) and \(a,n\)
grow polynomially in the interval parameter while \(k\)
grows at least as fast as \(x^5-2\sqrt{x}\).
Every bounded range of \(m\) has compact closure in \(G\),
by the simultaneous bounds on \(x,x^{-1}\), and the compact
averages in \eqref{eq:gramintegrand891} are bounded there.
This proves the absolute convergence assertion.
\end{proof}

For nonzero rational frequencies this tail has the decay required
in \eqref{eq:freqweight875}, since the weight there is at least
\(q^{-(k-n^7+6)t}\). The integral over \(J\) was calculated in
Theorem~\ref{thm:firstgram890}. A frequency-dependent compact
range of nonprincipal positions still lies between these two
estimates.
\clearpage

\subsection{Frequency decay on all integral positions}
\label{sec:integralfreq892}

Keep the relative wave and the groups of
Section~\ref{sec:gramtail891}. Put
\(c_\lambda=\max(0,-\operatorname{ord}\lambda)\), with \(c_0=0\),
and \(t=\max(c_\alpha,c_\beta)\). We will estimate the integral
of \(|b_H|\) over the whole of \(K\), including positions outside \(J\).

\begin{lemma}[The operator on pairs of end coordinates]
\leavevmode\label{lem:incidentoperator892}
For \(t\ge1\), write the distinct pairs indexing \(H_t\backslash H\)
as \((w_i,\mu_i)\in (O/\pi^t)^r\times ((O/\pi^t)^r)^*\).
On this finite set with probability measure the operator with
kernel
\[
             \Psi\big(\alpha\mu_i(w_j)+\beta\mu_j(w_i)\big)
\]
has norm at most \(\min(1,4q^{t+1-r(c_\alpha+c_\beta)/2})\)
on complex square-integrable functions.
\end{lemma}
\begin{proof}
Use counting probability on the full product of the vector and
covector spaces over \(O/\pi^t\). The ambient operator is a
product of two finite Fourier operators, interchanging the order
of the vector and covector coordinates. For one such operator
with coefficient \(\lambda\), its adjoint times itself averages
against the kernel \({\bf1}_{u-u'\in\pi^{c_\lambda}(O/\pi^t)^r}\),
by orthogonality of additive characters. This operator is
\(q^{-rc_\lambda}\) times the orthogonal projection onto functions
constant on these cosets. This includes \(c_\lambda=0\).
Consequently the norm of the product is
\(q^{-r(c_\alpha+c_\beta)/2}\).

By the basis calculation in Lemma~\ref{lem:flagwave890},
the set of pairs has uniform measure after conditioning on
membership. For \(K\) its ambient probability equals
\((1-q^{-r})q^{-t}(1-q^{-(r-1)})\). Indeed \(w\) is primitive,
and on the hyperplane \(\mu(w)=0\) the covector is primitive.
For \(Z\) the probability equals
\((1-q^{1-r})q^{-t-1}(1-q^{-(r-2)})\). Here \(\bar w\) is
independent of \(\bar z\), and inside \(\mu(w)=0\) the additional
residue equation \(\bar\mu(\bar z)=0\) imposes one independent
condition. Both probabilities are at least \(\tfrac14 q^{-t-1}\)
at these ranks. Extending a function by zero, applying the ambient
operator, and restricting back gives the desired operator after
division by the probability of the set. The norm bound follows
with this factor. The bound by one also follows directly from
Cauchy--Schwarz and the modulus of the kernel.
\end{proof}

\begin{lemma}[Mean square after fixing the two ends in the residue field]
\leavevmode\label{lem:residueends892}
For \(H=K,Z\) and \(t\ge1\), at the original split marks,
\[
 \frac{v_H}{v_K}\frac1{v_Jv_{H_t}^2}
             \int_K |C_{H_t}(x)|^2\,dx \ \le\ q^{100n^2}.
\]
\end{lemma}
\begin{proof}
For integral \(x\) the value of \(\phi(hxh^{-1})\) depends only
on the reductions of \(h,x\), by Lemma~\ref{lem:finiteposition881}.
The residue image \(D\) of \(H_t\) equals that of \(H_1\). To lift
a residue matrix in \(H_1\), prescribe its first column to be \(e\)
and its last row to be \(\ell\) exactly, and lift all remaining
entries arbitrarily. These prescriptions agree at their intersection
and the determinant is a unit. In the case \(Z\) the terminal
condition is already imposed modulo \(\pi\). Thus the normalized
averages for all \(t\ge1\) coincide on \(K\).

Write \(\bar H\) for the residue image of \(H\). The required
quantity is
\[
 \frac{|\bar H|}{|\bar G|}
       \sum_{\mathcal C}\frac{|S_{\mathcal C}|^2}{|\mathcal C|},
 \qquad
 S_{\mathcal C}=\sum_{X\in\mathfrak n\cap\mathcal C}\psi(\Lambda(X)),
\]
where the orbits are now for \(D\). This follows by averaging on
each orbit as in \eqref{eq:finiteposition881}. For \(X\) of
rank \(j>0\) let \(E_X\) be its centralizer in \(\bar G\) and
\(\mathcal C_K\) its full conjugacy orbit. Let \(T_X\) be the
size of the \(E_X\)-orbit of the list of fixed data, namely
\((e,\ell)\) for \(K\) and \((e,\ell,z)\) for \(Z\), with the
dual action on \(\ell\). Orbit--stabilizer gives
\[
 \frac{|\bar H|}{|\bar G|\cdot|\mathcal C|}
   =\frac{[\bar H:D]}{|\mathcal C_K| T_X}
   \le\frac{4q^{2r}}{q^{2rj-2j^2} T_X}.
\]
Here \([\bar H:D]\le q^{2r}\) counts the pairs of end coordinates.
The lower bound on \(|\mathcal C_K|\) follows from the commuting
algebra calculation in Theorem~\ref{thm:integralmass881}.
Always \(j\le2a\). Put \(M=n^2q^{4a^2}\).

For \(K\), Proposition~\ref{prop:middlecancellation882} gives
\(\sum_{\mathcal C}|S_{\mathcal C}|\le10qM q^{d(j-1)}\)
also on these refined orbits. To see this, the middle changes
of basis fix both ends. Each shear there has the form \(1+w\xi\),
with \(\xi\) the covector of a left coordinate strictly after
the first one and \(w\) in the middle space. It fixes \(e,\ell\).
The transposed argument exchanges these two conditions. All
averages therefore preserve each refined orbit, and its proof
uses counts on the whole stratum for the remaining exceptions.

We give the corresponding argument for \(Z\), writing \(D_X\)
for the right-to-middle block of \(X\), to distinguish it from
the group \(D\). On \(D_X\ne0\) transpose and reverse the flags.
In the new coordinates the left middle block \(B'\) is nonzero,
and the new terminal covector \(\eta'\) is the old fixed \(z\)
under duality, up to a nonzero scalar. Middle changes preserving
\(\eta'\) belong to the transformed group \(D\). They send the
new terminal vector \(z'\) uniformly over the affine hyperplane
\(\eta'(w)=\eta'(z')\ne0\). This follows by extending a basis of
\(\ker\eta'\) by \(w\). The term involving \(\eta'\) in the
phase is constant under this average. If the last row of \(B'\)
is nonzero, orthogonality on this affine space cancels its
contribution unless that row is proportional to \(\eta'\).
In the latter case the row space of \(B'\) contains \(\eta'\).

If the last nonzero row has index \(i<a\), use the shear in
Proposition~\ref{prop:middlecancellation882}, with \(w\in\ker\eta'\).
It fixes the two ends and the covector \(\eta'\), extended by
zero on the arms. Thus the cancellation in each class again applies.
The only exceptions in that proof have either the sum of the
two middle ranks strictly below \(j\), or the row space of \(B'\)
containing \(\eta'\). Each has size at most \(M q^{d(j-1)}\).
Indeed the first follows by counting the two rank factorizations
as in that proposition, and in the second one row-space basis
vector is prescribed. On \(D_X=0\) with \(\operatorname{rank}B<j\)
the first smaller count applies directly. Consequently the sum
of the absolute class contributions from all the parts just
considered for \(Z\) is at most \(12qM q^{d(j-1)}\).
Here the nonvanishing-block condition survives the conjugations
because \(B'\) remains nonzero under the middle changes and is
fixed by the interior-row shear.

It remains to count \(D_X=0,\ \operatorname{rank}B=j>0\) for
\(Z\). Now \(XV=B(W)\subset L\). Choose a covector \(\lambda\)
supported on \(W\) with \(\lambda(z)=1\). For each
\(v\in W\cap\ker B\cap\ker\lambda\), the transformation
\(1+v\lambda\) commutes with \(X\) and sends \(z\) to \(z+v\).
It is invertible since \(\lambda(v)=0\). This gives at least
\(q^{d-a-1}\) images of \(z\). There is a vector \(u\in\ker X\)
with \(\ell(u)=1\). Indeed subtract from the last right basis
vector a middle vector having the same image under \(X\), which
exists since \(XV=B(W)\). For any covector \(\nu\) vanishing on
\(XV+\langle e,z,u\rangle\), the map \(1+u\nu\) is invertible,
commutes with \(X\), fixes \(e,z\), and takes \(\ell\) under
the dual action to \(\ell-\nu\). There are at least
\(q^{r-j-3}\) choices. Taking the orbit of \(z\) first and
then the orbit under its stabilizer gives
\[
                 T_X\ge q^{d-a-1+r-j-3}\ge q^{2r-5a-4}.
\]
The total count on this part is at most \(M q^{dj}\) by the
full rank count in Proposition~\ref{prop:middlecancellation882}.
The orbit lower bound applies to every class meeting the part.

To sum, use \((u+v)^2\le2(u^2+v^2)\) for the two absolute
contributions within a class. For the signed smaller sums the
bound per rank is at most \(4q^{4a}(12qM)^2\) before this factor
of two, since \(r=d+2a\) and \(2j^2-4aj\le0\).
For the last part the count and its bound on \(T_X\) give at
most \(4q^{5a+4}M^2\) per rank. The zero orbit contributes
\(|\bar H|/|\bar G|\le1\). Summing at most \(2a\) positive ranks
proves the displayed inequality, using \(n=a+1\) and \(q\ge2\).
\end{proof}

\begin{theorem}[Every integral position at nonintegral frequency]
\leavevmode\label{thm:integralfreq892}
For \(H=K,Z\), uniformly at the original split marks of the common
interval and at all sufficiently large nodes, one has, for \(t\ge1\),
\[
 \int_K |b_H(x;\alpha,\beta)|\,dx
        \le 4q^{100n^2+t+1-r(c_\alpha+c_\beta)/2}.
\]
In particular for a nonzero rational frequency with local
indices \(I,D\) as in Theorem~\ref{thm:firstgram890},
\[
 \int_K |b_H(x;\alpha,\beta)|\,dx
                     \le\min\{1,I^{(k-n^7)/2}D^{-3}\}.
\]
The latter statement is asserted for \(t\ge1\). At \(t=0\) the
integral is exactly \(A_H v_H/v_K\).
\end{theorem}
\begin{proof}
In the finite sum for \eqref{eq:gramintegrand891} apply
Lemma~\ref{lem:incidentoperator892} to the two coefficient vectors
with entries \(A_i(x)\) of Theorem~\ref{thm:gramtail891}.
The resulting bound is the operator norm times
\((v_H/(v_Jv_K))N_t^{-1}\sum_i |A_i(x)|^2\).
Conjugating \(x\) by each \(k_i\) in its integral over \(K\)
gives the mean square in Lemma~\ref{lem:residueends892}.
This proves the first estimate.

Put \(s=c_\alpha+c_\beta\), so \(1\le t\le s\), \(D=q^s\),
and \(I\ge q^{-s}\). On this interval \(n\) tends to infinity,
and \(r=k+1+2a\). Therefore
\(4q^{100n^2+t+1}\le q^{(r-k+n^7-6)s/2}\) eventually,
uniformly in \(s\). The first bound is then at most
\(q^{-(k-n^7+6)s/2}\), which is at most both one and
\(I^{(k-n^7)/2}D^{-3}\) since \(k>n^7\). At \(t=0\) the
wave is identically one, and the last formula follows from
the definition of \(A_H\).
\end{proof}

The estimate here gives the frequency weight over all integral
positions. At selected marks it does not provide a separate
terminal mass factor multiplying this weight. For positions
outside \(K\), Theorem~\ref{thm:gramtail891} controls the
distant shells. The intervening nonintegral shells, the later
Gram coordinates, and the comparison with the arithmetic sum
still require estimates.
\clearpage

\subsection{Unequal conductors and small nonintegral heights}
\label{sec:shallowwave893}

Fix a position of height \(m\) in the sense of
Section~\ref{sec:nonspheretail884}, and put \(s=2m+1\).
For \(j\ge1\) let \(J(j)=1+\pi^j M_r(O)\). For every \(h\in H\)
the function \(f_x(h)=\phi(hxh^{-1})\) depends only on
\(h\bmod\pi^s\). Indeed, put \(y=hxh^{-1}\). For \(j_0\in J(s)\)
one has \(y^{-1}(j_0-1)y\in\pi M_r(O)\) since
\(\|y\|,\|y^{-1}\|\le q^m\). Hence \(j_0 y j_0^{-1}\in yJ\)
and the right \(J\)-invariance of \(\phi\) gives the assertion.

\begin{proposition}[The derivative of the wave on a congruence subgroup]
\leavevmode\label{prop:congruencewave893}
Suppose \(t=\max(c_\alpha,c_\beta)>s\), set
\[
 j_0=\max\{s,\lceil t/2\rceil\},\qquad h_0=t-j_0\ge1,
 \qquad A=\pi^t\alpha,\quad B=\pi^t\beta .
\]
In the compact double integral \eqref{eq:gramintegrand891} one
may insert the indicator of
\begin{equation}
      A(ge)\ell\equiv B e(\ell g^{-1})\pmod{\pi^{h_0}},
                \qquad g=h_1h_2^{-1}.
 \label{eq:stationarywave893}
\end{equation}
If \(c_\alpha\ne c_\beta\), then \(b_H(x;\alpha,\beta)=0\).
If \(c_\alpha=c_\beta\), the condition requires
\[
 (ge,\ell g^{-1})\equiv
          (\lambda e,(A/B)\lambda\ell)\pmod{\pi^{h_0}}
\]
for a unit \(\lambda\). For a uniform \(g\in H\) the probability
of the condition is at most \(4q^{1-(2r-2)h_0}\).
\end{proposition}
\begin{proof}
Replace \(h_1\) by \((1+\pi^{j_0}Y)h_1\), and average in the
integral matrix \(Y\) with additive probability. Each left
multiplier belongs to \(H\) and leaves \(f_x(h_1)\) fixed by the
height bound. The inverse multiplier is congruent to
\(1-\pi^{j_0}Y\) modulo \(\pi^t\), since \(2j_0\ge t\).
For fixed \(g\) the ratio of the new wave to the old wave is
therefore
\[
        \Psi\!\left(\pi^{j_0-t}
                \operatorname{tr}\big((A(ge)\ell-B e(\ell g^{-1}))Y\big)
              \right).
\]
Entrywise orthogonality gives precisely
\eqref{eq:stationarywave893}. If the conductors differ, exactly
one of \(A,B\) is a unit. Both unscaled rank-one matrices are
nonzero modulo \(\pi\), so the condition is impossible.

For equal conductors read the matrix congruence in a basis
starting with \(e\) in which \(\ell\) is the last coordinate
covector. In its last column all entries of \(ge\) below the
first one must vanish, and in its first row every entry of
\(\ell g^{-1}\) before the last must vanish. The corner entry
gives the stated relation of the remaining coefficients, which
are units by primitivity. Conversely these conditions suffice.
There are at most \(q^{h_0}\) admissible pairs. The number of
pairs occurring uniformly at this precision is at least
\(\tfrac14 q^{(2r-1)h_0-1}\) by the counts in
Lemma~\ref{lem:incidentoperator892}. Division proves the bound.
\end{proof}

\begin{theorem}[A terminal saving on small positive heights]
\leavevmode\label{thm:shallowwave893}
Let \(H=K,Z\) at one of the marked split places and let a nonzero
rational frequency have \(u=c_\alpha+c_\beta\ge1\).
Set \(K_0=k-n^7\) and \(W=\min\{1,I^{K_0/2}D^{-3}\}\),
with the local indices in Theorem~\ref{thm:firstgram890}.
At all sufficiently large nodes,
\begin{equation}
  \int_{\{m\ge1:\ 64r(m+1)\le n^7u\}}
                   |b_H(x;\alpha,\beta)|\,dx \ \le\ q^{-r} W .
 \label{eq:shallowwave893}
\end{equation}
If \(c_\alpha\ne c_\beta\), all the heights with
\(2m+1<t\) have identically zero integrand, independently of
the inequality in \eqref{eq:shallowwave893}.
\end{theorem}
\begin{proof}
Put \(l=\min(t,s)\). The reductions of \(H_t\) and \(H_l\)
modulo \(\pi^s\) coincide. This is immediate when \(t\le s\).
If \(t>s\), a matrix of \(H_s\) modulo \(\pi^s\) has first
column \(e\) and last row \(\ell\). Lift these entries exactly
and lift the others arbitrarily. The resulting unit matrix
belongs to \(H_t\), retaining the terminal condition modulo
\(\pi\) when present. Thus the normalized conjugacy average on
the shell is the same for \(H_t\) and \(H_l\), since \(f_x\)
factors through this reduction.

Use the operator of Lemma~\ref{lem:incidentoperator892} in the
coset expression \eqref{eq:gramintegrand891}. After integrating
the sum of squares of its coefficients on a shell \(\mathcal S\),
the preceding observation gives the upper bound
\[
  4q^{t+1-ru/2}\,
   \frac{v_H}{v_K}\frac{1}{v_J v_{H_l}^2}
                       \int_{\mathcal S}|C_{H_l}(x)|^2\,dx .
\]
To bound the last two factors use the orbit sum in the proof of
Theorem~\ref{thm:gramtail891}. All contributing shells have
\(P=Q\), and on each such shell the sum of the absolute signed
intersection counts is at most the unsigned count
\(q^{4a^2(m+1)+dP}\) of Lemma~\ref{lem:shellorbits884}.
Every contributing \(H_l\)-orbit has size at least
\(q^{2(d-2a)P}/(4[K:H_l])\). Since \([H:H_l]\le q^{2rl}\)
and \(P\le2a(m+1)\), this proves the explicit bound
\[
        \int_{\mathcal S}|b_H(x;\alpha,\beta)|\,dx
        \le 16q^{\,t+1-ru/2+2rl+16a^2(m+1)}.
\]
We now have \(2rl\le4r(m+1)\) and \(16a^2\le r\) on the given
interval at large nodes. Under the height condition in the
theorem, \(8r(m+1)\le n^7u/8\). Also
\(t+1+\log_q16\le6u\), and \(n\) tends to infinity. In particular
\[
 t+1+\log_q16+5r(m+1)+2r+2(m+1)
          \le (r-k+n^7-6)u/2
\]
for all sufficiently large nodes, using \(r\ge k\), \(r\ge2\).
It follows that the last shell bound is at most
\(q^{-2r-2(m+1)}q^{-(K_0+6)u/2}\).
The last factor is at most \(W\) by \(D=q^u\) and \(I\ge D^{-1}\).
At each \(m\) there are at most \(2a(m+1)+1\) possible values
of \(P\). The sum of \((2a(m+1)+1)q^{-2(m+1)}\) over all
\(m\ge1\) is \(O(a+1)\), which is at most \(q^r\) eventually.
This gives \eqref{eq:shallowwave893}. The exact vanishing is
Proposition~\ref{prop:congruencewave893}.
\end{proof}

Here \(q^{-r}\le v_Z/v_K=(q^r-1)^{-1}\), so the last positive-height
estimate retains a terminal mass factor together with the frequency
weight. For equal conductors the rank-one congruence controls the
relative position. Its unweighted probability by itself gives
no estimate for a coefficient concentrated on the same congruence.
The region between the small-height estimate and
Theorem~\ref{thm:gramtail891} remains in the first-coordinate
comparison.
\clearpage

\subsection{The two end ideals in the first relative transform}
\label{sec:endtransform894}

There is a part of the first relative transform for which the
frequency estimate, including the terminal mass factor, holds
at every matrix height. It is determined by two ideals of \(O\).
Let \(\mathcal L=O^r\), with the original ends \(e,\ell\). For
\(y\in G\), define
\begin{align*}
 I_e(y)&=\{\lambda(e):\lambda\in\mathcal L^*,\
                             \lambda(y-1)\mathcal L\subset\pi O\},\\
 I_\ell(y)&=\{\ell(v):v\in\mathcal L,\
                             (y^{-1}-1)v\in\pi\mathcal L\}.
\end{align*}
Each is the image of an \(O\)-module under a linear form. At height
\(m\) each contains \(\pi^{m+1}O\), by taking all rows or columns
divisible by \(\pi^{m+1}\) and using primitivity. Write
\(\delta_e,\delta_\ell\) for their valuations truncated above at \(t\).
The ideals depend only on \(yJ\). For \(j\in J\) one has
\(yj-1=(y-1)j+(j-1)\), and
\((yj)^{-1}-1=j^{-1}(y^{-1}-1)+(j^{-1}-1)\). These identities
preserve the row and column conditions. The truncated valuations
are constant on \(H_t\)-conjugacy orbits on \(G/J\). Indeed the
row and column modules are transported under conjugation, and
the end coordinates on which they are evaluated are preserved
modulo \(\pi^t\) by \(H_t\).

\begin{lemma}[Movement of both ends by a coset stabilizer]
\leavevmode\label{lem:endorbit894}
Let \(y\in U\), let \(\mathcal D_y\) be the stabilizer of \(yJ\)
in \(K\), and put \(d_*=d-2a\). The size of its orbit on the
pair \((e,\ell)\bmod\pi^t\) is at least
\[
                  q^{d_*(2t-\delta_e-\delta_\ell)} .
\]
If \(\bar z\notin\overline{T_D}\), with \(T_D\) the saturated
middle sublattice of Lemma~\ref{lem:terminalorbit886}, its orbit
on the pair together with \(\bar z\) has size at least
\[
           (q^{d_*}-1)q^{d_*(2t-\delta_e-\delta_\ell)} .
\]
\end{lemma}
\begin{proof}
Use the decomposition \(y=Y\oplus1\) on an integral complement
\(V_1\) and a middle summand \(W_0\) of rank \(d_0\ge d_*\) from
Lemma~\ref{lem:shellorbits884}. Thus \(e\in V_1\) and \(\ell\)
vanishes on \(W_0\). The element with diagonal blocks one and
a single lower block \(A:V_1\to W_0\) stabilizes \(yJ\) if
\(A(Y-1)\) is divisible by \(\pi\). This follows on multiplying
out \(y^{-1}g y-g\), which must lie in \(\pi M_r(O)\).
Each row of \(A\) can be chosen independently, and its values
on \(e\) fill \(I_e(y)\). Indeed adjoining the \(W_0\)-entries
of a row does not change the test or its value on \(e\).
Consequently these transformations move \(e\) through at least
\(q^{d_0(t-\delta_e)}\) residues.

Likewise a single upper block \(C:W_0\to V_1\) may be chosen with
\((Y^{-1}-1)C\) divisible by \(\pi\). It fixes \(e\) and moves
\(\ell\), in the dual action, by \(-\ell C\) on \(W_0\).
The entries of this row fill \(I_\ell(y)\) independently, giving
\(q^{d_0(t-\delta_\ell)}\) possible residues. Taking first the
vector orbit and then the covector orbit under its stabilizer
proves the first bound. Under the extra hypothesis we may choose
\(W_0\) and its complement as in Lemma~\ref{lem:terminalorbit886},
so that \(\bar z\) has nonzero \(W_0\)-component. The full group
\(\mathrm{GL}(W_0)(O)\) acts trivially on \(V_1\), commutes with
\(y\), and fixes both ends exactly. It moves that component
through \(q^{d_0}-1\) values. Taking the remaining orbit in the
stabilizer of the pair proves the second bound.
\end{proof}

For \(t\ge1\) and \(0\le M\le2t\) put
\[
 \phi_M(y)=\phi(y){\bf1}_{\delta_e(y)+\delta_\ell(y)\le M}.
\]
Let \(b_{H,M}\) denote \eqref{eq:gramintegrand891} with
\(\phi_M\) in both factors. This is a term in the original
integral obtained by writing each coefficient as
\(\phi_M+(\phi-\phi_M)\).

\begin{theorem}[Frequency and terminal mass at all heights with an end-ideal bound]
\leavevmode\label{thm:endtransform894}
For a nonzero rational frequency put \(u=c_\alpha+c_\beta\ge1\),
\(t=\max(c_\alpha,c_\beta)\), and
\[
 M=\min\{2t,\lfloor n^7u/(16d)\rfloor\},\qquad
 W=\min\{1,I^{(k-n^7)/2}D^{-3}\}.
\]
At every original split mark of the shorter interval, uniformly
for all sufficiently large nodes,
\[
       \int_G |b_{K,M}(x;\alpha,\beta)|\,dx\le W,\qquad
       \int_G |b_{Z,M}(x;\alpha,\beta)|\,dx\le q^{-r}W .
\]
\end{theorem}
\begin{proof}
On a shell \(\mathcal S(m,P,P)\) let \(E_H\) be the normalized
mean square of the \(H_t\)-conjugacy averages of \(\phi_M\),
namely \(v_H/v_K\) times their squared integral divided by
\(v_Jv_{H_t}^2\), with the compact averages taken using ordinary
Haar measure. The folded coset identity in the proof of
Theorem~\ref{thm:gramtail891} expresses it as
\[
 E_H=\frac{v_H}{v_K}\sum_{\mathcal C}
        \frac{|S_{\mathcal C}|^2}{|\mathcal C|},
\]
where only the \(H_t\)-orbits passing the ideal test are used and
\(S_{\mathcal C}\) is the original signed class sum. For each
contributing orbit choose \(y\in U\) in its preimage. Write
\(T_y=[\mathcal D_y:\mathcal D_y\cap H_t]\). This is the orbit
size in the lemma for the pair when \(H=K\), and for the pair and
\(\bar z\) when \(H=Z\). The orbit--stabilizer identity is
\[
        \frac{v_H}{v_K|\mathcal C|}
            =\frac{[H:H_t]}{|\mathcal C_K|T_y}.
\]
The full orbit has size \(|\mathcal C_K|\ge
q^{2(d-2a)P}/4\), and \([H:H_t]\le q^{2rt}\). Put
\(T_0=q^{d_*(2t-M)}\) and \(N=q^{4a^2(m+1)+dP}\).
Every \(T_y\) is at least \(T_0\), and the unsigned unipotent
measure of the entire shell is at most \(N\).

For \(Z\), split the tested orbits at \(T_y=q^{d_*}T_0/2\).
Every unipotent representative of an orbit below this threshold
must have \(\bar z\in\overline{T_D}\) by the lemma. The total
unipotent measure of these orbits is at most \(2q^{a-d}N\)
by Lemma~\ref{lem:terminalorbit886} applied on the whole shell.
This argument only requires the orbit size to be constant on an
\(H_t\)-orbit, which follows from its index formula. The subspace
\(\overline{T_D}\) may depend on the representative.
For all larger orbits use their denominator bound and the count
\(N\). Squaring these two counts separately shows that the
usual bound with \(T_y\ge T_0\) gains a factor at most
\(2q^{-d_*}+4q^{2a-2d}\le6q^{-d_*}\). For \(K\) use \(T_0\)
alone. Since \(P\le2a(m+1)\), we obtain in both cases
\[
      E_H\le24q^{16a^2(m+1)+8at+d_*M-\epsilon d_*},
      \qquad \epsilon={\bf1}_{H=Z}.
\]
Here \(2(r-d_*)=8a\). Lemma~\ref{lem:incidentoperator892} and
integration of its coefficient norm over the shell, as in
Theorem~\ref{thm:integralfreq892}, multiply this bound by at
most \(4q^{t+1-ru/2}\).

Sum first for \(m<1200n^3(t+1)\). There are at most
\(1200n^3(t+1)(2400an^3(t+1)+1)\) choices of \(m,P\).
Its base-\(q\) logarithm is \(O(\log(2n)+\log(2t))\) uniformly
for \(q\ge2\). At large nodes \(16a^2(m+1)\le n^7u/16\),
since \(a<n\), \(m+1\le2400n^3t\), \(t\le u\), and \(n\)
tends to infinity. Also \(d_*M\le n^7u/16\). The remaining
positive exponents \(8at+t+1\), the logarithm of the shell
count, \(\log_q192\), and the cost
\(\epsilon(r-d_*)\le4a\), together are at most \(n^7u/16\)
eventually, uniformly for \(t\ge1\). Since \(r\ge k\),
\(3n^7/16\le(r-k+n^7-6)/2\) at these nodes. It follows that
the entire bounded range contributes at most
\(\tfrac12 q^{-\epsilon r}q^{-(k-n^7+6)u/2}
\le\tfrac12 q^{-\epsilon r}W\).

Finally, on the distant shells the proof of
Theorem~\ref{thm:gramtail891} applies to \(\phi_M\) without
an increased constant. The estimate of
Lemma~\ref{lem:flagclass891} sums absolute values over all
\(H_t\)-orbits, and the truncation retains whole orbits, as
proved above. The sum on the tail is therefore at most
\(q^{-6k(t+1)}\). This is at most \(\tfrac12 q^{-r}W\), since
\(u\le2t\), \(W\ge q^{-(k-n^7+6)u/2}\), and \(r/k\to1\).
Combining the two ranges proves the theorem.
\end{proof}

The test on both factors is substantive. For example take
\(y=1+\pi^{-m}e f\), where \(f\) is a primitive middle covector
extended by zero on both arms. Then \(I_e(y)=\pi^{m+1}O\)
and \(I_\ell(y)=O\): the row condition forces divisibility of
\(\lambda(e)\), whereas the last right basis vector satisfies
the column condition and evaluates to one. Thus mixed terms
and terms in which both coefficients exceed the ideal bound
remain to be estimated.
\clearpage

\subsection{Successive row integration at a deep end}
\label{sec:rowdescent895}

Work at an original split mark with \(U,J,K,e,\ell,z,\eta\) as in
Sections~\ref{sec:allpositions883} and~\ref{sec:endtransform894}.
The basis vectors of the left arm \(L\) are \(e_1=e,\ldots,e_a\),
and \(L_j\) is the span of its first \(j\) vectors. Write \(y=1+X\in U\),
use \(B:W\to L\) for its middle columns, and let \(\beta_j\) be row
\(j\) of \(B\). Fix a bounded shell of height \(m\).

For \(1\le j\le a\) let \(p_j\) be the least valuation of
\(\lambda(e_j)\) over integral rows \(\lambda\) vanishing on
\(L_{j-1}\) with \(\lambda X\in\pi\mathcal L^*\). Define \(p_{a+1}\)
in the same way with \(\lambda|_L=0\) and with \(z\) in place of
\(e_j\). Each ideal contains \(\pi^{m+1}O\). In particular all
these integers lie between zero and \(m+1\). We have
\(\min(t,p_1)=\delta_e(y)\).

\begin{lemma}[A row change preserving all subsequent rows]
\leavevmode\label{lem:rowshear895}
Let \(H=K\), or let \(H\) be the subgroup of \(K\) fixing
the middle covector \(\eta\) modulo \(\pi\), extended by zero
on the arms. Let \(H_t\) impose the two end conditions of
Section~\ref{sec:gramtail891}. In an \(H_t\)-orbit of right
\(J\)-cosets, integrate the original \(\chi(y)\) on \(U\).
For each \(j=a,a-1,\ldots,1\) one may impose
\begin{equation}
        \beta_j(V_j)\subseteq\pi^{1-p_{j+1}}O
        \label{eq:rowtest895}
\end{equation}
without changing the integral. Here \(V_j\) is an integral direct
summand of \(W(O)\) of codimension at most \(a-j+2\) determined
by the rows strictly below \(j\).
\end{lemma}
\begin{proof}
Fix those lower rows. Take a row \(\kappa\) vanishing on \(L_j\)
with \(\kappa X\in\pi\mathcal L^*\), for which
\(\kappa(e_{j+1})\) has valuation \(p_{j+1}\) if \(j<a\), or
\(\kappa(z)\) has valuation \(p_{a+1}\) if \(j=a\).
It depends only on lower rows because \(\kappa\) annihilates
\(L_j\). To make a measurable choice, use a fixed list of
representatives of integral rows modulo \(\pi^{m+2}\) with the
prescribed leading zero entries. The test on \(\kappa X\) and a
valuation at most \(m+1\) are constant on these residue classes.
Choose the first row attaining the least valuation.
Let \(V_j\) be the common kernel on \(W(O)\) of \(\beta_i\)
for \(i>j\), of \(\kappa|_W\), and, when \(H\) fixes the
covector, of \(\eta\). This kernel is saturated since its
quotient embeds in a vector space over \(F_w\). It is therefore
an integral summand of the asserted codimension.

For \(w\in V_j\), extended by zero on the arms, put \(b=Xw=Bw\).
Then \(b\in L_j\) and \(\kappa w=\kappa b=0\). The exact identity
\begin{equation}
 (1-w\kappa)y(1+w\kappa y)=(1+b\kappa)y=:y'
 \label{eq:rowshear895}
\end{equation}
follows by multiplication using \(yw=w+b\). The left multiplier
belongs to \(H_t\). It fixes both ends exactly since
\(\kappa(e)=0\) and \(\ell(w)=0\), and in the covector case
\(\eta(w)=0\) as well. The product of the two multipliers in
left-to-right order is \(1+w\kappa X\in J\). Consequently \(y'J\)
is conjugate to \(yJ\) by the left multiplier. Also \(y'\in U\),
because \(b\) is supported in \(L_j\) and \(\kappa y\) vanishes
there. It has precisely the original rows below \(j\).

For a fixed lower part this is a measure-preserving substitution
on the upper rows. Indeed put \(\gamma=(\kappa y)|_W\), determined
by the lower part. Each upper middle row is transformed by
\(\beta_i\mapsto\beta_i(1+w\gamma)\), and \(\gamma(w)=0\).
The determinant of this linear map is one. The other upper
entries undergo translations depending on the old middle rows
and the fixed lower part. The upper block coordinates give
product additive Haar measure on \(U\) by
Lemma~\ref{lem:shellorbits884}. Equivalently the inverse in
\eqref{eq:rowshear895} uses \(-w\), since \(X'w=b\) and
\(\kappa y'=\kappa y\), where \(X'=y'-1\). This substitution
thus preserves the shell and the chosen orbit in both directions.

The only change in the character \eqref{eq:actualphase882} is
the \(j\)-th simple entry. If \(j<a\), it changes before division
by \(\pi\) by the unit \(c_j\) times
\(\beta_j(w)\kappa(e_{j+1})\), since
\(\kappa y e_{j+1}=\kappa(e_{j+1})\).
If \(j=a\), the change is \(c_a\beta_a(w)\kappa(z)\), since
\(yz-z\in L\). All earlier simple entries are fixed because
\(\kappa y\) vanishes on \(L_j\). All right-arm and middle-to-right
entries are fixed. Averaging in \(w\) on \(V_j\) with probability
Haar gives zero unless \eqref{eq:rowtest895} holds, by the
conductor of \(\Psi\). The substitutions fix all tests previously
imposed at indices greater than \(j\), as those tests only read
rows below or at their own index. This proves the assertion
successively in the stated order.
\end{proof}

\begin{theorem}[Signed sums at full first-end depth]
\leavevmode\label{thm:rowdescent895}
Use either group \(H\) of Lemma~\ref{lem:rowshear895} and any
\(t\ge1\). On \(\mathcal S(m,P,P)\) and for \(0\le h\le t\),
\[
 \sum_{\mathcal C:\,\delta_e\ge h}
      \left|\int_{\mathcal C\cap U}\chi(y)\,dy\right|
             \le 2 q^{16n^2(m+1)+d(P-h)} .
\]
The condition on \(\delta_e\) is constant on each of the indicated
\(H_t\)-orbits.
\end{theorem}
\begin{proof}
The orbit invariance follows from the definition of \(I_e\) in
Section~\ref{sec:endtransform894}. Apply the row lemma within
each orbit and bound the sum of absolute values by the unipotent
measure surviving all the tests and \(p_1\ge h\).

Consider the module generated by the rows of \(X\) modulo
\(\pi\mathcal L^*\) using integral coefficients. Its length is
\(P\), by the Smith entries of \(X\) and its transpose.
Let \(\nu\) be the length generated by the rows below the left
arm. Adding rows \(a,a-1,\ldots,1\) increases the length by
\(p_a,p_{a-1},\ldots,p_1\), respectively. In fact the new
quotient is cyclic, and the least integer \(b\) such that the
row times \(\pi^b\) belongs to the earlier module is exactly
\(p_j\). This follows by writing the linear relation as
\(\lambda X\in\pi\mathcal L^*\) and scaling its leading unit.
Thus
\[
                     P=\nu+\sum_{j=1}^a p_j .
\]
This also implies
\(\pi^{p_j}\beta_j(V_j)\subseteq\pi O\), since all middle
components of lower rows are either zero or among the
\(\beta_i\) with \(i>j\). Together with \eqref{eq:rowtest895}
this restricts the values of \(\beta_j\) on \(V_j\) to
\(\pi^{1-\min(p_j,p_{j+1})}O\).

Fix a profile \(p_1,\ldots,p_{a+1}\). For fixed lower rows,
choose a basis splitting \(V_j\) inside \(W(O)\). The basis
change preserves the maximum entry bound and additive measure.
Each of its complementary coordinates costs at most \(q^{m+1}\)
in measure for the row \(\beta_j\), since \(\pi O\) has measure
one. Each coordinate on \(V_j\) costs at most
\(q^{\min(p_j,p_{j+1})}\). The other variable entries of the
\(j\)-th row are in at most \(2a\) outer coordinates. Integrating
first row one and then proceeding downwards, these bounds give
a factor at most
\[
 q^{a(3a+1)(m+1)+d\sum_{j=1}^a\min(p_j,p_{j+1})}
\]
over the bottom rows. All conditional upper bounds used here
are uniform in the fixed lower rows.

The bottom rows of \(X\) have entries only in the last \(a\)
columns, comprising the blocks \(D: R\to W\) and \(E:R\to R\).
For fixed \(\nu\) their total measure within the entry bound
is at most \(q^{2a^2(m+1)+d\nu}\). Indeed their common row
module in \((\pi^{-m}O/\pi O)^a\) has at most \(a\) generators,
so there are at most \(q^{a^2(m+1)}\) possible modules.
There are at most \(a^2\) entries of \(E\), and each of the
\(d\) rows of \(D\) then lies in a module of size \(q^\nu\)
modulo \(\pi O^a\).

Requiring \(p_{a+1}=v\) reduces this last upper bound by a factor
at most \(2q^{-(d-a)v}\). To verify it let \(T_D=D(R)\cap W(O)\)
be the saturated lattice of rank at most \(a\) used in
Lemma~\ref{lem:terminalorbit886}. Every integral middle covector
annihilating it, extended by zero on the arms, annihilates \(X\).
Thus \(p_{a+1}\ge v\) implies \(z\in T_D+\pi^v W(O)\), by a
basis splitting this lattice. Integral middle changes of basis
preserve the entry bounds, bottom measure and \(\nu\).
For \(v\ge1\) a primitive vector is uniform under this action
on \(q^{dv}(1-q^{-d})\) residues, of which at most \(q^{av}\)
belong to the given submodule. Averaging the indicator gives
the bound, also valid trivially for \(v=0\).
This calculation is a count on all bottom rows after bounding
the upper ones, so it requires no change of basis inside a
restricted orbit.

Now
\[
 \sum_{j=1}^a\min(p_j,p_{j+1})
     =\sum_{j=1}^a p_j-\sum_{j=1}^a(p_j-p_{j+1})_+
     \le \sum_{j=1}^a p_j-p_1+p_{a+1}.
\]
After multiplication the exponent involving \(d\) is at most
\(d(P-p_1)+a p_{a+1}\). There are at most \((m+2)^{a+1}\)
profiles, and \(\nu\) is determined by \(P\) and the first
\(a\) entries. Since \(p_{a+1}\le m+1\), \(n=a+1\) and
\(m+2\le2^{m+1}\), all remaining exponents fit inside
\(16n^2(m+1)\). This proves the result.
\end{proof}

The additional terminal condition permitted in the row change
fixes a covector. The group fixing the original terminal vector
acts differently on \(1-w\kappa\). Its reduction sends \(\bar z\)
to \(\bar z-\bar w\,\bar\kappa(\bar z)\). This dependence matters
when using the same calculation on that side.
\clearpage

\subsection{Row cancellation with the terminal vector fixed}
\label{sec:twodepth896}

Keep \(p_1,\ldots,p_{a+1}\) and the unipotent coordinates of
Section~\ref{sec:rowdescent895}. At the original vector stabilizer
\(Z\) the shear used there must also preserve \(\bar z\). We account
for this condition by adjoining the column \(z\) to \(X=y-1\).
All row modules below are taken with integral coefficients modulo
\(\pi\) times the corresponding ambient lattice, also when some
entries have denominators.

\begin{lemma}[The extra column and the terminal orbit]
\leavevmode\label{lem:augmentedrows896}
For \(1\le i\le a+1\) let \(\sigma_i\) be the increase in row-module
length on adjoining the column \(z\) to the rows from \(i\)
downwards. At \(i=a+1\) this means all middle and right rows.
Each increase is zero or one. Moreover \(\sigma_i=1\) precisely
when some integral row \(\lambda\) vanishing on \(L_{i-1}\)
satisfies \(\lambda X\in\pi\mathcal L^*\) and
\(\lambda(z)\in O^\times\).

For \(1\le j\le a\) impose also \(\lambda(z)\in\pi O\) in
the definition of \(p_j\), calling the new value \(p_j^Z\).
Set \(p_{a+1}^Z=\max(1,p_{a+1})\). Then
\begin{align*}
 p_j^Z &=p_j+\sigma_j-\sigma_{j+1},\\
 p_{a+1}^Z &=p_{a+1}+\sigma_{a+1},\\
 \sum_{j=1}^a\min(p_j^Z,p_{j+1}^Z)
   &\le P-\nu-p_1+\sigma_2+p_{a+1}.
\end{align*}
If \(\sigma_2=1\), the stabilizer of \(yJ\) in \(K\) moves the
pair of ends at precision \(t\) together with \(\bar z\) through
at least \(\tfrac12 q^{d_*(2t-\delta_e-\delta_\ell+1)}\) positions,
where \(d_*=d-2a\).
\end{lemma}
\begin{proof}
Project the augmented row module to the old one. The kernel
is a submodule of \(O/\pi O\) in the last coordinate. Its
nonzero elements come exactly from linear combinations of the
specified rows satisfying the two tests on \(\lambda\).
The successive cyclic length increments for the augmented matrix
are \(p_j^Z\), just as for \(X\) in Theorem~\ref{thm:rowdescent895}.
Taking differences gives the first identity. The second follows
since \(\sigma_{a+1}=1\) precisely when \(p_{a+1}=0\). The
smallest positive valuation of the bottom evaluation at \(z\)
under the extra constraint is then one, by multiplying a unit
value by \(\pi\). The last inequality follows by summing the
first identities and subtracting at least \(p_1^Z-p_{a+1}^Z\)
in the positive successive drops.

If \(\sigma_2=1\), take a witnessing \(\lambda\) with
\(\lambda(e)=0\). In the middle summand \(W_0\) of
Lemma~\ref{lem:endorbit894}, take \(v\in W_0(O)\) with
\(1+\lambda(v)\) a unit. The transformation \(g=1+v\lambda\)
is integral and invertible. It fixes \(e,\ell\) exactly, and
\(y^{-1}gy-g=v\lambda X\) belongs to \(\pi M_r(O)\), so it
stabilizes \(yJ\). Since \(\lambda(z)\) is a unit, the distinct
residues of these \(v\) give distinct images of \(\bar z\).
At least half the residue vectors are allowed, because an
excluded affine hyperplane has fraction \(q^{-1}\), or is empty
if \(\lambda\) vanishes on the summand modulo \(\pi\).
Its rank is at least \(d_*\). Multiplying the resulting lower
bound in the pair stabilizer by the pair orbit size of
Lemma~\ref{lem:endorbit894} proves the assertion.
\end{proof}

\begin{theorem}[Two end depths in the refined orbital mean square]
\leavevmode\label{thm:twodepth896}
Let \(H=K\) or \(Z\), \(t\ge1\), and
\(\mathcal S=\mathcal S(m,P,P)\). With ordinary Haar measure in
\(C_{H_t}\), at the original split marks,
\begin{equation}
 \frac{v_H}{v_K}\frac1{v_Jv_{H_t}^2}
                \int_{\mathcal S}|C_{H_t}(x)|^2\,dx
 \le  256(t+1)^2q^{64n^2(m+1)+8at}.
 \label{eq:twodepth896}
\end{equation}
All end ideals and all mixed contributions to the coefficient
mean square are included.
\end{theorem}
\begin{proof}
Group the \(H_t\)-orbits meeting \(U\) by
\(p=\delta_e,\ s=\delta_\ell\). Put \(T_y=[\mathcal D_y:
\mathcal D_y\cap H_t]\) for the full coset stabilizer
\(\mathcal D_y\). This index is constant on each refined orbit.
For \(K\) assign all orbits to type \(\epsilon=0\). For \(Z\)
assign type \(\epsilon=1\) exactly when
\[
               T_y\ge\tfrac12 q^{d_*(2t-p-s+1)} .
\]
In either case we have
\(T_y\ge\tfrac12q^{d_*(2t-p-s+\epsilon)}\) by
Lemma~\ref{lem:endorbit894}. For \(Z\) every unipotent
representative on a type-zero orbit has \(\sigma_2=0\), by
Lemma~\ref{lem:augmentedrows896}.

For \(K\), the sum of the absolute signed class sums in this
group of orbits has the bound of Theorem~\ref{thm:rowdescent895}
with \(h=p\). For \(Z\) perform the same successive shears using
the covectors chosen with \(\kappa(z)\in\pi O\). For \(j<a\)
their leading valuations are \(p_{j+1}^Z\); for \(j=a\) their
evaluation at \(z\) has valuation \(p_{a+1}^Z\).
The left multiplier \(1-w\kappa\) now fixes \(\bar z\) for every
integral \(w\). The transformation, measure, and phase identities
in Lemma~\ref{lem:rowshear895} apply verbatim. All the substitutions
remain in the original refined orbit. Thus the integral on each
orbit can be cut by all the resulting lower-row tests. The
middle entries of row \(j\), on the summand in that lemma, now
cost at most \(q^{d\min(p_j^Z,p_{j+1}^Z)}\): a relation in the
augmented module also gives the required relation on the middle
columns of \(X\).

For a fixed profile, the bottom coordinate count is still the
one in Theorem~\ref{thm:rowdescent895}, including the gain
\(2q^{-(d-a)p_{a+1}}\) with the original bottom depth. We can
use that bound because \(\nu\) here is the length for the
unaugmented bottom rows. The last inequality in the preceding
lemma shows that the total exponent involving \(d\), after the
bottom count, is at most \(d(P-p_1+\sigma_2)+a p_{a+1}\).
In the selected group of orbits \(\sigma_2\le\epsilon\) for all
surviving representatives. We may sum profiles of both the
original and augmented valuations. Each lies in \([0,m+1]\);
there are at most \((m+2)^{2a+2}\) choices, and \(\nu\) is
determined by the original profile. The other row and submodule
costs are as in that theorem. In both groups \(H\) this gives
\[
        \sum_{\mathcal C}|S_{\mathcal C}|
          \le 2q^{20n^2(m+1)+d(P-p+\epsilon)} .
\]

We also have \(2q^{20n^2(m+1)+d(P-s)}\) for the same sum.
Indeed pass to the dual space by inverse transposition and
reverse the arm order. This is a group isomorphism, preserving
\(J\), right cosets and the two end conditions. The first end
valuation there is the original \(\delta_\ell\), by transposing
the defining column test on \(y^{-1}-1\). The two Smith lengths
are interchanged and both equal \(P\). The unipotent Haar
normalization is preserved as \(U\cap J\) maps to its counterpart.
On consecutive blocks inversion negates the entries, and
transposition exchanges the terminal vector and covector up to
units. Hence the character has the same form with unit
coefficients and with terminal pairing still a unit. The
original \(Z\) now fixes the new terminal covector modulo \(\pi\).
This is precisely the second group in
Theorem~\ref{thm:rowdescent895}. Its estimate holds also on any
subset of the orbits at depth \(s\), since it sums absolute
values. This proves the second bound.

Take the smaller of the two bounds, giving
\(2q^{20n^2(m+1)+d(P-\max(p-\epsilon,s))}\).
In the normalized orbital sum for the left of
\eqref{eq:twodepth896}, the multiplier on
\(|S_{\mathcal C}|^2\) is
\[
 \frac{[H:H_t]}{|\mathcal C_K|T_y}
 \le8q^{2rt-2d_*P-d_*(2t-p-s+\epsilon)}
\]
by Lemma~\ref{lem:shellorbits884}. Bounding the sum of squares
by the square of the preceding bound leaves, besides
\(40n^2(m+1)+4aP+8at\), the exponent
\[
 d_*(p+s-\epsilon)-2d\max(p-\epsilon,s)
  =-d|p-s-\epsilon|-2a(p+s-\epsilon)\le2a.
\]
As \(P\le2a(m+1)\), these terms are covered by the exponent
in \eqref{eq:twodepth896}. There are at most \(2(t+1)^2\)
groups of orbits. The stated constant includes them and the
numerical factors above.
\end{proof}

The possible loss of one residue digit from the vector condition
has thus been included in the same estimate as the two end
depths. No bound on either depth relative to the Fourier conductor
is imposed in \eqref{eq:twodepth896}.
\clearpage

\subsection{The complete unrestricted mass and its principal term}
\label{sec:allkmass897}

At zero frequency the full integral for \(K\) has principal
contribution one in our normalization. We estimate its complement
over all heights. Retain \(B,D,S\) for the blocks of \(y\in U\)
in Section~\ref{sec:armshear887}. The right-to-middle block of
\(y^{-1}\) is \(D^-=-D S^{-1}\).

\begin{lemma}[A saving on every nonprincipal shell]
\leavevmode\label{lem:onesave897}
For the \(K\)-orbits of right \(J\)-cosets in
\(\mathcal S(m,P,P)\) with \(P>0\),
\[
       \sum_{\mathcal C}|S_{\mathcal C}|
          \le 6q^{16n^2(m+1)+d(P-1)} .
\]
\end{lemma}
\begin{proof}
Partition the unipotent integral into three regions. First take
\(B\notin\pi M_{a,d}(O)\). On this region use the shear
\eqref{eq:rowshear895}. At step \(j\) the entire middle matrix
changes by \(B\mapsto B(1+w\gamma)\), including the unchanged
lower rows since those rows vanish on \(w\). Here
\(\gamma=(\kappa y)|_W\) is integral because \(\kappa\) is integral
and \(\kappa X\in\pi\mathcal L^*\). Also \(\gamma(w)=0\).
Thus \(1+w\gamma\) and its inverse are integral, so that
the first region is preserved. Every shear stays in the
\(K\)-orbit of the original coset. The row proof of
Theorem~\ref{thm:rowdescent895} therefore applies within this
region and each orbit, without imposing a positive end depth.

Retain the actual positive drops in the last line of its count.
Writing \(v=p_{a+1}\), the saving in the exponent of \(d\) is
\[
                  v+\sum_{i=1}^a(p_i-p_{i+1})_+
                         \ge \max_{1\le i\le a}p_i.
\]
The residual cost \(av\) is already included in \(16n^2(m+1)\).
Here at least one \(p_i>0\): the bottom rows of \(X\) have zero
middle component, while at least one of the top rows has
nonzero middle component modulo \(\pi O\). Thus the full row
module has larger length than the bottom one. The first region
contributes at most \(2q^{16n^2(m+1)+d(P-1)}\).

In the second region \(B\in\pi M_{a,d}(O)\) and
\(D^-\notin\pi M_{d,a}(O)\). Use the inverse-transpose isomorphism
and reverse the arms as in Theorem~\ref{thm:twodepth896}. The
new left middle block is the transpose of \(D^-\), up to the
basis reversal. The old test on \(B\) becomes the test on the
right-to-middle block of the inverse matrix in the new
coordinates. This latter block stays fixed under every left row
shear. Indeed \(y'^{-1}=y^{-1}(1-b\kappa)\), and \(y^{-1}b\in L_j\),
so the inverse changes only in the top \(j\) rows. Thus both
tests of the second region survive the calculation. The same
bound as in the first region follows.

In the third region \(B,D^-\) both have entries in \(\pi O\).
For each right unitriangular \(S\), the substitution
\(D\mapsto-D S^{-1}\) preserves additive Haar on the rows of \(D\)
as \(\lvert\det S\rvert=1\). Thus the total conditional measure
allowed for \(B,D\) is at most one, using
\(\operatorname{vol}(\pi O)=1\) for each coordinate. All other
upper entries together cost at most \(q^{2a^2(m+1)}\).
This estimates the sum over disjoint classes by absolute measure
and is within the displayed bound since \(P\ge1\). Adding the
three contributions proves the lemma.
\end{proof}

\begin{theorem}[The unrestricted position mass on the full group]
\leavevmode\label{thm:allkmass897}
At every original split mark, uniformly for all sufficiently
large nodes of the shorter interval,
\[
 0\le\frac{\int_G|C_K(x)|^2\,dx}{v_Jv_K^2}-1
          =\frac{\int_{G\setminus J}|C_K(x)|^2\,dx}{v_Jv_K^2}
          \le q^{-2k+20000n^5}.
\]
\end{theorem}
\begin{proof}
The contribution from \(J\) is one because \(J\) is normal in
\(K\) and \(\phi=1\) there. A shell with positive mass has \(P=Q\).
Also \(P=0\) exactly on \(J\), since \(P=0\) means
\((x-1)\mathcal L\subset\pi\mathcal L\).
For \(P>0\), apply the full-orbit formula
\eqref{eq:cosetnorm884} and the lower orbit bound
\(q^{2(d-2a)P}/4\) from Lemma~\ref{lem:shellorbits884}.
The preceding lemma yields at most
\[
       144q^{32n^2(m+1)+4aP-2d}
           \le 144q^{40n^2(m+1)-2d}
\]
per shell, using \(P\le2a(m+1)\). When \(m<400n^3\), there
are at most \(400n^3(800an^3+1)\) such shells. Their sum is
at most \(q^{-2k+17000n^5}\) eventually, since \(d=k+1\)
and \(m+1\le400n^3\).

For \(m\ge400n^3\) apply the signed large-shell bound in the
proof of Theorem~\ref{thm:allfixermass888}, namely
\(400q^{-31kL/16}\) per shell with
\(L=\lfloor m/(200n^3)\rfloor\). Now \(L\ge2\). There are at most
\(200n^3(800an^3L+1)\) pairs of shell indices at each \(L\).
Summing the geometric series with its linear factor gives at
most \(q^{-3k}\) at large nodes, since \(62/16>3\) and the
polynomial in \(a,n\) is absorbed by the difference of exponents.
Adding proves the stated inequality. On our interval
\(n^5=o(k)\), so the bound tends uniformly to zero.
\end{proof}

In particular this calculation determines the zero-frequency
principal term before any restriction to a terminal-vector fixer.
The exact relative-residue factor \(\rho_w\) at an unselected
mark is still given by Corollary~\ref{cor:allmarks888}.
\clearpage

\subsection{The full first relative Fourier bound at the marked places}
\label{sec:fullfirst898}

For the actual first relative coordinate of
\eqref{eq:gramintegrand891} write
\(\mathcal B_H(\alpha,\beta)=\int_G|b_H(x;\alpha,\beta)|\,dx\).
The coefficients in this definition are the original \(\phi\) in
both variables.

\begin{theorem}[All positions and both unrestricted end depths]
\leavevmode\label{thm:fullfirst898}
At every original split mark, for both \(H=K\) and \(H=Z\), put
\(u=c_\alpha+c_\beta\), \(t=\max(c_\alpha,c_\beta)\),
and \(K_0=k-n^7\). For a nonzero rational frequency with \(u\ge1\),
at all sufficiently large nodes of the shorter interval,
\begin{equation}
 \mathcal B_H(\alpha,\beta)
       \le q^{-n^7u/4}\min\{1,I^{K_0/2}D^{-3}\}.
 \label{eq:fullfirst898}
\end{equation}
In particular the factor \(q^{-r}\) in addition to the indicated
frequency weight is available whenever \(n^7u\ge4r\).
If \(u=0\), the values are the zero-frequency masses. They satisfy
\[
 1\le\mathcal B_K\le1+q^{-2k+20000n^5},\qquad
             \mathcal B_Z\le2q^{-9k/10}.
\]
\end{theorem}
\begin{proof}
On a fixed shell, use the coset-pair expression for
\(b_H(x;\alpha,\beta)\) in Theorem~\ref{thm:gramtail891}.
Apply the finite Fourier operator bound of
Lemma~\ref{lem:incidentoperator892} to its two coefficient
vectors. Conjugating \(x\) by each coset representative preserves
the shell and its measure. Thus the integral on the shell is
at most \(4q^{t+1-ru/2}\) times the left side of
\eqref{eq:twodepth896}. The latter includes the original
coefficients for all end valuations. In particular mixed
contributions between different depths in the two factors
are controlled by the same operator estimate.

For \(m<1200n^3(t+1)\), sum this bound over at most
\(1200n^3(t+1)(2400an^3(t+1)+1)\) shells. The total, for all large
nodes, is at most
\[
                       \tfrac12 q^{-ru/2+n^7u/8}.
\]
Indeed \(64n^2(m+1)\le153600 n^5u\) since \(t\le u\).
The linear terms \(8at+t+1\), the logarithms to base \(q\)
of the shell count and of \(2048(t+1)^2\), and this last
multiple of \(n^5u\), together are at most \(n^7u/8\)
eventually. This is uniform for \(u\ge t\ge1\), since
\(\log(2t)\ll t\), \(q\ge2\), and \(n\to\infty\).

For the other shells Theorem~\ref{thm:gramtail891} gives
\(q^{-6k(t+1)}\). This is also at most
\(\tfrac12 q^{-ru/2+n^7u/8}\), since \(u\le2t\) and \(r/k\to1\).
Hence the full integral is at most \(q^{-ru/2+n^7u/8}\).
As \(r\ge k\), at large nodes this is bounded by
\[
 q^{-n^7u/4}q^{-(k-n^7+6)u/2}
       \le q^{-n^7u/4}\min\{1,I^{K_0/2}D^{-3}\}.
\]
The last inequality uses \(D=q^u\), \(I\ge D^{-1}\) and \(K_0>0\).
At \(u=0\) the relative wave is identically one on \(H^2\), so
\(\mathcal B_H=\int_G|C_H|^2/(v_Jv_Hv_K)\). Apply
Theorems~\ref{thm:allkmass897} and~\ref{thm:allfixermass888}.
\end{proof}

\begin{corollary}[The selected factors in the first rational frequency sum]
\leavevmode\label{cor:firstsample898}
Use the field \(E\), shorter interval, selected set \(\Xi\),
and indices \(j_w\) of Theorem~\ref{thm:bothmarks873}.
Let \(T\) contain both primes of \(E\) at every \(w\in\Xi\)
and exclude the unselected marks as in
Proposition~\ref{prop:sampling889}. In the common first-coordinate
scale of Theorem~\ref{thm:freqnode875}, set
\[
 V_w(y)=I_w(y)^{K_0/2}D_w(y)^{-3},\qquad
 Z_w^0=\frac{v_{Z_w}}{v_{K_w}}=\frac{1}{q_w^r-1}.
\]
Write \(\mathcal B_{Z_w}(y)\) for the full first transform at the
two local components of the rational frequency. Then
\begin{equation}
 \sum_{0\ne y\in\mathcal O_{E,T}}
       W_{T,c,K_0}(y)
       \prod_{w\in\Xi}
         \frac{B_0q_w^{n-j_w}\mathcal B_{Z_w}(y)}
                        {Z_w^0 V_w(y)}
             \le b^{-7kf}
 \label{eq:firstsample898}
\end{equation}
at all sufficiently large nodes. The terminal volume in the
denominator permits comparison at its original principal scale.
\end{corollary}
\begin{proof}
Both components at a marked split place use an additive
character of conductor \(O\); the trace normalization changes
the frequencies there only by units. Thus their valuations give
\(I_w,D_w\) in the rational weight. If \(u_w\ge1\), the theorem
gives \(\mathcal B_{Z_w}(y)\le V_w(y)\). If \(u_w=0\), it gives
\(\mathcal B_{Z_w}(y)\le2q_w^{-9k/10}\le1\le V_w(y)\).
The remaining multiplier per selected place is at most
\[
        B_0 q_w^{a+r}\le b^{6k+11(a+r)}\le b^{18k},
\]
using \(B_0=b^{6k}\), \(q_w\le b^{11}\), \(n-j_w\le a\), and
\(r=k+1+2a\) with \(a=o(k)\).
There are fewer than \(f/4\) selected places. Their product
is at most \(b^{5kf}\). The sum of \(W_{T,c,K_0}\) itself
is at most \(b^{-12kf}\) by Theorem~\ref{thm:freqnode875}.
This proves \eqref{eq:firstsample898}, including the empty
product.
\end{proof}

The factors at the selected split places in \eqref{eq:firstsample898}
are now estimates of the actual first relative transform with no
test on end ideals. At the other places the formula still uses
the explicit arithmetic weight. Replacing those factors by their
actual transforms, and passing through later pivots and unequal
rational heads, are additional steps in the projected-period
comparison.
\clearpage

\subsection{The first transform at standard unmarked split places}
\label{sec:standardfirst899}

We change the level as well as the residue cardinality. At a split
place with the standard integral lattice take \(G=\mathrm{GL}_r(F_w)\),
\(K=\mathrm{GL}_r(O)\) and now \(J=K\). Both \(K\) and \(U\cap K\)
have volume one. Assume the simple character has additive conductor
\(O\), unit coefficients on the arms, and primitive integral
terminal vector and covector. Define
\(\phi_0(uy)=\chi_0(u)\) for \(u\in U,\ y\in K\), and extend it
by zero off \(UK\). It is well defined since \(\chi_0=1\) on
\(U\cap K\). In this normalization \(\mathcal B^0_w\) is the
integral of the absolute value of \eqref{eq:gramintegrand891}
with \(H=J=K\) and \(\phi=\phi_0\). Indeed the image of \(K\) in
\(U\backslash G\) carries its quotient measure from \(K\) with
fiber \(U\cap K\) of volume one. The relative wave is unchanged
on both \(U\)-orbits as in that equation.

These hypotheses hold in the integral split model of
Theorem~\ref{thm:orbitfactors871} at all split places outside a
finite set at the node. The fixed flag basis and its dual, the
form and its inverse, and the class representatives are integral
there. The level \(M\) is a unit, and the normalization by the
different makes the trace character have integer-ring conductor.
A fixed nonzero terminal vector and its paired covector are
primitive away from finitely many places. The arguments below
are uniform over the places satisfying the hypotheses.

Set \(P=\operatorname{length}\big(((x-1)\mathcal L+\mathcal L)/\mathcal L\big)\),
where length is taken on the quotient, and use the corresponding
integer for \(x^{-1}\). In particular \(P\) is the length of
\((x\mathcal L+\mathcal L)/\mathcal L\). Both lengths and \(m\)
are invariant under multiplication by \(K\) on either side.
For a contributing coset, \(\operatorname{ord}\det x=0\), so the
lengths agree by the elementary divisors of \(x\). Write
\(\mathcal S(m,P)\) for such a shell. Except on \(K\) one has
\(1\le m\le P\): a largest pole occurs in \(x\) or \(x^{-1}\)
and contributes at least its order to the respective length.
On a shell meeting \(U\), \(P\le2am\) by applying the Smith
formula to \(x-1\) at a unipotent representative of rank at
most \(2a\) after subtracting the identity.

\begin{lemma}[Signed shell estimates at level one]
\leavevmode\label{lem:standardrows899}
Let \(K_t\) fix both ends modulo \(\pi^t\), with \(K_0=K\) in
this lemma. Use the valuations \(\delta_e,\delta_\ell\) of
Section~\ref{sec:endtransform894} truncated at \(t\), now with
the row and column conditions taking values in the integral
lattices. On the \(K_t\)-orbits of right \(K\)-cosets in
\(\mathcal S(m,P)\), at prescribed end depths \(p,s\),
\[
      \sum_{\mathcal C}|S_{\mathcal C}|
            \le 2q^{16n^2(m+1)+d(P-\max(p,s))}.
\]
For \(P>0\) the unrestricted sum over all these orbits is at most
\begin{equation}
     6q^{16n^2(m+1)+d(P-L)},\qquad
                L=\left\lceil\frac{P}{2a^2+1}\right\rceil .
 \label{eq:standardshell899}
\end{equation}
\end{lemma}
\begin{proof}
Here the class sum is \(\int_{U\cap\mathcal C}\chi_0(y)\,dy\);
the intersection with a right coset has volume one if nonempty
and the character is constant there. In the row proof of
Theorem~\ref{thm:rowdescent895} take the row module modulo
\(\mathcal L^*\) and impose \(\kappa X\in\mathcal L^*\).
With \(p_j\) the corresponding least valuations, the row length
identity \(P=\nu+\sum p_j\) is unchanged. In
\eqref{eq:rowshear895} the product of the multipliers is
\(1+w\kappa X\in K\): it is integral with integral inverse
\(1-w\kappa X\), since \(\kappa Xw=0\). The left multiplier
still fixes both ends. The phase has no division by \(\pi\),
so its row test is \(\beta_j(V_j)\subseteq\pi^{-p_{j+1}}O\).
The relation modulo the integral lattice gives the second bound
\(\beta_j(V_j)\subseteq\pi^{-p_j}O\). The measure of their
intersection in one coordinate is precisely
\(q^{\min(p_j,p_{j+1})}\), since \(O\) now has volume one.
Every discarded outer coordinate costs at most \(q^m\).
The row substitution, its inverse, and its Jacobian are the
same algebraic identities as before. The bottom module uses
\((\pi^{-m}O/O)^a\), and \(p_{a+1}=v\) still forces
\(z\in T_D+\pi^vW(O)\), with the same probability bound.
Thus the full count of Theorem~\ref{thm:rowdescent895}, in this
normalization, costs no more than
\(2q^{16n^2(m+1)+d(P-p)}\). Inverse transposition gives the
bound with \(s\), hence the first assertion.

For the second assertion define \(b\) as the least nonnegative
integer with \(\pi^bB\) integral, and define \(b^-\) the same
way for \(D^-=-DS^{-1}\).
On \(b>a(L-1)\), the projection of the row module onto its middle
coordinates has length at least \(b\) by a largest Smith pole of
\(B\). The bottom module projects to zero. Hence
\(\sum_{j=1}^a p_j\ge b\) and some \(p_j\ge L\).
The calculation in Lemma~\ref{lem:onesave897} saves this maximum
in the exponent of \(d\), since
\(v+\sum(p_j-p_{j+1})_+\ge\max p_j\).
On \(b\le a(L-1),\ b^->a(L-1)\) use inverse transposition.
These two regions are preserved by their respective shears:
the tested left block is multiplied on the right by an integral
unit matrix, and the opposite right-to-middle inverse block is
unchanged. On the remaining region the measure of \(B,D^-\)
is at most \(q^{2a^2d(L-1)}\), using the substitution
\(D\mapsto-D S^{-1}\) at fixed \(S\). The other entries cost at
most \(q^{2a^2m}\). Since
\(2a^2(L-1)\le P-L\) by the definition of \(L\), summing the
three regions proves \eqref{eq:standardshell899}.
\end{proof}

\begin{theorem}[The full first transform at a standard split place]
\leavevmode\label{thm:standardfirst899}
Write \(u=c_\alpha+c_\beta\), \(t=\max(c_\alpha,c_\beta)\) and
\(K_0=k-n^7\). Uniformly at the places just described and at all
sufficiently large nodes,
\[
 \begin{aligned}
 1\le\mathcal B^0_w(0,0)&\le1+q^{-2k+1000n^4},\\
 \mathcal B^0_w(\alpha,\beta)&\le
 q^{-n^7u/4}\min(1,I^{K_0/2}D^{-3})\qquad(u\ge1)
 \end{aligned}
\]
for nonzero rational frequencies in the second line. Integral
local frequencies give the same value as \((0,0)\).
\end{theorem}
\begin{proof}
We detail the orbit denominator at this level. For a unipotent
representative write \(y=Y\oplus1_{d_0}\) as in
Lemma~\ref{lem:shellorbits884}, with \(d_0\ge d_*=d-2a\).
A matrix \(g\in K\) stabilizes \(yK\) precisely when
\(y^{-1}gy-g\) is integral. Indeed \(y^{-1}gy\) has unit
determinant, so integrality also gives an integral inverse.
In the two off-diagonal rectangles of \(g\), the individual
rows must take \(Y-1\) to integral rows, and the columns must
be taken to integral columns by \(Y^{-1}-1\).
Each condition has index \(q^P\) by the Smith formula. The
\(2d_0\) conditions are independent in a uniform integral matrix
before conditioning on invertibility. Its probability of
invertibility is \(\prod_{i=1}^r(1-q^{-i})\ge1/4\).
Thus the full orbit size is at least \(q^{2d_*P}/4\).
In its stabilizer the independent lower and upper rectangular
shears of Lemma~\ref{lem:endorbit894} are valid with integral
row and column conditions. The same orbit argument gives at
least \(q^{d_*(2t-p-s)}\) images of the end pair.

The coefficient mean square is the sum of squared class sums
divided by orbit sizes, by Haar averaging on these right cosets.
At level \(K_t\) its inverse orbit sizes are bounded by
\([K:K_t]\) divided by the full orbit size and the number of end
images under its stabilizer. Here \([K:K_t]\le q^{2rt}\).
Using the first bound in the lemma at each \(p,s\), summing at
most \((t+1)^2\) pairs and using \(P\le2am\) gives at most
\[
     16(t+1)^2 q^{40n^2(m+1)+8at}
\]
on the shell. The depth terms satisfy
\(d_*(p+s)-2d\max(p,s)\le0\).
The estimate of Lemma~\ref{lem:incidentoperator892} for the
relative wave on end pairs involves only \(K,e,\ell\) and its
additive conductor. It therefore multiplies the last bound by
at most \(4q^{t+1-ru/2}\) for \(t\ge1\).

Alternatively use \eqref{eq:standardshell899} and the operator
norm bound one, discarding the extra end images in the orbit
denominator. The resulting shell estimate is
\[
                  144q^{2rt+40n^2(m+1)-2dL}.
\]
For fixed \(L\ge1\) there are at most \((2a^2+2)^2 L^2\)
possible shells and \(m+1\le(2a^2+2)L\). At \(t=0\) their sum
is at most \(q^{-2k+1000n^4}\) for large nodes: after bounding
\(40n^2(2a^2+2)\le160n^4\), sum the geometrically decreasing
series with its quadratic factor. The contribution on \(K\)
is exactly one. This proves the first assertion.

For \(t\ge1\), at heights \(m\ge20(2a^2+1)(t+1)\) one has
\(L\ge20(t+1)\). For large nodes both \(2rt\) and
\(40n^2(m+1)\) are at most \(kL/8\). The series over these
\(L\) is therefore at most \(q^{-10k(t+1)}\), by the same
shell count. In the remaining range the end-pair operator
estimate gives a sum at most
\(\tfrac12 q^{-ru/2+n^7u/8}\). To see this, the exponent
\(40n^2(m+1)\) is \(O(n^4u)\), the other positive exponents
are \(O(nu)\), and the shell count and \((t+1)^2\) are polynomial
in \(n,t+1\), since \(P\le2am\). Their logarithms and numerical
constants are absorbed uniformly by \(n^7u/8\).
Adding the distant sum and applying the last exponent comparison
of Theorem~\ref{thm:fullfirst898} proves the second assertion.
\end{proof}

\begin{corollary}[Product over the standard unmarked split places]
\leavevmode\label{cor:standardeuler899}
In \(E=\mathbb Q(\sqrt{-d},\mu_{q_1})\) of
Corollary~\ref{cor:oddwindow677}, let \(\mathcal V_0\) be any set
of standard unmarked split places of \(F=E^+\) with residue
characteristic prime to \(q_1\). Here \(d\) in the field denotes
the quadratic prime. For a nonzero rational frequency, write
\(u_w,I_w,D_w\) for its local data. Then
\[
 \prod_{w\in\mathcal V_0}\mathcal B^0_w(y)
 \le \exp\!\big(O(q_1^{\,2-k})\big)
       \prod_{\substack{w\in\mathcal V_0\\u_w>0}}
                     \min(1,I_w^{K_0/2}D_w^{-3}).
\]
The product of the zero-frequency masses lies between one and
the same exponential. Its implied constant is absolute.
\end{corollary}
\begin{proof}
At a split place in this set, the residue field of \(E\) has the
same size \(q_w\) as that of \(F\). It contains all \(q_1\)-th
roots of unity since reduction is injective on roots of order
prime to the residue characteristic. For example, in the kernel
of reduction the first nonzero term of \((1+h)^{q_1}-1\)
has the same valuation as \(h\). Hence \(q_w\ge q_1+1\).
We have \(f=[F:\mathbb Q]\le\varphi(q_1)\le q_1\), and at most
\(f\) places can have any prescribed prime-power residue size.
At large nodes \(q_w^{-2k+1000n^4}\le q_w^{-k}\), so
\[
 \sum_{w\in\mathcal V_0}q_w^{-2k+1000n^4}
          \le f\sum_{j>q_1} j^{-k}\le q_1^{2-k}/(k-1).
\]
Apply the theorem at the finitely many nonintegral local frequencies
and use \(1+x\le\exp x\) at the others. This also proves convergence
and the assertion at zero frequency.
\end{proof}

This supplies the standard split factors without charging a
constant at each of the infinitely many good places. Nonsplit
places have a Hermitian relation between their two arms, so the
independent middle blocks used in the shell count above require
a different treatment there.
\clearpage

\section{The first relative Gram transform at inert places}\label{nav:inert-gram-transforms}

\subsection{The first coordinate on an unramified Hermitian lattice}
\label{sec:hermitiancone900}

Let \(A/F_w\) be an unramified quadratic field extension. Write
\(O_A,O\) for the rings of integers, \(\pi\in O\) for a common
uniformizer, and \(q=\#(O/\pi)\). Let \(\mathcal L\) be a self-dual
Hermitian \(O_A\)-lattice of rank \(r\), with form \(h\) linear in
the second variable. Write \(Q(v)=h(v,v)\). In the standard model
of Section~\ref{sec:orbitfactors871}, the first flag vector \(e\)
lies in an integral hyperbolic pair. We use probability Haar measure
on \(K=\mathrm U(\mathcal L)\), and take level \(J=K\). Fix an
additive character \(\psi\) of \(F_w\) of conductor \(O\), and put
\(\psi_A=\psi\circ\operatorname{Tr}_{A/F_w}\). The trace pairing
on \(O_A\) is perfect over \(O\), since its reduction is the
nondegenerate trace pairing on a separable quadratic extension.
Thus \(\psi_A\) has conductor \(O_A\).

For \(\alpha\in A\), let \(c=\max(0,-\operatorname{ord}_A\alpha)\),
taking \(c=0\) for \(\alpha=0\). The first relative Gram entry gives
the character
\[
                    \omega_\alpha(g)=\psi_A(\alpha h(e,ge)).
\]
Both copies of \(U\) preserve this function because \(U\) fixes \(e\).
Let \(\phi_0\) be the right \(K\)-invariant function supported on
\(UK\) with \(\phi_0(uk)=\chi_0(u)\), where the local Bessel
character \(\chi_0\) is trivial on \(U\cap K\).
In \eqref{eq:gramintegrand891} take both compact integration
groups equal to \(K\), and use \(\phi_0,\omega_\alpha\).
Denote the resulting integral by \(b^0_K(x;\alpha)\).
For \(x\in K\) both factors involving \(\phi_0\) are one, so the
contribution from \(K\) to \(\int |b^0_K(x;\alpha)|\,dx\) is
\(\big|\int_K\omega_\alpha(g)\,dg\big|\).

\begin{lemma}[The compact orbit of a primitive isotropic vector]
\leavevmode\label{lem:hermitianorbit900}
For \(c\ge1\) put \(V_c=\mathcal L/\pi^c\mathcal L\).
Let \(Y_c\) consist of the vectors of \(V_c\) nonzero modulo \(\pi\)
whose norm is zero modulo \(\pi^c O\).
For Haar-uniform \(g\in K\), the image of \(ge\) is uniform on
\(Y_c\). The counting probability \(\rho_c=|Y_c|/|V_c|\) satisfies
\[
 \rho_c=q^{-c}\big(1+(q-1)(-1)^r q^{-r}-q^{1-2r}\big),
                  \qquad \rho_c\ge (1-q^{-2})q^{-c}.
\]
\end{lemma}
\begin{proof}
For primitive \(v\) the map \(u\mapsto\operatorname{Tr}h(v,u)\)
is onto modulo \(\pi\). Indeed self-duality makes \(h(v,\cdot)\)
onto \(O_A\), and the residue trace is onto. In passing from
precision \(j\) to \(j+1\) the new norm digit of \(v+\pi^j u\)
changes by this trace, for \(j\ge1\). Each primitive solution has
exactly \(q^{2r-1}\) lifts one digit higher. Successive lifts and
completeness give an exactly isotropic primitive lift in \(\mathcal L\).

Given such a vector \(v\), choose \(f\) with \(h(v,f)=1\).
Choose \(b\in O_A\) with \(b+\bar b=Q(f)\), and replace \(f\) by
\(f-bv\). The trace equation is soluble integrally because
the residue trace has a unit in its image. The new \(f\) gives
a hyperbolic partner of \(v\). The plane splits off orthogonally
since its Gram determinant is a unit. Every unimodular Hermitian
lattice over \(O_A\) has an orthonormal basis. In the residue
space some vector has nonzero norm, since otherwise applying
the norm to \(v+\lambda w\) for all residue scalars \(\lambda\)
would give \(\operatorname{Tr}(\lambda h(v,w))=0\) for all
\(\lambda,v,w\), contradicting nondegeneracy. A lift has unit
norm. The norm on \(O_A^\times\) is onto \(O^\times\): it is
onto in the cyclic multiplicative groups of the residue fields,
and a correction \(1+\pi^j t\) changes its \(j\)-th digit by
\(\operatorname{Tr}t\). Thus the lift can be scaled to norm one.
Split off that line and induct. Applying this classification to
the two hyperbolic complements gives an integral isometry taking
\(e\) to \(v\). This proves the orbit assertion, including at
places above two.

For the count modulo \(\pi\), diagonalize the lattice as just
proved. For each nontrivial additive character of the residue
field, the sum of its values on the norms in one residue
coordinate equals \(1+(q+1)(-1)=-q\), by the fiber sizes of
the norm. Character orthogonality in the norm equation therefore
gives \(q^{2r-1}+(q-1)(-1)^r q^{r-1}\) zeros. Remove the zero
vector and use the one-digit lifting count above. For the lower
bound write a vector in the fixed hyperbolic plane and its
complement as \(x e+w+y f\). Choose the residue \(x\ne0\) and
choose \(w\) freely. The equation for \(y\) is one surjective
trace equation, with \(q\) solutions. This gives at least
\((1-q^{-2})q^{-1}\) at the first digit and the asserted lower
bound at every precision.
\end{proof}

\begin{theorem}[The principal first transform at an inert place]
\leavevmode\label{thm:hermitiancone900}
If \(c=0\), the average of \(\omega_\alpha\) on \(K\) is one.
If \(c\ge1\), its exact value is
\begin{equation}
 \int_K\omega_\alpha(g)\,dg
   =\frac{(1-q^{-1})(-1)^{rc}q^{-rc}
                   -{\bf1}_{c=1}q^{-2r}}{\rho_c}.
 \label{eq:hermitiancone900}
\end{equation}
In particular the absolute value is at most
\(3q^{-(r-1)c}\). At the nodes of Theorem~\ref{thm:fullfirst898},
for nonintegral rational frequency this contribution is at most
\(I_w^{(k-n^7)/2}D_w^{-3}\).
\end{theorem}
\begin{proof}
The assertion at \(c=0\) follows from integrality. For \(c\ge1\)
write \(\gamma=\pi^c\alpha\), a unit, and take counting
probability on \(V_c\). To compute the unnormalized average
over \(Q(x)=0\bmod\pi^c\), insert its indicator
\[
          q^{-c}\sum_{s\in O/\pi^c}\psi(\pi^{-c}s Q(x))
\]
alongside \(\psi_A(\pi^{-c}\gamma h(e,x))\).
If \(s\) is a nonunit, translation by \(\pi^{c-1}u\) fixes
the quadratic phase. For \(c=1\) this is because \(s=0\)
in the quotient. The translation multiplies the linear phase by
a nontrivial character of \(u\bmod\pi\), so this average vanishes.
For a unit \(s\), put \(v=s^{-1}\bar\gamma e\). Then
\[
          sQ(x)+\operatorname{Tr}(\gamma h(e,x))=sQ(x+v)
\]
because \(Q(v)=0\). In an orthonormal basis the shifted average
is the \(r\)-th power of the one-coordinate average
\(G_c(s)=\mathbb E_z\psi(\pi^{-c}s z\bar z)\).
The residue computation gives \(G_1(s)=-q^{-1}\).
At \(c\ge2\), averaging over the last digit of \(z\) forces
\(z=0\bmod\pi\) by trace duality. Substitution of \(z=\pi u\)
then gives \(G_c(s)=q^{-2}G_{c-2}(s)\), with \(G_0(s)=1\).
Consequently every unit contributes \((-1)^{rc}q^{-rc}\).

It remains to remove the imprimitive vectors from the average.
At \(c=1\) their only contribution is the mass \(q^{-2r}\) of
zero. At \(c\ge2\) put \(x=\pi y\). The norm condition on \(y\)
has precision \(c-2\), while the linear character has conductor
exponent \(c-1\). Translation of \(y\) by \(\pi^{c-2}u\)
preserves that norm condition and gives zero upon averaging over
\(u\bmod\pi\). The unit fraction among the \(s\) is \(1-q^{-1}\).
Dividing the resulting primitive average by \(\rho_c\) proves
\eqref{eq:hermitiancone900}. The bound follows from the lower
bound for \(\rho_c\). At an inert place \(I_w=q^{-2c}\) and
\(D_w=q^{2c}\) for a local pole of order \(c\).
Since \(r-1=k+2a\), the asserted arithmetic weight absorbs the
factor \(3\) for all sufficiently large nodes.
\end{proof}

\begin{proposition}[Compression of the first Fourier operator]
\leavevmode\label{prop:inertoperator900}
Let \(K_c=\{g\in K:ge\equiv e\bmod\pi^c\}\) and \(c\ge1\).
For \(v\in Y_c\) let \(A_x(v)\) be the conditional average of
\(\phi_0(h_1xh_1^{-1})\) over \(h_1^{-1}e=v\bmod\pi^c\)
with \(h_1\) Haar-uniform in \(K\). Then for every \(x\)
\[
  |b^0_K(x;\alpha)|\le
     \frac{q^{-rc}}{\rho_c}
           \frac{1}{|Y_c|}\sum_{v\in Y_c}|A_x(v)|^2 .
\]
These conditional averages are over the left cosets of \(K_c\).
\end{proposition}
\begin{proof}
For \(v=h_1^{-1}e\) and \(w=h_2^{-1}e\) the relative pairing is
\(h(e,h_1h_2^{-1}e)=h(v,w)\). The two compact variables are
independent; conditioning gives the corresponding quadratic form
on \(Y_c\) with counting probability and kernel
\(\psi_A(\alpha h(v,w))\). On \(V_c\) with counting probability
the squared operator norm of this kernel is \(q^{-2rc}\).
Indeed summation of the product of a kernel entry with the
conjugate entry in another row gives zero unless the two row
indices agree, by trace duality and self-duality of \(\mathcal L\).
The diagonal kernel is one, and the normalized diagonal operator
is \(1/|V_c|\) times the identity. Extending functions by zero
and then restricting the output compresses the ambient operator.
Changing the measure from ambient probability to conditional
probability introduces the factor \(\rho_c^{-1}\). Applying the
operator bound to \(A_x\) proves the estimate.
\end{proof}

The proposition retains the two original Bessel coefficients at
nonintegral \(x\). Integrating its right side outside \(K\)
requires a separate estimate of their conditional mean square.
\clearpage

\subsection{The first neighbors of an unramified Hermitian lattice}
\label{sec:inertneighbor901}

Continue in the standard inert model of Section~\ref{sec:hermitiancone900}.
We write \(\mathcal L=L\oplus W\oplus L^\vee\) for the integral
decomposition, with \(L,L^\vee\) dual totally isotropic summands
of rank \(a\) and \(W\) self-dual of rank \(d=r-2a\). Denote by
\(\mathcal N\) the set of positions \(x\in\mathrm U(\mathcal L
\otimes_{O_A}A)\) with
\[
       \operatorname{length}_{O_A}
       \big((x\mathcal L+\mathcal L)/\mathcal L\big)=1.
\]
Right \(K\)-cosets in \(\mathcal N\) are identified with their
lattices \(x\mathcal L\). The function \(\phi_0\) of the preceding
section takes a particularly simple form on these cosets.

\begin{lemma}[Neighbor lattices and the unipotent orbit]
\leavevmode\label{lem:inertneighbor901}
The set of cosets in \(\mathcal N\) is one \(K\)-orbit of size
\[
      N_1=\frac{q|Y_1|}{q^2-1}\ge q^{2r-2}.
\]
Its intersection with \(UK/K\) has size
\[
                       S_1=\frac{q^{2a}-1}{q+1}.
\]
The value of \(\phi_0\) on each of these \(S_1\) cosets is one.
\end{lemma}
\begin{proof}
For \(\mathcal L'=x\mathcal L\) set
\(\ell=(\pi\mathcal L'+\pi\mathcal L)/\pi\mathcal L\).
The length assumption makes \(\ell\) a line in
\(\mathcal L/\pi\mathcal L\). It is isotropic since the norm and
all pairings on \(\mathcal L'\) are integral. The intersection
\(\mathcal M=\mathcal L'\cap\mathcal L\) is the inverse image
of \(\ell^\perp\): for \(z\in\mathcal L\), orthogonality to \(\ell\)
modulo \(\pi\) means \(h(z,\mathcal L')\subset O_A\), which
characterizes membership in the self-dual \(\mathcal L'\).
Lift \(\ell\) to an isotropic vector \(v\) with hyperbolic
partner \(f\) as in Lemma~\ref{lem:hermitianorbit900}, and
let \(W_v\) be their orthogonal complement in \(\mathcal L\).
Then
\[
 \mathcal M=O_Av\oplus W_v\oplus\pi O_A f,\qquad
 \mathcal M^\#=\pi^{-1}O_Av\oplus W_v\oplus O_A f .
\]
Here the superscript denotes the dual lattice. The form \(\pi h\)
on \(\mathcal M^\#/\mathcal M\) has hyperbolic basis
\(\pi^{-1}v,f\) over the residue field of \(A\). Self-dual lattices
between \(\mathcal M\) and \(\mathcal M^\#\) correspond exactly
to self-orthogonal lines in this plane: integrality is isotropy
and duality takes a line to its perpendicular. In addition to
the original line of \(f\), these are the \(q\) lines spanned by
\(\pi^{-1}v+t f\) with \(t+\bar t=0\bmod\pi\).
All the resulting lattices are isometric to \(\mathcal L\)
by the orthonormal basis construction in the preceding lemma.

The group \(K\) is transitive on the choices of \(\ell\). Its
stabilizer permutes the \(q\) lines just found transitively.
Indeed the isometry fixing \(f,W_v\) and taking
\(v\) to \(v+\pi u f\), with \(u+\bar u=0\), adds the residue of
\(u\) to \(t\). Every trace-zero residue has a trace-zero lift:
subtract \(\pi\) times a solution of a trace equation from
any lift of that residue. Thus there is just one orbit.
Its size follows by dividing \(|Y_1|\) by the number \(q^2-1\)
of generators of a residue line and multiplying by \(q\).
The lower bound follows from Lemma~\ref{lem:hermitianorbit900}.

It remains to determine which neighbors occur with nonzero
coefficient. Write a representative of a contributing coset
in \(U\) as
\[
 y=\begin{pmatrix}T&B&C\\0&1&D\\0&0&S\end{pmatrix},
              \qquad S=T^{-*},\qquad HD=-B^*S ,
\]
where \(H\) is the Gram matrix on \(W\) and \(^{*}\) denotes
conjugate transpose in the paired bases. The matrix
\(\pi(y-1)\) is integral and has rank one modulo \(\pi\),
since its reduced image is \(\ell\). If the unipotent \(T\)
is nonintegral, then its inverse must also be nonintegral:
an integral inverse of determinant one has an integral inverse
by the adjugate identity. Both \(T,S\) would therefore have a
pole, and the reduction of \(\pi(y-1)\) would have a nonzero
image from \(L\) in \(L\) and a nonzero projection from
\(L^\vee\) to \(L^\vee\). This contradicts rank one.
Now \(T,S\) are integral and invertible over \(O_A\). A pole
in \(B\) would give one in \(D\) by the displayed relation,
producing a nonzero image from \(W\) in \(L\) and a nonzero
projection of the image to \(W\). Thus \(B,D\) are integral
as well.

Right multiplication by \(\operatorname{diag}(T^{-1},1,S^{-1})\)
makes both diagonal arms identity while keeping the same
lattice. Use primes for the resulting off-diagonal blocks.
Choose \(\theta\in O_A\) with \(\theta+\bar\theta=1\).
The integral unipotent with blocks \(B',D'\) and
\(C_0=-\theta (D')^*HD'\) is unitary: the remaining equation
is \(C_0+C_0^*=-(D')^*HD'\). Multiplication on the right
by its inverse leaves only the upper-right block \(C'-C_0\).
Consequently the supporting right cosets are indexed by the
nonzero rank-one skew-Hermitian matrices
\(\pi(C'-C_0)\bmod\pi\) on the arms. Every such residue lifts
to an integral skew-Hermitian matrix by choosing one triangular
half arbitrarily and lifting trace-zero diagonal entries as above.
Two such representatives give the same coset exactly when
their residues agree. The Bessel character reads the successive
flag entries and the terminal entry in \(W\); all those entries
vanish on this upper-right block. Right \(K\)-invariance thus
gives value one on each coset.

For each image line choose a generator \(v_0\) in the left
arm. A rank-one skew-Hermitian matrix with that image is
\(\lambda v_0v_0^*\), since its adjoint has the same image.
It has \(q-1\) choices for the nonzero trace-zero scalar
\(\lambda\). The number of lines in the left arm is
\((q^{2a}-1)/(q^2-1)\). Their product is \(S_1\).
\end{proof}

\begin{theorem}[The first inert shell at every conductor]
\leavevmode\label{thm:inertneighbor901}
Give \(K\) volume one as above. For integral \(\alpha\), one has
\[
             \int_{\mathcal N}|b^0_K(x;\alpha)|\,dx
                        = S_1^2/N_1 < q^{-2d}.
\]
For \(c=-\operatorname{ord}_A\alpha\ge1\), one has
\[
       \int_{\mathcal N}|b^0_K(x;\alpha)|\,dx
                     \le S_1\rho_c^{-1}q^{-rc}.
\]
In particular the contribution of \(K\cup\mathcal N\) is between
\(1\) and \(1+q^{-2d}\) at zero frequency. At a local pole of a
nonzero rational frequency its contribution is at most
\(q^{-n^7c/2}I_w^{(k-n^7)/2}D_w^{-3}\) at the large nodes of
Theorem~\ref{thm:fullfirst898}.
\end{theorem}
\begin{proof}
For integral \(\alpha\) the relative character on \(K^2\) is one,
so the compact double average is the squared absolute value
of the compact single average. Right \(K\)-invariance gives
\(\phi_0(h_1xh_1^{-1})=\phi_0(h_1x)\). The lattice of \(h_1x\)
is uniform on the \(N_1\) neighbors as \(h_1\) varies. The
single average is therefore \(S_1/N_1\) for every \(x\) in this
shell. Each neighbor has volume one, proving the equality.
Now \(S_1<q^{2a-1}\) and \(N_1\ge q^{2r-2}\), which give the
stated strict inequality.

For \(c\ge1\) integrate Proposition~\ref{prop:inertoperator900}
over \(\mathcal N\). The square of a conditional average is at
most the conditional average of the squared modulus. After
averaging over \(v\in Y_c\), the integral of the latter is
\[
     \int_K\int_{\mathcal N}|\phi_0(h_1 x)|^2\,dx\,dh_1=S_1,
\]
by invariance of \(\mathcal N\) on the left and the support
calculation of the lemma. This proves the second estimate.
Adding Theorem~\ref{thm:hermitiancone900} bounds the positive
conductor contribution on the two sets together by
\((3+2S_1)q^{-(r-1)c}\le5q^{-kc}\), since \(r-1=k+2a\)
and \(c\ge1\). For large \(n\), \(5\le q^{(n^7/2-6)c}\).
Use \(I_w=q^{-2c}\) and \(D_w=q^{2c}\) to obtain the assertion.
\end{proof}

This calculation uses the original character and accounts for every
lattice of length one. At length two an off-diagonal pair \(B,D\)
can both have poles, so the support on the middle block acquires
additional degrees of freedom.
\clearpage

\subsection{A lattice length bound for distant inert positions}
\label{sec:inertlength902}

Use the standard inert model and the measure conventions of
Section~\ref{sec:inertneighbor901}. For \(x\) in the unitary group
put \(P=\operatorname{length}_{O_A}((x\mathcal L+\mathcal L)/
\mathcal L)\) and let \(m\) be the least nonnegative integer
for which \(\pi^m x\) and \(\pi^m x^{-1}\) are integral.
Both invariants depend only on the double coset \(KxK\).
Write \(\mathcal S(m,P)\) for the positions with these invariants.
We give two counts, one for the supporting unipotents and one for
the size of their compact orbits.

For \(y\in U\) write \(T,B,D,S,C\) for its blocks in the notation
of Lemma~\ref{lem:inertneighbor901}. Define lattices in the left arm
and integers by
\[
  L'=L+T L,\qquad J_B=L'+BW,\qquad h_B=\operatorname{length}(J_B/L),
             \qquad \delta(y)=P-2h_B .
\]
For a right coset \(yK\) with \(y\in U\), these are independent
of the unipotent representative. Indeed two representatives
differ on the right by \(U\cap K\); this changes \(T\) to \(TT_0\)
for \(T_0\) integral upper unitriangular and changes \(B\) by
\(TB_0\) with \(B_0\) integral. Both lattices stay the same.

\begin{lemma}[Lengths of the two middle maps]
\leavevmode\label{lem:inertmiddle902}
One has \(2h_B\le P\). For \(m\ge1\) and an integer \(b\ge0\)
the number of supporting right cosets in \(\mathcal S(m,P)\) with
\(\delta(y)\ge b\) is at most
\begin{equation}
                         q^{5a^2m+d(P-b)}.
 \label{eq:inertmiddle902}
\end{equation}
Here and below a supporting coset is a coset on which \(\phi_0\)
is nonzero.
\end{lemma}
\begin{proof}
Let \(\tau=\operatorname{length}(L'/L)\) and
\(\beta=\operatorname{length}(J_B/L')\). The kernel of the
projection of \((y\mathcal L+\mathcal L)/\mathcal L\) to the
last two blocks contains \(J_B/L\). In the projected module
the last block has length \(\tau\), because \(S=T^{-*}\)
and \(\det T=1\). More explicitly, the positive and negative
parts of the elementary divisors of \(T\) have equal total
length, and dual inversion interchanges them. The kernel of
this last projection is the image of \(Du\) modulo \(W\)
for \(u\in L^\vee\) with \(Su\in L^\vee\). Since
\[
 S L^\vee\cap L^\vee=(T L+L)^\#,\qquad HD=-B^*S,
\]
its length is \(\beta\). Indeed the matrix of \(B\) in bases of
\(W,L'\) and its adjoint in the dual bases have identical Smith
exponents, and the sums of their pole orders compute the two
lengths. Adding the lengths gives
\(P\ge h_B+\tau+\beta=2h_B\).

For the count use product Haar measure in the entries of \(T\)
above its diagonal, all entries of \(B\), and a skew-Hermitian
matrix \(Z\) of size \(a\). To see the normalization, write
\(y=u_0\,\operatorname{diag}(T,1,S)\). The middle-to-dual block
of \(u_0\) is \(D_0=-H^{-1}B^*\), and its upper-right block is
\(Z-\theta D_0^*HD_0\), with \(\theta+\bar\theta=1\).
The product of two such \(u_0\)'s adds their \(B\)'s and adds
their \(Z\)'s with a translation depending on the \(B\)'s.
Conjugation by the diagonal sends \(B\) to \(TB\) and \(Z\)
to \(TZT^*\); both linear maps are unipotent and preserve volume.
For \(T\) multiplication is triangular in its entries with
diagonal one. Thus the product measure is Haar, and the integral
parameter lattices describe \(U\cap K\) exactly, of volume one.

The entries of \(T\) have poles of order at most \(m\), for a
cost at most \(q^{a(a-1)m}\). At fixed \(T\) the lattice \(J_B\)
lies between \(L\) and \(\pi^{-m}L\). The number of possibilities
is at most \(q^{2a^2m}\): each submodule of \(\pi^{-m}L/L\) is
generated by at most \(a\) elements, by taking a basis of its
inverse image. With \(J_B\) fixed, the \(d\) columns of \(B\)
are in \(J_B\), costing at most \(q^{2dh_B}\).
Finally \(CS^{-1}=Z-\theta D_0^*HD_0\) has entries in
\(\pi^{-2m}O_A\) since both \(C\) and \(S^{-1}\) have poles
of order at most \(m\). At fixed \(T,B\) any two possible \(Z\)'s
differ by a skew-Hermitian matrix in that ball. This costs
\(q^{2a^2m}\), as the integral skew-Hermitian matrices are a
full \(O\)-lattice of rank \(a^2\). On the specified set,
\(2h_B\le P-b\). Multiplication of the four factors proves
\eqref{eq:inertmiddle902}; Haar volume counts right cosets
because each intersection with \(U\) has volume one.
\end{proof}

\begin{lemma}[A compact orbit bound from a primitive frame]
\leavevmode\label{lem:inertframe902}
On any \(K\)-orbit of right cosets in \(\mathcal S(m,P)\)
meeting \(UK/K\), let \(N\) be the orbit size. For \(m\ge1\),
\[
                           N\ge \tfrac13 q^{2dP-4a^2m}.
\]
\end{lemma}
\begin{proof}
Write \(A_x=x\mathcal L+\mathcal L\).
Its elementary divisors give \(s\) independent primitive vectors
\(v_1,\ldots,v_s\), extendible to an integral basis, and positive
integers \(\lambda_i\le m\) with \(\sum_i\lambda_i=P\), such that
\(A_x=\mathcal L+\sum_i O_A\pi^{-\lambda_i}v_i\).
For an orbit meeting \(U\) one has \(s\le2a\).
Indeed \(A_y/\mathcal L\) is generated by the image of \(y-1\)
modulo \(\mathcal L\), and \(y-1\) has rank at most \(2a\):
the images of the left and middle blocks are in \(L\otimes A\)
and the remaining block has \(a\) columns. A rank-\(2a\)
lattice image has a basis with at most \(2a\) elements.

We record the compact counting argument for arbitrary such
frames, allowing degeneracy of their Gram matrix. In the
residue field \(|\mathrm U_j(q)|=q^{j^2}\prod_{i=1}^j
(1-(-1)^i q^{-i})\). Indeed the same additive-character count
as in Lemma~\ref{lem:hermitianorbit900} gives
\(q^{2j-1}-(-1)^j q^{j-1}\) vectors of norm one. Choosing
successively an orthonormal basis proves the formula. The
product is between \(2/3\) and \(2\); use
\(\sum_{i\ge1}q^{-2i}\le1/3\) for the lower bound and
\(\exp(\sum_{i\ge1}q^{-(2i-1)})<2\) for the upper bound.

Modulo \(\pi\), let the radical of the span of the \(s\)
frame vectors have dimension \(j\). A nondegenerate complement
in their span of dimension \(s-j\) is fixed pointwise, so one
can work in its perpendicular, with \(j\) fixed isotropic
vectors. Complete these to hyperbolic pairs. Concretely choose
dual partners with pairings the identity and subtract suitable
linear combinations of the fixed vectors to set the partners'
Gram matrix to zero. The required equation \(C+C^*=D\) for a
Hermitian \(D\) is soluble by choosing a triangular half and
solving the diagonal trace equations. In the resulting three
blocks of sizes \(j,r-s-j,j\), the isometries fixing the first
\(j\) vectors pointwise have identity diagonal arms, an arbitrary
unitary middle block, \(j(r-s-j)\) free entries over the
quadratic residue field, and \(j^2\) free scalar digits in the
skew-Hermitian upper-right part. These choices follow by first
pairing the middle and last blocks to determine the opposite
rectangle, then pairing the last block with itself. Hence the
residue pointwise stabilizer has size at most
\(2q^{(r-s-j)^2+2j(r-s-j)+j^2}=2q^{(r-s)^2}\).

An isometry modulo \(\pi^i\) lifts to one modulo \(\pi^{i+1}\),
in exactly \(q^{r^2}\) ways. Correct any lift by \(1+\pi^iX\);
the required correction in the Gram matrix is \(X^*H'+H'X\)
modulo \(\pi\), with \(H'\) the invertible Gram matrix of the
full lattice. This map is onto the Hermitian matrices by the
same trace equations and has kernel of dimension \(r^2\)
over the residue field of \(F_w\). Completeness lifts each
finite solution to \(K\). For lifts fixing the frame, two
choices differ by such a skew-adjoint \(X\) annihilating its
\(s\) independent residue vectors. The corresponding
skew-Hermitian form descends in both entries to the quotient
by their span, allowing at most \(q^{(r-s)^2}\) choices.
Orbit--stabilizer at precision \(m\) now gives at least
\(q^{(2rs-s^2)m}/3\) different images of the ordered frame.

Every \(g\in K\) fixing \(x\mathcal L\) also fixes \(A_x\).
In the chosen basis this forces the last \(r-s\) entries of
\(g v_i\) to be divisible by \(\pi^{\lambda_i}\), because
\(\pi^{-\lambda_i} g v_i\in A_x\). At most
\(q^{2rsm-2(r-s)P}\) matrices of frame images modulo \(\pi^m\)
meet those conditions, even when all other requirements are
dropped. The images are uniform on their orbit. The probability
of fixing \(x\mathcal L\) is thus at most
\(3q^{s^2m-2(r-s)P}\). Its reciprocal is \(N\). Use
\(s\le2a\) and \(r-s\ge d\) to conclude.
\end{proof}

\begin{theorem}[Unweighted inert shell energy]
\leavevmode\label{thm:inertlength902}
For \(m\ge1\), write \(\phi_b\) for \(\phi_0\) restricted to
supporting cosets with \(\delta(y)\ge b\), and give \(K\) volume
one. Then
\[
 \int_{\mathcal S(m,P)}
       \left|\int_K\phi_b(gx)\,dg\right|^2 dx
                       \le 3q^{14a^2m-2db}.
\]
For the original zero-frequency integral on an entire shell
one can take \(b=P\bmod 2\).
\end{theorem}
\begin{proof}
On an orbit of \(N\) right cosets the compact average is the
signed sum of the values on the orbit divided by \(N\); each
right coset has volume one. Its energy is the squared modulus
of the sum divided by \(N\). By the triangle inequality, the
sum of these squared moduli is bounded by the square of the
total number counted in Lemma~\ref{lem:inertmiddle902}.
Division by the lower bound of Lemma~\ref{lem:inertframe902}
proves the estimate. Since \(\delta=P-2h_B\) is nonnegative
and has the parity of \(P\), taking \(b=P\bmod 2\) includes
all supporting cosets. At zero frequency the square of the
compact average is exactly \(b^0_K(x;0)\).
\end{proof}

\begin{corollary}[A summable contribution at all heights]
\leavevmode\label{cor:inertsum902}
Restrict the original \(\phi_0\) outside \(K\) to the cosets
with \(\delta(y)\ge\lceil m/(100a^2)\rceil\). Call this
restriction \(f\), extended by zero elsewhere. At the large
nodes the full unweighted compact average satisfies
\[
         \int_G\left|\int_K f(gx)\,dg\right|^2dx
                       \le q^{-2d+1500a^4}.
\]
\end{corollary}
\begin{proof}
Put \(j=\lceil m/(100a^2)\rceil\ge1\).
Since \(P=\sum\lambda_i\le2am\), for each \(j\) there are at
most \((100a^2 j)(200a^3 j+1)\) possible pairs \(m,P\).
The theorem bounds each by \(3q^{-(2d-1400a^4)j}\).
Here \(d/a^4\to\infty\) along the specified nodes. The sum of
\(j^2 q^{-(2d-1400a^4)(j-1)}\) is bounded by an absolute constant
there. The remaining polynomial factor in \(a\) is absorbed by
\(q^{100a^4}\).
\end{proof}

The corollary concerns the compact average of the stated
restriction of the Bessel coefficient. Its interference with the
complement in the original square depends on the energy of that
complement. The two lengths can be equal, \(P=2h_B\), so the
length inequality alone supplies no decay in that case.
\clearpage

\subsection{Terminal row cancellation on inert lattice orbits}
\label{sec:inertterminal903}

Assume that the terminal vector \(z\) is primitive in the
self-dual middle lattice \(W\) and that its simple coefficient
is a unit. In the arm coordinates of Section~\ref{sec:inertlength902}
the terminal entry of the simple character is
\(\psi_A(\gamma\beta_a(z))\), where \(\beta_a\) is the last row
of \(B\) and \(\gamma\in O_A^\times\).
Let \(p=p(y)\) be the nonnegative integer for which the last
coordinate projection of \(J_B\) is \(\pi^{-p}O_A\).
The last row of \(T\) is the last coordinate row, so this
projection is \(O_A+\beta_a(W)\). In particular \(p\le m\).
We preserve the signs of the original Bessel character below.

Choose a totally isotropic direct summand \(Q\subset W\) of
rank \(g=\lfloor d/2\rfloor\) such that \(h(Q,z)=O_A\).
Here is an integral construction. In a Hermitian residue space
of dimension at least two, two norm-one orthogonal coordinates
give an isotropic vector by choosing a scalar of norm \(-1\).
Split off a hyperbolic plane and repeat. The isotropic vectors
span the residue space: the hyperbolic basis vectors span the
planes, and for any vector \(u\) in the remaining line, one can
make \(e'+u+t f'\) isotropic by a trace equation in a hyperbolic
pair \(e',f'\). Hence there is an isotropic vector pairing
nontrivially with the reduction of \(z\). Lift this vector and
split off its integral plane by Lemma~\ref{lem:hermitianorbit900}.
Continue in the complementary lattice to obtain \(Q\).
Put
\[
                         \eta=\min(g-a,d-g-1).
\]
This is positive at all sufficiently large nodes.

\begin{theorem}[Terminal cancellation in the signed class sums]
\leavevmode\label{thm:inertterminal903}
In \(\mathcal S(m,P)\), \(m\ge1\), restrict the supporting
cosets to \(\delta(y)\ge b,\ p(y)\ge s\), for integers \(b,s\ge0\).
Let \(\mathcal C\) run over the orbits on the right \(K\)-cosets
under \(K\), or under the subgroup \(K_c\) fixing the first flag
vector modulo \(\pi^c\), for any \(c\ge1\).
Write \(S_{\mathcal C}\) for the sum of \(\phi_0\) over the
restricted cosets in the orbit. Then
\begin{equation}
      \sum_{\mathcal C}|S_{\mathcal C}|
           \le 2(m+1)q^{5a^2m+d(P-b)-2\eta s}.
 \label{eq:inertterminal903}
\end{equation}
One may prescribe exact values of \(\delta\) and \(p\) in this
estimate.
\end{theorem}
\begin{proof}
We first give the change of variables before any absolute values
are taken. For \(v\) primitive in \(Q\cap\ker B\), write
\(v^\flat=h(v,\cdot)\). For \(w\in W\cap v^\perp\), define the
integral isometry
\[
       N_v(w)=1+w v^\flat-vw^\flat-\theta Q_W(w)v v^\flat,
       \qquad Q_W(w)=h(w,w),\quad \theta+\bar\theta=1 .
\]
Multiplying by the adjoint gives the identity, using \(Q_W(v)=0\)
and \(v^\flat(w)=0\). It fixes both arms. Since \(yv=v\) and
\(yw=w+b_w\) with \(b_w=Bw\in L\), conjugation by \(y\) sends
\(N_v(w)\) to \(N_v(w+b_w)\). The vectors \(w,b_w\) are
orthogonal and \(b_w\) is isotropic. The two translations along
them therefore commute and their product is \(N_v(w+b_w)\),
by multiplication of the displayed matrices. Consequently
\begin{equation}
     N_v(w)^{-1}yN_v(w)=N_v(b_w)y,\qquad
                 N_v(b_w)=1+b_w v^\flat-v b_w^\flat .
 \label{eq:inertconj903}
\end{equation}
The last factor on the left of \(y\) on the right-hand side
belongs to \(U\). Its entries on the simple arm roots are zero.
On the terminal root its contribution to the phase is
\(\gamma\beta_a(w)h(v,z)\). The conjugation preserves the
individual orbit \(\mathcal C\) even for \(K_c\), since the
conjugating isometry fixes the first flag vector.

We detail the measure argument when the vector \(v\) depends on
the data. Fix \(T\) and \(B|_Q\). Let \(\tau\) be the valuation
of \(h(Q\cap\ker B,z)\), with value \(+\infty\) for the zero
ideal. If \(\tau<p\), choose \(v\) in this kernel attaining
\(\tau\); it is primitive, since dividing a nonprimitive vector
by \(\pi\) would decrease the valuation. The choice can be made
measurably using a kernel basis obtained by elimination and the
first basis vector of minimum valuation. For fixed \(v,w\),
the new \(B\) equals \(B(1+w v^\flat|_W)\).
This integral change of basis has determinant one and is the
identity on \(Q\), so it preserves \(B|_Q,J_B\) and the product
measure in the other columns. In this formula \(T,S\) and
the upper-right block \(C\) stay fixed: the possible extra
term \(b_w v^\flat D\) is zero by \(HD=-B^*S\) and \(Bv=0\).
Moreover \(BH^{-1}B^*\) stays fixed on substituting the new
\(B\), again because \(Bv=0\) and \(v\) is isotropic. Thus the
skew-Hermitian parameter \(Z\) of Lemma~\ref{lem:inertmiddle902}
stays fixed as well. The inverse change in \(B\) uses
\(1-w v^\flat|_W\). This proves conditional Haar invariance
and preservation of every restriction in the statement.

As in Lemma~\ref{lem:inertmiddle902} a signed coset sum equals
the integral over the corresponding part of \(U\). On each
piece of fixed \(J_B,T,B|_Q\) with \(\tau<p\), average
\eqref{eq:inertconj903} over the compact lattice \(W\cap v^\perp\).
Orthogonality of its additive characters forces
\begin{equation}
               \beta_a(W\cap v^\perp)\subseteq\pi^{-\tau}O_A .
 \label{eq:inertrowtest903}
\end{equation}
Indeed the image is an \(O_A\)-module and \(\psi_A\) has
conductor \(O_A\). For \(\tau\ge p\) retain the whole piece.
The triangle inequality may now be applied separately in each
orbit and the surviving volumes added.

For this last count fix a candidate \(J=J_B\) and put
\(h_B=\operatorname{length}(J/L)\). Enlarge the column domain
to independent columns uniform in \(J\), of total measure
\(q^{2dh_B}\), retaining their constraints just obtained.
The number \(p\) is fixed by \(J\).
The law of the map \(B|_Q\) is invariant under
\(\mathrm{GL}_g(O_A)\) on its domain. Its kernel there is a
primitive submodule of rank at least \(g-a\). For \(j\ge1\),
\[
              \Pr(\tau\ge j)\le 2q^{-2(g-a)j}.
\]
In fact one may first fix the map and then precompose by an
independent uniform element of \(\mathrm{GL}_g(O_A)\). In a
basis extending the kernel, the coordinates of the resulting
primitive functional \(h(\cdot,z)|_Q\) have the distribution
of a uniform primitive row (with conjugation of coefficients).
The probability of \(g-a\) prescribed zero coordinates
modulo \(\pi^j\) is at most
\(q^{-2(g-a)j}/(1-q^{-2g})\). The bound also holds for \(j=0\).

Conditional on \(B|_Q\), complete \(Q\) by a fixed integral
complement of rank \(d-g\). Within it, \(v^\perp\) contains a
primitive summand of rank at least \(d-g-1\). Extend a basis
of that summand to a basis of the complement. In these columns
of \(B\) the last coordinates are still independent and uniform
in the projection \(\pi^{-p}O_A\) of \(J\). If \(\tau<p\),
condition~\eqref{eq:inertrowtest903} costs at most
\(q^{-2(d-g-1)(p-\tau)}\) of this probability. On the piece
\(\min(\tau,p)=j\), multiply this bound (one when \(j=p\)) by
the bound for \(\Pr(\tau\ge j)\). The product is at most
\(2q^{-2\eta p}\). Summing \(0\le j\le p\le m\) gives the
factor \(2(m+1)q^{-2\eta s}\).
The counts of \(T,J,Z\) in Lemma~\ref{lem:inertmiddle902} apply
unchanged, with \(2h_B\le P-b\). This proves the theorem.
\end{proof}

\begin{corollary}[An inert contribution controlled by the terminal row]
\leavevmode\label{cor:inertterminal903}
Outside \(K\) restrict \(\phi_0\) to supporting cosets satisfying
\(\delta(y)+p(y)\ge\lceil m/(100a^2)\rceil\), and extend this
restriction \(f\) by zero. Then at sufficiently large nodes
\[
           \int_G\left|\int_K f(gx)\,dg\right|^2dx
                              \le q^{-2d+1600(a+1)^4}.
\]
\end{corollary}
\begin{proof}
On a shell put \(j=\lceil m/(100a^2)\rceil\). Decompose each
orbit sum by the exact pair \(\delta,p\). For each admissible
pair, the theorem bounds the total absolute sum over orbits
by \(2(m+1)q^{5a^2m+dP-\min(d,2\eta)j}\).
There are at most \((P+1)(m+1)\) pairs. Bound the sum of
squared orbit sums by the square of their combined absolute
sum, then divide by the orbit lower bound from
Lemma~\ref{lem:inertframe902}. The shell energy is at most
\[
      12(m+1)^4(P+1)^2 q^{14a^2m-2\min(d,2\eta)j}.
\]
Since \(2\eta\ge d-2a-2\), \(m\le100a^2 j\) and
\(P\le2am\), the exponential factor is at most
\(q^{-(2d-1400a^4-4a-4)j}\).
The number of shells at each \(j\) is bounded as in
Corollary~\ref{cor:inertsum902}. All the other factors are
bounded by a fixed polynomial in \(a,j\). Summing the geometric
series times that polynomial gives the claimed bound, since
\(d/a^4\to\infty\) and \(q\ge2\).
\end{proof}

In particular this estimate includes cosets with \(\delta=0\)
when their terminal row has a sufficiently large pole. For
\(\delta=p=0\) the other arm rows can still have poles.
\clearpage

\subsection{Distant inert shells at every conductor}
\label{sec:inerttail904}

Retain the assumptions of Section~\ref{sec:inertterminal903}.
For an integral lattice \(J\supset L\) in \(L\otimes A\), write
\(b_i(J)\) for the integer such that the \(i\)-th coordinate of
\(J\cap (L_i\otimes A)\) ranges over \(\pi^{-b_i(J)}O_A\);
here \(L_i=\langle e_1,\ldots,e_i\rangle\).
Successive projections give \(\sum_i b_i(J)=\operatorname{length}(J/L)\).
We use \(J=J_B\), so \(b_a(J)=p\). To detect a pole situated
earlier in the flag, one can average over an isotropic vector
dual to the next row.

\begin{lemma}[An interior row on the compact orbits]
\leavevmode\label{lem:inertdrop904}
Fix \(1\le j<a\), a shell \(\mathcal S(m,P)\), \(m\ge1\),
and a candidate lattice \(J=J_B\). Write \(b_i=b_i(J)\),
\(\delta=P-2\operatorname{length}(J/L)\), and
\(D_j=(b_j-b_{j+1}-\delta-1)_+\), where \(s_+=\max(s,0)\).
On the cosets with this \(J\), the sum of the absolute values
of the signed orbit sums, for either group in
Theorem~\ref{thm:inertterminal903}, is at most
\[
              q^{C(a+1)^2(m+1)+d(P-\delta)-2(d-a)D_j}
\]
with an absolute \(C\). The simple character on the \(j\)-th
arm root is assumed to have a unit coefficient as at a standard
unmarked place.
\end{lemma}
\begin{proof}
Put \(A_y=y\mathcal L+\mathcal L\),
\(M_y=A_y^\#=\mathcal L\cap y\mathcal L\), and
\(J^+=A_y\cap (L\otimes A)\). The projection of
\(A_y/\mathcal L\) to the last two blocks has length \(h_B\):
its last projection has length \(\tau\) and the kernel has
length \(\beta\) in the proof of Lemma~\ref{lem:inertmiddle902}.
Its kernel before projection to those two blocks is \(J^+/L\).
Hence \(J\subseteq J^+\) and \(\operatorname{length}(J^+/J)=\delta\).

Set \(V=L_j^\perp\cap\mathcal L\).
Among the \(v\in M_y\cap V\) the least valuation of \(h(v,e_{j+1})\)
is \(t=b_{j+1}(J^+)\). In fact restriction from the dual of
\(A_y\) onto the dual of \(J^+\) is surjective because \(J^+\)
is saturated in \(A_y\). In the dual of \(J^+\) impose vanishing
on \(L_j\). The image on \(e_{j+1}\) is \(\pi^{b_{j+1}(J^+)}O_A\):
apply duality to the rank-one successive quotient of the
intersections with the flag, which is a direct summand in the
quotient by the preceding intersection. This also shows
\[
       0\le t\le m,\qquad b_{j+1}\le t\le b_{j+1}+\delta .
\]
Indeed \(\pi^{-m}L\supseteq J^+\), and the increases in all
the \(b_i\) under \(J\subseteq J^+\) are nonnegative and sum
to \(\delta\).

Give the smooth integral isotropic vectors in \(V\) their
invariant probability measure \(\mu\). More explicitly, keep
the norm-zero residues whose projection off the radical
\(L_j\bmod\pi\) is nonzero. At each higher precision condition
counting probability on these smooth solutions. The trace
derivative is surjective there, so all solutions have the
same number of lifts per digit. These finite distributions
give \(\mu\) on their inverse limit. Every integral isometry
preserving \(V\) and \(L_j\) preserves this measure. Let
\[
    H_y=\mu\{v\in M_y\cap V:
                    \operatorname{ord}_A h(v,e_{j+1})=t\}.
\]
Uniformly for the specified \(y\),
\begin{equation}
                   H_y\ge q^{-10(a+1)(m+1)} .
 \label{eq:inertvector904}
\end{equation}
We give the argument since a cost with \(d\) in this exponent
would spoil the shell estimate. The saturation of the adjoint
image of \(B\) in \(W\) has rank at most \(a\). Its orthogonal
complement in \(W\) reduces onto the perpendicular of its
residue image: the saturation is primitive, so restriction of
integral linear functionals to it is onto. In that residue
perpendicular, the radical has dimension at most \(a\).
Lift a basis of a nondegenerate complement, obtaining a
unimodular summand \(W_0\) of rank at least \(d-2a\).
The unipotent \(y\) fixes \(W_0\) pointwise because \(B W_0=0\),
and preserves its orthogonal complement. Thus all integrality
conditions on \(v\in M_y\) are confined to the orthogonal
complement of \(W_0\), of rank at most \(4a\).
In this complement fix the coordinates of \(v\), modulo
\(\pi^{m+1}\), equal to those of a vector attaining \(t\).
Take its component in \(V\cap W_0^\perp\), since projection
off \(W_0\) changes neither stability nor the leading coordinate.
This costs at most \(q^{8a(m+1)}\) in reciprocal probability.
For any choices of those coordinates at precision \(l\ge m+1\),
the norm equation in \(W_0\) has smooth solutions of probability
at least \((1-q^{-2})q^{-l}\): in an integral hyperbolic plane
take the first coefficient a unit and solve the resulting
surjective trace equation in its partner. The complementary
norm can have any prescribed value. Such a plane exists by the
construction preceding Theorem~\ref{thm:inertterminal903}.
The ambient probability of a smooth isotropic solution in \(V\)
is at most \(q^{1-l}\) by one-digit lifting from precision one.
Taking the ratio proves \eqref{eq:inertvector904}.

In a signed class sum insert
\[
            1=H_y^{-1}\int {\bf1}_{v\in M_y,\,
                   \operatorname{ord}h(v,e_{j+1})=t}\,d\mu(v).
\]
Partition first by \(t\) and hold the integration vector \(v\)
fixed. Put \(\kappa=h(v,\cdot)\), and let \(w\) range over
\(\pi(W\cap v^\perp\cap\ker B_{>j})\), where \(B_{>j}\) denotes
the rows after \(j\). Then \(b=Bw\in L_j\otimes A\), \(v\)
and \(y^{-1}v\) are integral isotropic vectors, and both are
orthogonal to \(w\), as \(h(v,yw)=h(v,w+b)=0\).
The translations defined in \eqref{eq:inertconj903} obey
\begin{equation}
       N_v(w)^{-1}yN_{y^{-1}v}(w)=N_v(b)y,\qquad
                     N_v(b)=1+b\kappa-vb^\flat .
 \label{eq:inertinner904}
\end{equation}
To verify it conjugate the right translation by \(y\), obtaining
\(N_v(yw)=N_v(w+b)=N_v(w)N_v(b)\), as \(w,b\) are orthogonal.
Here \(N_v(b)\in U\): \(b\) lies in the first \(j\) flag lines,
and \(v\) is orthogonal to them. On the left arm this adds only
earlier rows to later input coordinates, the dual action is the
inverse adjoint, and the change on the middle quotient is zero.
The only simple entry of \(N_v(b)\) is the \(j\)-th one,
\(\beta_j(w)\kappa(e_{j+1})\), since \(b\) has no coordinates
after \(j\) and \(\kappa\) vanishes on \(L_j\). The terminal
entry is zero since \(j<a\).

Write \(y'\) for the element on either side of
\eqref{eq:inertinner904}. The two exterior factors on the left
are integral isometries, and the left one fixes \(e_1\).
Thus \(y'\) belongs to the same refined orbit as \(y\).
The middle block with image in \(L\) changes from \(B\) to
\(B(1+w(\kappa y)|_W)\); this right multiplier is integral
with integral inverse since \(\kappa y\) is integral and
\(\kappa y(w)=0\). The left arm \(T\) changes by
\(b(\kappa y)|_L\). It follows in both directions that \(J\)
is unchanged, since \(b\in BW\) and the latter module is unchanged.
Also \(M_{y'}=N_v(w)^{-1}M_y\). This integral map fixes \(L_j\)
and preserves \(V\). For any \(v'\in V\) its change in the
\(h(v',e_{j+1})\) coordinate is divisible by \(\pi^{t+1}\):
the \(w\)-term has no dual-arm coordinate and the coefficient
of \(v\) in the change is divisible by \(\pi\), since
\(w\in\pi W\). Therefore \(t\) is still the least valuation
on the new stable lattice and \(H_{y'}=H_y\), by invariance of
\(\mu\). The chosen \(v\) is fixed by the left multiplier.

For clarity, this substitution preserves Haar measure also when
the allowed \(w\)'s depend on the lower rows of \(B\).
At fixed \(v\) use the closed additive subgroup of \(U\)
consisting of the \(N_v(u)=1+u\kappa-vu^\flat\) for
\(u\in L_j\otimes A\). Its parameters add since \(v,L_j\) are
mutually orthogonal and isotropic. On each left coset of this
subgroup the lower rows \(B_{>j}\) are fixed, by restriction
to \(W\). For any allowed \(w\), the value of \(Bw\) is also
fixed there, since \(N_v(u)\) fixes \(w+Bw\). Thus on the coset
the substitution is a translation by a fixed group element.
Disintegrating Haar measure along these cosets proves the
invariance. Equivalently, fix a representative in each coset
and use translation invariance in the additive coordinate \(u\);
both groups are unimodular. Its inverse uses \(-w\).
The orbit, integrality restrictions and the factor \(H_y^{-1}\)
have all been preserved. Averaging over the allowed \(w\)'s
with probability Haar now imposes
\[
        \beta_j(W\cap v^\perp\cap\ker B_{>j})
                                      \subseteq\pi^{-t-1}O_A
\]
by the unit simple coefficient and the conductor of \(\psi_A\).

Apply absolute values after this averaging. Bound \(H_y^{-1}\)
by \eqref{eq:inertvector904}, and use the column and central
parameter count of Lemma~\ref{lem:inertmiddle902}. At fixed
\(v,J\) enlarge the column domain to independent uniform
columns in \(J\), and condition on their coordinates after row
\(j\). Then their \(j\)-th coordinates are independently uniform
in cosets of \(\pi^{-b_j}O_A\). Indeed the kernel of that
projection on \(J\) is \(J\cap(L_j\otimes A)\). The tested
summand of \(W\) is primitive of rank at least \(d-a\), being
the kernel of at most \(a-j+1\) functionals. After extending
a basis of it to \(W\), the same independent Haar laws, up to
translations, hold in the new columns. The fraction surviving
the test is at most \(q^{-2(d-a)(b_j-t-1)_+}\), including when
some of the translated balls miss the test altogether.
Integrating over \(v\) costs at most one and the partition by
\(t\) at most \(m+1\). The volume of all columns is
\(q^{2d\operatorname{length}(J/L)}\); the remaining counts of
\(T,Z\) and the reciprocal probability in
\eqref{eq:inertvector904} are absorbed by
\(q^{C(a+1)^2(m+1)}\). Using \(t\le b_{j+1}+\delta\) proves
the lemma.
\end{proof}

\begin{theorem}[The distant part of the full inert transform]
\leavevmode\label{thm:inerttail904}
At the standard unmarked inert places above, let
\(c=\max(0,-\operatorname{ord}_A\alpha)\), with \(c=0\) at
\(\alpha=0\). For every sufficiently large node, uniformly in
the residue field and the conductor,
\[
    \int_{P(x)\ge1000(a+1)^2(c+1)}
                   |b^0_K(x;\alpha)|\,dx
                         \le q^{-6k(c+1)}.
\]
\end{theorem}
\begin{proof}
For fixed \(J\) define the integer
\(D=\max(\{\delta,b_a\}\cup\{D_j:1\le j<a\})\).
The inequalities \(b_j\le b_{j+1}+2D+1\) and \(b_a\le D\)
give \(P=\delta+2\sum b_i\le(4a^2+1)D+2a^2\).
For \(P\ge4a^2+2\) this implies
\(D\ge P/(10(a+1)^2)\).

For each \(J\) use the estimate corresponding to a maximum.
For \(\delta\) the unsigned count in
Lemma~\ref{lem:inertmiddle902} saves \(d\delta\) in the
exponent relative to \(dP\). For \(b_a\) the proof of
Theorem~\ref{thm:inertterminal903} at this fixed \(J\)
saves \(2\eta b_a\); for a \(D_j\) use the preceding lemma.
In the last two cases the count also saves \(d\delta\), which
we can discard. Each choice saves at least \(dD/2\) at large
nodes. There are at most \(q^{2a^2m}\) possible \(J\)'s by
Lemma~\ref{lem:inertmiddle902}. Consequently on the full shell
the sum of absolute signed class sums on the \(K_c\)-orbits
(with \(K_0=K\) here) is at most
\[
 q^{C_1(a+1)^2(m+1)+dP-dP/(20(a+1)^2)} .
\]
Here \(C_1\) is absolute.

For the final integration the operator of the relative character
on the compact orbit has norm at most one, by the modulus of
its kernel and Cauchy--Schwarz. The shell integral is therefore
bounded by the norm squared of the \(K_c\)-average of \(\phi_0\)
on that shell. More explicitly, the conditional averages of
Proposition~\ref{prop:inertoperator900} are left translates of
this average in the \(x\)-variable, whose squared integrals
agree on the \(K\)-invariant shell. At \(c=0\) this is the
ordinary compact average. For \(c\ge1\), \([K:K_c]=|Y_c|
\le q^{2rc}\); the same bound one applies for \(c=0\).
A refined orbit thus has size at least
\(q^{2dP-4a^2m-2rc}/3\), by Lemma~\ref{lem:inertframe902}
and orbit--stabilizer. Squaring the signed bound and dividing
by this size gives, after increasing the absolute constant,
\[
               q^{2rc+C_2(a+1)^2(m+1)-dP/(10(a+1)^2)} .
\]
We have \(m\le P\): the elementary divisors of a self-dual
lattice relative to \(\mathcal L\) occur in opposite pairs,
since duality changes their signs. Thus a pole of largest
order occurs in the length sum. Since \(d/(a+1)^4\to\infty\),
the displayed bound is at most \(q^{2rc-dP/(20(a+1)^2)}\)
at the stated large nodes. There are at most \(P\) choices
of \(m\) at fixed \(P\). For \(P\ge1000(a+1)^2(c+1)\),
these factors, summed against the decreasing geometric series,
give at most \(q^{-6k(c+1)}\): indeed \(r/d\to1\),
\(d=k+1\), and the first negative exponent before summation
is at least \(50d(c+1)\) in magnitude. The polynomial factor
is absorbed uniformly using \(\log_q P\le P\) and
\(d/(a+1)^2\to\infty\). This proves the assertion.
\end{proof}

The bound applies to the full Bessel coefficients, with both
compact integrations taken before the absolute value. Together
with Theorem~\ref{thm:inertneighbor901} it leaves a finite
range of intermediate lattice lengths at each conductor.
\clearpage

\subsection{The full unweighted mass at a standard inert place}
\label{sec:inertzero905}

For zero frequency Theorem~\ref{thm:inerttail904} leaves the
lengths \(P<1000(a+1)^2\). At these lengths a single
nonzero digit must still give decay in the middle rank. We
recover the digit lost by averaging over \(\pi W\) in
Lemma~\ref{lem:inertdrop904}. The necessary comparison
is between the ordinary measure of a lattice and its measure
on an isotropic cone.

\begin{lemma}[Isotropic density in the fixed middle summand]
\leavevmode\label{lem:inertdensity905}
With the notation \(H_y,M_y,V,t\) of
Lemma~\ref{lem:inertdrop904}, use additive probability on \(V\)
and put
\[
      H_y^{\rm lin}=(1-q^{-2})\operatorname{vol}_{V}(M_y\cap V).
\]
Uniformly in the shell and \(j\), if \(d-2a\ge4\), one has
\[
  H_y=H_y^{\rm lin}(1+O(q^{1-d+2a})),\qquad
               H_y^{\rm lin}\ge (1-q^{-2})q^{-2P}.
\]
All implied constants in this section are absolute.
\end{lemma}
\begin{proof}
Write \(D_0=\operatorname{rank}W_0\ge d-2a\) for the fixed
unimodular summand constructed in the proof of
Lemma~\ref{lem:inertdrop904}. We first count norms in \(W_0\).
For any target \(s\) and \(l\ge1\), the probability of norm
\(s\bmod\pi^l\) in \(W_0\) is
\[
                         q^{-l}(1+O(q^{1-D_0})).
\]
Indeed expand the norm indicator in additive characters modulo
\(\pi^l\). After the trivial character, the total at conductor
exponent \(i\) inside the normalized Fourier sum is at most
\(q^i q^{-D_0i}\), by the one-dimensional Gauss calculation
in Theorem~\ref{thm:hermitiancone900} and an orthonormal basis.
Summing over \(1\le i\le l\) proves the claim. The part with
all entries divisible by \(\pi\) costs at most
\(O(q^{-l+2-2D_0})\): for \(l\ge2\) divide the norm equation
by \(\pi^2\) if its target allows this, and use the same bound
at precision \(l-2\), including precision zero. For \(l=1\)
the imprimitive part has measure at most \(q^{-2D_0}\).

Now write \(V=(V\cap W_0^\perp)\oplus W_0\). Integrality of
\(y^{-1}v\) and the leading coordinate \(h(v,e_{j+1})\) depend
only on the first summand. Their joint additive probability is
\(H_y^{\rm lin}\), since the image of the coordinate on \(M_y\cap V\)
is precisely \(\pi^tO_A\), whose unit fraction at that valuation
is \(1-q^{-2}\). For each fixed value in the first summand,
the preceding norm count applies with its prescribed target.
A primitive \(W_0\)-component guarantees smoothness in \(V\),
since \(W_0\) is nondegenerate. The estimate of the imprimitive
part thus shows that the smooth norm-zero probability at
precision \(l\ge m+1\) is \(q^{-l}(1+O(q^{1-D_0}))\)
uniformly for each first component. Apply this first with the
two linear restrictions and then without them. Dividing the
two probabilities gives the claimed estimate of \(H_y\) at the
large ranks, in fact uniformly for \(D_0\ge4\): the denominator
is bounded below by \((1-q^{-2})q^{-l}\) using a hyperbolic
pair in \(W_0\). Finally the index of \(M_y\cap V\) in \(V\)
is at most \([\mathcal L:M_y]=q^{2P}\), by duality with
\(A_y/\mathcal L\).
\end{proof}

\begin{proposition}[Cancellation at an interior drop of one digit]
\leavevmode\label{prop:inertonedigit905}
Fix \(J\) on a shell as in Lemma~\ref{lem:inertdrop904}, and
keep either of its compact orbit decompositions. If \(\delta=0\)
and \(b_j>b_{j+1}\) for some \(j<a\), the sum of the absolute
signed class sums on the cosets with that \(J\) is at most
\[
                       q^{C(a+1)^2(m+1)+dP-d+2a+1}.
\]
\end{proposition}
\begin{proof}
Here \(t=b_{j+1}\). In the integral for a signed orbit sum
insert the isotropic vectors as in
Lemma~\ref{lem:inertdrop904}, now dividing by \(H_y^{\rm lin}\).
This multiplies the integrand after integration in \(v\) by
\(H_y/H_y^{\rm lin}\). Consequently the total absolute error
over all orbits is at most \(O(q^{1-d+2a})\) times the unsigned
count for this \(J\), by Lemma~\ref{lem:inertdensity905} and
\(|\chi_0|=1\). The count is at most
\(q^{5a^2m+dP}\) by Lemma~\ref{lem:inertmiddle902}.

For the modified integral allow all
\(w\in W\cap v^\perp\cap\ker B_{>j}\) without the factor
\(\pi\). Identity~\eqref{eq:inertinner904}, preservation of
the orbit and \(J\), and translation invariance on the additive
cosets in its proof still hold. The left multiplier \(N_v(w)^{-1}\)
fixes \(L_j\) pointwise, hence preserves \(V\) and its additive
Haar measure. It carries \(M_y\) onto \(M_{y'}\). Thus
\(H_y^{\rm lin}=H_{y'}^{\rm lin}\) exactly.
The vector \(v\) is fixed by the multiplier and remains in
\(M_{y'}\). Its leading valuation is still \(t\). For instance,
\(J,\delta=0\) have stayed fixed, so the formula
\(t=b_{j+1}(J^+)=b_{j+1}(J)\) applies on both sides.
These facts verify invariance of all factors used in the
modified averaging.

Character orthogonality now requires integrality after multiplying
\(\beta_j(W\cap v^\perp\cap\ker B_{>j})\) by \(\pi^t\).
The conditional column count from the preceding section saves
\(2(d-a)(b_j-t)\ge2(d-a)\) in the exponent.
The reciprocal of the linear normalizer costs at most
\(2q^{2P}\) by the lemma. Since \(P\le2am\), this enters
\(q^{C(a+1)^2(m+1)}\). The bound on this modified sum and
the error estimate above prove the proposition.
\end{proof}

\begin{theorem}[Complete unweighted mass at a standard inert place]
\leavevmode\label{thm:inertzero905}
Under the good inert hypotheses of Theorem~\ref{thm:inerttail904}
there is an absolute constant \(C\) such that at all sufficiently
large nodes
\[
          1\le\int_G |b^0_K(x;0)|\,dx
               \le 1+q^{-2k+C(a+1)^4}.
\]
Every integral local frequency has the same mass.
\end{theorem}
\begin{proof}
For each \(m\ge1,\ P>0\) and candidate \(J\) there are three
possibilities. If \(\delta\ge1\), the unsigned count saves
at least \(d\) in the exponent relative to \(dP\).
If \(\delta=0\) and \(b_a\ge1\), the terminal-row bound saves
\(2\eta b_a\ge d-2a-2\). Otherwise \(2\sum b_i=P>0\) and
\(b_a=0\); some consecutive pair has a strict drop. Apply
Proposition~\ref{prop:inertonedigit905}. In all cases, adding
over at most \(q^{2a^2m}\) candidate lattices and absorbing the
polynomial factors gives
\[
       \sum_{\mathcal C}|S_{\mathcal C}|
               \le q^{C_1(a+1)^2(m+1)+dP-d}
\]
for the original coefficient on the entire shell. Its zero-frequency
energy is at most the square of this bound divided by the
orbit lower bound in Lemma~\ref{lem:inertframe902}.
When \(0<P<1000(a+1)^2\), one has \(m\le P\). Summing the
result over these polynomially many shells costs at most
\(q^{-2d+C_2(a+1)^4}\). Theorem~\ref{thm:inerttail904}
bounds all greater lengths by \(q^{-6k}\).
For \(P=0\) the two lattices coincide: inclusion of one
self-dual lattice in another gives equality by taking duals.
The position is in \(K\), where the contribution is one.
This proves both inequalities, using \(d=k+1\).
At an integral frequency the character on \(K^2\) is
identically one.
\end{proof}

\begin{corollary}[The standard unmarked Euler factors at zero]
\leavevmode\label{cor:inertproduct905}
On the fields and shorter intervals of
Corollary~\ref{cor:oddwindow677}, let \(\mathcal V\) be any set of
unmarked places of \(F=E^+\), of residue characteristic prime
to \(q_1\), satisfying the standard split or inert hypotheses
of the respective theorems. Then
\[
                  1\le\prod_{w\in\mathcal V}\mathcal B^0_w(0)
                               \le\exp(O(q_1^{\,2-k})).
\]
The logarithms converge absolutely. The same bound applies
at all members of this set at which a given rational frequency
is integral.
\end{corollary}
\begin{proof}
Use Corollary~\ref{cor:standardeuler899} for the split factors.
At an inert place the automorphism of the residue field of
\(E\) given by raising to \(q_w\) is the nonidentity element of
\(\operatorname{Gal}(E/F)\), acting on the \(q_1\)-th roots of
unity by inversion. Reduction is injective on these roots as
the residue characteristic is prime to \(q_1\). Thus
\(q_1\mid q_w+1\). By the theorem and \(k/(a+1)^4\to\infty\),
each logarithm is between zero and \(q_w^{-3k/2}\) at large
nodes. There are at most \(f\le q_1\) places of a given residue
size, so their sum is at most
\[
             q_1\sum_{j\ge q_1-1}j^{-3k/2}
                 \le \frac{q_1(q_1-2)^{1-3k/2}}{3k/2-1}
                 \ll q_1^{2-k}.
\]
For the last inequality use \((q_1-2)^{3/2}\ge q_1\) at large
nodes. Exponentiating proves the assertion.
\end{proof}

For a nonintegral frequency the distant estimate in
Theorem~\ref{thm:inerttail904} still applies. The interior
range then requires control of the compact average after
fixing a vector modulo its full conductor.
\clearpage

\subsection{The first inert transform at prime denominator}
\label{sec:inertprime906}

We use the integral lattice \(M_x=\mathcal L\cap x\mathcal L\)
to measure how many end vectors remain indistinguishable in a
compact average at conductor \(\pi^c\). Put
\[
  h_c(x)=\min\{c,\operatorname{ord}_A h(M_x,e)\},
             \qquad d_1=d-2a-2,\qquad c\ge1.
\]
The valuation here is that of the ideal of all pairings with the
first flag vector \(e\). It is finite since \(M_x\) has full rank.
The integer \(h_c(x)\) is constant on the \(K_c\)-orbit of the
right coset \(xK\). Indeed a left multiplier \(g\in K_c\)
carries \(M_x\) to \(gM_x\) and \(g^{-1}e\equiv e\bmod\pi^c\).
For a supporting coset \(yK\) with \(y\in U\), write \(\delta,p\)
as in Sections~\ref{sec:inertlength902} and~\ref{sec:inertterminal903}.

\begin{lemma}[End vectors under the lattice stabilizer]
\leavevmode\label{lem:inertend906}
The stabilizer in \(K\) of a supporting coset \(yK\) has at
least \(q^{2d_1(c-h_c(y))}\) different images of \(e\bmod\pi^c\).
In particular the \(K_c\)-orbit of that coset in
\(\mathcal S(m,P)\), for \(m\ge1\), has size at least
\[
          \tfrac13 q^{2dP-4a^2m+2d_1(c-h_c(y))-2rc}.
\]
\end{lemma}
\begin{proof}
Split off the integral unimodular \(W_0\) of rank at least \(d-2a\)
fixed pointwise by \(y\), as in Lemma~\ref{lem:inertdrop904}.
Take an element of \(M_y\) minimizing the valuation of the
pairing with \(e\), and project it off \(W_0\). This preserves
the valuation and membership in \(M_y\). In a hyperbolic plane
of \(W_0\) add an integral vector of norm equal to the negative
of the norm of this projection. Every integral value is attained
there by taking coefficient one on the first vector and prescribing
the trace of the other coefficient. The resulting vector \(v\)
is isotropic, belongs to \(M_y\), and still attains the minimum.
Let \(W_1\) be the orthogonal complement of this plane in
\(W_0\), of rank at least \(d_1\).

For \(w\in W_1\) the translation \(N_v(w)\) of
\eqref{eq:inertconj903} is integral. Its conjugate by \(y^{-1}\)
is \(N_{y^{-1}v}(w)\), also integral, since \(y\) fixes \(w\)
and \(v\in M_y\). It therefore stabilizes \(yK\).
The \(W_1\)-component of \(N_v(w)e\) is exactly \(h(v,e)w\),
all other terms lying in \(W_1^\perp\). Reduction modulo
\(\pi^c\) gives the asserted number of images.
For a coset with stabilizer \(H\subset K\) the size of its
\(K_c\)-orbit is its \(K\)-orbit size multiplied by
\([H:H\cap K_c]/[K:K_c]\), by the index formula. The numerator
counts these images of \(e\) and the denominator is at most
\(q^{2rc}\). Apply Lemma~\ref{lem:inertframe902}.
\end{proof}

For a right \(K\)-invariant coefficient \(f\), write
\(\mathcal B(f;\alpha)\) for the integral of the absolute
value of the compact double average defining
\(b^0_K(x;\alpha)\), with \(f\) used at both positions.
Let \(f_c\) be \(\phi_0\) on the supporting cosets satisfying
\[
                 h_c\le 1\qquad\text{or}\qquad h_c\le\delta+p ,
\]
extended by zero elsewhere.

\begin{theorem}[The first transform at prime denominator]
\leavevmode\label{thm:inertprime906}
At the standard unmarked inert places of
Theorem~\ref{thm:inerttail904}, for \(\operatorname{ord}_A\alpha=-c<0\),
\[
          \mathcal B(f_c;\alpha)
             \le q^{-n^7c/2} I_w^{(k-n^7)/2}D_w^{-3}
\]
at all sufficiently large nodes. For \(c=1\) this gives the
bound for the original coefficient \(\phi_0\) throughout the
whole group.
\end{theorem}
\begin{proof}
First fix \(m\ge1,P>0\) and partition the \(K_c\)-orbits by
\(h=h_c\). On the orbits of each given \(h\), the absolute
signed class sums of \(f_c\) satisfy
\begin{equation}
      \sum_{\mathcal C}|S_{\mathcal C}(f_c)|
                \le q^{C(a+1)^2(m+1)+dP-d_1h}.
 \label{eq:inertendsum906}
\end{equation}
For \(h\le1\) use the proof of Theorem~\ref{thm:inertzero905}
before division by orbit sizes. All three estimates used there
apply on the refined orbits as well, by the statements of
Theorem~\ref{thm:inertterminal903} and
Proposition~\ref{prop:inertonedigit905} and the unsigned count.
The saving \(d\), with terms \(O(a)\) in the other exponent,
suffices for \(d_1h\). If \(h\ge2\), only \(J=J_B\) with
\(\delta+p\ge h\) are retained. On each such \(J\) the terminal
calculation in Theorem~\ref{thm:inertterminal903} saves
\(d\delta+2\eta p\ge d_1h\) relative to \(dP\).
For a fixed \(J\) the condition on \(h\) merely selects whole
orbits, as observed above. The number \(q^{2a^2m}\) of
possible lattices and the factor \(2(m+1)\) are absorbed by the
absolute constant \(C\).

Let \(P_{K_c}f_c(x)=\int_{K_c} f_c(gx)\,dg\) with probability
Haar on \(K_c\). Its squared integral on the shell is the sum
of squared signed class sums divided by their orbit sizes.
Use Lemma~\ref{lem:inertend906} on orbits of each depth in
\eqref{eq:inertendsum906}. The terms involving \(h\) cancel,
giving the bound
\[
        (c+1)q^{C_1(a+1)^2(m+1)+2(r-d_1)c}.
\]
For the compact double average, the inequality of
Proposition~\ref{prop:inertoperator900} applies to \(f_c\)
because its proof permits arbitrary coefficients. On integration
over a \(K\)-invariant shell the conditional mean squares are the
squared integral just bounded, by the translation argument in
Theorem~\ref{thm:inerttail904}. The operator norm supplies the
factor \(q^{-rc}/\rho_c\le2q^{-(r-1)c}\).

Sum first over \(0<P<1000(a+1)^2(c+1)\), where \(m\le P\).
Since \(r-d_1=4a+2\), the entire sum is bounded by
\(q^{-(r-1)c+C_2(a+1)^4c}\). The polynomial number of shells
and powers of \(c+1\) are absorbed since \(c\ge1\) and
\(\log_q(c+1)\le c\). The part on \(K\) has \(h=0\) and
is bounded by Theorem~\ref{thm:hermitiancone900}; enlarge
\(C_2\) to include it.

For completeness the distant estimate remains valid for this
particular restriction despite the signs of the sums.
In the proof of Theorem~\ref{thm:inerttail904} the signed
estimate was obtained first for each \(J\), then summed over
\(J\) by the triangle inequality. For fixed \(J\) the restriction
to \(h\le\max(1,\delta+p)\) discards whole \(K_c\)-orbits;
hence each absolute sum over orbits can only decrease at
this step. Use the same orbit denominator and the same distant
series to obtain \(q^{-6k(c+1)}\). Finally \(n=a+1\),
\(r-1=k+2a\), \(I_w=q^{-2c}\) and \(D_w=q^{2c}\).
For sufficiently large \(n\) the saving \(n^7c/2\) available
in the arithmetic weight absorbs \(C_2n^4c+6c\) and the sum
of the two contributions. If \(c=1\), every supporting coset
has \(h_c\le1\).
\end{proof}

\begin{corollary}[Good Euler factors for squarefree inert denominators]
\leavevmode\label{cor:primeproduct906}
Use the fields, intervals and set \(\mathcal V\) of
Corollary~\ref{cor:inertproduct905}. Let \(y\in E^\times\)
have pole order at most one at each inert member of \(\mathcal V\).
For the original first relative Gram transform,
\[
      \prod_{w\in\mathcal V}\mathcal B^0_w(y)
        \le \exp(O(q_1^{2-k}))
            \prod_{\substack{w\in\mathcal V\\ D_w(y)>1}}
                   \min\{1,I_w(y)^{(k-n^7)/2}D_w(y)^{-3}\}.
\]
\end{corollary}
\begin{proof}
Only finitely many local denominators exceed one. At the split
ones apply Theorem~\ref{thm:standardfirst899}; at the inert ones
apply the theorem with \(c=1\). Discard the extra factors at
most one in those bounds. At every other place the frequency
is integral in each component, so the relative character on
the compact pair is one. Apply
Corollary~\ref{cor:inertproduct905} to their product.
\end{proof}

For a deeper inert denominator the remaining part has
\(h_c\ge2\) and \(h_c>\delta+p\). Both this part and its
mixed terms with \(f_c\) enter the original double average.
\clearpage

\subsection{Simultaneous inert row tests in an isotropic summand}
\label{sec:inertjoint907}

Work in the unramified quadratic field and self-dual integral
decomposition \(\mathcal L=L\oplus W\oplus R\) of
Section~\ref{sec:inertneighbor901}, where \(R=L^\vee\).
The rank of \(W\) is \(d\); \(z\in W\) is primitive.
We use the standard unmarked level and the unit coefficients
of the simple character. Write \(M_y=\mathcal L\cap y\mathcal L\)
and \(J=J_B=L+TL+BW\) for a unipotent representative \(y\).
We keep \(m,P,\delta,b_i(J)\) from Section~\ref{sec:inerttail904}.
The two copies of the coefficient can sample very deep first
endpoints. We construct one family of admissible row tests
which can be applied successively inside a single signed sum.

There is a hyperbolic middle summand \(W_H\) of rank \(2a\)
and a totally isotropic direct summand \(Q\subset W_H^\perp\cap W\)
of rank \(g=\lfloor(d-2a)/2\rfloor\) with \(h(z,Q)=O_A\).
First split off a hyperbolic pair whose first vector pairs to
a unit with \(z\), as in Section~\ref{sec:inertterminal903}.
Take \(a\) hyperbolic planes for \(W_H\) in the complement.
Inside the perpendicular to \(W_H\), extend the first vector
of the original pair to \(Q\) by using further integral
hyperbolic pairs. Their existence follows from the norm and
trace constructions of Lemma~\ref{lem:hermitianorbit900}.
Choose an integral map \(\iota:L\oplus R\longrightarrow W_H\)
taking the form to its negative. In paired hyperbolic bases
one uses the first vectors and the negatives of their partners.
Define
\[
  \mathcal P=\{ s+\iota(s)+u:\ s\in L\oplus R,\ u\in Q\},
       \qquad W_U=W_H\oplus Q .
\]
The summand \(\mathcal P\) is totally isotropic by orthogonality
and the defining property of \(\iota\). Its projection onto \(W\)
is \(W_U\). An integral complement \(W_C\) of \(W_U\) in \(W\)
has rank \(d-2a-g\); it need not be orthogonal to \(W_U\).

Write \(f_i=e_i\) for \(1\le i\le a\) and \(f_{a+1}=z\).
For a full sublattice \(D\) of \(O_A^{a+1}\), let its \(i\)-th
integer be the least valuation of the \(i\)-th coordinate among
vectors whose earlier coordinates vanish. The sum of the integers
is \(\operatorname{length}(O_A^{a+1}/D)\) by successive coordinate
projections. Use the linear evaluation map
\(E(v)=(h(f_i,v))_{i=1}^{a+1}\). Define \(t_i,u_i\) as the
integers for \(E(M_y)\), \(E(M_y\cap\mathcal P)\), respectively.
Both lattices have full rank. In fact \(E(\mathcal P)=O_A^{a+1}\)
because the \(R\)-component prescribes the first \(a\) values
and \(h(z,Q)=O_A\) corrects the last one. Also
\(\pi^m\mathcal L\subset M_y\). Thus \(0\le t_i\le u_i\le m\).
For each \(i\) select an integral \(v_i\in M_y\cap\mathcal P\)
attaining \(u_i\), with all earlier evaluations zero. Make the
choice depend only on \(M_y\cap\mathcal P\) and the fixed flag.
On a bounded shell there are finitely many possible such
sublattices, since they contain \(\pi^m\mathcal P\).

\begin{theorem}[Simultaneous simple-root tests at an inert place]
\leavevmode\label{thm:inertjoint907}
For a \(K_c\)-orbit of right \(K\)-cosets on a fixed shell,
including \(K_0=K\), its signed unipotent integral is unchanged
on retaining just the \(y\) satisfying all \(a\) conditions
\begin{equation}
     \beta_j\big(W\cap v_{j+1}^\perp\cap\ker B_{>j}\big)
                               \subseteq \pi^{-u_{j+1}}O_A,
                 \qquad 1\le j\le a .
 \label{eq:inertjoint907}
\end{equation}
Here \(\beta_j\) is row \(j\) of \(B\), and \(B_{>a}\) has no
rows. The identity holds also after fixing \(J\) and
\(M_y\cap\mathcal P\) inside the integral.
\end{theorem}
\begin{proof}
Fix these two lattices, and start at \(j=a\), proceeding down
to \(j=1\). At the current row write \(v=v_{j+1}\).
For \(w\in W\cap v^\perp\cap\ker B_{>j}\), put \(b=Bw\).
One has \(b\in L_j\otimes A\), \(v\perp L_j\), and both \(v\)
and \(y^{-1}v\) are integral isotropic vectors orthogonal to \(w\).
Identity~\eqref{eq:inertinner904} gives
\[
              y'=N_v(w)^{-1} y N_{y^{-1}v}(w)=N_v(b)y .
\]
This identity allows arbitrary integral \(w\) in the specified
summand. The left exterior multiplier fixes \(e_1\) and the
right one is integral, so \(y'\) remains in the same orbit.
As in Lemma~\ref{lem:inertdrop904}, \(B'=B(1+w h(v,y\,\cdot)|_W)\)
and \(T'=T+b h(v,y\,\cdot)|_L\). The first right multiplier is
an integral automorphism, since \(h(v,yw)=h(v,w+b)=0\).
Hence \(J\) stays fixed. On \(W\), the change \(B'-B\) has
image in \(L_j\otimes A\), so all rows strictly after \(j\)
stay fixed.

The additional invariance is integral and exact. On
\(\mathcal P\) the terms containing \(h(v,p)\) in \(N_v(w)p\)
vanish; this map sends \(p\) to \(p-vh(w,p)\). Thus it is an
integral automorphism of \(\mathcal P\), acting by addition of
integral multiples of \(v\). It preserves \(M_y\cap\mathcal P\)
in both directions because \(v\) belongs to that sublattice.
Since \(M_{y'}=N_v(w)^{-1}M_y\), intersection with \(\mathcal P\)
gives \(M_{y'}\cap\mathcal P=M_y\cap\mathcal P\).
In particular all the selected \(v_i\)'s remain exactly the same.

For \(j<a\) the only possible simple entry of \(N_v(b)\)
is \(\beta_j(w)h(v,e_{j+1})\); for \(j=a\) it is
\(\beta_a(w)h(v,z)\). These statements follow by restricting
\(1+b h(v,\cdot)-v h(b,\cdot)\) to the flag and to \(z\);
\(b\perp L\oplus W\). Therefore the simple character changes
by \(\psi_A(\gamma_j\beta_j(w)h(v,f_{j+1}))\), with
\(\gamma_j\) a unit. This is the simple-root calculation
underlying the Eichler map in the supplied source.

The change preserves unipotent Haar measure. At fixed \(v\)
disintegrate it along left cosets of the additive subgroup
\(\{N_v(u):u\in L_j\otimes A\}\). The allowed \(w\)'s are the
same along such a coset, since the lower rows of \(B\) stay
fixed. For each such \(w\), \(N_v(u)\) fixes \(w+Bw\), so the
translating element \(N_v(Bw)\) is constant on the coset.
Translation preserves its Haar measure, with inverse obtained
using \(-w\), as in Lemma~\ref{lem:inertdrop904}.
For \(j=a\) the group is again additive and closed: the nonzero
pairing \(h(v,z)\) makes \(u\) recoverable from \(N_v(u)z-z\);
for the other \(j\) use the pairing with \(e_{j+1}\).
The restrictions at indices greater than \(j\) depend only on
their own and lower rows and on the unchanged vectors \(v_i\).
They are consequently preserved in the successive integrations.
Averaging in \(w\) gives exactly the current condition by
character orthogonality over \(O_A\). Induction proves the
identity. Finally sum over the disjoint finite choices of the
two lattices whenever they are not prescribed.
\end{proof}

\begin{proposition}[Loss in the flag of evaluations]
\leavevmode\label{prop:inertloss907}
Let \(\Gamma:W\to J\oplus O_A\) be
\(\Gamma(w)=(Bw,h(z,w))\) and put
\(s_Q=\operatorname{length}((J\oplus O_A)/\Gamma(Q))\),
with value \(+\infty\) when the quotient is not finite. Set
\[
               G_B=\sum_{j=1}^a (b_j(J)-u_{j+1})_+ .
\]
If \(s_Q\) is finite, then
\begin{equation}
                           G_B+\delta+s_Q\ge t_1 .
 \label{eq:inertloss907}
\end{equation}
\end{proposition}
\begin{proof}
For \(v=s+w\in M_y\), \(s\in L\oplus R\), form the class of
\(\Gamma(w-\iota(s))\) modulo \(\Gamma(Q)\).
If it vanishes, choose \(u\in Q\) representing this difference.
The vector \(v'=s+\iota(s)+u\) differs from \(v\) by an
integral middle vector in \(\ker B\) orthogonal to \(z\).
That vector is fixed by \(y\), hence belongs to \(M_y\).
Thus \(v'\in M_y\cap\mathcal P\) and \(E(v')=E(v)\).
The kernel of the indicated linear map from \(M_y\) therefore
maps to zero in \(E(M_y)/E(M_y\cap\mathcal P)\), giving
\[
         \sum_{i=1}^{a+1}(u_i-t_i)
              \le\operatorname{length}(\Gamma(W)/\Gamma(Q)).
\]
Also every vector in \(\ker B\cap W\) is in \(M_y\) and has
zero pairing with \(L\). Thus \(t_{a+1}\) is at most the
valuation of \(h(z,\ker B\cap W)\). The latter is at most
\(\operatorname{length}((J\oplus O_A)/\Gamma(W))\), since
the intersection of \(\Gamma(W)\) with the last coordinate
line consists precisely of these pairings. Adding gives
\(t_{a+1}+\sum_i(u_i-t_i)\le s_Q\).

For \(i\le a\), \(t_i=b_i(J^+)\) with
\(J^+=(y\mathcal L+\mathcal L)\cap (L\otimes A)\).
The lattice proof in Lemma~\ref{lem:inertdrop904} applies also
to \(i=1\), with no earlier coordinates to remove.
Consequently \(\sum_{i=1}^a t_i=\sum_{i=1}^a b_i(J)+\delta\).
Now sum \(b_j(J)-u_{j+1}\) before taking positive parts.
The result is \(t_1-\delta-t_{a+1}
-\sum_{i=2}^{a+1}(u_i-t_i)\), at least \(t_1-\delta-s_Q\).
Taking positive parts only increases the sum.
\end{proof}

The integer \(s_Q\) uses only the restriction of \(B\) to \(Q\).
The row conditions in \eqref{eq:inertjoint907} also test the
columns on the independent complement \(W_C\).
\clearpage

\subsection{The full first transform at standard inert places}
\label{sec:inertfull908}

Continue with the summands and lattices of
Section~\ref{sec:inertjoint907}. We estimate the cost of the
index \(s_Q\) using the columns on \(Q\); the independent
columns on \(W_C\) account for the remaining endpoint depth.
Write \(\lambda=d-6a-4\), which is positive at sufficiently
large nodes.

\begin{lemma}[Refined signed class sums at every endpoint depth]
\leavevmode\label{lem:inertdepth908}
On \(\mathcal S(m,P)\) with \(m\ge1\), let \(\mathcal C\) range
over the \(K_c\)-orbits of right cosets with \(h_c=h\), where
\(c\ge1\) and \(0\le h\le c\). For the entire original coefficient
let \(S_{\mathcal C}\) be its signed sum on the orbit. Then
\begin{equation}
       \sum_{\mathcal C:\,h_c=h}|S_{\mathcal C}|
           \le (h+1)q^{C(a+1)^2(m+1)+dP-\lambda h}
 \label{eq:inertdepth908}
\end{equation}
for an absolute \(C\).
\end{lemma}
\begin{proof}
Apply Theorem~\ref{thm:inertjoint907} in each orbit before taking
its absolute value. It suffices to count the volume of the
resulting simultaneous tests. We use the product coordinates
\(T,B,Z\) and the central-block notation \(C_y\) for the block
called \(C\) in Lemma~\ref{lem:inertmiddle902}. Fix \(T\), a
candidate \(J=J_B\), and the entries of \(C_y\) modulo \(O_A\).
The latter have at most \(q^{2a^2m}\) possibilities, since all
their pole orders are at most \(m\). For any fixed \(T,B\) the
allowed \(Z\)-volume is at most \(q^{2a^2m}\), as in that lemma.
Discarding other restrictions on it after this count is valid
even though \(C_y\) is related to \(B\) by unitarity.

Once \(B|_{W_U}\) is also given, the entire sublattice
\(M_y\cap\mathcal P\) is determined by the fixed data.
Indeed for an integral \(v=s+w\in\mathcal P\) its middle
component \(w\) belongs to \(W_U\).
The condition \(v\in M_y\) means that \(h(v,y\,\cdot)\) is
integral. On the left and middle input blocks this is precisely
integrality of its \(L\)-functional on \(J=L+TL+BW\), since the
middle component of \(v\) is already integral. On the last
input block the pairing is the sum of its \(L\)-functional on
\(C_y\), its \(R\)-functional on \(S\), and the middle term
\(w^*HD=-(Bw)^*S\). All are determined modulo integral
functionals by the prescribed data. Thus the selected \(v_i\),
their \(u_i\), and \(G_B\) are fixed before the columns on
\(W_C\) are counted. One may now enlarge the domain of \(B\)
to independent columns uniform in \(J\), keeping these fixed
data as parameters for the tests. Its unnormalized volume is
\(q^{2d\operatorname{length}(J/L)}=q^{d(P-\delta)}\).
Only data occurring for an actual coset of the indicated
depth need be retained.

For this column measure, choose an integral basis of \(Q\)
in which the row \(h(z,\cdot)\) is \((1,0,\ldots,0)\).
The index length \(s_Q\) is then the length of \(J\) modulo
the last \(g-1\) random columns from \(Q\). This follows by
use of the first column, whose last coordinate is one, to
eliminate the last coordinate in \(J\oplus O_A\).
For any integer \(s\ge0\),
\begin{equation}
                    \Pr(s_Q\ge s)\le q^{-(2g-3a)s}.
 \label{eq:inertindex908}
\end{equation}
If the length is at least \(s\), all \(g-1\) columns belong
to some sublattice of colength \(s\) in \(J\), also when the
quotient has a free part. A fixed such lattice costs
\(q^{-2(g-1)s}\). Their number is at most
\((s+1)^a q^{2(a-1)s}\). To see this in a basis of \(J\),
choose the successive diagonal valuations of a triangular
basis, whose sum is \(s\). In each pivot row the entries on
one side of that pivot reduce modulo its pivot ideal, and
there are at most \(a-1\) such entries. This yields the stated
count. Use \(s+1\le 2^s\) for \(s\ge1\) and \(q\ge2\).
For \(s=0\) the assertion is immediate.

Next fix the columns on \(W_U\). Reveal the rows of the
\(W_C\)-columns from the last row up to the first. Conditional
on their coordinates strictly after row \(j\), the row
\(\beta_j\) there has independent coordinates uniform in
translates of \(\pi^{-b_j(J)}O_A\): the projection kernel
in \(J\) is \(J\cap(L_j\otimes A)\). Within \(W_C\) the
submodule \(\ker B_{>j}\cap v_{j+1}^\perp\) is primitive and
has rank at least \(\operatorname{rank}W_C-a\). It is specified
by the fixed columns, \(v_{j+1}\), and the already revealed
rows. Completing a basis of that kernel inside \(W_C\)
preserves the conditional product measure up to translation.
The fraction meeting the \(j\)-th test is therefore at most
\[
           q^{-2(\operatorname{rank}W_C-a)(b_j(J)-u_{j+1})_+}.
\]
This conditional bound holds for every permitted set of later
rows. Iteration bounds the fraction meeting all the tests by
\(q^{-2(\operatorname{rank}W_C-a)G_B}\).

Partition the volume according to \(s=\min(h,s_Q)\).
This has at most \(h+1\) values and is decided on \(W_U\).
If \(s<h\), Proposition~\ref{prop:inertloss907} gives
\(G_B\ge(h-\delta-s)_+\), since \(t_1\ge h_c=h\).
When \(s=h\) the same inequality follows from \(G_B\ge0\).
It continues to bound the conditional count after enlargement:
if the fixed data occur at the required depth, then their
already determined \(G_B\) has this value or a greater one.
Combining the two counts, the total saving in the exponent
relative to \(dP\) is at least
\[
 d\delta+(2g-3a)s+
          2(\operatorname{rank}W_C-a)(h-\delta-s)_+
                                  \ \ge\ \lambda h .
\]
Here \(g=\lfloor(d-2a)/2\rfloor\) and
\(\operatorname{rank}W_C=d-2a-g\).
Finally \(T\) costs at most \(q^{a(a-1)m}\), the choice of
\(J\) at most \(q^{2a^2m}\), and the residue choices and
central volume have the bounds above.
After the signed identities we have counted a subset of
the original unipotents; there is no factor counting
\(M_y\cap\mathcal P\), since that lattice is already determined
by the conditioned data. This proves \eqref{eq:inertdepth908}.
\end{proof}

\begin{theorem}[The complete first transform at a standard inert place]
\leavevmode\label{thm:inertfull908}
Let \(w\) satisfy the standard unmarked inert hypotheses of
Theorem~\ref{thm:inerttail904}. With both original Bessel
coefficients on the full group, at every nonintegral frequency
of conductor exponent \(c\ge1\) one has
\begin{equation}
       \mathcal B^0_w(\alpha)=\int_G|b^0_K(x;\alpha)|\,dx
          \le q^{-n^7c/2} I_w^{(k-n^7)/2}D_w^{-3}
 \label{eq:inertfull908}
\end{equation}
at all sufficiently large nodes, uniformly in \(w,c\).
Integral frequencies have the zero-frequency mass of
Theorem~\ref{thm:inertzero905}.
\end{theorem}
\begin{proof}
For \(m\ge1\) and fixed depth \(h\), divide the square of
\eqref{eq:inertdepth908} by the lower bound for the refined
orbit sizes in Lemma~\ref{lem:inertend906}. The result bounds
the sum of the squared signed sums divided by their individual
orbit sizes. Adding over \(h\) bounds the squared integral of
the \(K_c\)-average of the full \(\phi_0\) on the shell by
\[
       (c+1)^3 q^{C_1(a+1)^2(m+1)+2(r-\lambda)c}.
\]
Indeed \(d_1=d-2a-2\ge\lambda\) and the depth contribution
after division is \(2(r-d_1)c+2(d_1-\lambda)h\).
In particular this calculation already includes all products
of distinct parts of the Bessel coefficient occurring in a
squared orbit sum.

The conditional Fourier operator of
Proposition~\ref{prop:inertoperator900} has norm at most
\(2q^{-(r-1)c}\). Its application to the original double
average, integrated on the shell, multiplies the last bound
by this factor, as in Theorem~\ref{thm:inertprime906}.
For \(0<P<1000(a+1)^2(c+1)\), one has \(m\le P\), and
\(r-\lambda=8a+4\). The sum over all these shells is thus
at most \(q^{-(r-1)c+C_2(a+1)^4c}\), including the polynomial
factors because \(\log_q(c+1)\le c\). On \(K\) use
Theorem~\ref{thm:hermitiancone900}; beyond that range of
\(P\) use Theorem~\ref{thm:inerttail904}. The latter contribution
is bounded by \(q^{-6k(c+1)}\). With \(n=a+1\), \(r-1=k+2a\),
\(I_w=q^{-2c}\) and \(D_w=q^{2c}\), the sum satisfies
\eqref{eq:inertfull908} once \(n^7/2>C_2 n^4+8\) and the
node is sufficiently large. At integral frequency the relative
character on the two compact variables is one.
\end{proof}

\begin{corollary}[Standard finite places at arbitrary rational frequency]
\leavevmode\label{cor:standardall908}
On the fields and intervals of
Corollary~\ref{cor:inertproduct905}, let \(\mathcal V\) be any
set of its standard unmarked split and inert places.
For every \(y\in E^\times\),
\[
       \prod_{w\in\mathcal V}\mathcal B_w^0(y)
          \le \exp(O(q_1^{2-k}))
               \prod_{\substack{w\in\mathcal V\\D_w(y)>1}}
                    \min\{1,I_w(y)^{(k-n^7)/2}D_w(y)^{-3}\}.
\]
\end{corollary}
\begin{proof}
Apply the theorem at the inert poles and
Theorem~\ref{thm:standardfirst899} at the split ones. There
are only finitely many poles. For an inert pole the displayed
arithmetic factor is less than one. All the integral masses
have product bounded by Corollary~\ref{cor:inertproduct905}.
\end{proof}

This controls the first relative coordinate at every place
satisfying the stated integral hypotheses. The remaining
coordinates involve the longer flag and the last vector of norm
\(-N\), and the original block also has real factors.
\clearpage

\section{Local transforms in the rational frequency sum}\label{nav:local-rational-frequency-sum}

\subsection{The first frequency sum on the period fields}
\label{sec:firstglobal909}

Take an actual node \(E=\mathbb Q(\sqrt{-d_0},\mu_{q_1})\),
\(F=E^+\), of the shorter interval of
Corollary~\ref{cor:oddwindow677}; \(d_0\) here denotes the
quadratic prime. Let \(f=[F:\mathbb Q]\), and write
\(\mathcal W=\{w:w\mid p_0p_1\}\) for the original marks,
\(\mathcal U\) for all the other finite places of \(F\).
Use the finite position factors after the class changes in
Theorem~\ref{thm:spheremean872}. We verify the integral
hypotheses also at the places above the separate prime \(N\).

\begin{theorem}[All unmarked factors on the specified fields]
\leavevmode\label{thm:allunmarked909}
In these coordinates each \(w\in\mathcal U\) satisfies the
hypotheses of Theorem~\ref{thm:standardfirst899} or of
Theorems~\ref{thm:inertzero905} and~\ref{thm:inertfull908}.
Put \(K_0=k-n^7\), \(V_w(y)=I_w(y)^{K_0/2}D_w(y)^{-3}\),
and
\[
             \epsilon=q_1^{\,2-k}
                         +f\log(2q_1)\,2^{-3k/2}.
\]
For every \(y\in E^\times\), at the same sufficiently large nodes,
\begin{equation}
     \prod_{w\in\mathcal U}\mathcal B^0_w(y)
       \le \exp(O(\epsilon))
            \prod_{\substack{w\in\mathcal U\\D_w(y)>1}}
                                  \min(1,V_w(y)).
 \label{eq:allunmarked909}
\end{equation}
The product at frequency zero lies between one and
\(\exp(O(\epsilon))\).
\end{theorem}
\begin{proof}
The two imaginary quadratic subfields \(\mathbb Q(i)\) and
\(\mathbb Q(\sqrt{-d_0})\) have coprime discriminants. At each
rational prime its inertia group in the abelian extension
\(E/\mathbb Q\) fixes at least one of these fields, whereas
complex conjugation acts nontrivially on both.
Thus no finite inertia group contains complex conjugation.
The quadratic extension \(E/F\) is unramified at every finite
place, as in Corollary~\ref{cor:fieldrelations712}. Every
quadratic local algebra under consideration is accordingly split
or an unramified field, even at a prime dividing \(q_1\).

In the rational middle Gram matrix the positive vectors \(x_i\)
are orthonormal, \(h(z,z)=-N\), and \(z\) is orthogonal to them.
The actual finite lattice uses the basis replacing a single
\(x_*\) by
\[
                 u=x_*/\bar\pi_N+z/N,
                         \qquad \pi_N\bar\pi_N=N .
\]
This is the lattice of Section~\ref{sec:discriminant681} and of
Lemma~\ref{lem:massscale869}. In the pair \(u,z\) its Gram
matrix is
\[
                   \begin{pmatrix}0&-1\\-1&-N\end{pmatrix}.
\]
Indeed the two terms in \(h(u,u)\) are \(1/N,-1/N\);
each cross pairing with \(z\) is \(-1\). This block has unit
determinant up to sign at every place. The remaining middle
lines have norm one and the paired arms are integral hyperbolic
planes. The entire lattice is self-dual; \(z\) is a basis
vector and its covector takes the unit value \(-1\) on \(u\).
Thus both terminal directions required in the split calculation
are primitive in the corresponding components. This argument
also applies over \(N\), where the original positive basis and
the actual lattice basis can differ integrally.

The local level is the full stabilizer of this lattice when
\(w\nmid M=p_0p_1\), since the level imposes only the reduction
modulo \(M\). The finite simple character reads the successive
entries in the frame \(e_1,\ldots,e_a,z\), with common scalar
\(M^{-1}\), a unit at \(w\). The finite additive pairing is
the trace divided by the different generator \(\delta_E\), whose
dual integer lattice is the integer ring itself. In a split
algebra this gives two conductor-integer-ring characters, each
a unit multiple of one such character of \(F_w\). For an unramified
quadratic algebra it is a conductor-integer-ring character of
that algebra; it is a unit multiple in its argument of a
conductor-integer-ring character of \(F_w\) composed with
relative trace. This follows from the perfect integral trace
pairing and additive duality, or by transitivity of the different
and the unit relative different. All the flag vectors and their
simple-coordinate functionals in the preceding lattice basis
are integral with the required primitive coefficients.
The same unit changes for the first frequency preserve its
local valuations.

In the class integrals the changes of variables by \(h_c^{-1}\)
and \(h_{c'}^{-1}\) in Theorem~\ref{thm:spheremean872} carry the
two normalized unipotent measures to those giving \(U\cap J_w\)
volume one. These middle isometries fix \(z\) and the two arms,
normalize \(U\), and preserve its character. The first relative
entry also satisfies
\(h(e,h_c^{-1}g h_{c'}e)=h(e,ge)\).
Thus the class-removed position calculation uses the specified
lattice even when a chosen class representative is nonintegral
at \(w\). The local hypotheses have now been checked.

Outside the residue characteristics dividing \(q_1\) apply
Corollary~\ref{cor:standardall908} and its zero-frequency
counterpart. There are at most \(f\log(2q_1)/\log 2\) other
places, since at most \(f\) lie over any rational prime.
At each of these, an integral-frequency mass lies in
\([1,1+q_w^{-3k/2}]\) at large nodes by the two local theorems:
their positive error exponents are \(O(n^4)\) and \(n^4=o(k)\).
At its poles the corresponding weighted theorem applies unchanged.
Using \(q_w\ge2\) on this last finite set proves the asserted
upper bounds. At zero all local masses are at least one,
and the absolute convergence of their logarithms follows from
the same estimates.
\end{proof}

\begin{corollary}[Rational first frequencies with the real power weight]
\leavevmode\label{cor:firstglobal909}
Let \(\Xi\subseteq\mathcal W\) be the selected set of
Theorem~\ref{thm:bothmarks873}, and let
\(\mathcal V=\mathcal W\setminus\Xi\).
Write \(c_v=2\pi i/(\Phi_v(\delta_E)b_{1,v})\) for the actual
first-coordinate scale. For \(y\ne0\) define the real weight
\[
       A_\infty(y)=\prod_{v\mid\infty}
           |c_v\Phi_v y|^{-K_0}\min\{1,|c_v\Phi_v y|^{-6}\}.
\]
Using \(B_0,j_w,Z_w^0\) of Corollary~\ref{cor:firstsample898},
one has
\begin{align}
 &\sum_{\substack{0\ne y\in E\\y\ {\rm integral\ at\ all\ primes\
                  over}\ \mathcal V}}
 A_\infty(y)
 \prod_{w\in\Xi}
       \frac{B_0 q_w^{n-j_w}\mathcal B_{Z_w}(y)}{Z_w^0}
 \prod_{w\in\mathcal U}\mathcal B^0_w(y)
                    \ \le\ \exp(O(\epsilon))\, b^{-7kf}.
 \label{eq:firstglobal909}
\end{align}
\end{corollary}
\begin{proof}
Take first a finite conjugation-stable set \(T\) of primes of
\(E\), containing both primes above each member of \(\Xi\),
disjoint from the primes above \(\mathcal V\). For
\(0\ne y\in\mathcal O_{E,T}\), the product of \(V_w(y)\) over
the places under \(T\) is \(I_T(y)^{K_0/2}D_T(y)^{-3}\).
At an integral place \(V_w(y)\ge1\). By
\eqref{eq:allunmarked909}, the product of all unmarked masses
divided by the product of \(V_w(y)\) at the unmarked places
under \(T\) is therefore at most \(\exp(O(\epsilon))\).
Multiplying the summand in \eqref{eq:firstsample898} by this
ratio removes exactly the remaining finite power weights and
inserts the unmarked masses. This proves the estimate with
\(y\) restricted to \(\mathcal O_{E,T}\). Increase \(T\)
through all finite conjugation-stable sets allowed above.
Every \(y\) integral at the unselected marks has only finitely
many further poles and hence occurs in the increasing union.
All summands in the displayed corollary are nonnegative, so
monotone convergence gives \eqref{eq:firstglobal909}.
\end{proof}

In \eqref{eq:firstglobal909} the real factor \(A_\infty\) is the
specified power weight. Applying this estimate to the original
theta-period integral requires a bound on its real Bessel
correlations at these scales, with the other frame entries and
the rational sampling hypotheses retained.
\clearpage

\subsection{The Fourier operator on the real frame quadrics}
\label{sec:realquadrics911}

At a real place write \(h(v,w)=v_+^*w_+-v_-^*w_-\) on
\(\mathbb C^p\oplus\mathbb C^q\), and \(r'=p+q\).
The signature at the first frame vector is
\((p,q)=(k+a,a+1)\). After removing \(j\) hyperbolic pairs it
is \((k+a-j,a+1-j)\). On the first \(a\) orbits the remaining
frame starts with a nonzero null vector. On the last orbit it
starts with the vector of norm \(-1\), and \(q=1\).
For \(\tau=0,1\) put
\[
 \mathcal C_\tau=\{v:h(v,v)=-\tau,\ v\ne0\},\qquad
 d\sigma_\tau(v)=\delta(h(v,v)+\tau)\,d^{2r'}v .
\]
Here the delta notation means the fiber density obtained by
using \(h(v,v)+\tau\) as one real coordinate. On the indicated
orbit its differential is nonzero. The determinant of an
isometry has modulus one, so this measure is invariant.
In the following formulas \(\sigma_\tau\) always has this
normalization. Write \(s=\|v_+\|\). Let \(\sigma_\tau^R\) be
its restriction to \(s\le R\). For \(\nu\ne0\) set
\[
 (T_{\nu,R}g)(v)
   =\int e^{2i\operatorname{Im}(\nu h(v,w))}
               g(w)\,d\sigma_\tau^R(w),\qquad s(v)\le R .
\]

\begin{theorem}[A uniform estimate for the real quadric operator]
\leavevmode\label{thm:realquadrics911}
Suppose \(p,q\) are positive integers, \(R>0\), and either
\(\tau=0\) or \((\tau,q)=(1,1)\). For every real
\(0\le\ell\le r'-2\),
\begin{equation}
 \|T_{\nu,R}\|_{L^2(\sigma_\tau^R)\to L^2(\sigma_\tau^R)}
 \le \frac{\pi^{r'}}{r'-1-\ell}\,
         R^{2(r'-1-\ell)}|\nu|^{-\ell}.
 \label{eq:realquadrics911}
\end{equation}
\end{theorem}
\begin{proof}
Use probability measure on each of the two unit spheres. Put
\(v=(s\omega,\sqrt{s^2+\tau}\,\eta)\). Integrating first in
the negative radius divides its radial measure by twice that
radius. Since the sphere in \(\mathbb C^m\) has area
\(2\pi^m/\Gamma(m)\), the resulting radial density is
\[
 d\lambda(s)=
 \frac{2\pi^{r'}}{\Gamma(p)\Gamma(q)}
          s^{2p-1}(s^2+\tau)^{q-1}\,ds.
\]
In both cases of the theorem its variable factor is \(s^{2r'-3}\).
The endpoints omitted from the polar charts have measure zero.

On the probability sphere \(S^{2m-1}\) the integral operator
with kernel \(e^{iu\omega\cdot\theta}\), \(u>0\), acts on the
harmonics of degree \(d\) by
\[
       i^d\Gamma(m)(u/2)^{1-m}J_{m+d-1}(u).
\]
For completeness let \(P_d\) be a homogeneous harmonic
polynomial on \(\mathbb R^{2m}\). Completing the square for real
\(z\) in the following Gaussian integral and taking spherical
means around \(z\omega/2\) gives
\[
 \int_{\mathbb R^{2m}}e^{-|x|^2+z x\cdot\omega}P_d(x)\,dx
      =\pi^m e^{z^2/4}(z/2)^d P_d(\omega).
\]
The mean-value property applies to each real and imaginary
part of the polynomial. Both sides are entire in \(z\).
Expanding the left side in polar coordinates, the coefficient
of \(z^l\) is the degree-\(l\) angular moment divided by \(l!\),
multiplied by \(\pi^m\Gamma(m+(d+l)/2)/\Gamma(m)\).
Thus the angular moments vanish unless \(l=d+2b\) with \(b\ge0\).
Inserting the remaining moments in the series for the spherical
kernel gives \(i^d\Gamma(m)P_d(\omega)\) times
\[
  (u/2)^d\sum_{b\ge0}
       \frac{(-u^2/4)^b}{b!\,\Gamma(m+d+b)}.
\]
This is the asserted Bessel expression. For integer \(N\ge0\)
the coefficient of \(z^N\) in
\(\exp((u/2)(z-z^{-1}))\), expanded on \(|z|=1\), has the
series for \(J_N(u)\). The integrand defining that Fourier
coefficient has modulus one. Hence \(|J_N(u)|\le1\).
These formulas also give the usual integral and series
normalizations of \(J_N\).
\href{https://dlmf.nist.gov/10.9.E2}{DLMF, Eq.~10.9.2} and
\href{https://dlmf.nist.gov/10.2.E2}{Eq.~10.2.2}.

Spherical harmonics of different degrees are orthogonal and
their span is dense, by the harmonic decomposition of
polynomials and density of polynomial restrictions on a sphere.
The angular operator therefore has norm at most
\(\Gamma(m)(u/2)^{1-m}\). Its norm is also at most one by
Cauchy--Schwarz for probability measure. At fixed positive
radii \(s,t\), the two angular operators for \(h\) form a
tensor product. The phase of \(\nu\), the imaginary part, and
the minus sign in \(h\) compose the factors with orthogonal
changes of variables. The positive factor has \(u=2|\nu|st\).
On the null orbit the negative factor has the same \(u\).
On the orbit with \(q=1\) we use only the probability bound
for the negative factor. In both cases the tensor norm is
at most one and at most
\(\Gamma(p)\Gamma(q)(|\nu|st)^{-(r'-2)}\).
Since the factorial product is at least one, it is at most
\(\Gamma(p)\Gamma(q)(|\nu|st)^{-\ell}\) in the stated range.

For two test functions, bound each angular pairing at radii
\(s,t\) by this norm times their two angular \(L^2\) norms.
The resulting two radial integrals separate. Cauchy--Schwarz
bounds their product by the product of the full \(L^2\) norms
times \(\int_0^R s^{-2\ell}d\lambda(s)\). Evaluation cancels
the two gamma factors and gives \eqref{eq:realquadrics911}.
\end{proof}

\begin{corollary}[The first Gram wave before position integration]
\leavevmode\label{cor:realgramposition911}
Let \(G=\mathrm U(k+a,a+1)\), let \(U\) be the Bessel
unipotent, and let \(e\) be the first frame vector, of
Euclidean length one. For \(a=0\) take \(e=z\). Set
\(\tau=-h(e,e)\). Use the quotient Haar measure of
Theorem~\ref{thm:normalizedcut699}. For the cut vector
\(F=F_v^\flat\) from that theorem and \(x\in G\), let
\(D_x(v)d\sigma_\tau(v)\) be the image of
\(\langle F(y),F(yx)\rangle\,d\dot y\) on \(U\backslash G\)
under \(\dot y\mapsto y^{-1}e\). Define
\[
 B(x;\nu)=\iint e^{2i\operatorname{Im}(\nu h(y^{-1}e,t^{-1}e))}
   \langle F(y),F(yx)\rangle\,
   \overline{\langle F(t),F(tx)\rangle}\,d\dot y\,d\dot t .
\]
Use \(n=a+1,\ r=k+2a+1,\ K_0=k-n^7\), at nodes with \(K_0\ge0\)
and \(n^7\ge7\). For \(R\ge\max(1,b^{6n+2})\),
\begin{equation}
 |B(x;\nu)|\le\pi^r R^{2(n^7+2a)}
       |\nu|^{-K_0}\min(1,|\nu|^{-6})
       \int_{\mathcal C_\tau}|D_x|^2\,d\sigma_\tau .
 \label{eq:realgramposition911}
\end{equation}
\end{corollary}
\begin{proof}
The subgroup \(U\) fixes \(e\). Its scalar character in
\(F(y)\) cancels that in \(F(yx)\), so the density descends
to the quotient. The integral is supported in the compact
image of the matrix norm ball of radius \(R\). Indeed the cut
is supported where \(\inf_{s\in S}(1+\|sy\|)\le b^{6n+2}\)
with \(S=U\rtimes\mathrm U(k)\), and the compact factor acts
unitarily. In the chosen Hermitian coordinates
\(\|g^{-1}\|=\|g\|\). Hence \(s(y^{-1}e)\le R\).
The map to the orbit is a smooth submersion. On a compact
support a bounded density has bounded image density relative
to smooth orbit measure, by integration in a finite collection
of local product charts and a partition of unity. This proves
\(D_x\in L^2(d\sigma_\tau)\), including for the continuous
cut. Also \(h(y^{-1}e,t^{-1}e)=h(e,yt^{-1}e)\), the first
Gram coordinate of the relative element. Pushing forward
both integrals gives a pairing for \(T_{\nu,R}\) or its
transpose, which has the same norm.

For \(|\nu|\le1\) take \(\ell=K_0\). For \(|\nu|>1\) take
\(\ell=K_0+6\). Both are at most \(r-2\). The denominator
in \eqref{eq:realquadrics911} is at least one and
\(r-1-K_0=n^7+2a\), proving the inequality. The same exponent
choices work on every residual quadric of signature
\((k+a-j,a+1-j)\), since its dimension parameter is at least
\(k+1\).
\end{proof}

The right side of \eqref{eq:realgramposition911} still
depends on the position through the conditional correlation
\(D_x\). Estimating its integral in \(x\) is necessary before
inserting the real factor in Corollary~\ref{cor:firstglobal909}.
\clearpage

\subsection{The first real transform on a ball of positions}
\label{sec:realpositionball912}

We retain the normalized real frame of Section~\ref{sec:realgram682}.
Let \(G_j=\mathrm U(k+a-j,a+1-j)\) and \(r_j=k+2(a-j)+1\)
for \(0\le j\le a\). Let \(f_j=e_{j+1}\) if \(j<a\) and
\(f_a=z\). Write \(\sigma_j\) for the measure of
Section~\ref{sec:realquadrics911} on the orbit of \(f_j\).
Put \(r=r_0,\ n=a+1\). All matrix norms are Euclidean operator
norms in coordinates in which the two vectors of a hyperbolic
pair are orthonormal for the Euclidean metric and have
Hermitian pairing one.

\begin{lemma}[Truncated Plancherel estimate on one quadric]
\leavevmode\label{lem:truncquad912}
Let \(m=p+q\), \(R\ge1\), \(T>0\), and let \(g\in L^2(\sigma_\tau)\)
be supported in \(R^{-1}\le\|v\|\le R\). Then
\[
 \int_{\|t\|\le T}
   \left|\int g(v)e^{2i\operatorname{Im} h(t,v)}
                                      d\sigma_\tau(v)\right|^2 dt
 \le 2(2m)^{3/2}RT\,\pi^{2m-1}\int |g|^2d\sigma_\tau .
\]
\end{lemma}
\begin{proof}
Diagonalize the form by a Euclidean unitary change of coordinates.
The imaginary part of \(h(t,v)\) is the ordinary real dot product
of \(v\) with an orthogonal transform of \(t\). In \(D=2m\)
real coordinates partition the support according to the first
coordinate of largest absolute value and its sign. On each
of the resulting \(2D\) sets, deletion of this coordinate is
one-to-one. Its absolute value is at least \(1/(R\sqrt D)\).
The graph density over the other coordinates \(v'\) is
\(J(v')=(2|v_i|)^{-1}\le R\sqrt D/2\). Fix the dual variable
\(\xi_i\). Plancherel in the other \(D-1\) coordinates, with
kernel \(e^{2i v'\cdot\xi'}\), gives
\(\pi^{D-1}\int |g|^2J^2\,dv'\) on this patch. The graph phase
in \(\xi_i\) has modulus one. Integrating \(\xi_i\) over an
interval of length \(2T\) and using
\(J^2dv'\le(R\sqrt D/2)d\sigma_\tau\) proves the estimate for
each patch before the factor \(2D\). Cauchy--Schwarz for the
sum of the transforms contributes that factor, and the
squared input norms add because the patches are disjoint.
\end{proof}

We describe the measures in the recursion. Let \(H_j\) fix
\(f_j\) in \(G_j\), let \(U_j\) be the unipotent for the remaining
frame, and put \(C=\mathrm U(k)\). Thus \(H_a=C,\ U_a=1\).
For \(j<a\) use the order \(e,V,e^*\) with \(e=f_j\) and
\(V\) the residual space of dimension \(r_{j+1}\). There is a
subgroup \(N_j\) consisting of the transformations
\[
 \begin{split}
 e&\longmapsto e,\qquad w\longmapsto w-h(t,w)e\quad(w\in V),\\
 e^*&\longmapsto e^*+t+\big(-h(t,t)/2+i u\big)e
                    \qquad(t\in V,\ u\in\mathbb R).
 \end{split}
\]
Write this element as \(n(t,u)\). Pairing the images of
\(e,w,e^*\) gives the displayed coefficients. Conversely any
isometry fixing \(e\) acts on \(e^\perp/\mathbb C e\) by some
\(l\in G_{j+1}\). Removing \(l\) leaves exactly one \(n(t,u)\).
The condition on the remaining flag gives
\(H_j=N_j\rtimes G_{j+1}\) and \(U_j=N_j\rtimes U_{j+1}\).
Multiplication adds the \(t\) coordinates and adds the \(u\)
coordinates with a translation depending on the two \(t\)'s.
Thus \(dt\,du\) is Haar. The conjugation by \(l\) sends \(t\)
to \(lt\) and preserves this Haar measure. On \(N_j\) the
simple-root character evaluates the coefficient of \(e\) in
the image of \(f_{j+1}\). It is therefore
\(\exp(2i\operatorname{Im}(-h(t,f_{j+1})))\), up to complex
conjugation according to the character convention. The
real frame normalization makes the simple coefficients one.

Give \(C\) probability measure. Choose Haar measures \(dg_j^*\)
on \(G_j\) by integration along \(g\mapsto g^{-1}f_j\), with
base \(\sigma_j\) and fiber \(dt\,du\,dg_{j+1}^*\) on \(H_j\),
or probability on \(H_a\). These are Haar measures since the
orbit measure is invariant and the fiber group is unimodular.
Give \(U_j\) the successive product of \(dt\,du\). Write
\(d\dot y_j^*\) for the quotient measures. In particular,
at \(j<a\) one can integrate on \(U_j\backslash G_j\) by
integrating first on \(U_{j+1}\backslash G_{j+1}\) and then
on the orbit with measure \(\sigma_j\).

Let \(A,B\) be bounded vector-valued functions, compactly
supported modulo \(U_j\), with Bessel equivariance for \(U_j\)
and the same irreducible unitary representation of \(C\) of
dimension \(d\). We say the functions have radius \(R\) if
their support modulo \(U_j\) is in the image of
\(\{y:\|y\|\le R\}\). Define \(D_{j,x}^{A,B}\,d\sigma_j\) as
the image under \(\dot y\mapsto y^{-1}f_j\) of
\(\langle A(y),B(yx)\rangle d\dot y_j^*\). Put
\[
 E_j(X;A,B)=\int_{\|x\|\le X}\int |D_{j,x}^{A,B}(v)|^2\,
                      d\sigma_j(v)\,dg_j^*(x).
\]

\begin{proposition}[Conditional correlations on a position ball]
\leavevmode\label{prop:recenergy912}
For \(X,R\ge1\) and the functions above,
\begin{equation}
 E_j(X;A,B)\le
    (16r^{3/2}R^5X^2)^{a-j}
    \pi^{\,\sum_{i=j+1}^{a}(2r_i-1)}
    \frac{\|A\|_{2,*}^2\,\|B\|_{2,*}^2}{d},
 \label{eq:recenergy912}
\end{equation}
where both norms use \(d\dot y_j^*\).
\end{proposition}
\begin{proof}
Choose a measurable section \(s_v^{-1}f_j=v\). At \(j=a\),
equivariance under \(C\) makes the integrand constant in the
fiber, and \(D_{a,x}(v)=\langle A(s_v),B(s_v x)\rangle\).
For fixed \(v\), extend the \(x\) integral to \(G_a\), and
translate \(s_v x\) to a Haar variable. For a fixed vector
\(A(s_v)\), averaging the square of the inner product over
\(C\) introduces the factor \(1/d\). Indeed the \(C\)-average
of the rank-one operator \(bb^*\) commutes with the irreducible
representation and has trace \(\|b\|^2\), so it is
\(\|b\|^2 I/d\). The remaining \(v\)-integral is
\(\|A\|_{2,*}^2\) by disintegration and equivariance.
This proves the base case.

For \(j<a\) choose \(s_v=a_vk_v\), with \(k_v\) unitary for
the Euclidean metric, and \(a_v\) scaling \(e,e^*\) by
\(\|v\|^{-1},\|v\|\) and fixing \(V\). Such \(k_v\) exists
because the positive and negative parts of a null vector
have equal length and can be rotated independently.
It can be chosen measurably by extending unit vectors to
orthonormal bases on finitely many coordinate patches. Define
\(A_v(m)=A(m s_v)\) and define \(B_{v'}\) in the same way.
Both have the required residual equivariance. For nonzero
slices, \(R^{-1}\le\|v\|,\|v'\|\le R\), since the isometries
in the norm ball and their inverses have norm at most \(R\).
The residual support also has radius \(R\). In fact, if
\(\|u_0 m a_vk_v\|\le R\) for \(u_0\in U_j\), discard the
right unitary factor. The middle block of \(u_0 m a_v\)
is \(u_1m\) with \(u_1\in U_{j+1}\), and its norm is at most
\(R\). The slice identity
\(\int\|A_v\|_{2,*}^2\,d\sigma_j(v)=\|A\|_{2,*}^2\)
follows directly from the quotient measures.

Hold \(v\) fixed. For \(v'=x^{-1}v\) write
\(s_v x=n(t,u)l s_{v'}\). The measure in \(x\) is
\(dt\,du\,dg_{j+1}^*(l)\,d\sigma_j(v')\). If \(\|x\|\le X\)
and the two slices are nonzero, we have
\[
        \|l\|\le X,\qquad \|t\|\le R^2X,\qquad |u|\le R^2X .
\]
Indeed \(l\) is the middle block of
\(a_v k_v x k_{v'}^{-1}a_{v'}^{-1}\), where the diagonal
factors act identically on the middle block. The full matrix
has norm at most \(R^2X\), and its last column has middle
component \(t\) and top entry with imaginary part \(u\).

At this \(x\), the fiber integral for \(D_{j,x}(v)\) is
\[
 \int_{U_{j+1}\backslash G_{j+1}}
       \langle A_v(m),B(m n(t,u)l s_{v'})\rangle\,d\dot m^* .
\]
Move \(n(t,u)\) to the left of \(m\) by conjugation. Its
character is the wave at \(h(t,m^{-1}f_{j+1})\), with a fixed
sign. Disintegration at \(w=m^{-1}f_{j+1}\) now gives the
Fourier transform of \(D_{j+1,l}^{A_v,B_{v'}}(w)d\sigma_{j+1}(w)\)
at \(t\) or \(-t\). The result is independent of \(u\).
Its support as a density in \(w\) satisfies the annulus
condition of Lemma~\ref{lem:truncquad912} by the radius of
\(A_v\). Each density on compact support is bounded, as in
Corollary~\ref{cor:realgramposition911}.

Enlarge the \(t,u,l\) domain to the three bounds just proved.
The central interval has length \(2R^2X\). Apply the lemma in
dimension \(r_{j+1}\) with \(T=R^2X\), and integrate the
result in \(l,v,v'\). We obtain
\[
 E_j(X;A,B)\le
 16r^{3/2}R^5X^2\,\pi^{2r_{j+1}-1}
 \iint E_{j+1}(X;A_v,B_{v'})\,d\sigma_j(v)\,d\sigma_j(v').
\]
All integrals of squares used in enlarging the domains are
nonnegative. By induction and the slice identities for the
two norms, this proves \eqref{eq:recenergy912}.
\end{proof}

\begin{theorem}[The normalized first real transform on a position ball]
\leavevmode\label{thm:realpositionball912}
Let \(F=F_v^\flat\), \(d=d_v\), and let \(B(x;\nu)\) and \(R\)
be as in Corollary~\ref{cor:realgramposition911}.
For the Haar \(dx\) of Theorem~\ref{thm:realgram682}, put
\[
 S_{k,a}=\pi^{(2a+1)k} V_k/V_r,\qquad
 L_{k,a}(R,X)=2^n\pi^{2a^2+a+1}
        (16r^{3/2})^a R^{2n^7+9a}X^{2a}.
\]
For \(X\ge1\) the actual double correlation satisfies
\begin{equation}
 \frac1d\int_{\|x\|\le X}|B(x;\nu)|\,dx
 \le S_{k,a}\, L_{k,a}(R,X)
              |\nu|^{-K_0}\min(1,|\nu|^{-6}).
 \label{eq:realpositionball912}
\end{equation}
On the shorter interval one may take \(R=b^{6n+2}\)
and \(X=b^{100n^4}\). Then \(L_{k,a}(R,X)\le b^{2000n^9}\)
for all sufficiently large nodes.
\end{theorem}
\begin{proof}
Use the recursive measures first. The angular bound
\eqref{eq:realgramposition911} and Proposition~\ref{prop:recenergy912}
apply with \(j=0\) and \(A=B=F\), keeping the two occurrences
of \(F\) in the Hermitian pairing. The sum of the powers of
\(\pi\) is
\[
    r+\sum_{i=1}^a(2r_i-1)=(2a+1)k+2a^2+a+1 .
\]
The power of \(R\) is \(2(n^7+2a)+5a=2n^7+9a\).

We compute the change of group Haar. For the matrix Gram
density \(d\mu_j=\delta(g^*J_jg-J_j)\,dg_{\rm Leb}\) use real
diagonal and complex upper coordinates for Hermitian matrices.
Over the first column \(e\), its pairing with each middle
column sets that column's bottom entry to zero. Its pairing
with the dual column sets the latter's bottom entry to one.
The zero norm of that dual column determines the real part
of its top entry with Jacobian two, leaving the variables
\(t,u\) above. Each pairing of a middle column with the dual
then fixes the top complex entry of that middle column
with Jacobian one. The middle Gram equations have density
\(d\mu_{j+1}\). These fiber computations hold over every
first column by translation invariance. Consequently at each
null column the product disintegration of \(d\mu_j\) acquires
a factor \(1/2\) relative to \(\sigma_j\,dt\,du\,d\mu_{j+1}\).
At the last column, of norm \(-1\), the orthogonal complement
is compact positive of dimension \(k\). Its fiber has mass
\(V_k/2^k\), as computed in Theorem~\ref{thm:realgram682}.
Haar is unimodular, so the first-column disintegration and
the inverse-column disintegration give the same constants.
That theorem normalizes \(dx=2^r d\mu_0/V_r\). Thus
\(dx=2^{r-a-k}(V_k/V_r)dg_0^*=2^n(V_k/V_r)dg_0^*\).

Any conversion factor between the two measures on
\(U\backslash G\) appears twice in the double correlation and
twice in the square of \(\|F\|_2^2\). It therefore cancels in
the homogeneous inequality. In the original quotient Haar we
have \(\|F\|_2^2\le d\) by Theorem~\ref{thm:normalizedcut699}.
Together with the \(1/d\) from \eqref{eq:recenergy912}, this
proves \eqref{eq:realpositionball912}. For the last assertion,
\(r\ll x^6\) and \(\log b\ge(\log x)^6\) on the shorter
interval, so \(16r^{3/2}\le b\) at large nodes. Also
\(6n+2\le 8n\), \(a<n\), and the logarithm to base \(b\)
of \(2^n\pi^{2a^2+a+1}\) is at most \(10n^2\) there.
Inserting these inequalities in the definition of \(L_{k,a}\)
gives the stated exponent.
\end{proof}

Inequality~\eqref{eq:realpositionball912} uses a bounded
position domain. Its factor \(X^{2a}\) records the lengths
of the central and Fourier variables in the recursion.
A bound for the complementary positions is required for the
full archimedean integral.
\clearpage

\subsection{Common dual vectors and large real positions}
\label{sec:commondirection913}

In the notation of Proposition~\ref{prop:recenergy912}, the
interval and finite places are already those of
Corollary~\ref{cor:firstglobal909}. For example, for each real
place insert
\[
 \frac{1}{d_v S_{k,a}}
 \int_{\|x\|\le b^{100n^4}}
             |B_v(x;c_v\Phi_v y)|\,dx
\]
in place of the corresponding factor of \(A_\infty(y)\) in
\eqref{eq:firstglobal909}. The resulting nonnegative sum is
at most \(\exp(O(\epsilon)) b^{-(7-o(1))kf}\). Indeed
Theorem~\ref{thm:realpositionball912} bounds the product of
the new factors by \(b^{2000n^9f}A_\infty(y)\).
At a node of size \(u\) on the shorter interval we have
\(n\le\sqrt u\) and \(k\asymp u^5\), whence \(n^9=o(k)\).
The real integral over larger positions requires an estimate
for the mass that can pass to the next vector of the flag.

At a null step \(j<a\), and with \(h\) linear in its second
variable, define the positive number
\[
  \Delta_x(v)=
     \max_{\|w\|\le1,\ \|x^{-1}w\|\le1}|h(v,w)|,
       \qquad 0\ne v\in\mathcal C_0,\quad x\in G_j .
\]
The maximum exists by compactness. Scaling down an arbitrary
vector on which the pairing is nonzero shows positivity.
Equivalently one maximizes
\(|h(v,w)|/\max(1,\|x^{-1}w\|)\) on the Euclidean unit sphere.
This last expression gives continuity in \(v,x\).
Throughout this subsection suppose \(k\ge6\) and put
\(z_0=k-5\), \(d_0=k-1\).

\begin{lemma}[A negative moment of the common dual pairing]
\leavevmode\label{lem:commondirection913}
For \(p=k+a-j,\ q=a+1-j,\ r_j=p+q,\ R\ge1\), uniformly in \(x\),
\begin{align}
 \int_{\|v\|\le R}\Delta_x(v)^{-2z_0}\,d\sigma_j(v)
 &\le\frac{\pi^{r_j}\Gamma(d_0-z_0)R^{2(r_j-1-z_0)}}
 {(r_j-1-z_0)\Gamma(q)\Gamma(d_0)\Gamma(p-z_0)} \notag\\
 &\le 6\pi^{r_j}R^{4n+10}.
 \label{eq:commondirection913}
\end{align}
\end{lemma}
\begin{proof}
In positive and negative coordinates there is a complex
\(d_0\)-plane \(P_x\) in the positive space which \(x^{-1}\)
takes to the positive space. In fact the negative projection
of \(x^{-1}\) on the positive space has rank at most \(q\),
and \(p-q=d_0\). On \(P_x\) the isometry \(x^{-1}\) preserves
Euclidean length because the Hermitian and Euclidean norms
agree on both positive spaces. Testing on its unit ball gives
\(\Delta_x(v)\ge\|\operatorname{pr}_{P_x}v\|\).

Write \(v=(s\omega,s\eta)\) in the polar coordinates of
Theorem~\ref{thm:realquadrics911}. The squared norm of the
projection of a uniform \(\omega\in S^{2p-1}\) onto \(P_x\)
has beta density with parameters \(d_0,p-d_0\).
To see this, choose \(p\) independent complex Gaussians with
density \(\pi^{-1}e^{-|z|^2}\) and divide by their vector
norm. The squared norms of the first \(d_0\) and the other
\(p-d_0\) coordinates are independent gamma variables, since
each squared modulus has exponential density. Changing
variables to their sum and their ratio to the sum gives the
beta density. Its moment of order \(-z_0\) equals
\(\Gamma(p)\Gamma(d_0-z_0)/(\Gamma(d_0)\Gamma(p-z_0))\)
by the beta integral. The radial integral up to \(s=R\)
bounds the indicated ball and contributes
\(R^{2(r_j-1-z_0)}/(2(r_j-1-z_0))\). The radial density
constant from Theorem~\ref{thm:realquadrics911} gives the
first inequality. The numerator gamma factor is \(6\), all
denominator gamma arguments are positive integers, and
\(r_j-1-z_0=2(a-j)+5\). This proves the second.
\end{proof}

\begin{proposition}[Size of the residual position and a weighted fiber]
\leavevmode\label{prop:positiondrop913}
Let \(v'=x^{-1}v\) with \(R^{-1}\le\|v\|,\|v'\|\le R\).
Use the same sections as in Proposition~\ref{prop:recenergy912}
and write
\[
        s_v x s_{v'}^{-1}=n(t,u)l,\qquad l\in G_{j+1}.
\]
Then
\begin{equation}
     \|x\|\le 12R^4\,\Delta_x(v)^{-2}\|l\| .
 \label{eq:positiondrop913}
\end{equation}
For fixed \(v,v',l\), put \(r'=r_{j+1}\). For \(Y\ge1\)
and \(T=R^2Y\),
\begin{equation}
 \int_{\|x\|\le Y}\Delta_x(v)^{2z_0}\,dt\,du
       \ \le\ 2(2R)^{2z_0}\pi^{r'}T^{\,1+2(r'-z_0)} .
 \label{eq:weightedfiber913}
\end{equation}
Here \(x=s_v^{-1}n(t,u)l s_{v'}\) varies with the fiber
coordinates, and \(1+2(r'-z_0)\le 4n+5\).
\end{proposition}
\begin{proof}
Choose a common vector \(w\) attaining \(\Delta=\Delta_x(v)\).
The vectors \(w_1=s_{v'}x^{-1}w\) and \(w_2=s_v w\) have
norm at most \(R\), and \(n(t,u)l\) takes \(w_1\) to \(w_2\).
Their bottom coefficients in the order \(e,V,e^*\) coincide
with \(\gamma=h(v,w)\), so \(|\gamma|=\Delta\le R\).
The middle row gives
\[
             (w_2)_V=l(w_1)_V+\gamma t .
\]
It follows that \(\|t\|\) and \(\|l^{-1}t\|\) are both at most
\(2R\|l\|/\Delta\), using \(\|l^{-1}\|=\|l\|\).
The middle-to-top row of \(n(t,u)l\) has norm
\(\|l^{-1}t\|\), since it is the functional \(-h(t,l\,\cdot)\).
Using the top row on \(w_1\) bounds the remaining top right entry
by \(2R/\Delta+2R^2\|l\|/\Delta^2\).
The three diagonal blocks have norms \(1,\|l\|,1\).
Adding the bounds on all blocks gives a norm at most
\(12R^2\|l\|/\Delta^2\). Each outer section has norm and
inverse norm at most \(R\), proving \eqref{eq:positiondrop913}.

For the integral, choose a positive \(d_0\)-plane \(P_l\) for
\(l^{-1}\) in the residual space as in the lemma.
For every \(b\) in that space,
\(\|\operatorname{pr}_{P_l}lb\|\le\|b\|\): pair with a unit
vector of \(P_l\), move \(l^{-1}\) to that vector, and use
Cauchy--Schwarz. Projecting the middle row therefore gives
\(\Delta\|\operatorname{pr}_{P_l}t\|\le2R\). The conditions
\(\|t\|,|u|\le T\) follow from the norm of \(s_v x s_{v'}^{-1}\).
Use the weight bound
\(\Delta^{2z_0}\le(2R)^{2z_0}\|\operatorname{pr}_{P_l}t\|^{-2z_0}\)
almost everywhere. Integrating over the product of two balls
of radius \(T\) in \(P_l\) and its Euclidean complement gives
\[
 \int_{\|t\|\le T}\|\operatorname{pr}_{P_l}t\|^{-2z_0}dt
 \le \frac{\pi^{r'} T^{2(r'-z_0)}}
 {(d_0-z_0)\Gamma(d_0)\Gamma(r'-d_0+1)}.
\]
All factors in the denominator are at least one. The
\(u\)-interval has length \(2T\). Since
\(r'-z_0=2(a-j)+4\) and \(a=n-1\), this proves
\eqref{eq:weightedfiber913}.
\end{proof}
\clearpage

\subsection{The tail of the real position integral}
\label{sec:realremotetail914}

Write \(M_j(x;A,B)=(\int |D_{j,x}^{A,B}|\,d\sigma_j)^2\)
for the functions and recursive measures of
Proposition~\ref{prop:recenergy912}. Absolute values here are
taken after the integration over the stabilizer of the current
frame vector. In particular \(|B(x;\nu)|\le M_0(x;F,F)\)
when both are computed with the recursive quotient measure.

\begin{theorem}[Large positions in the first real Gram transform]
\leavevmode\label{thm:realremotetail914}
Suppose \(a\ge1,\ k\ge 40n^3,\ R\ge\max(16,2r)\), and that
the real cut has radius \(R\). With the original measures of
Theorem~\ref{thm:realpositionball912}, for every \(Y\ge R^2\)
and every complex frequency \(\nu\) one has
\begin{equation}
  \frac1{d_v S_{k,a}}\int_{\|x\|>Y}|B(x;\nu)|\,dx
       \le R^{80nk}\,Y^{-k/(2a)} .
 \label{eq:realremotetail914}
\end{equation}
\end{theorem}
\begin{proof}
We first work on \(Y<\|x\|\le2Y\) in the recursive measures.
Put \(z_0=k-5\) and define
\[
  \lambda=(Y/R^2)^{1/a},\quad Y_j=Y\lambda^{-j},\quad
  \delta_*^2=12 R^4/\lambda,\quad T=2R^2Y.
\]
For \(0\le j\le a\) let \(I_j(A,B)\) be the integral of
\(M_j(x;A,B)\) over \(Y_j<\|x\|\le2Y\) in \(G_j\).
At \(j=a\), \(U_a=1\), so nonzero integrand in a correlation
requires both \(y,yx\) to be of norm at most \(R\).
Then \(\|x\|\le R^2=Y_a\), giving \(I_a=0\).

For \(j<a\) put
\[
 A_j^0=6\pi^{r_j}R^{4n+10},\qquad
 H_j^0=2(2R)^{2z_0}\pi^{r_{j+1}}T^{4n+5},\qquad
 P=16r^{3/2}R^5(2Y)^2 .
\]
Cauchy--Schwarz with \(\Delta_x(v)^{\pm z_0}\) in the integral
defining \(M_j^{1/2}\), followed by Lemma~\ref{lem:commondirection913},
gives
\[
       M_j(x;A,B)\le A_j^0
          \int \Delta_x(v)^{2z_0}|D_{j,x}^{A,B}(v)|^2\,d\sigma_j(v).
\]
The integrand is supported where both \(v\) and \(v'=x^{-1}v\)
have norm in \([1/R,R]\). Integrate in \(x\) and use the
coordinates \(v',t,u,l\) of Proposition~\ref{prop:recenergy912}.
On the part with \(\Delta_x(v)>\delta_*\), the inequality
\eqref{eq:positiondrop913} gives \(\|l\|>Y_j/\lambda=Y_{j+1}\).
Also \(\|l\|\le2Y\). The Fourier identity in that proposition
gives \(|D_{j,x}(v)|^2\le M_{j+1}(l;A_v,B_{v'})\) by the
triangle inequality. This upper bound is independent of \(t,u\).
By \eqref{eq:weightedfiber913} the integral of the weight in
these two variables is at most \(H_j^0\). The contribution to
\(I_j\) is therefore at most
\[
    A_j^0 H_j^0
       \iint I_{j+1}(A_v,B_{v'})\,d\sigma_j(v)\,d\sigma_j(v').
\]

On the part with \(\Delta_x(v)\le\delta_*\), the weight is at
most \(\delta_*^{2z_0}\). Retain the Fourier transform in \(t\).
The integral of its square over \(\|t\|\le T\), integrated also
in \(|u|\le T\), is at most
\[
 P\,\pi^{2r_{j+1}-1}
     \int |D_{j+1,l}^{A_v,B_{v'}}(w)|^2d\sigma_{j+1}(w)
\]
by Lemma~\ref{lem:truncquad912}. Enlarge the residual position
domain to \(\|l\|\le 2Y\), apply \eqref{eq:recenergy912}, and
integrate the slice norms. This bounds the small-pairing
contribution to \(I_j(A,B)\) by
\[
 A_j^0\delta_*^{2z_0} P^{a-j}
    \pi^{\sum_{i=j+1}^a(2r_i-1)}
                \|A\|_{2,*}^2\|B\|_{2,*}^2/d_v .
\]
Here \(A,B\) and all their slices have the given compact type.
In the other contribution we iterate down the flag. At each
iteration the integral of \(\|A_v\|_{2,*}^2\|B_{v'}\|_{2,*}^2\)
is the product of the two unsliced squared norms. Since the
recursion ends with \(I_a=0\), we have proved
\begin{align}
 I_0(A,B)
 &\le \frac{\|A\|_{2,*}^2\|B\|_{2,*}^2}{d_v}\,
 \delta_*^{2z_0}
 \sum_{j=0}^{a-1}
 \left(\prod_{i<j}A_i^0 H_i^0\right)
 A_j^0 P^{a-j}
 \pi^{\sum_{i=j+1}^a(2r_i-1)} .
 \label{eq:realshell914}
\end{align}
The small-pairing variable occurs at just one stage in
each summand. Integrating its Fourier transform before taking
a bound in the remaining variables is what yields the factor
\(\delta_*^{2z_0}\).

For clarity we bound the constants in \eqref{eq:realshell914}.
The exponent of \(\pi\) in the \(j\)-th summand is
\(r+2\sum_{i=1}^a r_i-(a-j)\le (2a+1)k+2n^2\).
The power of \(Y\) outside \(\delta_*^{2z_0}\) is
\((4n+5)j+2(a-j)\le10n^2\).
After extracting their displayed powers of \(\pi,Y\),
the three constants \(A_i^0,H_i^0,P\) are at most
\(R^{4n+11},R^{4k+12n+16},R^{10}\), respectively.
These follow from \(R\ge\max(16,2r)\) and \(T=2R^2Y\).
Finally
\(\delta_*^{2z_0}\le R^{7k}Y^{-z_0/a}\).
For \(n\ge2,\ a=n-1,\ k\ge40n^3\), these bounds, including
the number \(a\) of summands, the factor \(2^n\) in the group
measure and the remaining power \(\pi^{2n^2}\), are bounded by
\(R^{79nk}Y^{-k/(2a)}\) after removal of \(\pi^{(2a+1)k}\).
Indeed the exponent of \(Y\) is at most \(10n^2-(k-5)/a\le-k/(2a)\).
The exponent of \(R\) is at most
\(a(4k+16n+27)+(4n+11)+10a+7k+2n^2+2n\), which is below \(79nk\).

Take \(A=B=F\). The quotient-measure conversion cancels
against the two squared norms exactly as in
Theorem~\ref{thm:realpositionball912}, and
\(dx=2^n(V_k/V_r)dg_0^*\). In the original quotient
\(\|F\|_2^2\le d_v\). Thus the bound on the normalized shell
is \(R^{79nk}Y^{-k/(2a)}\). Apply it to the disjoint shells
with lower endpoints \(2^mY,\ m=0,1,\ldots\). The geometric
sum costs at most two, absorbed in \(R^{80nk}\).
This proves \eqref{eq:realremotetail914}.
\end{proof}
\clearpage

\subsection{The first real transform on the whole position group}
\label{sec:fullrealfirst915}

At a real place \(v\) define the normalized mass at a nonzero
first Gram frequency by
\[
       \mathcal B^\infty_v(\nu)
          =\frac{1}{d_v S_{k,a}}
                     \int_{G_0}|B_v(x;\nu)|\,dx .
\]
Here \(B_v\) uses the cut vector, its original quotient Haar
and the first frame vector in Corollary~\ref{cor:realgramposition911}.
We use group Haar from Theorem~\ref{thm:realgram682}. In particular
\(S_{k,a}=\pi^{(2a+1)k}V_k/V_r\) is displayed separately from
the contribution of the compact type.

\begin{theorem}[All positions at a nonzero real Gram frequency]
\leavevmode\label{thm:fullrealfirst915}
At all sufficiently large nodes on the shorter interval, for
every real place and every \(\nu\in\mathbb C^\times\),
\begin{equation}
       \mathcal B^\infty_v(\nu)
       \le b^{2000n^9}|\nu|^{-K_0}\min(1,|\nu|^{-6}),
       \qquad K_0=k-n^7 .
 \label{eq:fullrealfirst915}
\end{equation}
\end{theorem}
\begin{proof}
For any \(0\le\ell\le r-2\), the proof of
Theorem~\ref{thm:realpositionball912}, using
\eqref{eq:realquadrics911} with this \(\ell\), gives
\begin{align}
 \frac1{d_v S_{k,a}}\int_{\|x\|\le X}|B_v(x;\nu)|\,dx
 \le
  2^n\pi^{2a^2+a+1}(16r^{3/2})^a
  R^{\,2(r-1-\ell)+5a} X^{2a}|\nu|^{-\ell}.
 \label{eq:variableball915}
\end{align}
Indeed only the norm of the first angular operator uses
\(\ell\), and \(r-1-\ell\ge1\) in its denominator.
Put \(R=b^{6n+2}\). On the shorter interval we have \(a\ge1\),
\(n\le\sqrt u,\ k\asymp u^5\) at a node \(u\), and
\(\log b\ge(\log u)^6\). These imply the size hypotheses of
Theorem~\ref{thm:realremotetail914} uniformly. Take
\[
        X=R^{200an}\max(1,|\nu|)^{2a}.
\]
The contribution from \(\|x\|>X\) is at most
\(R^{-20nk}\max(1,|\nu|)^{-k}\) by \eqref{eq:realremotetail914}.
It is therefore at most
\(|\nu|^{-K_0}\min(1,|\nu|^{-6})\), since \(n^7\ge6\).

If \(0<|\nu|\le1\), choose \(\ell=K_0\) in
\eqref{eq:variableball915}. If \(|\nu|>1\), choose
\(\ell=K_0+6+4a^2\). In the latter case \(4a^2\) pays the
power of \(|\nu|\) in \(X^{2a}\).
Both choices lie in \([0,r-2]\), since \(n^7\ge6+4a^2+1\)
for \(n=a+1\ge2\). The constant left from the ball is at most
\[
    2^n\pi^{2a^2+a+1}(16r^{3/2})^a
                   R^{2n^7+9a+400a^2n}.
\]
At large nodes \(16r^{3/2}\le b\), \(6n+2\le 8n\), and
\(\log_b(2^n\pi^{2a^2+a+1})\le10n^2\).
Thus the logarithm to base \(b\) of the last display is at
most \(10n^2+a+16n^8+72n^2+3200n^4<1999n^9\) for \(n\ge2\).
Adding the outer contribution proves \eqref{eq:fullrealfirst915}.
\end{proof}

\begin{corollary}[The rational first-frequency sum with its real integrals]
\leavevmode\label{cor:adelicfirst915}
Use the fields, the marks \(\mathcal W=\Xi\sqcup\mathcal V\),
the unmarked finite places \(\mathcal U\), and the scalars
\(c_v=2\pi i/(\Phi_v(\delta_E)b_{1,v})\) of
Corollary~\ref{cor:firstglobal909}. Then
\begin{align}
 \sum_{\substack{0\ne y\in E\\ y\ {\rm integral\ above}\ \mathcal V}}
  &\left(\prod_{v\mid\infty}\mathcal B^\infty_v(c_v\Phi_v y)\right)
   \left(\prod_{w\in\Xi}
     \frac{B_0 q_w^{n-j_w}\mathcal B_{Z_w}(y)}{Z_w^0}\right)
   \prod_{w\in\mathcal U}\mathcal B_w^0(y)
   \ \le\ \exp(O(\epsilon))\,b^{-7kf+2000n^9 f}.
 \label{eq:adelicfirst915}
\end{align}
In particular the exponent of \(b\) on the right is
\(-(7-o(1))kf\) along the shorter interval.
\end{corollary}
\begin{proof}
For each \(y\ne0\) all real frequencies \(c_v\Phi_v y\)
are nonzero. Multiply \eqref{eq:fullrealfirst915} over the
\(f\) real places of the totally real base. The product is
bounded by \(b^{2000n^9f}A_\infty(y)\), with exactly the
archimedean arguments of \eqref{eq:firstglobal909}. Multiplying
by its nonnegative finite factors and summing proves
\eqref{eq:adelicfirst915}. The unmarked Euler products and the
sum over denominators are the convergent ones justified in
that corollary. Since \(n^9/k\ll u^{-1/2}\), the last
assertion follows uniformly.
\end{proof}

Formula~\eqref{eq:adelicfirst915} controls the first relative
entry with no earlier entries prescribed. In the integral
theta-period problem one must also treat subsequent entries
with the previous relative frame products prescribed, sum
their coupled rational contributions, and compare the resulting
signed integral period to the coefficient lattice. The
strict cohomology estimate in Section~\ref{sec:missing} is a
further requirement for the deduction of Dirichlet GRH.
\clearpage

\subsection{The real identity mass and the first pivot density}
\label{sec:realcutmass916}

At a real place let \(F=F_v\) be the full test and \(W=F_v^\flat\)
its cut. Write \(d=d_v,\ D=d_{\pi,v}\) and
\(C_Z(x)=\langle Z,R(x)Z\rangle\), where \(R(x)\) acts by
right translation. At frequency zero the double integral is
\(B_v(x;0)=|C_W(x)|^2\). The Haar measures are those of
Theorem~\ref{thm:fullrealfirst915}.

\begin{theorem}[The zero-frequency mass with both tests cut]
\leavevmode\label{thm:realcutmass916}
Uniformly at all real places on the shorter interval, for
sufficiently large nodes,
\begin{equation}
 (1-b^{-r^3 f})^4
 \le \frac D{d^2}\int_G |C_W(x)|^2\,dx
 \le 1+b^{-20}.
 \label{eq:realcutmass916}
\end{equation}
Consequently the product of the real masses equals
\((d_B^2/D_\pi)(1+O(f b^{-20}))\), where
\(d_B=\prod_v d_v\) and \(D_\pi=\prod_vd_{\pi,v}\).
\end{theorem}
\begin{proof}
Put \(R=b^{6n+2}\) and \(Y=R^{400an}\).
First \(S_{k,a}D/d\le b^{12nr}\) at large nodes.
Indeed in the compact-dual formal degree formula of
Theorem~\ref{thm:norm91} the first \(k-1\) consecutive
entries of Proposition~\ref{prop:stationaryweights695}
cancel the corresponding Vandermonde factors. At most
\(2nr\) pairs have an entry among the remaining \(2n\).
Their numerator gaps are at most \(b^3\), since the entries
are \(-\kappa-2jr^8\) for \(1\le j\le2n\), the other entries
have size at most \(k\), and \(\kappa\le ub\) at a node of
size \(u\). Each remaining denominator is at least one.
Thus \(D\le b^{6nr}\). Also
\(V_k/V_r=\prod_{j=k+1}^r(j-1)!/(2\pi^j)
 \le r^{(2a+1)r}\). With \(r,\pi\le b\) and \(d\ge1\)
this proves the claimed bound. Applying
Theorem~\ref{thm:realremotetail914} with \(\nu=0\) gives
\[
 \frac D{d^2}\int_{\|x\|>Y}|C_W(x)|^2\,dx
       \le b^{12nr}R^{-120nk}\le b^{-40}
\]
eventually, uniformly in the place, since \(r\le2k\).

On the ball the \(L^2\) error for the test vectors gives the
pointwise bound \(|C_W-C_F|\le2d\,b^{-r^3f}\):
expand the difference in its two arguments and apply
Cauchy--Schwarz and Theorem~\ref{thm:normalizedcut699}.
Lemma~\ref{lem:cosetpacking699} gives
\(D\operatorname{vol}(\|x\|\le Y)\le b^{38r^2}Y^{9r^2}\)
by bounding the integrand of its second inequality from below.
As \(\log_bY=400an(6n+2)=O(n^3)\) and \(n^3=o(r)\),
\[
 \frac{\sqrt D}{d}
 \|C_W-C_F\|_{L^2(\|x\|\le Y)}
        \le 2b^{-r^3 f+(38r^2+9r^2\log_bY)/2}\le b^{-40}.
\]
By Proposition~\ref{prop:averagedschur684} the full
coefficient is \(d\) times a unit discrete-series coefficient,
so \(\|C_F\|_2=d/\sqrt D\) by Schur orthogonality. The
triangle inequality on the ball followed by addition of the
outer squared norm now gives the stated upper bound.

For the lower bound the one-place identity in the proof of
Theorem~\ref{thm:mixedmass876} gives
\(\int\overline{C_F}C_W=(d/D)\|P_v W\|^2\), where \(P_v\)
is the coefficient projection fixing \(F\).
It is a contraction, so
\(\|P_v W\|\ge (1-b^{-r^3 f})\sqrt d\). Applying
Cauchy--Schwarz to that identity gives the lower bound in
\eqref{eq:realcutmass916}. The product assertion follows by
multiplication, since \(f\ll u^{12}\) and \(f b^{-20}=o(1)\).
\end{proof}

\begin{proposition}[Inversion in the first scalar Gram entry]
\leavevmode\label{prop:realpivotdensity916}
Use the density \(D_x\,d\sigma_0\) for \(W\) from
Corollary~\ref{cor:realgramposition911}, with its original
quotient normalization. There is a compactly supported
\(C^{K_0+3}\) function \(g_v\) on \(\mathbb C\), supported on
\(|c|\le R^2\), such that for every bounded Borel function
\(\varphi\),
\begin{align}
 \int_{\mathbb C}\varphi(c)g_v(c)\,d^2c
  &=\int_G\iint
     \varphi(h(v_1,v_2))D_x(v_1)\overline{D_x(v_2)}
                      \,d\sigma_0(v_1)\,d\sigma_0(v_2)\,dx .
 \label{eq:realpivotdensity916}
\end{align}
Here \(d^2c\) is real planar Lebesgue measure. At every \(c\)
one has the absolutely convergent formula
\[
 g_v(c)=\pi^{-2}\int_{\mathbb C}
     e^{-2i\operatorname{Im}(\nu c)}
       \left(\int_G B_v(x;\nu)\,dx\right)\,d^2\nu.
\]
Its integral is the mass in \eqref{eq:realcutmass916} before
normalization by \(D/d^2\).
\end{proposition}
\begin{proof}
On a fixed ball in \(G\), the integral of
\((\int |D_x|\,d\sigma_0)^2\) is finite by Cauchy--Schwarz
on the finite measure part \(\|v_1\|\le R\) of the cone
and Proposition~\ref{prop:recenergy912}. On the complement
the proof of Theorem~\ref{thm:realremotetail914} bounds this
very same squared integral before insertion of the character.
Changing quotient Haar multiplies the bound by a finite
constant. The right side of \eqref{eq:realpivotdensity916}
therefore defines a finite complex measure \(\mu\). It is
supported on the claimed disk since
\(|h(v_1,v_2)|\le\|v_1\|\|v_2\|\le R^2\).

With Fourier convention \(e^{2i\operatorname{Im}(\nu c)}\),
Fubini gives \(\widehat\mu(\nu)=\int_G B_v(x;\nu)\,dx\).
It is bounded on \(|\nu|\le1\) by the total variation of
\(\mu\). For \(|\nu|>1\), Theorem~\ref{thm:fullrealfirst915}
gives decay of order \(|\nu|^{-K_0-6}\). The transform
multiplied by \((1+|\nu|)^{K_0+3}\) is integrable in two
real dimensions. Its inverse Fourier integral therefore
defines a \(C^{K_0+3}\) function by differentiation under the
integral sign. The inverse factor is \(\pi^{-2}\) because the
bilinear form \(\operatorname{Im}(\nu c)\) differs from the
real dot product by an orthogonal change of coordinates and
the exponent has coefficient two. Applying Euclidean Fourier
inversion to compactly supported smooth test functions and
using Fubini identifies this function's density with \(\mu\).
It vanishes off the disk by continuity. Finally at \(\nu=0\)
the transform equals \(\int |C_W|^2\), as asserted.
\end{proof}
\clearpage

\subsection{The second Gram entry at a prescribed first pivot}
\label{sec:conditionedpivot917}

Use \(e=f_0\), \(f=f_1\), the residual space \(V\) of
dimension \(r_1=r-2\), and the cut \(W=F_v^\flat\).
Here \(f\) is the next null frame vector if \(a>1\), and
the terminal vector of norm \(-1\) if \(a=1\). For two
quotient points with representatives \(y_1,y_2\) write
\[
 v=y_1^{-1}e,\quad w=y_2^{-1}e,\quad
 A=y_1^{-1}f,\quad B=y_2^{-1}f,\quad c=h(v,w).
\]
For \(c\ne0\) the second successive Gram pivot is
\[
       P_2=h(A,B)-\frac{h(A,w)h(v,B)}{c}.
\]
This is the determinant of the leading two-by-two relative
Gram block divided by its first entry. Adding any multiple of
\(v\) to \(A\), or of \(w\) to \(B\), preserves \(P_2\).
Consequently \(P_2\) descends to the two quotient points:
the Bessel unipotent fixes \(e\) and fixes \(f\) modulo \(e\).

For \(x\in G\) put \(K_x(y)=\langle W(y),W(yx)\rangle\).
Keep the sections \(s_v^{-1}e=v\) from
Proposition~\ref{prop:recenergy912}. Let \(\eta\) be defined
by \(d\dot y=\eta\,d\dot y_0^*\), where the star denotes the
recursive measures there. Given any section value \(s\in G\),
write \(D_{s,x}(z)\,d\sigma_1(z)\) for the image of
\(\eta K_x(ms)\,d\dot m_1^*\) under \(m\mapsto m^{-1}f\)
on \(U_1\backslash G_1\). This convention includes the full
conversion of the original quotient Haar in each branch.

\begin{theorem}[Disintegration with one nonzero real pivot prescribed]
\leavevmode\label{thm:conditionedpivot917}
For \(c\in\mathbb C^\times\), let \(w_c\) act by
\(e\mapsto c e^*\), \(e^*\mapsto e/\bar c\) and identically
on \(V\). For \(t\in V,\ u\in\mathbb R\) write \(n=n(t,u)\)
as in Section~\ref{sec:realpositionball912} and put
\(s'=w_c^{-1}n^{-1}s_v\).
For fixed \(x,\nu\), the image under the first pivot of
\[
  e^{2i\operatorname{Im}(\nu P_2)}
     K_x(y_1)\overline{K_x(y_2)}\,
                        d\dot y_1\,d\dot y_2
\]
on \(\mathbb C^\times\) has the following continuous density
relative to \(d^2c\):
\begin{align}
 J_{x,c}(\nu)
 &=\frac{|c|^{2r_1}}{2}\int d\sigma_0(v)
    \int_{V\times\mathbb R}dt\,du
    \iint e^{2i\operatorname{Im}(\nu h(z,z'))}
       D_{s_v,x}(z)\overline{D_{s',x}(z')}
                              \,d\sigma_1(z)\,d\sigma_1(z').
 \label{eq:conditionedpivot917}
\end{align}
For each \(c\ne0\) the integrals displayed at fixed \(x\)
converge absolutely. At \(\nu=0\) one has
\begin{equation}
       g_v(c)=\int_G J_{x,c}(0)\,dx,\qquad
                   \int_G|J_{x,c}(0)|\,dx<\infty ,
 \label{eq:conditionedzero917}
\end{equation}
where \(g_v\) is the real density of
Proposition~\ref{prop:realpivotdensity916}.
\end{theorem}
\begin{proof}
Disintegrate both copies of \(U\backslash G\) over their first
vectors. At fixed \(v\), integrate \(s_vw\) in coordinates
\(\alpha e+t_0+c e^*\) on its null orbit. The condition is
\(2\operatorname{Re}(\bar\alpha c)+h(t_0,t_0)=0\).
For \(c\ne0\) write \(\alpha=c(s+iu)\).
Lebesgue measure in \(\alpha\) is \(|c|^2ds\,du\) and
integration of this equation in \(s\) divides by \(2|c|^2\).
Setting \(t_0=ct\) gives exactly
\[
 s_vw=n(t,u)c e^*,\qquad
     d\sigma_0(w)=\tfrac12 |c|^{2r_1}\,d^2c\,dt\,du .
\]
The section \(s'\) sends this \(w\) to \(e\), since \(w_c\)
is an isometry with \(w_c e=c e^*\).
Using \(s'\) instead of any other section gives the same
conditional Haar: the change is a left factor in \(H_0\),
whose unipotent part cancels in \(K_x\), and whose residual
part acts by right translation of \(d\dot m_1^*\).

Take the quotient representatives \(y_1=m s_v,\ y_2=m's'\)
and put \(z=m^{-1}f,\ z'=(m')^{-1}f\). Upon applying \(s_v\)
to all four columns they become
\(e,\ n c e^*,\ z,\ n z'\), respectively in the order
\(v,w,A,B\). We have \(nz'=z'-h(t,z')e\), so \(h(e,nz')=0\)
and \(h(z,nz')=h(z,z')\). The formula for \(P_2\) therefore
reduces to \(h(z,z')\). Further disintegration in \(z,z'\)
with the prescribed two quotient constants proves
\eqref{eq:conditionedpivot917} as a density.

Here are the support and regularity details needed for pointwise
use. A nonzero integrand requires \(R^{-1}\le\|v\|,\|w\|\le R\).
Since \(\|s_v\|\le R\), the middle and top coordinates of
\(s_vw=c n(t,u)e^*\) then give
\(\|t\|,|u|\le R^2/|c|\). Also \(|c|\le R^2\).
On compact subsets of \(0<|c|<\infty\) the sections under
consideration and their inverses are uniformly bounded.
For such a section \(s\), nonzero \(W(ms)\) requires
\(\|u_0 m s\|\le R\) for some \(u_0\in U\). Taking the
residual middle block of \(u_0m\) shows that \(m\) has
residual quotient radius at most \(R\|s^{-1}\|\).
Thus all the integrals run over compact sets in the orbit,
fiber variables and residual quotients, locally uniformly in
\(c\ne0\). The integrands in the residual quotient before
disintegrating \(z,z'\) are continuous in \(c\) for each
choice of \(v,t,u,m,m'\). Boundedness and compact support
there give the claimed absolute convergence and continuity
by dominated convergence.

At fixed \(x\) and \(\nu=0\) the measure in the first pivot
is the one with Fourier transform \(B_v(x;\xi)\). Its
Fourier inverse is continuous, since \eqref{eq:realgramposition911}
is integrable at \(|\xi|>1\), and at \(|\xi|\le1\) its
modulus is at most \((\int |D_x|\,d\sigma_0)^2\).
Here \(D_x\) is the first-vector density defined in
Corollary~\ref{cor:realgramposition911}. The fiber \(c=0\)
has zero measure for fixed \(v\): off \(t_0=0\) one may
solve the null equation in a nonzero residual real coordinate
with nonzero derivative, leaving \(c\) free, and the remaining
slice at \(c=t_0=0\) lies on a complex line of zero orbit
measure. This proves the Fourier identification across the
whole plane. Continuity identifies the inverse with
\(J_{x,c}(0)\) at every nonzero \(c\). Finally
\(\int_G\int_{\mathbb C}|B_v(x;\xi)|\,d^2\xi\,dx<\infty\):
for \(|\xi|\le1\) use the integrable squared bound in the
proof of Proposition~\ref{prop:realpivotdensity916}, and for
\(|\xi|>1\) apply \eqref{eq:fullrealfirst915}. Inversion and
Fubini give \eqref{eq:conditionedzero917} with the asserted
absolute integrability in \(x\).
\end{proof}

\begin{corollary}[The second real wave at a fixed position]
\leavevmode\label{cor:conditionedwave917}
For \(0<|c|\le R^2\) set \(R_c=8R^8\max(1,|c|^{-1})^3\).
Define \(E_x(c)\) to be \(|c|^{2r_1}/2\) times the integral
in \(v,t,u\) in \eqref{eq:conditionedpivot917} with the inner
pairing replaced by
\(\|D_{s_v,x}\|_{L^2(\sigma_1)}\|D_{s',x}\|_{L^2(\sigma_1)}\).
For every \(x\), \(E_x(c)<\infty\), and for \(\nu\ne0\) and
\(0\le\ell\le r_1-2\) one has
\[
       |J_{x,c}(\nu)|\le
       \frac{\pi^{r_1}R_c^{\,2(r_1-1-\ell)}}{r_1-1-\ell}
                    |\nu|^{-\ell}E_x(c).
\]
\end{corollary}
\begin{proof}
Put \(H=\max(1,|c|^{-1})\). On the relevant support,
\(\|n^{\pm1}\|\le 3+2\|t\|+\|t\|^2/2+|u|
                  \le 8R^4H^2\)
by its block matrix. The Weyl element and its inverse have
norm at most \(R^2H\). Thus
\(R\max(\|s_v^{-1}\|,\|(s')^{-1}\|)\le R_c\).
The residual radius observation in the theorem bounds both
orbit supports by \(R_c\). Pushing forward bounded continuous
densities on the indicated compact subsets by the residual
orbit submersion gives bounded densities there, uniformly
for the bounded sections at fixed \(c,x\). The residual
first vector has norm at least \(R_c^{-1}\) on the support
by the same inverse norm bound. This ensures a compact
subset of the nonzero orbit even for the null quadric.
These bounds also show directly that \(E_x(c)<\infty\).
Apply \eqref{eq:realquadrics911} on this residual orbit to
the inner pairing and integrate the resulting inequality.
\end{proof}

For a nonzero second frequency the conditional energy
\(E_x(c)\) still has to be controlled as \(x\) varies, with
sufficient dependence on the prescribed pivot \(c\) for
the rational sum.
\clearpage

\subsection{All positions with one real pivot prescribed}
\label{sec:secondallpositions918}

At a fixed nonzero first pivot the stabilizer of the two null
vectors acts on their common perpendicular space. Using this
stabilizer lets us integrate the position variable in the
second wave of \eqref{eq:conditionedpivot917}. Write
\(W_0(\nu)=|\nu|^{-K_0}\min(1,|\nu|^{-6})\).

\begin{lemma}[A mixed estimate on the residual group]
\leavevmode\label{lem:mixedresidual918}
Let \(A,B,C,D'\) on \(G_1\) be bounded functions with quotient
radius \(L\ge\max(16,2r)\), with the residual Bessel character
and the same irreducible compact type of dimension \(d\).
Use recursive measures, and for \(l\in G_1\) let
\[
 Q(l;\nu)=\iint e^{2i\operatorname{Im}(\nu h(z,z'))}
      D_{1,l}^{A,B}(z)\overline{D_{1,l}^{C,D'}(z')}\,
                                    d\sigma_1(z)d\sigma_1(z').
\]
At large nodes of the shorter interval,
\begin{equation}
 \int_{G_1}|Q(l;\nu)|\,dg_1^*(l)
 \le \pi^{(2a-1)k}L^{40n^7} W_0(\nu)
        \frac{\|A\|_{2,*}\|B\|_{2,*}\|C\|_{2,*}\|D'\|_{2,*}}{d}
 \quad(\nu\ne0).
 \label{eq:mixedresidual918}
\end{equation}
\end{lemma}
\begin{proof}
Set \(b'=a-1,\ r'=k+2b'+1\). On a ball of radius \(X\)
apply \eqref{eq:realquadrics911} and then Cauchy--Schwarz
in \(l\) to the two \(L^2(\sigma_1)\) norms. Apply
\eqref{eq:recenergy912} to each of the squared integrals.
Apart from the product of the four norms divided by \(d\)
this gives
\[
 \pi^{(2b'+1)k+2b'^2+b'+1}(16r^{3/2})^{b'}
    L^{2(r'-1-\ell)+5b'}X^{2b'}|\nu|^{-\ell}.
\]
For \(b'>0\), on the complement use
\(|Q|\le\sqrt{M_1(l;A,B)M_1(l;C,D')}\). Cauchy--Schwarz
and the shell estimate \eqref{eq:realshell914}, starting at
\(G_1\) with \(b'\) null steps, give at most
\(\pi^{(2b'+1)k}L^{80(b'+1)k}X^{-k/(2b')}\) times the
same norm factor. Here the proof of
Theorem~\ref{thm:realremotetail914} through the dyadic
summation uses only radius, equivariance and squared norms.
With recursive Haar there is no final group-measure conversion.
Take \(X=L^{200 b'n}\max(1,|\nu|)^{2b'}\).
The outer bound is at most the norm factor times
\(\pi^{(2b'+1)k}W_0(\nu)\). If \(b'=0\), positions of
norm greater than \(L^2\) contribute zero since \(U_1=1\),
and the ball bound is independent of \(X\).

Take \(\ell=K_0\) for \(|\nu|\le1\) and
\(\ell=K_0+6+4b'^2\) for \(|\nu|>1\).
Both are at most \(r'-2\), even for \(b'=0\), since
\(n=a+1\ge2\). As in \eqref{eq:variableball915}, the
power \(4b'^2\) pays the frequency contribution of \(X^{2b'}\).
The remaining power of \(L\) before the numerical constants
is at most \(2n^7+9b'+400b'^2 n\). Since
\(\pi\le L,\ 16r^{3/2}\le L^3,\ b'<n\) and \(n\ge2\),
this exponent plus \(2b'^2+4b'+2\) is at most \(40n^7\).
The additional unit in the exponent absorbs addition of the
outer bound. This proves the lemma.
\end{proof}

\begin{theorem}[The second frequency at a prescribed nonzero pivot]
\leavevmode\label{thm:secondallpositions918}
With \(R=b^{6n+2}\) and \(R_c=8R^8\max(1,|c|^{-1})^3\), the
conditional real transform of Theorem~\ref{thm:conditionedpivot917}
satisfies, for \(0<|c|\le R^2,\ \nu\ne0\),
\begin{equation}
 \frac{1}{d_vS_{k,a}}\int_G |J_{x,c}(\nu)|\,dx
 \le
 \frac{2^{n+1}\pi^{4a-2}}{(r_1!)^2}
    R_c^{40n^7} R^{8r_1+4}|c|^{-2r_1-2} W_0(\nu).
 \label{eq:secondallpositions918}
\end{equation}
For \(|c|>R^2\) the integrand vanishes.
\end{theorem}
\begin{proof}
For an ordered pair \(P=(v,w)\) in the fiber \(h(v,w)=c\),
put
\[
 d\beta_c(P)=d\sigma_0(v)\,dt\,du,\qquad
 s_P=n(t,u)^{-1}s_v .
\]
The section sends the pair to \((e,c e^*)\). Its stabilizer
fixes both \(e,e^*\), hence acts by an arbitrary isometry
of \(V\), giving precisely \(G_1\). At a second pair
\(P'=x^{-1}P\) write \(s_Px=l s_{P'}\) with \(l\in G_1\).
For fixed \(P\) the position Haar is
\[
                 dg_0^*(x)=dg_1^*(l)\,d\beta_c(P').
\]
Indeed disintegrate \(s_P x=n(t_0,u_0)l s_{v'}\) over
its inverse first vector \(v'\). The coordinates of \(P'\)
are then \(t'=-l^{-1}t_0,\ u'=-u_0\), since
\(n(t_0,u_0)^{-1}=n(-t_0,-u_0)\). This change of variables
preserves Lebesgue measure for fixed \(l\). It also gives
\(s_P x=l n(t',u')^{-1}s_{v'}=l s_{P'}\).

Define two functions on \(G_1\) by
\(A_P(m)=W(m s_P)\) and \(C_P(m)=W(m w_c^{-1}s_P)\).
Use the same definitions at \(P'\). For both branches the
residual position is \(l\), because \(w_c\) commutes with
\(G_1\). We can use \(s_P\) as the section for the first
branch in \eqref{eq:conditionedpivot917}: it differs from
\(s_v\) by \(n^{-1}\) on the left, and conjugation of this
element by \(m\) contributes the same unit character to the
two vectors in \(K_x\). In particular the first residual
pairing remains the one for \(h(z,z')\). Both residual
functions have radius at most \(R_c\) whenever both
branches are nonzero. This follows by the middle-block
argument for sections in Corollary~\ref{cor:conditionedwave917},
using \(\|s_P\|\le8R^5\max(1,|c|^{-1})^2\) and the bound
there for \(\|w_c^{-1}s_P\|\).

Let \(H(v)=\|m\mapsto W(m s_v)\|_{2,*}\). This norm is
independent of the chosen first-vector section: a change on
the left by \(n_0 l_0\in H_0\) contributes a character of
modulus one and the right translation \(m\mapsto m l_0\)
of quotient Haar. Hence \(\|A_P\|_{2,*}=H(v)\) and
\(\|C_P\|_{2,*}=H(w)\).
The disintegration of \(U\backslash G\) gives
\(\int H(v)^2d\sigma_0(v)=\|W\|_{2,*}^2\).
Let \(\mathcal A=\{v:R^{-1}\le\|v\|\le R\}\), containing
the support of \(H\). At each \(v\in\mathcal A\),
\[
 \int_{w\in\mathcal A}dt\,du\le
  D_c:=\frac{2\pi^{r_1}}{r_1!}
                  (R^2/|c|)^{2r_1+1}
\]
by the ball and central interval bounds in
Theorem~\ref{thm:conditionedpivot917}. The same bound applies
with the two vectors exchanged. To see the equality of the
two pair measures under exchange, the preceding Haar formula
identifies \(\beta_c\) with the quotient by the pair stabilizer
\(G_1\) using its fixed Haar. Exchanging the pair changes the
base value \(c\) to \(\bar c\), and \(w_c\) maps
\((e,\bar c e^*)\) to \((c e^*,e)\). It centralizes
\(G_1\), so left Haar invariance and disintegration identify
the exchanged measure with \(\beta_{\bar c}\).
Cauchy--Schwarz on \(\mathcal A\times\mathcal A\) in this
fiber therefore gives
\[
         \int H(v)H(w)d\beta_c(P)\le D_c\|W\|_{2,*}^2 .
\]

Apply the triangle inequality to the pair integral, and
make the position change of variables above. In recursive
quotient measures the result is bounded, by Tonelli and
Lemma~\ref{lem:mixedresidual918}, by
\[
 \frac{|c|^{2r_1}}{2}
 \,\pi^{(2a-1)k} R_c^{40n^7} W_0(\nu)
 \,\frac1{d_v}
     \left(\int H(v)H(w)d\beta_c(P)\right)^2.
\]
Here the four residual functions are
\(A_P,A_{P'},C_P,C_{P'}\), in that order. For original
quotient measures the two factors \(\eta\) give
\(\eta^2\|W\|_{2,*}^4=\|W\|_2^4\le d_v^2\).
Position Haar contributes \(2^n V_k/V_r\). Dividing by
\(d_v S_{k,a}\), inserting \(D_c^2\), and using
\(r_1-k=2a-1\) gives \eqref{eq:secondallpositions918}.
The support assertion follows from
\(|h(v,w)|\le\|v\|\|w\|\).
\end{proof}

The estimate is finite for each prescribed \(c\ne0\).
The factor \(|c|^{-2r_1-2}\), with \(r_1=k+2a-1\), is a loss
in summing over small rational pivots. Controlling that sum
requires more than the pointwise bound above.
\clearpage

\subsection{Common dual pairings on a prescribed Gram fiber}
\label{sec:crossfiber919}

Let \(e,V,e^*\) be the decomposition at the first null vector.
Set \(d_1=\dim_{\mathbb C}V=r_1=k+2a-1\),
\(z_0=k-5\), \(m=2(d_1-z_0)+1=4a+9\), and
\(\mathcal A_R=\{v\in\mathcal C_0:R^{-1}\le\|v\|\le R\}\).
Assume \(a\ge1,\ k\ge 10,\ R\ge1\). For \(p\ne0\) let
\(d\mu_p(v,w)\) be the measure on \(h(v,w)=p\) for which
\[
  d\sigma_0(v)d\sigma_0(w)=d^2p\,d\mu_p(v,w),\qquad
  d\mu_p(v,w)=\tfrac12 |p|^{2d_1}\,d\beta_p(v,w)
\]
in the explicit coordinates of Theorem~\ref{thm:conditionedpivot917}.
The formula on the right specifies the measure at every nonzero
\(p\). Write
\(\Delta_x(v)=\max_{\|b\|,\|x^{-1}b\|\le1}|h(v,b)|\).

\begin{theorem}[Crossed weights on the first Gram fiber]
\leavevmode\label{thm:crossfiber919}
For nonnegative measurable functions \(A,B\) on the null cone
supported in \(\mathcal A_R\), and any position \(x\),
\begin{align}
 \int A(v)B(w)\,d\mu_p(v,w)
 &\le
 \frac{\pi^{d_1}R^{m+1}}{|p|}
 \left(\int \Delta_x(v)^{2z_0} A(v)^2\,d\sigma_0(v)\right)^{1/2}
 \left(\int \Delta_x(w)^{2z_0} B(w)^2\,d\sigma_0(w)\right)^{1/2}.
 \label{eq:crossfiber919}
\end{align}
If the left integral is restricted to a set of pairs, each
function on the right can also be restricted to any measurable
set containing the corresponding projection of those pairs.
\end{theorem}
\begin{proof}
Fix \(v\in\mathcal A_R\). Write \(s_v=a_vk_v\) for the section
of Proposition~\ref{prop:recenergy912}. In
\(s_vw=\alpha e+y+p e^*\) one has \(y=pt\) and
\(\alpha=p(-h(t,t)/2+iu)\). Conditional on \(v\), the
measure \(\mu_p\) is \(\tfrac12\,dy\,du\). On the support
\(\|w\|\le R\), we have \(\|y\|\le R\): the middle block of
\(a_v\) is the identity and \(k_v\) is Euclidean unitary.
Also \(|\alpha|\le R^2\), whence \(|u|\le R^2/|p|\).

Let \(P_x\) be a positive \((k-1)\)-plane taken into the
positive space by \(x^{-1}\), as constructed in
Lemma~\ref{lem:commondirection913}. Since \(k_v\) preserves
the positive and negative spaces, \(k_vP_x\) has intersection
of dimension at least \(k-2\) with the positive part of
\(V\): removal of the hyperbolic pair removes only one positive
direction. Choose a \((k-2)\)-plane \(S\) in that intersection.
For every unit vector in \(k_v^{-1}S\), both constraints in
the maximum defining \(\Delta_x\) hold. Pairing \(w\) with
these vectors gives
\(\Delta_x(w)\ge\|\operatorname{pr}_S y\|\).
With \(s=k-2>z_0\), polar integration yields
\[
 \int_{\|y\|\le R}\|\operatorname{pr}_S y\|^{-2z_0}\,dy
 \le\frac{\pi^{d_1}R^{2(d_1-z_0)}}
 {(s-z_0)\Gamma(s)\Gamma(d_1-s+1)} .
\]
Indeed enlarge the ball to the product of the radius \(R\)
balls in \(S\) and \(S^\perp\), and integrate the \(S\)-radius
with exponent \(2s-1-2z_0\). The denominator is at least one.
Including the \(u\)-interval proves that the conditional
integral of \(\Delta_x(w)^{-2z_0}\) over \(w\in\mathcal A_R\)
is at most \(\pi^{d_1}R^{m+1}/|p|\), since \(m+1=2(d_1-z_0)+2\).
The exchanged estimate holds at each \(p\), by the symmetry
of \(\beta_p\) with \(\beta_{\bar p}\) in the proof of
Theorem~\ref{thm:secondallpositions918}.

In the left side of \eqref{eq:crossfiber919} multiply \(A(v)\)
by \(\Delta_x(v)^{z_0}/\Delta_x(w)^{z_0}\) and \(B(w)\)
by its reciprocal. Cauchy--Schwarz and the two conditional
estimates give the result. In the restricted case one may
first insert the two projection cutoffs and then enlarge the
domain, since the integrands are nonnegative.
\end{proof}

To allow different first entries at the two endpoints, let
\(d_\lambda=\operatorname{diag}(\lambda,I,\bar\lambda^{-1})\)
and \(M=\max(|\lambda|,|\lambda|^{-1})\), for
\(\lambda\in\mathbb C^\times\). If \(v\) and
\(v'=\lambda^{-1}x^{-1}v\) are in \(\mathcal A_R\), write
\[
             d_\lambda s_v x s_{v'}^{-1}=n(t,u)l,
                      \qquad l\in G_1 .
\]
The matrix on the left fixes \(e\), which proves that these
coordinates exist. For fixed \(v\), left invariance of Haar
and the first-vector disintegration give
\(dg_0^*(x)=dt\,du\,dg_1^*(l)\,d\sigma_0(v')\).

\begin{proposition}[The weighted position fiber with two endpoint scales]
\leavevmode\label{prop:twistedfiber919}
In these coordinates \(\|l\|\le\|x\|\) and
\begin{align}
  \|x\|&\le12 R^4 M^5 \Delta_x(v)^{-2}\|l\|,
                       \label{eq:twisteddrop919}\\
  \int_{\|x\|\le Y}\Delta_x(v)^{2z_0}\,dt\,du
       &\le 2\pi^{d_1}|\lambda|^{2z_0}(R^2 M Y)^m
            \qquad(Y\ge1)
                       \label{eq:twistedfiber919}.
\end{align}
In the second formula \(v,v',l\) are fixed.
\end{proposition}
\begin{proof}
The middle block of \(d_\lambda a_v k_v x
k_{v'}^{-1}a_{v'}^{-1}\) is the middle block of the matrix
\(k_v x k_{v'}^{-1}\). It is \(l\), giving the first bound.
The full matrix with \(\|x\|\le Y\) has norm at most
\(T=R^2 M Y\). Its last column bounds both \(\|t\|\) and
\(|u|\) by \(T\).

Let \(b\) attain \(\Delta=\Delta_x(v)\), and put
\(b_1=s_{v'}x^{-1}b,\ b_2=d_\lambda s_v b\). The matrix
\(n(t,u)l\) takes \(b_1\) to \(b_2\). The common bottom
entry is \(\gamma=\bar\lambda^{-1}h(v,b)\).
The two vectors have norm at most \(R M\). More precisely,
their middle components both have norm at most one since
the diagonal factors act identically there. With \(P_l\)
the positive \((k-1)\)-plane of Lemma~\ref{lem:commondirection913}
in the residual space, the middle equation
\((b_2)_V=l(b_1)_V+\gamma t\) yields
\[
        \Delta\,\|\operatorname{pr}_{P_l}t\|\le 2|\lambda|.
\]
Here \(\|\operatorname{pr}_{P_l}l(b_1)_V\|\le\|(b_1)_V\|\)
by pairing with positive unit vectors of \(P_l\).
Integrate \((2|\lambda|)^{2z_0}\|\operatorname{pr}_{P_l}t\|^{-2z_0}\)
over the ball of radius \(T\), then integrate over the central
interval. The same polar computation as above, now with
\(s=k-1\), gives the right side of \eqref{eq:twistedfiber919}
times
\[
       \frac{2^{2z_0}}
              {(s-z_0)\Gamma(s)\Gamma(d_1-s+1)}\le1 .
\]
Indeed \((k-2)!\ge 4^{k-5}\) for \(k\ge10\), by checking
\(k=10\) and increasing \(k\) one at a time.

For the size transfer, the middle equation and its inverse
bound \(\|t\|,\|l^{-1}t\|\) by \(2 R M\|l\|/|\gamma|\).
The top row then bounds the top right entry by
\(2RM/|\gamma|+2(RM)^2\|l\|/|\gamma|^2\). Adding the block
bounds as in Proposition~\ref{prop:positiondrop913}, and
using \(|\gamma|\le RM\), gives
\(\|n(t,u)l\|\le12(RM)^2 |\gamma|^{-2}\|l\|\).
The two sections and \(d_\lambda^{-1}\) cost at most
\(R^2M\) on returning to \(x\). Substituting
\(|\gamma|=\Delta/|\lambda|\) proves \eqref{eq:twisteddrop919}.
\end{proof}
\clearpage

\subsection{The second real frequency with two prescribed endpoints}
\label{sec:secondsharp920}

Use \(d_1,z_0,m\) from Section~\ref{sec:crossfiber919}.
Let \(p,q\in\mathbb C^\times\) be the first pivots at the two
endpoints, put \(\lambda=p/q\), and define
\[
 H=\max(1,|p|^{-1},|q|^{-1}),\quad M=\max(|\lambda|,|\lambda|^{-1}),
 \qquad R=b^{6n+2}.
\]
We specify the mixed form, including its measure. At a pair
\(P=(v,w)\) with \(h(v,w)=p\), choose \(s_P\) taking it to
\((e,p e^*)\), as in Theorem~\ref{thm:secondallpositions918}.
In recursive quotient Haar, let \(D^\alpha_{s,x}\,d\sigma_1\)
be the image at \(z=m_1^{-1} f_1\) of
\[
   \langle W(m_1s),W(m_1d_\alpha s x)\rangle\,d\dot m_1^*,
                \qquad m_1\in U_1\backslash G_1 .
\]
For \(\alpha=1\) this is the untwisted density. Define
\(J_{p,q}(x;\nu)\) by integrating
\[
 \iint e^{2i\operatorname{Im}(\nu h(z,z'))}
        D^1_{s_P,x}(z)\overline{D^\lambda_{w_p^{-1}s_P,x}(z')}
                                      \,d\sigma_1(z)d\sigma_1(z')
\]
against \(\eta^2|\lambda|^{-d_1}d\mu_p(P)\).
Thus the two second-branch sections are \(w_p^{-1}s_P\) and
\(d_\lambda w_p^{-1}s_Px=w_q^{-1}s_Px\). The density factor is
\(\tfrac12|pq|^{d_1}d\beta_p(P)\), the geometric mean of the
two endpoint Jacobian factors in the pair coordinates.
The remaining freedom in \(s_P\) is left multiplication by
\(G_1\). It transports both residual orbits by the same isometry,
so the integral is independent of that choice. For \(p=q=c\)
this is precisely \(J_{x,c}(\nu)\) in \eqref{eq:conditionedpivot917}.
Write \(D_{\rm type}=d_v\) for the compact-type dimension.

\begin{theorem}[The whole second real wave at prescribed first entries]
\leavevmode\label{thm:secondsharp920}
At sufficiently large nodes of the shorter interval, for
\(\nu\ne0\) and \(0<|p|,|q|\le R^2\),
\begin{align}
 \frac{\int_G |J_{p,q}(x;\nu)|\,dx}{D_{\rm type}S_{k,a}}
  &\le b^{2000n^9}\,
      \max(1,|p|^{-1})^{n^6}\max(1,|q|^{-1})^{n^6}
      |\nu|^{-K_0}\min(1,|\nu|^{-6}).
 \label{eq:secondsharp920}
\end{align}
The integral vanishes if either pivot has modulus greater than
\(R^2\). In particular the small-pivot loss
\(|c|^{-2r_1-2}\) of \eqref{eq:secondallpositions918} can be
removed in the equal-pivot case.
\end{theorem}
\begin{proof}
We work in recursive measures first. Put \(W_v(m_1)=W(m_1s_v)\).
For \(\alpha\in\{1,\lambda\}\) and \(j=1,2\) let
\(A_{\alpha,j}(v,x)=\|D^\alpha_{s_v,x}\|_{L^j(\sigma_1)}\).
These norms are independent of the first-vector section.
Indeed replacement by \(n(t,u)l_0s_v\) changes the first
amplitude by a character and a right translation by \(l_0\).
In the second amplitude the unipotent is
\(d_\alpha n(t,u)d_\alpha^{-1}=n(\bar\alpha t,|\alpha|^2 u)\).
After moving the two unipotents to the left of \(m_1\), the ratio
of the characters is a unit-modulus function of \(m_1^{-1}f_1\).
The right translation transports this vector by \(l_0\), which
preserves \(\sigma_1\). This proves the norm assertion.

At fixed \(v\) put \(v'=\alpha^{-1}x^{-1}v\) and use
\(d_\alpha s_vx=n(t,u)l s_{v'}\). Moving \(n(t,u)\) past
\(m_1\) likewise gives
\begin{equation}
  A_{\alpha,j}(v,x)=
       \|D_{1,l}^{W_v,W_{v'}}\|_{L^j(\sigma_1)}.
 \label{eq:twistedrow920}
\end{equation}
The slices have radius \(R\), and their squared \(L^2\) norms
integrate to \(\|W\|_{2,*}^2\) in the first-vector variables.
Define \(F_{\alpha,j}(Y)\) to be the integral over
\(\|x\|\le Y\) and \(v\in\mathcal A_R\) of
\(\Delta_x(v)^{2z_0} A_{\alpha,j}(v,x)^2\) in
\(dg_0^*(x)d\sigma_0(v)\). Only \(v'\in\mathcal A_R\) contribute.
Set \(N_\alpha=\max(|\alpha|,|\alpha|^{-1})\).
Integrating \eqref{eq:twistedrow920} in \(t,u\) using
\eqref{eq:twistedfiber919}, then using \(\|l\|\le Y\) in
Proposition~\ref{prop:recenergy912}, gives
\begin{align}
 F_{\alpha,2}(Y)\le {}&
 2\pi^{d_1}|\alpha|^{2z_0}(R^2 N_\alpha Y)^m
 (16r^{3/2}R^5Y^2)^{a-1}
 \pi^{\sum_{i=2}^a(2r_i-1)}\,\frac{\|W\|_{2,*}^4}{D_{\rm type}} .
 \label{eq:weightedendpoint920}
\end{align}
All residual norms on the right of \eqref{eq:twistedrow920}
are independent of these \(t,u\).

The four first vectors on the support are
\(v,w,v'=x^{-1}v,w'=\lambda^{-1}x^{-1}w\), all in
\(\mathcal A_R\). Since \(h(v',w')=q\), this proves the
support assertion and yields \(M\le R^2H\).
At fixed \(x\) the absolute value of the inner wave
in the definition of the form is bounded by
\(A_{1,2}(v,x)A_{\lambda,2}(w,x)\) times
\[
     \pi^{d_1}\min\{R^{2(d_1-1)},\
       (8R^8H^3)^{2(d_1-1-z_0)}|\nu|^{-z_0}\}.
\]
For the oscillatory bound use \eqref{eq:realquadrics911}
with \(\ell=z_0\) in the common pair section. The radius
\(8R^8H^3\) follows from the first amplitude of each branch
as in Corollary~\ref{cor:conditionedwave917}. The second
amplitudes may only reduce these supports. For the other
bound take the product of the two \(L^1\) norms. Each
\(L^1/L^2\) ratio is at most the square root of the orbit
volume of radius \(R\) in its own canonical section, by
\eqref{eq:twistedrow920}. The polar formula in
Theorem~\ref{thm:realquadrics911} bounds this volume by
\(\pi^{d_1}R^{2(d_1-1)}\).

Now apply \eqref{eq:crossfiber919} and Cauchy--Schwarz in
\(x\). The two resulting square integrals are
\(F_{1,2}(Y),F_{\lambda,2}(Y)\). Their geometric mean supplies
\(|\lambda|^{z_0} M^{m/2}\), while the form carries
\(|\lambda|^{-d_1}\). The remaining ratio is at most
\(M^{d_1-z_0+m/2}\le M^m\). Passing to original Haar
multiplies by \(2^n V_k/V_r\) in \(x\) and by \(\eta^2\)
in the two quotient measures. We also have \(\eta^2\|W\|_{2,*}^4\le D_{\rm type}^2\).
Put \(s=m+2(a-1)=6a+7\). Denote the normalized mass on a
set \(B\subset G\) by \(I_B=(D_{\rm type}S_{k,a})^{-1}
\int_B |J_{p,q}(x;\nu)|dx\). We have proved
\begin{equation}
 I_{\|x\|\le Y}
 \le (RH)^{100n^2}Y^s
                \min(R^{2z_0},|\nu|^{-z_0}).
 \label{eq:secondball920}
\end{equation}
Here and below the minimum at \(\nu=0\) is \(R^{2z_0}\).
To check the constant, the uncanceled power of \(\pi\) is
\(3d_1+\sum_{i=2}^a(2r_i-1)-(2a+1)k=2a^2+a\).
The other factors outside \(Y^s\) and the minimum are at most
\[
  2^{n+1}(16r^{3/2})^{a-1}\,
     |p|^{-1}R^{3m+1+5(a-1)}
           M^m(8R^8H^3)^{2(d_1-1-z_0)} .
\]
Since \(d_1-1-z_0=2a+3\), \(R\ge\max(16,2r)\) and
\(M\le R^2H\), all these factors, including \(\pi^{2a^2+a}\),
fit in \((RH)^{100n^2}\) at large nodes.

For the distant positions use the product of the two \(L^1\)
norms before the pair integration. The support implies
\(\Delta_x(v)\ge\min(|p|,|q|)/R\), by testing on
\(w/(R\max(1,|\lambda|))\). Likewise
\(\Delta_x(w)\ge |p|/R\) by testing on \(v/R\).
Thus in each use of \eqref{eq:twistedrow920} the residual
position satisfies \(\|x\|\le D_*\|l\|\), with
\(D_*=R^{17}H^7\), by \eqref{eq:twisteddrop919}. The
projection cutoffs in Theorem~\ref{thm:crossfiber919} retain
these lower bounds on \(\Delta_x\) in both square integrals.
On \(Y<\|x\|\le2Y\) integrate the weighted positions as
before, with \(2Y\) as the upper bound. This time the
residual squared \(L^1\) norm is \(M_1(l;W_v,W_{v'})\).
For \(a\ge2,\ Y\ge D_*R^2\), its integral over \(\|l\|>Y/D_*\)
is at most
\[
 \pi^{(2a-1)k} R^{80ak}
    (Y/D_*)^{-k/(2(a-1))}
   \|W_v\|_{2,*}^2\|W_{v'}\|_{2,*}^2/D_{\rm type}.
\]
This is the recursive-measure shell bound proved in
\eqref{eq:realshell914}, with \(a-1\) null rows and its
dyadic summation. For \(a=1\), such residual positions give
zero. Taking the square root of the product for the two
branches cancels the large power of \(|\lambda|\) as above,
and the remaining power of \(\pi\) is \(4a-2\). It follows that
\begin{equation}
 I_{Y<\|x\|\le2Y}\le (RH)^{100n^2}
     R^{80ak} Y^m (Y/D_*)^{-k/(2(a-1))}
                     \qquad (a\ge2).
 \label{eq:secondshell920}
\end{equation}
Indeed the factors from the conditional fiber and the two
position fibers, outside \(Y^m\) and the residual tail, are
bounded by \(2^{n+1+m}\pi^{4a-2}|p|^{-1}R^{3m+1}M^m\).

Choose \(X=R^{400n^2}H^{16}\max(1,|\nu|)^{4a}\).
Since \(k\ge40n^3\) eventually,
\(k/(2(a-1))-m\ge k/(4a)\) for \(a\ge2\). Summing
\eqref{eq:secondshell920} over \(Y=2^iX\) costs at most
\[
   2(RH)^{100n^2} D_*^m R^{80ak}
                  (X/D_*)^{-k/(4a)}
       \le 2(RH)^{200n^2}\max(1,|\nu|)^{-k}.
\]
Here \((400n^2-17)/(4a)>80a\).
In the ball, for \(|\nu|\le1\) use
\(\min(R^{2z_0},|\nu|^{-z_0})\le
 R^{2(z_0-K_0)}|\nu|^{-K_0}\), by splitting at
\(|\nu|=R^{-2}\). For \(|\nu|>1\), the \(4as\) frequency
power from \(X^s\) is paid by
\(z_0-K_0-6=n^7-11\ge4as\) at large nodes.
The total coefficient is bounded by
\(R^{2n^7+O(n^3)}H^{O(n^2)}\) with absolute constants,
by \eqref{eq:secondball920} and the tail just proved.
Since \(R=b^{6n+2}\) and \(n\to\infty\) along the shorter
interval, the first factor is at most \(b^{2000n^9}\) and
the second is at most
\(\max(1,|p|^{-1})^{n^6}\max(1,|q|^{-1})^{n^6}\).
This gives \eqref{eq:secondsharp920}.
\end{proof}
\clearpage

\subsection{Reciprocal differences at the first real pivot}
\label{sec:firstmismatch921}

For the comparison of two rational heads there is also a
frequency when the weight on the next Gram entry is constant.
Let \(p,q\ne0\), \(p\ne q\), and set
\(\xi=p^{-1}-q^{-1}\). Keep the form \(J_{p,q}\), its Haar
measures and \(H,M,R,d_1,m,z_0\) of the preceding subsection.

\begin{theorem}[A reciprocal-difference estimate on the whole real group]
\leavevmode\label{thm:firstmismatch921}
On the same shorter interval as Theorem~\ref{thm:secondsharp920},
\begin{align}
 \frac1{d_v S_{k,a}}\int_G |J_{p,q}(x;0)|\,dx
 &\le b^{2000n^9}
 \max(1,|p|^{-1})^{n^6}\max(1,|q|^{-1})^{n^6}
 |\xi|^{-K_0}\min(1,|\xi|^{-6}).
 \label{eq:firstmismatch921}
\end{align}
\end{theorem}
\begin{proof}
By support we may suppose \(|p|,|q|\le R^2\).
Write \(\lambda=p/q\). At zero next frequency the residual
pairing factors into the product of the two ordinary
conditional correlations. The first is \(D_{0,x}^{W,W}(v)\),
independent of the pair section, since its two arguments use
the same simple-root character. Denote it by \(C_x(v)\) in
recursive measures. For the other correlation fix \(w\) and
parametrize the first vector and pair section by
\[
      v=s_w^{-1}n(-t,-u)\bar p e^*,\qquad
      w_p^{-1}s_P=n(t,u)s_w .
\]
These formulas satisfy \(s_Pv=e,\ s_P w=p e^*\), since
\(w_p(\bar p e^*)=e\). The pair measure is
\(\tfrac12 |p|^{2d_1}d\sigma_0(w)\,dt\,du\) by
Theorem~\ref{thm:conditionedpivot917} and exchange of the
two vectors. Moving \(n(t,u)\) to the left of the residual
variable \(m_1\) gives the character at \(-h(t,m_1^{-1}f_1)\)
in the first amplitude. In the second it gives the character
at \(-\lambda h(t,m_1^{-1}f_1)\), by conjugation by \(d_\lambda\).
As the inner product is conjugate linear in its first
argument, the remaining factor is, up to the choice of sign
for the Bessel character,
\[
      T_w(t)=\int e^{2i\operatorname{Im}((1-\lambda)h(t,z))}
                        D^\lambda_{s_w,x}(z)\,d\sigma_1(z).
\]
The central variable does not occur in this transform.

At fixed \(x\) take absolute values and apply Cauchy--Schwarz
in \(d\mu_p\) to
\(|C_x(v)|\Delta_x(w)^{-z_0}\) and
\(\Delta_x(w)^{z_0}|T_w(t)|\), keeping the annular supports
of both first vectors. The first squared integral is at most
\[
    \pi^{d_1}R^{m+1}|p|^{-1}
                         \int |C_x(v)|^2d\sigma_0(v)
\]
by the conditional negative moment in
Theorem~\ref{thm:crossfiber919}. For the second one we have
\(\|t\|,|u|\le T=R^2/|p|\). The support of
\(D^\lambda_{s_w,x}\) in its canonical coordinates satisfies
\(R^{-1}\le\|z\|\le R\) by the first amplitude \(W(m_1s_w)\).
Apply Lemma~\ref{lem:truncquad912} with the frequency
variable scaled by \(\overline{1-\lambda}\), and integrate
the \(u\)-interval. Conditional on \(w\), this gives
\begin{align}
 \tfrac12 |p|^{2d_1}\int_{|u|,\|t\|\le T}|T_w(t)|^2dt\,du
 &\le 2(2d_1)^{3/2}\pi^{2d_1-1}R^5
          |p|^{-1}|\xi|^{1-2d_1} A_{\lambda,2}(w,x)^2 .
 \label{eq:reciprocalplancherel921}
\end{align}
In fact complex scalar multiplication has real Jacobian
\(|1-\lambda|^{2d_1}\), the Fourier ball has radius
\(|1-\lambda|T\), and \(|1-\lambda|=|p|\,|\xi|\).

Integrate the geometric mean of these two bounds over
\(\|x\|\le Y\) and use Cauchy--Schwarz in the position.
The two remaining squared integrals are respectively
\(E_0(Y;W,W)\) from \eqref{eq:recenergy912} and
\(F_{\lambda,2}(Y)\) from \eqref{eq:weightedendpoint920}.
In the second integral the factor \(\Delta_x(w)^{2z_0}\)
is the one supplied by the preceding application of
Cauchy--Schwarz. The large scale in the latter bound becomes
\(|\lambda|^{z_0}\) on taking the square root. Including the
pair normalization leaves
\(|\lambda|^{z_0-d_1}M^{m/2}\le M^m\).
The total power of \(Y\) is
\((2a+m+2(a-1))/2\le s=6a+7\). The uncanceled power of \(\pi\)
after the Haar conversion is \(2a^2+a-1\), by adding the
two Plancherel exponents in the square roots. All remaining
constant factors other than \(|\xi|^{1/2-d_1}\) are bounded
by \((RH)^{100n^2}\) at large nodes: the powers of \(R\)
outside \(M^m\) are at most \(3m+5a+10\), the polynomial
factors in \(r\) are at most \((32r^2)^{a+1}\), and
\(|p|^{-1}\le H,\ M\le R^2H\). With
\(I_B=(d_vS_{k,a})^{-1}\int_B|J_{p,q}(x;0)|dx\) we have
\[
 I_{\|x\|\le Y}\le (RH)^{100n^2}Y^s
           \min\bigl(R^{2z_0},\,|\xi|^{1/2-d_1}\bigr).
\]
Here the bound by \(R^{2z_0}\) is \eqref{eq:secondball920}
at zero next frequency.

For the complementary positions \eqref{eq:secondshell920}
holds with the constant weight. Take
\(X=R^{400n^2}H^{16}\max(1,|\xi|)^{4a}\) and sum its
dyadic shells as in Theorem~\ref{thm:secondsharp920}.
The tail is at most \(2(RH)^{200n^2}\max(1,|\xi|)^{-k}\).
For \(0<|\xi|\le R^{-2}\) use the radius term in the minimum.
For \(R^{-2}\le|\xi|\le1\) the excess exponent
\(d_1-\tfrac12-K_0\) costs at most
\(R^{2(d_1-1/2-K_0)}\). For \(|\xi|>1\) that excess is
greater than \(4as+6\) eventually, which pays the growth
from \(X^s\) while retaining the sixth power of decay.
The resulting coefficient is bounded by
\(R^{2n^7+O(n^3)}H^{O(n^2)}\). Conversion to the powers
of \(b\) and of the two inverse first entries is the same as
in \eqref{eq:secondsharp920}.
\end{proof}
\clearpage

\subsection{Ordered null flags with prescribed Gram entries}
\label{sec:orderedflags922}

Let \(1\le j\le a\). Write \(e_i=f_i\) for \(0\le i<j\),
and \(e_i^*\) for its dual hyperbolic partner. In this subsection
\(J\) denotes the matrix of the Hermitian form in the fixed Euclidean
coordinates, so \(J^*=J=J^{-1}\). The space \(X_j\) consists of
ordered, linearly independent, mutually orthogonal null vectors
\(A_0,\ldots,A_{j-1}\), modulo addition to each column of a linear
combination of earlier columns. We retain the scale of every
successive vector. Put \(N_{<j}=N_0\rtimes(N_1\rtimes\cdots
\rtimes N_{j-1})\), embedded in the successive middle blocks.
The stabilizer of the standard class is \(N_{<j}\rtimes G_j\).
Indeed an isometry preserving the class fixes \(e_0\). Removing
its \(N_0\) component reduces the assertion to the same assertion
on \(e_0^\perp/\mathbb C e_0\). This also proves transitivity,
by applying the null-vector transitivity at each step.

If \(b_i\in\mathcal C_i=G_i f_i\), choose the sections
\(s_{b_i}=a_{b_i}k_{b_i}\) of
Proposition~\ref{prop:recenergy912}. Set \(S_0=1\),
\(S_{i+1}=s_{b_i}S_i\). A point of \(X_j\) has unique coordinates
\(b_0,\ldots,b_{j-1}\) in these sections: \(b_i\) is the middle
projection of \(S_i A_i\). The earlier components of \(S_i A_i\)
along \(e_0,\ldots,e_{i-1}\) are precisely the freedom in
the representative. Its dual components vanish by orthogonality.
The invariant flag measure \(d\Sigma_j\) is
\(\prod_{i<j}d\sigma_i(b_i)\), by repeated use of the quotient
measures of Section~\ref{sec:realpositionball912}. At a nonzero
slice of the cut \(W\), all these coordinates have norm between
\(R^{-1}\) and \(R\), and \(W_{S_j}(m)=W(m S_j)\) has residual
quotient radius \(R\). This follows at each step by taking the
middle block as in Proposition~\ref{prop:recenergy912}. Also the
flag then has a simultaneous representative of column-operator
norm at most \(R\): take the inverse of a representative of the
support of \(W\) of norm at most \(R\) and apply it to the first
\(j\) standard vectors. Representatives differing by a left
Bessel unipotent give the same flag class. We use the closed set
of flags admitting such a representative when estimating supports.

For a pair of flags with column matrices \(A,B\), put
\(K=A^*JB\) and let \(\delta_i\) be its leading determinant
of size \(i+1\), with \(\delta_{-1}=1\). All these determinants
are invariant under the allowed column changes. On the open set
where none vanish let \(p_i=\delta_i/\delta_{i-1}\).
Write \(\mathbf p=(p_0,\ldots,p_{j-1})\). By successive
elimination there are unique representatives \(\widetilde A_i,
\widetilde B_i\) satisfying
\[
             h(\widetilde A_i,\widetilde B_l)
                    ={\boldsymbol 1}_{i=l}p_i .
\]
In fact the earlier \(i\)-by-\(i\) block is invertible, so we
can subtract a unique combination of the earlier \(A\)'s
from \(A_i\) to make it orthogonal to all earlier \(B\)'s.
The same operation with \(A,B\) exchanged fixes \(B_i\).
The new diagonal pairing is the quotient of consecutive
determinants, because these triangular operations have
determinant one.

\begin{theorem}[Coarea for two ordered null flags]
\leavevmode\label{thm:flagmeasure922}
In the open set just described,
\[
 d\Sigma_j(A)d\Sigma_j(B)=
 \left(\prod_{i<j}d^2p_i\right)d\mu_{\mathbf p}(A,B),\qquad
 d\mu_{\mathbf p}=2^{-j}
       \left(\prod_{i<j}|p_i|^{2r_{i+1}}\right)d\beta_{\mathbf p}.
\]
Here \(d\beta_{\mathbf p}\) is the Haar quotient by \(G_j\) in
the pair orbit, with the group and residual group in recursive
measures. The formula specifies the conditional measure at
every \(\mathbf p\in(\mathbb C^\times)^j\). Exchange of the
two flags takes this measure to \(d\mu_{\overline{\mathbf p}}\).

There is also a formula keeping the first flag fixed and exposing
the second in its own canonical coordinates. Given its first
\(i\) coordinates, let \(u_i\) be the residual middle component
of \(S_i\widetilde A_i\), where \(S_i\) belongs to the second
flag. The conditional measure for \(b_i\) is the coarea of
\(d\sigma_i(b_i)\) at \(h(u_i,b_i)=p_i\), relative to planar
Lebesgue measure. These successive conditional measures, together
with \(d\Sigma_j(A)\), give \(d\mu_{\mathbf p}\).
\end{theorem}

\begin{proof}
First fix \(A_0=v\). At \(p_0\ne0\) write
\(B_0=s_v^{-1}n(t,u)p_0e_0^*\). The one-vector calculation
of Theorem~\ref{thm:conditionedpivot917} gives
\(\tfrac12 |p_0|^{2r_1}d^2p_0\,dt\,du\).
With \(\tau=n(t,u)^{-1}s_v\), the pair is now \(e_0,p_0e_0^*\).
Changing either first-vector section to \(\tau\) or to
\(w_{p_0}^{-1}\tau\), respectively, preserves its conditional
residual flag measure. Indeed a change by \(N_0\) acts trivially
on the quotient \(e_0^\perp/\mathbb C e_0\), and the remaining
change is a \(G_1\)-isometry preserving that measure.
Represent each remaining column in the common middle space
by subtracting a multiple of its own first column to make it
orthogonal to the other first column. The resulting cross Gram
matrix is the Schur complement of \(p_0\), since a diagonal
first row and column have been removed. Its successive pivots
are therefore \(p_1,\ldots,p_{j-1}\). In this chart the two
remaining flags have the product of their invariant measures:
the two sections have the same middle identification because
\(w_{p_0}\) acts identically there. Induction gives the asserted
Jacobian.

After both columns of each pair have been fixed, the remaining
stabilizer is \(G_j\), acting on their common orthogonal
complement. The other coordinates in the above calculation,
at each step \(d\sigma_i(v)\,dt\,du\), are exactly the quotient
of \(dg_i^*\) by \(dg_{i+1}^*\). This identifies the stated
\(\beta_{\mathbf p}\) and proves the pointwise normalization.
For exchange, the standard pair for \(\mathbf p\) is carried
after swapping to the standard pair for \(\overline{\mathbf p}\)
by the direct sum of the maps \(w_{p_i}^{-1}\); it commutes
with the residual group. Left invariance of group Haar thus
preserves the quotient measure, and the Jacobian factors
are unchanged.

For the other description, \(\widetilde A_i\) depends only on
\(A\) and the earlier \(B\)'s. It is orthogonal to all earlier
\(B\)'s. Thus \(S_i\widetilde A_i\) has no earlier dual components.
Its middle projection \(u_i\) is null, and its pairing with
\(S_iB_i\) is \(h(u_i,b_i)\). Elimination proved above identifies
this with \(p_i\). If \(u_i=0\) there is no extension with the
required nonzero value. Otherwise take a section carrying \(u_i\)
to \(f_i\) and use the same one-vector coarea in this middle
space. Iteration gives the conditional formula. To see pointwise
agreement of the two descriptions, one may choose smooth local
sections in these calculations. At nonzero pivot the coordinates
of the other vector are its complex pivot, unrestricted middle
coordinates and the imaginary central coordinate. Thus every
successive pivot map is a submersion on the indicated open set.
Both calculations are changes of variables for the same smooth
density with the same ordered base differentials. The resulting
fiber density agrees in each local chart at every point.
\end{proof}

\begin{proposition}[A bounded common section for a prescribed pair]
\leavevmode\label{prop:flagsection922}
Suppose both flags admit simultaneous representatives of column
norm at most \(R\), with \(R\ge\max(16,2r)\). Put
\[
 H_{\mathbf p}=\max(1,\max_{i<j}|\delta_i|^{-1}),\qquad
 L_{\mathbf p}=R^{10n}H_{\mathbf p}^3.
\]
There is an isometry \(\tau_P\) taking the normalized columns
\(\widetilde A_i,\widetilde B_i\) to \(e_i,p_i e_i^*\),
respectively. Write \(w_{\mathbf p}\) for the simultaneous
swaps \(w_{p_i}\), equal to the identity on the middle space.
Then
\[
 R\max(\|\tau_P\|,\|w_{\mathbf p}^{-1}\tau_P\|)\le L_{\mathbf p}.
\]
The same bound applies to the norms of the inverse maps.
Also \(|\delta_i|\le R^{2(i+1)}\) for all \(i<j\).
\end{proposition}

\begin{proof}
For the bounded representatives, every cross block has norm at
most \(R^2\). The determinant bound follows from its singular values.
The inverse of the leading block of size \(i\ge1\) has norm at
most \(R^{2(i-1)}/|\delta_{i-1}|\), again by singular values.
The earlier columns subtracted in elimination have coefficient
vector of norm at most \(R^{2i}/|\delta_{i-1}|\). In particular
both normalized columns of index \(i\) have norm at most
\(R(1+R^{2i}H_{\mathbf p})\), also for \(i=0\).
Both \(|p_i|\) and \(|p_i|^{-1}\) are at most \(R^{2j}H_{\mathbf p}\)
by their determinant formulas.

Let \(T\) send \(e_i,e_i^*\) to
\(\widetilde A_i,\widetilde B_i/p_i\), respectively. Thus
\(T^*JT=J_0\), the standard form on these \(2j\) vectors, and
\(\|T\|\le4\sqrt j\,R^{4j+1}H_{\mathbf p}^2\).
Put \(Q=I-TJ_0T^*J\). It is the projection onto the Hermitian
orthogonal complement \(L\) along the range of \(T\). Let
\(\Pi_L\) be the Euclidean orthogonal projection. The inverse of
the compression \(\Pi_L J|_L\) is \(QJ|_L\): their product on
\(L\) is \(\Pi_L Q^*|_L=I\), since \(JQJ=Q^*\) and \(Q|_L=I\).
Thus the inverse compression has norm at most \(1+\|T\|^2\).
Diagonalize the Hermitian compression in an orthonormal basis
of \(L\) and rescale each eigenvector of eigenvalue \(\gamma\)
by \(|\gamma|^{-1/2}\). The resulting extension on the middle
space has norm at most \(\sqrt{1+\|T\|^2}\). Its signs match
the middle signature by additivity on the orthogonal direct
sum. Joining it to \(T\) gives an isometry of norm at most
\(3(1+\|T\|)\), whose inverse is \(\tau_P\). Every isometry
has the same norm as its inverse since \(g^{-1}=Jg^*J\).
The swap has norm at most \(R^{2j}H_{\mathbf p}\).
The stated \(L_{\mathbf p}\) absorbs these bounds, since
\(j<n\) and \(R\ge\max(16,2r)\).
The construction can be made measurable by choosing
orthonormal bases in successive eigenspaces by Gram--Schmidt
from the coordinate vectors. Any two common sections differ by
left multiplication by \(G_j\).
\end{proof}
\clearpage

\subsection{Weighted correlations on prescribed null flags}
\label{sec:flagenergy923}

Keep \(X_j,d\Sigma_j,d\mu_{\mathbf p}\) from
Theorem~\ref{thm:flagmeasure922} and \(z=k-5\). At a flag with
canonical coordinates \(b_0,\ldots,b_{j-1}\), let
\(\boldsymbol\alpha=(\alpha_0,\ldots,\alpha_{j-1})\) be nonzero
complex scales. For \(x_0=x\) define \(b_i'\) and \(x_{i+1}\) by
\[
 b_i'=\alpha_i^{-1}x_i^{-1}b_i,\qquad
 d_{\alpha_i}s_{b_i}x_i s_{b_i'}^{-1}=n(t_i,u_i')x_{i+1}.
\]
In this formula \(u_i'\) is a real central coordinate. Write
\(\Delta_y(b)=\max_{\|v\|,\|y^{-1}v\|\le1}|h(b,v)|\) in
the corresponding residual space. Put
\[
 D_{\boldsymbol\alpha,x}(B)=
 \prod_{i<j}\Delta_{x_i}(b_i),\quad M_i=\max(|\alpha_i|,|\alpha_i|^{-1}),
 \quad d_i'=r_{i+1},\quad m_i=2(d_i'-z)+1 .
\]
These weights are finite and strictly positive. Each \(x_i\)
depends only on \(x\), the scales, and the part of the flag
strictly preceding \(b_i\). Call a flag admissible here if
\(R^{-1}\le\|b_i\|\le R\) at each step in the fixed canonical
sections. Let \(\mathbf1\) denote the list of scales equal to one.

\begin{theorem}[Crossed conditional moments for a whole flag]
\leavevmode\label{thm:crossflags923}
If \(f,g\) are nonnegative measurable functions supported on
admissible flags, then, for each \(x\) and any scales
\(\boldsymbol\alpha,\boldsymbol\gamma\),
\begin{align}
 \int f(A)g(B)\,d\mu_{\mathbf p}
 &\le C_{\mathbf p}
 \left(\int D_{\boldsymbol\alpha,x}^{2z} f^2\,d\Sigma_j\right)^{1/2}
 \left(\int D_{\boldsymbol\gamma,x}^{2z} g^2\,d\Sigma_j\right)^{1/2},
 \label{eq:crossflags923}\\
 C_{\mathbf p}&=\prod_{i<j}
        \pi^{d_i'} R^{m_i+1}|p_i|^{-2}. \notag
\end{align}
For an integral restricted to a subset of pairs, one may impose
any measurable one-flag restrictions containing its projections
in the corresponding right-hand integrals.
\end{theorem}

\begin{proof}
Fix \(A\) and the first \(i\) coordinates of \(B\). Use \(u_i\)
for the null constraint vector of Theorem~\ref{thm:flagmeasure922}.
If its fiber is nonempty in \(\|b_i\|\le R\), then
\(\|u_i\|\ge |p_i|/R>0\). Carry \(u_i\) to the standard null
vector by \(s_{u_i}=a_{u_i}k_{u_i}\) as before. In the coordinates
\[
 s_{u_i}b_i=c e+y+p_i e^*,\qquad
           s=\operatorname{Im}(c/p_i),
\]
the conditional measure is \(\tfrac12dy\,ds\), since integration
in \(\operatorname{Re}(c/p_i)\) cancels the complex scaling
Jacobian as in Theorem~\ref{thm:conditionedpivot917}. We have
\(\|y\|\le R\) and
\(|s|\le R/(\|u_i\|\,|p_i|)\le R^2/|p_i|^2\); here \(a_{u_i}\)
acts as the identity on the middle space and as
\(\|u_i\|^{-1}\) on \(e\).

In the positive space there is a \((k-1)\)-plane on which
\(x_i^{-1}\) is a Euclidean isometry into the positive space.
Indeed the kernel of its negative projection on the positive
space has dimension at least the positive index minus the
negative index, which is \(k-1\); preservation of \(h\)
gives the isometry there. After applying \(k_{u_i}\) this
plane meets the middle positive space in dimension at least
\(k-2\). For a \((k-2)\)-plane \(E\) in that intersection we
have \(\Delta_{x_i}(b_i)\ge\|\operatorname{pr}_E y\|\), by
testing on its unit vectors transported back by \(k_{u_i}^{-1}\).
Integrating the negative power in the two orthogonal balls in
\(E\) and its complement as in Theorem~\ref{thm:crossfiber919}
gives at most \(\pi^{d_i'}R^{2(d_i'-z)}\). Here \(k-2>z\).
Including the central interval bounds this conditional integral
by \(\pi^{d_i'}R^{m_i+1}|p_i|^{-2}\).
This bound is independent of all earlier choices. Integrating
backwards from \(i=j-1\) gives
\(\int D_{\boldsymbol\gamma,x}(B)^{-2z}\le C_{\mathbf p}\)
conditionally on each \(A\), with the admissibility restrictions.
The exchanged statement follows at every prescribed pivot by the
exchange formula of Theorem~\ref{thm:flagmeasure922} and the same
calculation for \(\boldsymbol\alpha\).
Multiply \(f(A)\) by \(D_{\boldsymbol\alpha,x}(A)^z/
D_{\boldsymbol\gamma,x}(B)^z\), multiply \(g(B)\) by its
reciprocal, and apply Cauchy--Schwarz. The two conditional
bounds give \eqref{eq:crossflags923}. The same proof works
after inserting the two one-flag restrictions and enlarging
the pair domain.
\end{proof}

Write \(d_{\boldsymbol\alpha}\) for the simultaneous diagonal
scales on the first \(j\) hyperbolic pairs. For a flag section
\(S\), let \(\mathcal D^{\boldsymbol\alpha}_{S,x}(z')d\sigma_j(z')\)
denote the pushforward under \(m\mapsto m^{-1}f_j\) of
\[
     \langle W(mS),W(m d_{\boldsymbol\alpha} Sx)\rangle\,d\dot m_j^* .
\]
Use \(\mathcal A_{\boldsymbol\alpha,h}(B,x)\) for the
\(L^h(\sigma_j)\) norm of this density in the canonical section
of \(B\), for \(h=1,2\). Write \(d=d_v\).

\begin{proposition}[Position integral for one flag with prescribed scales]
\leavevmode\label{prop:flagenergy923}
The norm \(\mathcal A_{\boldsymbol\alpha,h}\) is independent of
the choice of flag section. In canonical coordinates it equals
\(\|D_{j,x_j}^{W_{S_j},W_{S_j'}}\|_h\), where \(S_j'\) is the
canonical product for the successive \(b_i'\). In particular
\begin{align}
 &\int_{\|x\|\le Y}\int_{X_j}
       D_{\boldsymbol\alpha,x}(B)^{2z}
       \mathcal A_{\boldsymbol\alpha,2}(B,x)^2
                 d\Sigma_j(B)dg_0^*(x) \notag\\
 &\ \le
 \left[\prod_{i<j}2\pi^{d_i'}|\alpha_i|^{2z}(R^2 M_iY)^{m_i}\right]
 (16r^{3/2}R^5Y^2)^{a-j}
 \pi^{\sum_{i=j+1}^a(2r_i-1)}\,\frac{\|W\|_{2,*}^4}{d}.
 \label{eq:flagenergy923}
\end{align}
For \(h=1\) on the further restriction \(\|x_j\|>Z\ge R^2\),
the same bound holds with the last line's residual factor after
the square bracket replaced by
\[
 \pi^{(2(a-j)+1)k}R^{80(a-j+1)k}Z^{-k/(2(a-j))}
                          \frac{\|W\|_{2,*}^4}{d}
                       \qquad(j<a).
\]
For \(j=a\) that restricted integral is zero.
\end{proposition}

\begin{proof}
Applying the successive position decompositions gives
\(d_{\boldsymbol\alpha}S_j x=n x_j S_j'\) with \(n\in N_{<j}\).
To see this inductively, the new residual section and the new
scale normalize the previously removed radical. After putting
that radical on the left the remaining factor is exactly the
one-step position decomposition defining \(x_{i+1}\).

For \(m\in G_j\) the Bessel character on \(m n m^{-1}\) is a
function only of \(m^{-1}f_j\). Indeed the simple entries
strictly before row \(j-1\) are unchanged: the corresponding
basis vectors and dual coordinate functionals are fixed by
\(m\). In row \(j-1\) one evaluates the earlier-coordinate
functional of \(n\) on \(m^{-1}f_j\). All later simple entries
vanish because the induced residual action of \(n\) is the
identity. This argument includes \(f_a\) at the terminal row.
Thus moving \(n\) past \(m\) multiplies the pushed-forward
density by a function of its orbit variable of absolute value
one. This gives the claimed equality to the canonical
correlation norm. A change of section by \(n_0l_0\), with
\(n_0\in N_{<j}\), \(l_0\in G_j\), similarly contributes
two such phase factors and a common right translation of the
residual variable by \(l_0\). Conjugation by the scales preserves
\(N_{<j}\) and fixes \(l_0\). The induced residual isometry
preserves \(\sigma_j\), giving section independence.

Fix the initial flag. Repeated Haar disintegration in \(x\)
gives the product of the \(dt_i\,du_i'\), the measures
\(d\sigma_i(b_i')\), and \(dg_j^*(x_j)\). The fixed left
translations by the scales preserve Haar at each step.
The two nonzero slices have radius \(R\) at every stage by the
middle-block argument. In particular both endpoints' successive
vector coordinates satisfy the required annular bounds.
Also \(\|x_{i+1}\|\le\|x_i\|\le Y\), by taking the middle block
between the compact parts of the two canonical sections.
Integrate first in \(t_0,u_0'\), then in order through \(j-1\).
At each step later factors depend only on the residual position
and the remaining flag coordinates. Proposition~\ref{prop:twistedfiber919}
in \(G_i\) bounds that integral of \(\Delta_{x_i}(b_i)^{2z}\)
by \(2\pi^{d_i'}|\alpha_i|^{2z}(R^2M_iY)^{m_i}\) at fixed
\(b_i,b_i',x_{i+1}\). Enlarge the remaining position domain to
\(\|x_{i+1}\|\le Y\) after each integration. Finally apply
\eqref{eq:recenergy912} at \(j\); the squared norms of the
two slices integrate separately to \(\|W\|_{2,*}^2\) in
their respective flag coordinates. This proves \eqref{eq:flagenergy923}.
For the restricted squared \(L^1\) norm use instead the
dyadic tail estimate obtained from \eqref{eq:realshell914}
with \(a-j\) null steps and recursive group Haar. It is uniform
for any two residual slices of radius \(R\) and type dimension
\(d\), as its proof uses only these properties and their squared
norms. At \(j=a\) a nonzero correlation requires
\(x_a=y^{-1}(yx_a)\) with both factors of norm at most \(R\),
since \(U_a=1\); thus \(\|x_a\|\le R^2\).
\end{proof}

\begin{proposition}[Transfer of a large position to either residual group]
\leavevmode\label{prop:flagtransfer923}
Let the initial flag pair \(P=(A,B)\) have pivots \(\mathbf p\).
Choose \(\lambda_i=p_i/q_i\), where \(q_i\ne0\). Suppose each
flag of \(P\), and each flag at the final pair
\[
       P'=(x^{-1}A,\ (x^{-1}B_i/\lambda_i)_{i<j}),
\]
is admissible and admits a representative of column norm at
most \(R\). These two requirements hold for flags coming from
nonzero slices of the cut \(W\). For the
residual position \(x_j\) defined on either \(A\) with scales
\(\mathbf1\), or \(B\) with scales \(\boldsymbol\lambda\),
\[
      \|x\|\le D_*\|x_j\|,\qquad
 D_*=R^{2j}L_{\mathbf p}^2L_{\mathbf q}^2
       \le R^{42n}(H_{\mathbf p}H_{\mathbf q})^6 .
\]
\end{proposition}

\begin{proof}
The final normalized representatives are exactly
\(x^{-1}\widetilde A_i\) and \(x^{-1}\widetilde B_i/\lambda_i\)
by uniqueness of diagonal elimination: their Gram matrix is
diagonal with entries \(q_i\), and the column changes are still
unitriangular after the coordinate scalings. Hence
\(l=\tau_P x\tau_{P'}^{-1}\) fixes the \(2j\) standard columns
and lies in \(G_j\). In particular
\(\|x\|\le L_{\mathbf p}L_{\mathbf q}\|l\|\).
Take \(T,T'\) to be the pair sections on the side in question,
so they are \(\tau_P,\tau_{P'}\) on the first side and
\(w_{\mathbf p}^{-1}\tau_P,w_{\mathbf q}^{-1}\tau_{P'}\)
on the other. In both cases \(d_{\boldsymbol\alpha}T x(T')^{-1}=l\),
using \(d_{\boldsymbol\lambda}w_{\mathbf p}^{-1}=w_{\mathbf q}^{-1}\).
The successive final one-sided coordinates really give the
indicated final flag: \(d_{\boldsymbol\alpha}S_jx=n x_j S_j'\)
fixes its standard flag after removal of \(S_j'\), with the
column scales \(\alpha_i^{-1}\) on the original transported flag.

Write \(TS_j^{-1}=n_0 h\), \(T'(S_j')^{-1}=n_1 h'\) in the
flag stabilizers. Taking middle blocks in the identity for \(l\)
gives \(l=h x_j (h')^{-1}\), since conjugation by the scales
fixes the residual block. By admissibility each elementary section
and its inverse have norm at most \(R\), so \(S_j,S_j'\) and
their inverses have norm at most \(R^j\). It follows that
\(\|h\|\le R^j L_{\mathbf p}\), \(\|(h')^{-1}\|=\|h'\|\le
R^j L_{\mathbf q}\), by the bounds on the pair sections.
Combining the two estimates for \(\|l\|\) and \(\|x\|\) proves
the assertion on the stated support.
\end{proof}
\clearpage

\subsection{Real Gram waves after an arbitrary number of prescribed pivots}
\label{sec:allrealwaves924}

Use the shorter interval, the cut \(W\), and the original
normalizations \(d=d_v\), \(S_{k,a}=\pi^{(2a+1)k}V_k/V_r\)
of Theorem~\ref{thm:fullrealfirst915}. For lists \(\mathbf p,\mathbf q\)
of \(j\) nonzero complex pivots, let \(\lambda_i=p_i/q_i\).
At a pair \(P=(A,B)\) with pivots \(\mathbf p\) take
\(\tau=\tau_P\) from Proposition~\ref{prop:flagsection922}.
Form the two residual densities
\[
 g=\mathcal D^{\mathbf1}_{\tau,x},\qquad
 g'=\mathcal D^{\boldsymbol\lambda}_{w_{\mathbf p}^{-1}\tau,x}.
\]
At \(j=0\) both sections and the scale matrix are the identity.
For a bounded measurable function \(\phi\) on \(\mathbb C\) the
inner integral is
\[
 C_P(x;\phi)=\iint
    \phi(h(v,w))g(v)\overline{g'(w)}\,d\sigma_j(v)d\sigma_j(w).
\]
Changing the common section by \(l\in G_j\) transports both
residual variables by the same isometry, leaving \(C_P\) unchanged.
Use \(\phi_\nu(c)=\exp(2i\operatorname{Im}(\nu c))\) and put
\[
 J_{\mathbf p,\mathbf q}(x;\nu)=
    \eta^2 \prod_{i<j}|\lambda_i|^{-r_{i+1}}
                  \int C_P(x;\phi_\nu)\,d\mu_{\mathbf p}(P).
\]
This is the real position form on the two prescribed heads.
Indeed in common coordinates the remaining frame vectors are
middle vectors, and the relative Gram entry is \(h(v,w)\).
The two endpoint sections for the second flag are
\(w_{\mathbf p}^{-1}\tau\) and \(w_{\mathbf q}^{-1}\tau x\),
since \(d_{\boldsymbol\lambda}w_{\mathbf p}^{-1}=w_{\mathbf q}^{-1}\).
The outside measure in recursive coordinates is
\(2^{-j}\prod_{i<j}|p_iq_i|^{r_{i+1}}d\beta_{\mathbf p}\), by
Theorem~\ref{thm:flagmeasure922}. We keep the factors \(1/2\)
from real coarea in this convention.

For a list \(\mathbf p\) set \(\delta_l(\mathbf p)=\prod_{i=0}^l p_i\)
and define the weight
\[
 \mathcal H(\mathbf p)=\prod_{l<j}
 {\boldsymbol 1}_{|\delta_l(\mathbf p)|\le b^{60n^2}}\,
                  \max(1,|\delta_l(\mathbf p)|^{-1})^{n^6}.
\]
An empty product is one. Write \(W_0(\nu)=|\nu|^{-K_0}
\min(1,|\nu|^{-6})\) with \(K_0=k-n^7\).

\begin{theorem}[Uniform real wave bound for prescribed heads]
\leavevmode\label{thm:allrealwaves924}
At sufficiently large nodes on the shorter interval, uniformly
in the real place, \(0\le j\le a\), the two nonzero pivot lists,
and \(\nu\ne0\),
\begin{equation}
 \frac1{d S_{k,a}}\int_G |J_{\mathbf p,\mathbf q}(x;\nu)|\,dx
 \le b^{2000n^9}\mathcal H(\mathbf p)\mathcal H(\mathbf q) W_0(\nu).
 \label{eq:allrealwaves924}
\end{equation}
For \(j\ge1\) the same bound holds with the absolute value taken
inside the flag integral, immediately after the integration
in the two residual orbit variables. Also the integral over all
positions and flag pairs of the product \(\|g\|_1\|g'\|_1\)
with the defining outside measures is finite at every pair
of prescribed heads.
\end{theorem}

\begin{proof}
For \(j=0\) the wave assertion is Theorem~\ref{thm:fullrealfirst915};
integrability of the product of the \(L^1\) norms follows from
the position-ball and tail bounds in
Sections~\ref{sec:realpositionball912} and~\ref{sec:realremotetail914}.
Let \(j\ge1\), \(z=k-5\), \(d_i'=r_{i+1}\),
\(m_i=2(d_i'-z)+1\), and
\(H=H_{\mathbf p}H_{\mathbf q}\) with the maxima of
Proposition~\ref{prop:flagsection922}. Every pair contributing
has both initial and final flags on the four cut projections:
a nonzero residual density requires a nonzero value of each
amplitude on its corresponding flag. Thus the support conditions
of Proposition~\ref{prop:flagtransfer923} hold. In particular
every prefix determinant at either endpoint is at most \(R^{2j}\)
in modulus. Since \(R=b^{6n+2}\) and \(j<n\), a violated
indicator in \(\mathcal H\) makes the form zero. We may henceforth
assume all prefix determinants are at most \(R^{2j}\).
Consequently \(M_i=\max(|\lambda_i|,|\lambda_i|^{-1})\le
R^{4j}H\).

Let \(I(E)\) be the normalized mass on \(x\in E\), with the
absolute value inside the flag integral. On the support the
first amplitude bounds both residual orbit radii in the common
pair section by \(L_{\mathbf p}\). In detail, if
\(\|u mT\|\le R\) for \(T=\tau\) or \(w_{\mathbf p}^{-1}\tau\),
the middle block of \(um\) is \(u_jm\) with \(u_j\in U_j\),
of norm at most \(R\|T^{-1}\|\). The inverse has the same
norm and takes \(f_j\) to \(m^{-1}f_j\). Apply
\eqref{eq:realquadrics911} with \(\ell=z\); it is in range
even at \(r_a=k+1\). Alternatively use \(\|g\|_1\|g'\|_1\),
taking their individual canonical sections. In these sections
the radii are \(R\) and the restricted orbit volume is at most
\(\pi^{r_j}R^{2(r_j-1)}\) by polar integration. The norm
invariance in Proposition~\ref{prop:flagenergy923} therefore gives
\[
 |C_P(x;\phi_\nu)|\le
 \pi^{r_j}L_{\mathbf p}^{2(r_j-1-z)}
 \min(R^{2z},|\nu|^{-z})
       \mathcal A_{\mathbf1,2}(A,x)\mathcal A_{\boldsymbol\lambda,2}(B,x).
\]
With the minimum replaced by \(R^{2z}\) the same right side
also bounds \(\|g\|_1\|g'\|_1\).

Apply \eqref{eq:crossflags923}, Cauchy--Schwarz in \(x\), and
\eqref{eq:flagenergy923} to the two squared integrals.
The geometric mean of the two position estimates contains
\(\prod_{i<j}|\lambda_i|^z M_i^{m_i/2}\). Multiplication by the
measure of the form leaves a factor at most \(\prod_{i<j}M_i^{m_i}\),
since \(d_i'-z=(m_i-1)/2\).
The original group conversion is \(2^n V_k/V_r\) and
\(\eta^2\|W\|_{2,*}^4\le d^2\), canceling the type and volume
factors upon normalization. Write
\(\sigma=\sum_{i<j}m_i\), \(s=\sigma+2(a-j)\le10n^2\).
There is an absolute \(C>0\) for which, with \(B_*=R^{Cn^3}H^{Cn^2}\),
\begin{equation}
 I(\|x\|\le Y)\le B_* Y^s\min(R^{2z},|\nu|^{-z}),\qquad Y\ge1 .
 \label{eq:allwaveball924}
\end{equation}
For verification of the dimension dependence, the exponent
of \(\pi\) before division by \(\pi^{(2a+1)k}\) is
\[
       r_j+2\sum_{i<j}r_{i+1}+\sum_{i=j+1}^a(2r_i-1).
\]
The coefficient of \(k\) is \(2a+1\), and the remainder
is \(O(n^2)\) since \(|r_i-k|\le 2n\). Apart from that power
of \(\pi\), the constants outside the minimum and \(Y^s\)
are bounded by
\[
 2^{n+j}(16r^{3/2})^{a-j}
 R^{3\sigma+j+5(a-j)}
 |\delta_{j-1}(\mathbf p)|^{-2}
 (R^{4j}H)^\sigma
 L_{\mathbf p}^{\,2(r_j-1-z)}.
\]
Here \(\prod|p_i|^{-2}\) telescopes. Now \(m_i\le4n+5\),
\(r_j-1-z\le2n+5\), \(L_{\mathbf p}=R^{10n}H_{\mathbf p}^3\)
and \(R\ge\max(16,2r)\). These estimates give the asserted
absolute exponents in \(B_*\).

On \(Y<\|x\|\le2Y\) use the product of the \(L^1\) norms.
By Proposition~\ref{prop:flagtransfer923}, its projection onto
either one-flag domain satisfies \(\|x_j\|>Y/D_*\) with
\(D_*\le R^{42n}H^6\). Keep that restriction in
\eqref{eq:crossflags923}. For \(Y\ge R^2 D_*\) the \(L^1\)
part of Proposition~\ref{prop:flagenergy923} now gives
\begin{equation}
 I(Y<\|x\|\le2Y)
 \le B_* R^{80(a-j+1)k} Y^\sigma
           (Y/D_*)^{-k/(2(a-j))} \qquad(j<a).
 \label{eq:allwaveshell924}
\end{equation}
The same inequality bounds the normalized integral of the
product of the two \(L^1\) norms on the shell. Here the exponent
of \(\pi\) before normalization is
\(2\sum_{i<j}r_{i+1}+(2(a-j)+1)k\); its excess is again
\(O(n^2)\). The extra factor \(2^\sigma\) from the upper
shell radius fits in \(B_*\) after enlarging its absolute
constant. For \(j=a\) these shells are empty on the support.

Take \(X=R^{500n^2}H^{16}\max(1,|\nu|)^{4a}\).
For \(j<a\) and large nodes,
\(k/(2(a-j))-\sigma\ge k/(4a)\), since
\(\sigma\le10n^2\) and \(k/n^3\) tends to infinity.
Sum the shells \(Y=2^l X\), \(l\ge0\). Their total is at most
\[
 2B_*D_*^\sigma R^{80(a-j+1)k}(X/D_*)^{-k/(4a)}
 \le R^{O(n^3)}H^{O(n^2)}\max(1,|\nu|)^{-k}.
\]
Indeed \((500n^2-42n)/(4a)>80(a-j+1)\) eventually.
This also proves convergence at fixed heads for the \(L^1\)
majorant, by combining the shells with the finite ball bound.

For \(|\nu|\le1\) the minimum in \eqref{eq:allwaveball924} is at
most \(R^{2(z-K_0)}|\nu|^{-K_0}\), by splitting at \(R^{-2}\).
When \(|\nu|>1\), the power \(4as\) from \(X^s\) is absorbed
by \(z-K_0-6=n^7-11\) at large nodes. Since \(X^s\) contributes
at most \(R^{O(n^4)}H^{O(n^2)}\) apart from this frequency
power, the total coefficient is
\(R^{2n^7+O(n^4)}H^{O(n^2)}\), with absolute constants.
The first factor is at most \(b^{2000n^9}\) eventually.
The second fits in the product of prefix weights since the
product of the maxima of one and the individual inverse
determinants dominates their maximum on each list, and
\(O(n^2)\le n^6\) eventually. These comparisons have enough
slack also to allow the factor \(2^j\) if the convention
without the coarea halves is used. This proves the theorem.
\end{proof}
\clearpage

\subsection{Unequal last entries of arbitrary real heads}
\label{sec:allmismatch925}

Let \(j=i+1\le a\), let \(\mathbf p,\mathbf q\) have length \(j\),
and keep \(J_{\mathbf p,\mathbf q},\mathcal H,W_0,\boldsymbol\lambda\)
from Section~\ref{sec:allrealwaves924}. Superscript minus denotes
the prefix of length \(i\); the prefix at \(i=0\) is a point
with measure one. Write \(p=p_i,\ q=q_i,\ \lambda=p/q\).

\begin{theorem}[Uniform reciprocal-difference bound for real heads]
\leavevmode\label{thm:allmismatch925}
On the same shorter interval, at sufficiently large nodes,
if \(p\ne q\) and \(\xi=p^{-1}-q^{-1}\), then
\[
 \frac1{d_v S_{k,a}}\int_G |J_{\mathbf p,\mathbf q}(x;0)|\,dx
 \le b^{2000n^9}\mathcal H(\mathbf p)\mathcal H(\mathbf q)W_0(\xi).
\]
\end{theorem}

\begin{proof}
At zero frequency the inner integral is the product of the
two ordinary conditional correlations, the second conjugated.
On the first side all scales are one. The scalar correlation
at a full flag \(A\) is independent of the full flag section:
both amplitudes acquire the same character under its unipotent
part and a common change of residual integration variable under
its residual isometry. In canonical coordinates let \(v_0\)
be the last vector of \(A\) and \(S_A\) the section of \(A^-\).
The modulus of this scalar is
\[
       F(A,x)=|\mathcal D^{\mathbf1}_{S_A,x}(v_0)|.
\]
Here the density on the right is the one at depth \(i\).
Indeed disintegrating \(U_i\backslash G_i\) over \(m^{-1}f_i=v_0\)
gives the full-flag scalar with the canonical last section.
In particular its squared integral in \(v_0\) is
\(\mathcal A_{\mathbf1,2}(A^-,x)^2\).

Fix the prefix pair \(P^-\) and choose a common section \(\tau\).
Write \(T_A=\tau=n_Ah_A S_A\) and
\(T_B=w_{\mathbf p^-}^{-1}\tau=n_Bh_B S_B\), where \(S_B\)
is the canonical section of \(B^-\). The \(n\)'s are in
\(N_{<i}\) and \(h_A,h_B\in G_i\). For \(i=0\) these matrices
are identities, and \(L_{\mathbf p^-}\) may be taken equal to one.
In the last residual space use the coordinates of \(S_B\) for
both vectors. If \(w_0\) is the last canonical vector of \(B\),
the two vectors in these coordinates are
\[
       v=h_B^{-1}h_A v_0,\qquad w=w_0,\qquad h(v,w)=p.
\]
This follows because the two common sections identify the same
middle space. The change there by \(h_B^{-1}\) preserves its
orbit and coarea measures. On the four cut projections
\(\|v_0\|,\|w\|\le R\), both are at least \(R^{-1}\), and
\(\|h_A\|,\|h_B^{-1}\|\le R^iL_{\mathbf p^-}\), as in
Proposition~\ref{prop:flagtransfer923}. Put
\(U=R^{2i+1}L_{\mathbf p^-}^2\); then \(\|v\|\le U\).

At fixed \(w\), choose \(t,u\) by
\(v=s_w^{-1} n(-t,-u)\bar p e_i^*\). The resulting full pair
section is \(w_p n(t,u)s_w h_B^{-1}\tau\), with the four
residual factors embedded at depth \(i\) (the prefix
section \(\tau\) acts on the original space). It sends the
initial last pair to \(e_i,p e_i^*\) and has the prescribed
action on all previous pairs. Its full second-flag section is
\[
             n(t,u)\,\widetilde n\, S,\qquad
 S=s_w S_B,\qquad
 \widetilde n=(s_w h_B^{-1})n_B(s_w h_B^{-1})^{-1}\in N_{<i}.
\]
Here \(\widetilde n\) is fixed when the two prefixes and \(w\)
are fixed. We must account for this earlier unipotent before
taking a Fourier transform in \(t\).

Let \(m\in G_j\), and abbreviate \(n(t,u)\) by \(n'\).
When \(\widetilde n\) is moved left through \(m n'\), its
character is constant in \(m,t,u\): the earlier simple entries
are fixed and the boundary entry depends only on
\((mn')^{-1}e_i=e_i\), by the simple-entry calculation in
Proposition~\ref{prop:flagenergy923}. In the second amplitude
the conjugating matrix is \(m d_{\boldsymbol\lambda} n'\).
Move \(d_{\boldsymbol\lambda^-}\) into \(\widetilde n\) first,
which preserves the earlier radical since this scale commutes
with \(G_i\); the remaining boundary vector is
\((m d_\lambda n')^{-1}e_i=e_i/\lambda\), also fixed.
At \(i=0\) the earlier radical is trivial. Removing these two
constant phases, the characters of \(n'\) itself differ by the
same factor as in Theorem~\ref{thm:firstmismatch921}. Thus the
modulus of the second scalar correlation is
\[
 |T_B(t)|=\left|\int
 e^{2i\operatorname{Im}((1-\lambda)h(t,z'))}
      \mathcal D^{\boldsymbol\lambda}_{S,x}(z')\,d\sigma_j(z')\right|
\]
up to changing the sign of the Fourier variable, which leaves
the following ball estimate unchanged. In particular the density
in this formula is the full canonical one, independent of \(t,u\).

The last conditional measure at fixed \(w\) is
\(\tfrac12|p|^{2r_j}dt\,du\).
The bounds on \(v,w\) give
\(\|t\|,|u|\le T=RU/|p|\) by applying \(s_w\) to \(v\).
The residual orbit in the last display has canonical annular
radius \(R\). Lemma~\ref{lem:truncquad912} applied after the
complex change of Fourier variable by \(\overline{1-\lambda}\)
therefore gives
\begin{align}
 \tfrac12|p|^{2r_j}\int_{\|t\|,|u|\le T}|T_B(t)|^2\,dt\,du
 &\le P_i|\xi|^{1-2r_j}
           \mathcal A_{\boldsymbol\lambda,2}(B,x)^2, \notag\\
 P_i&=2(2r_j)^{3/2}\pi^{2r_j-1}R(RU)^2|p|^{-1}.
 \label{eq:lastflagplancherel925}
\end{align}
The scaled ball has radius \(|1-\lambda|T\), its Jacobian
is \(|1-\lambda|^{2r_j}\), and the central interval has length
\(2T\). Substituting \(|1-\lambda|=|p|\,|\xi|\) gives the
displayed single inverse power of \(|p|\).

We describe the weighted integrations, since the first squared
integral uses one fewer flag vector than the second. Put
\(D_A=D_{\mathbf1,x}(A^-)\) at depth \(i\) and
\(D_B=D_{\boldsymbol\lambda,x}(B)\) at depth \(j\).
Apply Cauchy--Schwarz to
\(F(A,x)D_A^z D_B^{-z}\) and
\(|T_B(t)|D_B^z D_A^{-z}\) on the supported full pairs. The
first squared integral is bounded by
\[
  C_{\mathbf p}\int_{X_i}D_A^{2z}
               \mathcal A_{\mathbf1,2}(A^-,x)^2\,d\Sigma_i(A^-).
\]
Indeed first integrate \(D_B^{-2z}\) at fixed full \(A\)
by the conditional negative moment in the proof of
Theorem~\ref{thm:crossflags923}, then integrate \(F^2\) in
\(v_0\). For the second integral disintegrate first at the
prefix pair and \(w\), using Theorem~\ref{thm:flagmeasure922}.
The measure in \(w\) is \(d\sigma_i(w)\); the change from common
to canonical middle coordinates preserves it for each prefix
pair. Apply \eqref{eq:lastflagplancherel925} in \(t,u\), keeping
the initial support until the uniform value of \(U\) has
been inserted. The remaining dependence on \(A^-\) is only
\(D_A^{-2z}\) and admissibility. Integrating over the prefix
fiber at fixed \(B^-\) costs \(C_{\mathbf p^-}\) by the exchanged
negative moment. Thus the second squared integral is at most
\[
 P_i C_{\mathbf p^-}|\xi|^{1-2r_j}\int_{X_j}
        D_B^{2z}\mathcal A_{\boldsymbol\lambda,2}(B,x)^2\,d\Sigma_j(B).
\]

Integrate the geometric mean in \(\|x\|\le Y\) by
Cauchy--Schwarz and \eqref{eq:flagenergy923} at depths \(i,j\).
At \(i=0\) its first factor is just \eqref{eq:recenergy912}.
The position normalization and factors
\(\prod|\lambda_l|^{z-r_{l+1}}M_l^{m_l/2}\) are the same as in
Theorem~\ref{thm:allrealwaves924}. The total exponent of \(\pi\)
before normalization is
\[
 2\sum_{l<i}r_{l+1}+3r_j-1+\sum_{l=j+1}^a(2r_l-1)
                 =(2a+1)k+O(n^2).
\]
With \(H=H_{\mathbf p}H_{\mathbf q}\), all constants outside
the frequency power are \(R^{O(n^3)}H^{O(n^2)}\). To check the
new factors, \(C_{\mathbf p^-}\) has the same product bound as
\(C_{\mathbf p}\) with one fewer factor, \(U\le R^{22n}H^6\),
and \(|p|^{-1}\le R^{2j}H\). The power of \(Y\) is the
average of \(s_i=\sum_{l<i}m_l+2(a-i)\) and \(s_j\);
both are at most \(10n^2\). Together with the constant-weight
ball bound in the proof of \eqref{eq:allwaveball924}, the
normalized mass is at most
\[
   R^{O(n^3)}H^{O(n^2)}Y^{10n^2}
           \min(R^{2z},|\xi|^{1/2-r_j}).
\]
For the tail the product of the two \(L^1\) norms from
\eqref{eq:allwaveshell924} still applies. Take
\(X=R^{500n^2}H^{16}\max(1,|\xi|)^{4a}\) as before. For
\(|\xi|\le R^{-2}\) the radius bound costs at most
\(R^{2(z-K_0)}|\xi|^{-K_0}\). For \(R^{-2}\le|\xi|\le1\)
the frequency bound costs at most
\(R^{2(r_j-1/2-K_0)}|\xi|^{-K_0}\). For \(|\xi|>1\) the
difference \(r_j-1/2-K_0\) exceeds \(4a\cdot10n^2+6\) eventually.
The shell sum contributes at most
\(R^{O(n^3)}H^{O(n^2)}\max(1,|\xi|)^{-k}\).
All constants are absolute, so the same final conversion as in
\eqref{eq:allrealwaves924} gives the asserted inequality,
including its determinant support restrictions.
\end{proof}
\clearpage

\subsection{Fourier descent through the real null rows}
\label{sec:realnulldescent926}

For \(0\le j\le a\), let \(\mathcal J_{\mathbf p,\mathbf q}(x;\phi)\)
be the form of Section~\ref{sec:allrealwaves924} with \(C_P(x;\phi)\)
in place of \(C_P(x;\phi_\nu)\); thus
\(J_{\mathbf p,\mathbf q}(x;\nu)=\mathcal J_{\mathbf p,\mathbf q}(x;\phi_\nu)\).
Use \(J_{\mathbf p,\mathbf q}(x;0)\) for its constant-weight value.
Appending \(c\ne0\) to both lists will be denoted by
\(\mathbf p c,\mathbf q c\). Use planar Lebesgue measure at this
real place and the Fourier kernel \(\phi_\xi(c)=e^{2i\operatorname{Im}(\xi c)}\).

\begin{theorem}[Pointwise Fourier descent at the null rows]
\leavevmode\label{thm:realnulldescent926}
At the nodes of Theorem~\ref{thm:allrealwaves924}, let \(0\le j<a\)
and fix any two nonzero pivot lists of length \(j\), including
the empty lists. Then
\[
  \int_{\mathbb C^\times}\int_G
       |J_{\mathbf p c,\mathbf q c}(x;0)|\,dx\,d^2c<\infty .
\]
Define \(f(c)=\int_G J_{\mathbf p c,\mathbf q c}(x;0)\,dx\)
for \(c\ne0\). For every bounded measurable \(\phi\),
\[
  \int_{\mathbb C^\times}\phi(c)f(c)\,d^2c
       =\int_G\mathcal J_{\mathbf p,\mathbf q}(x;\phi)\,dx.
\]
For every \(c\ne0\), with absolutely convergent integrals,
\begin{equation}
 f(c)=\pi^{-2}\int_{\mathbb C} e^{-2i\operatorname{Im}(\xi c)}
                 \left[\int_G J_{\mathbf p,\mathbf q}(x;\xi)\,dx\right]
                       d^2\xi .
 \label{eq:realnulldescent926}
\end{equation}
The right-hand side extends \(f\) to a compactly supported
\(C^h\) function of the two real coordinates for every nonnegative
integer \(h<K_0+4\).
\end{theorem}

\begin{proof}
Fix the prefix pair \(P\) of length \(j\) and its common section
\(\tau\). After integrating the later residual stabilizers, the
two orbit variables \(v,w\) at this step have densities
\(g=\mathcal D^{\mathbf1}_{\tau,x}\) and
\(g'=\mathcal D^{\boldsymbol\lambda}_{w_{\mathbf p}^{-1}\tau,x}\).
The step is null because \(j<a\).
For fixed \(v\) choose a residual section \(s_v\), and at
\(c=h(v,w)\ne0\) write \(w=s_v^{-1} n(t,u)c e_j^*\).
The measure in \(w\) is
\(\tfrac12|c|^{2r_{j+1}}d^2c\,dt\,du\).
The excluded fiber \(c=0\) has zero orbit measure at fixed \(v\).
Indeed in coordinates \(\alpha e+y+c e^*\) its nonzero middle
points allow a middle real coordinate to solve the null equation,
leaving \(c\) free. The remaining points on the complex line
\(\mathbb C e\) have zero measure in the smooth nonzero orbit.

Take the full pair section \(\tau_+=n(t,u)^{-1}s_v\tau\).
Its two flag sections are this one and
\(w_{\mathbf p c}^{-1}\tau_+\). At fixed \(v,w\) the first
scalar \(g(v)\) is the integral of its next density in
the section \(\tau_+\). The analogous integral in the other
section gives \(g'(w)\). To see that the scalar is unaffected
by either section choice, at the residual step any two sections
for one fixed vector differ by \(n_0 l\in N_j\rtimes G_{j+1}\).
The preceding diagonal scales commute with this whole residual
group, and the new scale is \(c/c=1\). Thus the two amplitudes
in each scalar acquire the same unipotent character and a common
right translation in the residual variable. Their scalar
integral is unchanged. These statements follow directly by
disintegrating the quotient \(U_j\backslash G_j\) at the fixed
vector, including at \(j+1=a\).

Thus the product \(g(v)\overline{g'(w)}\) is the constant-weight
inner integral at the extended pair. Its coarea and the old
prefix measure combine to \(d^2c\,d\mu_{\mathbf p c}\) by
Theorem~\ref{thm:flagmeasure922}. The outside scale factor is
unchanged since the new scale is one. All quotient constants
are therefore exactly those of \(J_{\mathbf p c,\mathbf q c}\).
Taking absolute values after the two scalar integrations and
then integrating in \(c\) costs at most the old integral of
\(\|g\|_1\|g'\|_1\), with the old outside flag and position
measures. This majorant is integrable by
Theorem~\ref{thm:allrealwaves924}. Fubini now proves both the
first assertion and the identity for bounded \(\phi\).

For pointwise inversion we show that \(f\) is continuous away
from zero. First restrict positions to \(\|x\|\le Y\) and \(c\)
to a compact subannulus of \(\mathbb C^\times\).
All contributing full flags are projections of the original
compact cut support in \(U_0\backslash G_0\). Thus the initial
flag pairs lie in a compact set. All their prescribed pivots,
including the varying last one, stay bounded away from zero.
On this set choose a finite covering by smooth local pair
sections and a subordinate partition of unity. The orbit and
coarea calculations in Theorem~\ref{thm:flagmeasure922} supply
local smooth fiber coordinates at every such pair. The two
amplitudes are continuous in the parameters; in each residual
quotient the first amplitude restricts integration to a common
compact set. Indeed its residual radius is bounded by \(R\)
times the maximum norm of the inverse local section, by taking
the middle block. The flag integrations also have common compact
support after localization. Dominated convergence in these
coordinates gives continuity of the bounded-position integral.
The part \(\|x\|>Y\) tends to zero uniformly on the subannulus:
the weights \(H_{\mathbf p c},H_{\mathbf q c}\), and hence \(D_*\)
in \eqref{eq:allwaveshell924}, are uniformly bounded there.
The dyadic shell powers are summable. If the new depth is \(a\)
the tail vanishes for all sufficiently large \(Y\).
This proves the asserted continuity.

The bounded-weight identity identifies the Fourier transform of
\(f\) with the bracket in \eqref{eq:realnulldescent926}. For
\(|\xi|\le1\) its absolute value is bounded by the integral
of the old \(L^1\) majorant. For \(|\xi|>1\) it is at most a
constant depending on the fixed heads and node times
\(|\xi|^{-K_0-6}\), by Theorem~\ref{thm:allrealwaves924}.
The Fourier transform is integrable in the plane, also after
multiplication by \(|\xi|^h\) for \(h<K_0+4\).
Ordinary Fourier inversion identifies the displayed integral
with \(f\) almost everywhere. The inversion constant is
\(\pi^{-2}\), since \(2\operatorname{Im}(\xi c)\) is the real
two-dimensional dot product with an orthogonally transformed
frequency scaled by two. Continuity gives equality at every
nonzero \(c\); differentiation under the integrable transform
gives the specified regularity at all \(c\).
For compact support, \(|c\,\delta_{j-1}(\mathbf p)|\) and
\(|c\,\delta_{j-1}(\mathbf q)|\) are both at most
\(R^{2(j+1)}\) on the support, also for \(j=0\) with
\(\delta_{-1}=1\). These inequalities bound \(|c|\) for
the fixed heads.
\end{proof}
\clearpage

\subsection{The terminal real density and the orbital tensor}
\label{sec:terminaltensor927}

At depth \(a\) the residual space is \(\mathbb C^k\oplus\mathbb C z\)
with form \(\|u\|^2-|t|^2\). We use \(z=(0,1)\),
\(C=\mathrm U(k)\) with probability Haar and
\(S_k=2\pi^k/(k-1)!\). Write \(dx\) for the original group
Haar from Theorem~\ref{thm:realgram682} and put
\[
  c_G=V_k S_k/V_r,\qquad
  \kappa_G=2^n V_k/V_r,\qquad dx=\kappa_G\,dg_0^*(x).
\]
For \(|c|>1\) set \(s=(|c|^2-1)^{1/2}\); let \(\sigma_c\)
fix all but the first positive coordinate and \(z\), on which
its matrix in that order is
\(\left(\begin{smallmatrix}-\bar c&s\\s&-c\end{smallmatrix}\right)\).
Its columns are Hermitian orthogonal of norms \(1,-1\), and
\(h(z,\sigma_c z)=c\). For a null prefix \(\mathbf p\) of
length \(a\) put \(\gamma_p=w_{\mathbf p}\sigma_c\) and define
the matrix
\begin{align}
 O_p(c)&=\int_G W(y) W(\gamma_p^{-1}y)^*\,dy,\notag\\
 \Xi_p(c)&=\frac{\sqrt{c_G}}{d_v}(|c|^2-1)^{(k-1)/2}
          \prod_{l<a}|p_l|^{r_{l+1}}\,O_p(c).
 \label{eq:terminaltensor927}
\end{align}
Set \(\Xi_p(c)=0\) for \(|c|\le1\). Matrix inner products below
are Hilbert--Schmidt inner products, conjugate linear in their
first argument.

\begin{theorem}[The terminal Gram density in the original Haar measures]
\leavevmode\label{thm:terminaltensor927}
At every real place of the shorter family the tensor in
\eqref{eq:terminaltensor927} is continuous jointly in
\(\mathbf p\in(\mathbb C^\times)^a\) and \(c\).
For any two such prefixes and any bounded measurable \(\phi\),
\begin{equation}
 \int_{\mathbb C}\phi(c)\langle\Xi_p(c),\Xi_q(c)\rangle_{\rm HS}\,d^2c
 =\frac1{d_v^2}\int_G\mathcal J_{\mathbf p,\mathbf q}(x;\phi)\,dx .
 \label{eq:terminaldescent927}
\end{equation}
Both integrals converge absolutely. The common-entry Gram value
is unchanged if \(\sigma_c\) in both orbital tensors is replaced
by any residual isometry \(\sigma\) with \(h(z,\sigma z)=c\)
and \(|c|>1\).
\end{theorem}

\begin{proof}
We first identify the quotient constant
\(d\dot y=\eta\,d\dot y_0^*\) used in the preceding formulas.
The original Haar on \(U\) is the product of Lebesgue measure on
the upper entries of the frame group \(N_n\) and the measures
\(dB\,dK\) on its kernel, as specified in
Section~\ref{sec:gramtrace683} and Lemma~\ref{lem:unipotentvolume703}.
This agrees exactly with the successive \(dt_l\,du_l\) on the
Heisenberg groups \(N_l\). It suffices to check the two Haar
densities at the identity. In the tangent map of \(n_l(t_l,u_l)\),
the component of \(t_l\) along a later \(e_h^*\) yields the
upper entry in row \(l\), column \(h\) of the frame group by
negative conjugation. Its \(z\) component yields the entry of the
terminal column by conjugation, since \(h(z,z)=-1\).
The \(k\) positive components form column \(l\) of \(B\).
For each \(h>l\) the component along \(e_h\) gives the skew
entry \(K_{lh}\) by negative conjugation, and \(u_l\) gives
\(\operatorname{Im} K_{ll}\). These are all coordinates:
the quadratic terms in the Heisenberg matrix have zero
differential at the identity. The resulting real Jacobian has
absolute value one. Uniqueness of Haar gives equality globally.
Dividing \(dx=\kappa_G dg_0^*\) by this common unipotent
measure shows \(\eta=\kappa_G\).

For fixed \(\mathbf p\) and \(|c|>1\) the orbital integral in
\eqref{eq:terminaltensor927} is an ordinary absolutely convergent
integral. In fact if its integrand is nonzero, the two flags
represented by \(y\) and \(\gamma_p^{-1}y\) both occur on the
cut. Their already normalized columns are \(y^{-1}e_l\) and
\(p_l y^{-1}e_l^*\), since \(\sigma_c\) fixes the earlier
hyperbolic pairs. Thus \(y=m\tau_P\) for the bounded pair
section of Proposition~\ref{prop:flagsection922} and some
\(m\in G_a\). From \(W(y)\ne0\) we obtain
\(\|u m\tau_P\|\le R\) for some \(u\in U\). Its residual
middle block after removal of \(\tau_P\) is \(m\), since
\(U_a=1\), so \(\|m\|\le R\|\tau_P^{-1}\|\le L_{\mathbf p}\).
In particular \(\|y\|\le L_{\mathbf p}^2\). This gives a
common compact set when the prefixes range in a compact set
of nonzero entries. Dominated convergence for the continuous
cut proves continuity on \(|c|>1\). The orbital integral stays
bounded as \(|c|\downarrow1\); its factor
\((|c|^2-1)^{(k-1)/2}\) tends to zero since \(k>1\), proving
continuity there too.

For \(v,w\) on the negative unit orbit and a section \(m(v)\)
with \(m(v)v=z\), write \(m(v)w=(u,-c)\) where \(c=h(v,w)\).
The orbit equation is \(\|u\|^2=|c|^2-1\). Polar integration
in the \(k\) complex coordinates of \(u\) gives the measure
\[
  d\sigma_a(w)=\kappa(c)d^2c\,d\omega,\qquad
  \kappa(c)=\tfrac12 S_k (|c|^2-1)^{k-1}{\boldsymbol 1}_{|c|>1}
                         \quad (|c|\ge1)
\]
with zero density on \(|c|<1\); \(d\omega\) is probability
sphere measure. The boundary has measure zero: near a point
with \(u=0\) the nonzero terminal coordinate permits solving
the orbit equation in its radial variable, leaving \(u\) free.
For \(|c|>1\) write \(w=m(v)^{-1}h\sigma_c z\). As \(h\)
varies in \(C\) with probability Haar its positive direction
has law \(d\omega\), by transitivity and invariance.

Fix a prefix pair \(P\), with section \(\tau\). At the terminal
level a scalar density at \(v\) is the inner product at any
section \(m(v)\), since left translation by \(C\) preserves
the inner product of the two amplitudes. Similarly use
\(\sigma_c^{-1}h^{-1}m(v)\) for the second orbit vector \(w\).
In the first inner product replace \(m(v)\) by \(h^{-1}m(v)\)
as well; it is unchanged by compact equivariance. Integrating
\(v,h\) now integrates \(m=h^{-1}m(v)\) over \(G_a\) in
\(dg_a^*\). The residual group commutes with the earlier Weyl
swaps. It follows that after extracting \(\kappa(c)d^2c\)
the two inner products become
\[
  \langle W(y),W(yx)\rangle
       \langle W(\gamma_q^{-1}yx),W(\gamma_p^{-1}y)\rangle,
                  \qquad y=m\tau .
\]
The order of the second pair is reversed because the second
scalar was conjugated. The \(P\)-measure after its endpoint
scaling is \(2^{-a}\prod_{l<a}|p_l q_l|^{r_{l+1}}\,d\beta_{\mathbf p}\).
By its quotient normalization \(d\beta_{\mathbf p}\,dg_a^*(m)\)
integrates \(y\) in \(dg_0^*(y)\).
Theorem~\ref{thm:allrealwaves924} bounds the integral of the
product of the terminal scalar moduli, over the two orbits,
the prefix pairs and all positions. The above rewriting
preserves those moduli, so Fubini applies before identifying
the integrand at \(c\).

Finally put \(y'=yx\). The product \(dg_0^*(y)\,dx\) becomes
\(\kappa_G^{-1}dy\,dy'\) in original Haar. Coordinate expansion
of the displayed pair of inner products now gives exactly
\(\langle O_p(c),O_q(c)\rangle_{\rm HS}\) after integration in
\(dy\,dy'\). The remaining prefactor on the right side of
\eqref{eq:terminaldescent927} is
\[
  \frac{\eta^2}{d_v^2\kappa_G}\,2^{-a}\kappa(c)
                              \prod_{l<a}|p_l q_l|^{r_{l+1}}.
\]
Since \(n=a+1\) and \(\eta=\kappa_G=2^n V_k/V_r\), this is
\(c_G (|c|^2-1)^{k-1}\prod_{l<a}|p_l q_l|^{r_{l+1}}/d_v^2\)
on \(|c|>1\), as in \eqref{eq:terminaltensor927}.
If \(\sigma\) has the same \(c\), choose \(h_1\in C\) to rotate
the first positive direction of \(\sigma_c z\) to that of
\(\sigma z\). Then \(h_2=(h_1\sigma_c)^{-1}\sigma\) fixes
\(z\), so belongs to \(C\). In the orbital integral with
\(\sigma=h_1\sigma_c h_2\), set \(y=h_1 y_0\). Compact
equivariance and commutation with \(w_{\mathbf p}\) multiply
every tensor, for fixed \(\sigma\), by the same unitary matrices
on its two sides. Its Gram products are therefore unchanged.
\end{proof}

\begin{corollary}[The complete real-place family of Gram forms]
\leavevmode\label{cor:realsystem927}
For \(0\le j\le a\) take
\(L_j(\mathbf p,\mathbf q)=d_v^{-2}\int_G
J_{\mathbf p,\mathbf q}(x;0)\,dx\). For lists of length \(n\),
with respective null prefixes \(\mathbf p,\mathbf q\) and
terminal entries \(c,c'\), take
\(L_n=\langle\Xi_p(c),\Xi_q(c')\rangle_{\rm HS}\).
Then on the shorter interval, at each real place, the bottom
value satisfies
\[
 (1-b^{-r^3 f})^4/d_{\pi,v}\ \le L_0\le
                       (1+b^{-20})/d_{\pi,v}.
\]
At every \(j<n\), integration of \(L_{j+1}\) in a common next
entry yields \(L_j\), absolutely. At a nonzero Fourier frequency
\(\xi\) the absolute value of that integral with weight
\(\phi_\xi\) is bounded by
\[
     d_v^{-1} S_{k,a} b^{2000n^9}\,
                     \mathcal H(\mathbf p)\mathcal H(\mathbf q)W_0(\xi).
\]
At each \(1\le j\le a\) the same bound holds for \(L_j\) itself
when \(\xi\) is the nonzero difference of the reciprocal last
entries. The top form is continuous and positive semidefinite
on any finite set of lists.
\end{corollary}
\begin{proof}
The bottom inequality is \eqref{eq:realcutmass916}. The common
entry identities are Theorem~\ref{thm:realnulldescent926} up to
the last null row and \eqref{eq:terminaldescent927} at the
terminal row. The absolute estimates are
\eqref{eq:allrealwaves924} and Theorem~\ref{thm:allmismatch925}
after division by \(d_v^2\). Continuity of the orbital tensors
was proved above. Every finite quadratic form for the top
kernel is the squared Hilbert--Schmidt norm of the corresponding
linear combination of those tensors.
\end{proof}
\clearpage

\subsection{Packing the rational prefix determinants}
\label{sec:prefixpack928}

Use the shorter period fields \(E=\mathbb Q(\sqrt{-d_0},\mu_{q_1})\),
\(F=E^+\) and \(f=[F:\mathbb Q]\) of
Theorem~\ref{thm:allunmarked909}. Number the entries in a rational
list starting at one in the formulas for the field scales.
Write \(\Phi_v\) for one embedding at each conjugate pair. Put
\[
 B=b^{60n^2},\qquad
 s_{l,v}=\prod_{i=1}^{l} b_{i,v},\qquad
 b_{i,v}=\kappa_v^{2(n-i)}/N,\qquad
 S_l=\left(\prod_v s_{l,v}\right)^{1/f}.
\]
These are the scales of Theorem~\ref{thm:freqnode875}.
For a list \(\mathbf p=(p_1,\ldots,p_j)\) in \(E^\times\) let
\(\delta_l=\prod_{i=1}^l p_i\). Its real prefixes at \(v\)
are \(\mathbf p_v=(b_{i,v}\Phi_v p_i)_{1\le i\le j}\).
For a finite set \(T\) of primes of \(E\), set
\(I_T(y)=\prod_{\mathfrak p\in T}{\rm N}\mathfrak p^{\operatorname{ord}_{\mathfrak p}y}\)
and define
\[
 H_T(\mathbf p)=\prod_{l=1}^j
    {\boldsymbol 1}_{\delta_l\ {\rm integral\ at}\ T}
                     I_T(\delta_l)^{n^6/2},\qquad
 H(\mathbf p)=H_T(\mathbf p)\prod_v \mathcal H(\mathbf p_v).
\]
The local real weight \(\mathcal H\) is that in
Theorem~\ref{thm:allrealwaves924}. These definitions keep the
finite weights at their original values even when a prefix has
zeros at primes in \(T\).

\begin{theorem}[A field-uniform sum of the coupled prefix weights]
\leavevmode\label{thm:prefixpack928}
For \(0\le j\le a\) sum over all lists of length \(j\) whose
prefix products are integral outside \(T\). Then
\begin{align}
 \sum_{\mathbf p} H(\mathbf p)
 &\le \prod_{l=1}^{j}
       (1+2B/S_l)^{2f}\max(1,B/S_l)^{n^6 f}
 \le b^{130n^9 f}
 \label{eq:prefixpack928}
\end{align}
at sufficiently large nodes, uniformly in \(T\) and \(j\).
Only finitely many terms have positive weight. In particular the
double sum of the product of two such weights is at most
\(b^{260n^9f}=b^{o(kf)}\), also with arbitrary further restrictions
on the pair of heads.
\end{theorem}
\begin{proof}
A positive finite weight makes every \(\delta_l\) integral at
the primes of \(T\); the hypothesis gives integrality elsewhere.
Thus each \(\delta_l\) is a nonzero algebraic integer.
For any such \(\delta\),
\[
 I_T(\delta)^{1/2}\le |{\rm N}_{E/\mathbb Q}\delta|^{1/2}
                         =\prod_v|\Phi_v\delta|.
\]
For \(\delta=\delta_l\), multiply this inequality by the
product over all real places of
\(\max(1,|s_{l,v}\Phi_v\delta|^{-1})\). The right side becomes
\(\prod_v s_{l,v}^{-1}\max(1,|s_{l,v}\Phi_v\delta|)\), which is
at most \((B/S_l)^f\) on the support. Raising to \(n^6\) bounds
the combined weight of this prefix. This uses the product of
the actual local scales and allows them to differ at the
various embeddings.

It remains to count the possible \(\delta_l\). Map
\(\delta\in\mathcal O_E\) into \(\mathbb C^f\) by the coordinates
\((s_{l,v}/S_l)\Phi_v\delta\). Distinct images have distance at
least one in the maximum of the complex absolute values, since
the product of the squared absolute differences is the absolute
norm of a nonzero algebraic integer. Positive weight places
every coordinate in the disc of radius \(B/S_l\). The products
of open discs of radius \(1/2\) centered at these images are
disjoint and lie in the product of discs of radius \(B/S_l+1/2\).
Comparison of ordinary Lebesgue volumes gives at most
\((1+2B/S_l)^{2f}\) images. This argument on finite subsets
also proves finiteness. The tuple of the \(j\) successive
determinants determines the list, using \(p_l=\delta_l/\delta_{l-1}\)
with \(\delta_0=1\). Counting all the tuples and multiplying the
weight bounds gives the first inequality and the assertion for
the empty list.

On our fields \(Q_E=(\prod_v\kappa_v)^{1/f}=M\rho_E/(2\pi)\ge1\),
where \(\rho_E=|\operatorname{disc}E|^{1/(2f)}\). Hence
\[
 S_l=N^{-l}Q_E^{\,2nl-l(l+1)}\ge N^{-l}.
\]
The calibrated choice in Lemma~\ref{lem:calibratedlevel703}
satisfies \(N\le b^{60n}\) at large nodes; its upper endpoint
there is at most \(b^{24(2n)}e^{\sqrt{\log b}}\). Consequently
\(B/S_l\le b^{120n^2}\) for \(l\le a\). For \(b\ge3\) we have
\(1+2b^{120n^2}\le b^{120n^2+1}\). The logarithm to base \(b\)
of the first bound in \eqref{eq:prefixpack928} is at most
\(jf(120n^8+240n^2+2)\le130 n^9 f\) for large \(n\).
Finally \(n^9/k\to0\) on this same interval, as in
Corollary~\ref{cor:adelicfirst915}. Squaring the bound proves
the assertion for pairs, since all weights are nonnegative.
\end{proof}

Here is the geometric reason for integrality on the finite
support. At a split or unramified place with the integral
lattice of Theorem~\ref{thm:allunmarked909}, two flags arising
from points of \(UK\), with \(K\) the integral isometry group,
have simultaneous integral representatives. Indeed each point
has a representative in \(K\) after a left unipotent translation,
and that translation changes each of the first columns only by
earlier columns. The cross Gram determinant of the first \(l\)
columns is thus integral; it is invariant under both column
changes. The same fact holds when \(K\) is replaced by any
subgroup of \(K\) on a selected split mark. This proves the
integrality assertion for the determinants even if some of the
individual ratios of successive determinants have poles.
Theorem~\ref{thm:prefixpack928} is an estimate for weights of
these rational lists. The finite correlation at a prescribed
list is a separate integral.
\clearpage

\subsection{Rational frequency errors in units of the identity mass}
\label{sec:identityfrequency929}

The normalized real wave bound contains the factor
\(S_{k,a}^f/d_B\), where \(d_B=\prod_v d_v\).
The product of the real bottom forms is of order
\(1/D_\pi\), with \(D_\pi=\prod_v d_{\pi,v}\).
We compare these scales before summing the rational frequencies.
Use \(m=2n\), \(h=r-k=m-1\), \(K=k-n^7\), the terminal
prime \(N\) and the geometric mean \(w\) of the \(w_v\) in
Lemma~\ref{lem:calibratedlevel703}.

\begin{lemma}[A formal-degree bound matched to the frequency scale]
\leavevmode\label{lem:identityscale929}
On these fields and this interval,
\[
          D_\pi S_{k,a}^f
           \le w^{m(k-1)f}\,b^{o(kf)}.
\]
Consequently, for each of the two choices of \(c_v\) in
Theorem~\ref{thm:freqnode875}, at any one of the \(n\) rows and
for any finite set \(T\) of primes of \(E\),
\[
 D_\pi\, S_{k,a}^f
       \sum_{0\ne y\in\mathcal O_{E,T}} W_{T,c,K}(y)
          \le b^{-(25/2-o(1))kf}.
\]
The terms tending to zero are uniform in the row and \(T\).
\end{lemma}
\begin{proof}
By Theorem~\ref{thm:norm91}, the formal degree is the compact
Weyl dimension at the same regular parameter. Its first \(k-1\)
entries are consecutive and the other \(m\) entries are
\(-w_v-2lr^8\), \(1\le l\le m\), by
Proposition~\ref{prop:stationaryweights695}. In the Vandermonde
formula the first block cancels the denominator factorials
\(1!,\ldots,(k-2)!\). The remaining denominator is
\(\prod_{t=k-1}^{r-1}t!\). Each of the \(m(k-1)\) mixed gaps
is at most \(w_v+2mr^8+k\), and the product of the gaps in
the last block is
\((2r^8)^{m(m-1)/2}\prod_{l=1}^{m-1}l!\).
On the other hand the compact volume formula gives exactly
\[
 S_{k,a}=\pi^{hk}V_k/V_r
           =2^{-h}\pi^{-h(h+1)/2}\prod_{t=k}^{r-1}t!.
\]
All but \((k-1)!\) in the remaining denominator cancel.
Using \(\log(1+t)\le t\), summing over \(v\), and discarding
the remaining denominator for an upper bound gives
\[
 \log(D_\pi S_{k,a}^f)
 \le fm(k-1)\log w+
 O\!\left(fm^2\log(2r)+\frac{fm k(2mr^8+k)}{\min_v w_v}\right).
\]
Here \(\min_v w_v\ge b^{1-10^{-6}}\), \(m=O(\sqrt u)\),
\(r=O(u^6)\), \(k\asymp u^5\) and \(\log b\ge(\log u)^6\)
at the large nodes of size \(u\); the same weight estimates
were used in Lemma~\ref{lem:massscale869}. Both error terms are
\(o(kf\log b)\), giving the first assertion.

Use the first, scale-dependent bound in \eqref{eq:actualfreq875}.
After multiplication by the first assertion, the factors involving
\(M,w,b\), apart from \(b^{o(kf)}\), are
\[
          C^f (b^{15/2}/M)^{Kf}w^{m(n^7-1)f}.
\]
Since \(M\ge b^{20}\), \(1\le w\le b\), \(m n^7=o(k)\)
and \(n^7=o(k)\), their product is at most
\(b^{-(25/2-o(1))kf}\). The frequency bound was uniform
in \(T\) and the row, so this proves the claim.
\end{proof}

For the scaled heads \(\mathbf p_v,\mathbf q_v\) of
Section~\ref{sec:prefixpack928} abbreviate the actual real mass by
\[
 A_{j,v}(\mathbf p,\mathbf q;\nu)=
     d_v^{-2}\int_G |J_{\mathbf p_v,\mathbf q_v}(x;\nu)|\,dx.
\]
Write \(V_T(y)=I_T(y)^{K/2}D_T(y)^{-3}\), using the denominator
of \eqref{eq:freqweight875}. Every rational list below has its
prefix products integral off \(T\). For \(j=0\) use the empty list.

\begin{theorem}[Summed rational head and frequency bounds at the identity scale]
\leavevmode\label{thm:identityfrequency929}
For all sufficiently large nodes, uniformly in \(0\le j\le a\)
and finite \(T\), put \(c_v=2\pi i/(\Phi_v(\delta_E)b_{j+1,v})\).
Then
\begin{align}
 D_\pi \sum_{\mathbf p,\mathbf q} H_T(\mathbf p)H_T(\mathbf q)
     \sum_{0\ne y\in\mathcal O_{E,T}} V_T(y)
              \prod_v A_{j,v}(\mathbf p,\mathbf q;c_v\Phi_v y)
       &\le b^{-12kf}.
 \label{eq:identityfrequency929}
\end{align}
For \(1\le j\le a\), restrict to lists with unequal last entries
and \(y=(p_j^{-1}-q_j^{-1})/M\) integral off \(T\). Then the
same bound holds for
\[
 D_\pi\sum_{\mathbf p,\mathbf q} H_T(\mathbf p)H_T(\mathbf q)
                    V_T(y)\prod_v A_{j,v}(\mathbf p,\mathbf q;0).
\]
Both conclusions allow arbitrary further restrictions on the
summation sets.
\end{theorem}
\begin{proof}
At nonzero frequency,
Theorem~\ref{thm:allrealwaves924} bounds the product of real
masses by
\[
 \frac{S_{k,a}^f}{d_B}\,b^{2000n^9 f}
   \prod_v\mathcal H(\mathbf p_v)\mathcal H(\mathbf q_v)
              \prod_v W_0(c_v\Phi_v y).
\]
Together with \(V_T(y)\) the frequency factors form precisely
\(W_{T,c,K}(y)\). The remaining head factors, including the
two finite weights in the theorem, sum to at most \(b^{260n^9f}\)
by Theorem~\ref{thm:prefixpack928}. For unequal last entries
Theorem~\ref{thm:allmismatch925} gives the same estimate for
the constant-weight masses, since the reciprocal difference at
\(v\) equals \(c_v\Phi_v y\) with \(c_v=M/b_{j,v}\). For each
pair in this case the single nonnegative \(W_{T,c,K}(y)\) is
bounded by its sum over all the nonzero \(T\)-integers. Thus
both cases are bounded by
\(d_B^{-1}b^{2260 n^9 f}b^{-(25/2-o(1))kf}\)
using Lemma~\ref{lem:identityscale929}. Now \(d_B\ge1\) and
\(n^9/k\to0\); this proves the stated inequalities. No
multiplicity assumption on the rational prefixes is required.
\end{proof}

\begin{corollary}[The actual first-frequency error at the original mass scale]
\leavevmode\label{cor:actualidentity929}
For the original selected and unselected marks \(\Xi,\mathcal V\)
of Corollary~\ref{cor:adelicfirst915}, write
\(c_v=2\pi i/(\Phi_v(\delta_E)b_{1,v})\). On the shorter interval,
at sufficiently large nodes,
\begin{align*}
 D_\pi\sum_{\substack{0\ne y\in E\\ y\ {\rm integral\ above}\ \mathcal V}}
 &\prod_v A_{0,v}(\varnothing,\varnothing;c_v\Phi_v y)
  \prod_{t\in\Xi}\frac{B_0 q_t^{\,n-j_t}\mathcal B_{Z_t}(y)}{Z_t^0}
  \prod_{t\in\mathcal U}\mathcal B_t^0(y)
       \ \le\ b^{-7kf}.
\end{align*}
\end{corollary}
\begin{proof}
First take \(T\) conjugation-stable, containing all selected
places and no primes above the unselected marks. For rational
\(T\)-integers the product of finite factors divided by \(V_T(y)\)
is at most \(\exp(O(\epsilon)) b^{5kf}\). For the unmarked factors
this follows from \eqref{eq:allunmarked909}: at a pole use
\(\min(1,V_t)\le V_t\), and at an integral prime under \(T\)
one has \(V_t\ge1\). Outside \(T\) the frequency is integral.
At a selected mark the proof of \eqref{eq:firstsample898}
bounds \(\mathcal B_{Z_t}(y)\le V_t(y)\), and its remaining
normalizing factors are at most \(b^{18k}\) per place. There
are fewer than \(f/4\) such places. In fact this gives the
sharper exponent \( (9/2) kf\), which absorbs the bounded
\(\exp(O(\epsilon))\) within \(b^{5kf}\) at large nodes.
Apply \eqref{eq:identityfrequency929} at \(j=0\) and increase
\(T\) through all allowed finite sets. The actual summands are
nonnegative, independent of the auxiliary set, and every
allowed rational frequency occurs. Monotone convergence proves
the inequality. By Corollary~\ref{cor:realsystem927} the product
of the real bottom forms is \((1+o(1))/D_\pi\).
Thus this estimates the nonzero first frequencies relative to
the mass used by the period comparison. At larger depths
\eqref{eq:identityfrequency929} still uses the finite weights
\(V_T,H_T\); identifying them as majorants for the original
finite integrals at those depths is an additional requirement.
\end{proof}
\clearpage

\subsection{Removing unit prefixes at the standard finite places}
\label{sec:finiteunitreduce930}

Fix one unmarked place \(w\) of the period fields in
Theorem~\ref{thm:allunmarked909}. Write \(O\) for the integers of
\(F_w\), \(q\) for their residue cardinality, and \(A=E\otimes_F F_w\)
with integers \(O_A\). Thus \(A\) is split or unramified quadratic;
use the same uniformizer \(\pi\) as in \(O\). Give \(O_A\) additive
volume one, also in the split case. Let \(G_i,K_i,U_i\) be the
residual isometry group, integral lattice stabilizer and Bessel
unipotent after the first \(i\) hyperbolic pairs, and
\(r_i=k+2(a-i)+1\). Normalize \(K_i\) and \(U_i\cap K_i\) to
have volume one in their respective groups. The finite test
at this level is the scalar
\[
    W_i(uk)=\chi_i(u)\quad (u\in U_i,\ k\in K_i),\qquad
                    W_i=0\ \hbox{off }U_iK_i ,
\]
where \(\chi_i\) is the restriction of the original character.
It is trivial on \(U_i\cap K_i\). For \(\ell\in G_i\), set
\(K_{i,\ell}(y)=\overline{W_i(y)}W_i(y\ell)\). The unconditioned
first form of the residual group is
\[
 b_i(\ell;\nu)=\iint \psi\!\left(\nu\,h(y^{-1}f_i,t^{-1}f_i)\right)
          K_{i,\ell}(y)\overline{K_{i,\ell}(t)}\,d\dot y\,d\dot t ,
\]
using the two quotients by \(U_i\), with \(f_i=e_i\) for \(i<a\)
and \(f_a=z\). Here \(\psi\) is the additive character of \(A\)
with integer-ring annihilator used in Section~\ref{sec:firstglobal909}.
At \(i=0\) this agrees with \eqref{eq:gramintegrand891} at
\(H=J=K_0\): the first amplitudes restrict the two representatives
to \(K_0\), and the last right factor in each conjugation may
be removed by right \(K_0\)-invariance.

\begin{lemma}[Haar density of a list of unit pivots]
\leavevmode\label{lem:finiteunitdensity930}
Take two independent length \(j\le a\) flags from Haar on \(K_0\).
Use the ordered quotient vectors and the ratios of leading
Gram determinants of Section~\ref{sec:orderedflags922}, over
\(A\). On \((O_A^\times)^j\) their list of pivots has constant
density, relative to product additive Haar, given by
\[
 c_{j,w}=\prod_{i<j}\theta_{r_i,w}^{-1},\qquad
 \theta_{d,w}=
 \begin{cases}
 (1-q^{-d})(1-q^{1-d})&A\text{ split},\\
 1+(q-1)(-1)^d q^{-d}-q^{1-2d}&A\text{ unramified field}.
 \end{cases}
\]
At each prescribed unit list the conditional measure on the
compact pair orbit is \(c_{j,w}\) times invariant probability.
\end{lemma}

\begin{proof}
Fix the first vector of one of the flags and take it by an
integral isometry to \(e\). At a null row of dimension \(d\),
work modulo \(\pi^t\), \(t\ge1\). In the unramified field case
the second vector is uniform on the primitive solutions of the
null equation. Their number is
\(q^{t(2d-1)}\theta_{d,w}\) by Lemma~\ref{lem:hermitianorbit900},
whose proof uses the surjective residue trace also in residue
characteristic two. For a prescribed unit \(p\bmod\pi^t\), write
the vector as \(\alpha e+y+p e^*\). Each of the
\(q^{2t(d-2)}\) middle vectors \(y\) leaves one trace equation
for \(\alpha\), with \(q^t\) solutions. All these vectors are
primitive.

In the split case reduction of \(\mathrm{GL}_d(O)\) is onto.
The orbit consists of primitive pairs \((u,\ell)\) of a vector
and a covector with \(\ell(u)=0\) in \(O/\pi^t\). Indeed take
\(u\) to the first coordinate vector, and use the general
linear group of the remaining coordinates on the primitive
covector annihilating it. There are
\(q^{td}(1-q^{-d})q^{t(d-1)}(1-q^{1-d})\) choices. With the
fixed pair equal to the first vector and the last covector, the
two prescribed unit components of \(p\) fix \(u_d\) and \(\ell_1\).
We can choose \(u_1\) and both middle lists freely and solve
for \(\ell_d\) using the unit \(u_d\). This again gives
\(q^{t(2d-3)}\) solutions, all primitive in both components.
In either case the probability for the specified class is
\(q^{-2t}/\theta_{d,w}\). As every such additive class has
measure \(q^{-2t}\), this gives the one-row density on all units.

For the recursion with \(p\) a unit, \(t_0=y/p\) and
\(s=\alpha/p\) are integral and satisfy
\(s+\bar s=-h(t_0,t_0)\). The integral Heisenberg element
fixing \(e\), taking \(e^*\) to \(e^*+t_0+s e\) and taking
\(v\) in the middle space to \(v-h(t_0,v)e\), has integral
inverse. It carries \(p e^*\) to the second vector. Thus one
integral pair section identifies their common orthogonal
complement with the original residual lattice. The integral
stabilizer of one fixed vector projects onto \(K_{i+1}\)
by its middle block, with integral Heisenberg kernel, by the
block formula just used. Conditional on the two first vectors,
each remaining flag therefore has the Haar law of its residual
compact group. They are independent, and changing the two
sections to common integral coordinates preserves each law.
The next relative pivots are those of the Schur complement
after the first column, by successive elimination. Induction
proves the product formula and transitivity of \(K_0\) on
the integral pair orbit at the prescribed list. Invariance
identifies its conditional probability with Haar probability.
The unit-fiber normalization can be taken pointwise: in the
coordinates above the middle entries are free, the trace
equation has a fixed integral linear section and a free
trace-zero coordinate, and divisions by \(p\) are by units.
These are local product charts over the unit \(p\); in the
split case one may instead solve directly for \(\ell_d\).
Consequently the coarea densities are locally constant in
these charts and the finite-precision computation fixes their
normalization at every unit list.
\end{proof}

For clarity, use the two-endpoint form with the same coarea
convention as in Section~\ref{sec:allrealwaves924}. On
\(X_j=(N_{<j}\rtimes G_j)\backslash G_0\) the Haar quotient
of the flag stabilizer is normalized to give its compact
\(K_0\)-orbit mass one. This is the measure obtained by
disintegrating \(U_0\backslash G_0\) using the above measures
on \(U_j\backslash G_j\): all the subgroup intersections with
the indicated compact groups have mass one in the successive
Heisenberg and residual decompositions. Let
\(\mu_{\mathbf p}\) be the coarea of the product of these flag
measures at \(\mathbf p\). For initial pair \(P\) write
\(\tau_P\) for a section taking its normalized columns to
\(e_l,p_l e_l^*\). All pivots in the following theorem are
units. Let \(T=w_{\mathbf p}^{-1}\tau_P\) and
\(d_{\boldsymbol\lambda}=w_{\mathbf q}^{-1}w_{\mathbf p}\) with
\(\lambda_l=p_l/q_l\). At fixed \(P,x\) integrate the residual
wave \(\psi(\nu h(m^{-1}f_j,t^{-1}f_j))\) against the product
\[
 \overline{W_0(m\tau_P)}W_0(m\tau_Px)\,
                W_0(tT)\overline{W_0(t d_{\boldsymbol\lambda}Tx)}
\]
on the two residual quotients \(U_j\backslash G_j\).
Denote that integral by \(C(P,x;\nu)\) and its \(\mu_{\mathbf p}\)
integral by \(J_w(x;\nu)\). Here \(W_0\) denotes the finite
test at depth zero. The endpoint coarea factor relative to
\(\mu_{\mathbf p}\) is one, since every \(p_l,q_l\) has
multiplicative modulus one. Thus the equal-head case is the
original condition on the Gram pivots. The second flag in
the mixed case uses sections \(w_{\mathbf p}^{-1}\tau_P\)
and \(w_{\mathbf q}^{-1}\tau_Px\), retaining its original
two endpoint tests. Right invariance of the quotient measures
and common transport of the residual wave make the integral
independent of the residual choice of pair section.

\begin{theorem}[Reduction of the unit-prefix integral on the whole position group]
\leavevmode\label{thm:finiteunitreduce930}
For prescribed unit lists \(\mathbf p,\mathbf q\) of equal
length \(j\le a\),
\[
 \int_{G_0}\int |C(P,x;\nu)|\,d\mu_{\mathbf p}(P)\,dx
       =c_{j,w}\int_{G_j}|b_j(\ell;\nu)|\,d\ell .
\]
This bounds \(\int|J_w(x;\nu)|\,dx\). If the right side is
finite the signed complex integrals also satisfy
\(\int J_w(x;\nu)\,dx=c_{j,w}\int b_j(\ell;\nu)\,d\ell\).
Both identities hold with the respective positions restricted
to the integral groups on the two sides or to their complements.
\end{theorem}

\begin{proof}
The two initial flags on a contributing pair are in the image
of \(K_0\). This follows from \(W_0\) being supported on
\(U_0 K_0\); the leading scaled columns of the quotient
depend only on the flag. Lemma~\ref{lem:finiteunitdensity930}
allows a common section \(\tau_P\in K_0\) for that pair.
The two final flags are represented by \(x^{-1}A\) and the
columns \(x^{-1}B_l/\lambda_l\), with their earlier column
changes understood as in Proposition~\ref{prop:flagtransfer923}.
They too must lie in the compact flag orbit or the amplitude
is zero. Their pivots are \(\mathbf q\). Choose an integral
section \(\tau_{P'}\) at this final pair. Then
\[
       \ell=\tau_P x\tau_{P'}^{-1}\in G_j,\qquad
 d_{\boldsymbol\lambda}w_{\mathbf p}^{-1}\tau_P x
                    =\ell w_{\mathbf q}^{-1}\tau_{P'} .
\]
Indeed \(\ell\) fixes each of the \(2j\) earlier normalized
columns, and the swaps commute with its middle action.
In particular \(x\in K_0\) if and only if \(\ell\in K_j\).

In restriction to \(G_j\) the original function \(W_0\) equals
\(W_j\). If \(u m\in K_0\) for \(m\in G_j\) and \(u\in U_0\),
its middle block is \(u_jm\) integral, with \(u_j\in U_j\).
This block is itself an isometry of the residual self-dual
lattice; its inverse is integral as well. Conversely
\(m\in U_jK_j\) is already in \(U_0K_0\), and its character
value is \(\chi_j\) by restriction. Every section and Weyl swap
appearing after \(\ell\) in the two branch formulas is integral.
Right \(K_0\)-invariance therefore reduces the entire residual
integral to \(b_j(\ell;\nu)\). This retains the oscillation in
the two quotient variables even on nonintegral positions.

On the initial pair orbit \(d\mu_{\mathbf p}\) restricts to
\(c_{j,w}\) times the probability measure. For fixed \(P\),
disintegrate the position by the Haar quotient \(G_j\backslash
G_0\) of \(\tau_Px\); the residual measure is \(d\ell\).
The coset corresponding to \(P'\) is represented by
\(\tau_{P'}\). The condition of an integral final flag pair is
exactly that this coset meet \(K_0\), by the integral-section
construction. Its quotient measure is one, since \(K_0\)
and \(K_j=K_0\cap G_j\) both have mass one. Integration gives
the asserted formula using Tonelli for absolute values, and
Fubini for the complex integral when finite. The same argument
with the compact position cutoff uses the equivalence
\(x\in K_0\Longleftrightarrow\ell\in K_j\) above.
\end{proof}
\clearpage

\section{References}\label{nav:bibliography-XXVI}
The local factors use the marked Bessel character of Section~\ref{nav:two-mark-flags} and the Haar measures of Section~\ref{nav:marked-open-closed}. The construction and normalization of the real test vectors are in Section~\ref{nav:stationary-cutoff}. The global second moment and its terminal norm constraint are in Sections~\ref{nav:poincare-spanning} and~\ref{nav:spanning-and-sphere}.
\end{document}
